\documentclass[11pt,oneside]{book}
\usepackage[utf8]{inputenc}
\usepackage[T2A]{fontenc}
\usepackage[russian]{babel}
\usepackage{graphicx}
\graphicspath{{target/img/},{target/img/tall/},{target/img/wide/}}

\usepackage{wrapfig}
\usepackage{indentfirst}
\usepackage{url}
\usepackage{ifpdf}
\usepackage{lettrine}

\renewcommand\sectionmark[1]{}
\newcommand{\partornament}{cross}
\newcommand{\ornament}{uzor_begin_10}
\newcommand{\ornamentending}{uzor_end_3}

% calculate header offset from text body
\setlength\headheight{14pt}
\newlength\headoff
\addtolength\headoff{\headheight}
\addtolength\headoff{\headsep}
\addtolength\headoff{14pt}

\title{Молитвослов на всякую потребу}
\author{www.molitvoslov.org}
\date{2011-10-13}
\setcounter{secnumdepth}{-3}
%\setcounter{tocdepth}{-1}

\newcommand{\mypart}[1]{
    \part{#1}
}

\newcommand{\mychapter}[1]{
    \renewcommand{\thechapter}{}
    \chapter{#1}
}

\newcommand{\mychapterz}[2]{
    \renewcommand{\thechapter}{}
    \chapter[#1]{#2}
}

\newcommand\mysubtitle[1]{
    \subsubsection{#1}
}

\newcommand\mysubsection[1]{
    \subsection{#1}
}

\newcommand\myparagraph[1]{
    \paragraph{#1}
}

\newenvironment{mymulticols}[1]{}{}

% the idea of identity environment is to provide a mechanism
% to assign some id to the element it surrounds. Its main use is
% with htlatex configs
\newcommand\myid{}
\newenvironment{identity}{}{}

\newenvironment{Parallel}{}{}

\newenvironment{centericon}{}{}

\newcommand{\myfig}[2][0.43]{}

\newcommand{\myfigh}[3][0.43]{}

\newcommand{\myfigr}[2][0.43]{}

\newcommand{\myfigrh}[3][0.43]{}

\newcommand{\myfigure}[2][0.43]{\myfig[#1]{#2}}

\newcommand{\mychapterending}[1][2.5]{}

\newcommand{\bukvaending}{\mychapterending}

\newcommand\longpage[1][1]{}
\newcommand\shortpage[1][1]{}

% save the value of current parindent to be used
% later in minipage environments
\newlength\myparindent
\setlength\myparindent{\parindent}
\newcommand\noparindent{\setlength\parindent{0pt}}
\newcommand\restoreparindent{\setlength\parindent{\myparindent}}

\long\def\symbolfootnote[#1]#2{\begingroup%
\def\thefootnote{\fnsymbol{footnote}}\footnote[#1]{#2}\endgroup}

\newcommand{\firstletter}[1]{#1}%{\myheadingcolor #1}}

\newcommand\pripev[2][Припев:]{\mysmall{\myemph{#1} #2}}
\newcommand\irmos[1]{\pripev[Ирм\'{о}с:]{#1}}
\newcommand\Bogorodichen[1]{\myemph{Богородичен:} #1}
\newcommand\slavan{Сл\'{а}ва Отц\'{у} и С\'{ы}ну и Свят\'{о}му Д\'{у}ху.}
\newcommand\slava{\mysmall{\slavan}}
\newcommand\inynen{И н\'{ы}не и пр\'{и}сно и во в\'{е}ки век\'{о}в. Ам\'{и}нь.}
\newcommand\inyne{\mysmall{\inynen}}
\newcommand\slavainynen{Сл\'{а}ва Отц\'{у} и С\'{ы}ну и Свят\'{о}му Д\'{у}ху, и н\'{ы}не и пр\'{и}сно и во в\'{е}ки век\'{о}в. Ам\'{и}нь.}
\newcommand\slavainyne{\mysmall{\slavainynen}}

\newcommand{\pripevc}[1]{\mysmall{ \centerline{#1} \nopagebreak}}
\newcommand{\pripevmskipc}[1]{\medskip\pripevc{#1}}
\newcommand{\pripevpomiluj}{\pripevmskipc{\pripev{\firstletter{П}омилуй мя, Боже, помилуй мя.}}}
\newcommand{\slavac}{\pripevmskipc{\slavan}}
\newcommand{\inynec}{\pripevmskipc{\inynen}}

\newcommand{\TsariuNebesnyj}{%
Цар\'{ю} Неб\'{е}сный, Ут\'{е}шителю, Д\'{у}ше \'{и}стины, \'{И}же везд\'{е} сый и вся исполн\'{я}яй, Сокр\'{о}вище благ\'{и}х и ж\'{и}зни Под\'{а}телю, приид\'{и} и всел\'{и}ся в ны, и оч\'{и}сти ны от вс\'{я}кия скв\'{е}рны, и спас\'{и}, Бл\'{а}же, д\'{у}ши н\'{а}ша.}

\newcommand{\TrisviatoePoOtcheNash}{%
Свят\'{ы}й Б\'{о}же, Свят\'{ы}й Кр\'{е}пкий, Свят\'{ы}й Безсм\'{е}ртный, пом\'{и}луй нас. \myemph{(Tрижды)}

Сл\'{а}ва Отц\'{у} и С\'{ы}ну и Свят\'{о}му Д\'{у}ху, и н\'{ы}не и пр\'{и}сно и во в\'{е}ки век\'{о}в. Ам\'{и}нь.

Пресвят\'{а}я Тр\'{о}ице, пом\'{и}луй нас; Г\'{о}споди, оч\'{и}сти грех\'{и} н\'{а}ша; Влад\'{ы}ко, прост\'{и} беззак\'{о}ния н\'{а}ша; Свят\'{ы}й, посет\'{и} и исцел\'{и} н\'{е}мощи н\'{а}ша, \'{и}мене Твоег\'{о} р\'{а}ди.

Г\'{о}споди, пом\'{и}луй. \myemph{(Трижды)}

Сл\'{а}ва Отц\'{у} и С\'{ы}ну и Свят\'{о}му Д\'{у}ху, и н\'{ы}не и пр\'{и}сно и во в\'{е}ки век\'{о}в. Ам\'{и}нь.

\'{О}тче наш, \'{И}же ес\'{и} на небес\'{е}х! Да свят\'{и}тся \'{и}мя Тво\'{е}, да при\'{и}дет Ц\'{а}рствие Тво\'{е}, да б\'{у}дет в\'{о}ля Тво\'{я}, \'{я}ко на небес\'{и} и на земл\'{и}. Хлеб наш нас\'{у}щный даждь нам днесь; и ост\'{а}ви нам д\'{о}лги н\'{а}ша, \'{я}коже и мы оставл\'{я}ем должник\'{о}м н\'{а}шим; и не введ\'{и} нас во искуш\'{е}ние, но изб\'{а}ви нас от лук\'{а}ваго.
}

\newcommand{\priiditepoklonimsia}{%
Приид\'{и}те, поклон\'{и}мся Цар\'{е}ви н\'{а}\-ше\-му Б\'{о}гу. \myemph{(Поклон)}

Приид\'{и}те, поклон\'{и}мся и припад\'{е}м Христ\'{у}, Цар\'{е}ви н\'{а}шему Б\'{о}гу. \myemph{(Поклон)}

Приид\'{и}те, поклон\'{и}мся и припад\'{е}м Самом\'{у} Христ\'{у}, Цар\'{е}ви и Б\'{о}гу н\'{а}шему. \myemph{(Поклон)}}


\newcommand{\PsalmFifty}{%
Пом\'{и}луй мя, Б\'{о}же, по вел\'{и}цей м\'{и}лости Тво\'{е}й, и по мн\'{о}жеству щедр\'{о}т Тво\'{и}х оч\'{и}сти беззак\'{о}ние мо\'{е}. Наип\'{а}че ом\'{ы}й мя от беззак\'{о}ния моег\'{о}, и от грех\'{а} моег\'{о} оч\'{и}сти мя; яко беззак\'{о}ние мо\'{е} аз зн\'{а}ю, и грех мой пр\'{е}до мн\'{о}ю есть в\'{ы}ну. Теб\'{е} Ед\'{и}ному согреш\'{и}х и лук\'{а}вое пред Тоб\'{о}ю сотвор\'{и}х, \'{я}ко да оправд\'{и}шися во словес\'{е}х Тво\'{и}х, и побед\'{и}ши внегд\'{а} суд\'{и}ти Ти. Се бо, в беззак\'{о}ниих зач\'{а}т есмь, и во грес\'{е}х род\'{и} мя м\'{а}ти мо\'{я}. Се бо, \'{и}стину возлюб\'{и}л ес\'{и}; безв\'{е}стная и т\'{а}йная прем\'{у}дрости Твое\'{я} яв\'{и}л ми ес\'{и}. Окроп\'{и}ши мя исс\'{о}пом, и оч\'{и}щуся; ом\'{ы}еши мя, и п\'{а}че сн\'{е}га убел\'{ю}ся. Сл\'{у}ху моем\'{у} д\'{а}си р\'{а}дость и вес\'{е}лие; возр\'{а}дуются к\'{о}сти смир\'{е}нныя. Отврат\'{и} лиц\'{е} Тво\'{е} от грех мо\'{и}х и вся беззак\'{о}ния мо\'{я} оч\'{и}сти. С\'{е}рдце ч\'{и}сто соз\'{и}жди во мне, Б\'{о}же, и дух прав обнов\'{и} во утр\'{о}бе мо\'{е}й. Не отв\'{е}ржи мен\'{е} от лиц\'{а} Твоег\'{о} и Д\'{у}ха Твоег\'{о} Свят\'{а}го не отым\'{и} от мен\'{е}. Возд\'{а}ждь ми р\'{а}дость спас\'{е}ния Тво\'{е}го и Д\'{у}хом влад\'{ы}чним утверд\'{и} мя. Науч\'{у} беззак\'{о}ныя пут\'{е}м Тво\'{и}м, и нечест\'{и}вии к Теб\'{е} обрат\'{я}тся. Изб\'{а}ви мя от кров\'{е}й, Б\'{о}же, Б\'{о}же спас\'{е}ния моег\'{о}; возр\'{а}дуется яз\'{ы}к мой пр\'{а}вде Тво\'{е}й. Г\'{о}споди, устн\'{е} мои отв\'{е}рзеши, и уст\'{а} мо\'{я} возвест\'{я}т хвал\'{у} Тво\'{ю}. \'{Я}ко \'{а}ще бы восхот\'{е}л ес\'{и} ж\'{е}ртвы, дал бых \'{у}бо: всесожж\'{е}ния не благовол\'{и}ши. Ж\'{е}ртва Б\'{о}гу дух сокруш\'{е}н; с\'{е}рдце сокруш\'{е}нно и смир\'{е}нно Бог не уничиж\'{и}т. Ублаж\'{и}, Г\'{о}споди, благовол\'{е}нием Тво\'{и}м Си\'{о}на, и да соз\'{и}ждутся ст\'{е}ны Иерусал\'{и}мския. Тогд\'{а} благовол\'{и}ши ж\'{е}ртву пр\'{а}вды, вознош\'{е}ние и всесожег\'{а}емая; тогд\'{а} возлож\'{а}т на oлт\'{а}рь Твой тельц\'{ы}.\par}

\newcommand{\PsalmNinety}{%
Жив\'{ы}й в п\'{о}мощи В\'{ы}шняго, в кр\'{о}ве Б\'{о}га Неб\'{е}снаго водвор\'{и}тся. Реч\'{е}т Г\'{о}сподеви: Заст\'{у}пник мой ес\'{и} и Приб\'{е}жище мо\'{е}, Бог мой, и упов\'{а}ю на Нег\'{о}. \'{Я}ко Той изб\'{а}вит тя от с\'{е}ти л\'{о}вчи и от словес\'{е} мят\'{е}жна, плещм\'{а} Сво\'{и}ма осен\'{и}т тя, и под крил\'{е} Ег\'{о} над\'{е}ешися: ор\'{у}жием об\'{ы}дет тя \'{и}стина Ег\'{о}. Не убо\'{и}шися от стр\'{а}ха нощн\'{а}го, от стрел\'{ы} лет\'{я}щия во дни, от в\'{е}щи во тме преход\'{я}щия, от ср\'{я}ща и б\'{е}са пол\'{у}деннаго. Пад\'{е}т от стран\'{ы} твое\'{я} т\'{ы}сяща, и тма одесн\'{у}ю теб\'{е}, к теб\'{е} же не прибл\'{и}жится, об\'{а}че оч\'{и}ма тво\'{и}ма см\'{о}триши, и возда\'{я}ние гр\'{е}шников \'{у}зриши. \'{Я}ко Ты, Г\'{о}споди, упов\'{а}ние мо\'{е}, В\'{ы}шняго полож\'{и}л ес\'{и} приб\'{е}жище тво\'{е}. Не при\'{и}дет к теб\'{е} зло и р\'{а}на не прибл\'{и}жится телес\'{и} твоем\'{у}, \'{я}ко \'{А}нгелом Сво\'{и}м запов\'{е}сть о теб\'{е}, сохран\'{и}ти тя во всех пут\'{е}х тво\'{и}х. На рук\'{а}х в\'{о}змут тя, да не когд\'{а} преткн\'{е}ши о к\'{а}мень н\'{о}гу тво\'{ю}, на \'{а}спида и васил\'{и}ска наст\'{у}пиши, и попер\'{е}ши льва и зм\'{и}я. \'{Я}ко на Мя упов\'{а} и изб\'{а}влю \'{и}, покр\'{ы}ю \'{и}, \'{я}ко позн\'{а} \'{и}мя Мо\'{е}. Воззов\'{е}т ко Мне и усл\'{ы}шу eг\'{о}, с ним есмь в ск\'{о}рби, изм\'{у} eг\'{о}, и просл\'{а}влю eг\'{о}, долгот\'{о}ю дний исп\'{о}лню eг\'{о}, и явл\'{ю} eм\'{у} спас\'{е}ние Мо\'{е}.\par}

\newcommand{\Predstatelstvo}{%
Предст\'{а}тельство христи\'{а}н непост\'{ы}дное, ход\'{а}тайство ко Творц\'{у} непрел\'{о}жное, не пр\'{е}зри гр\'{е}шных мол\'{е}ний гл\'{а}сы, но предвар\'{и}, \'{я}ко Благ\'{а}я, на п\'{о}мощь нас, в\'{е}рно зов\'{у}щих Ти: ускор\'{и} на мол\'{и}тву, и потщ\'{и}ся на умол\'{е}ние, предст\'{а}тельствующи пр\'{и}сно, Богор\'{о}дице, чт\'{у}щих Тя.}

\newcommand{\Chestneyshuyu}{%
Дост\'{о}йно \'{е}сть \'{я}ко во\'{и}стинну бла\-ж\'{и}\-ти Тя, Богор\'{о}дицу, Присноблаж\'{е}нную и Пренепор\'{о}чную и М\'{а}терь Б\'{о}га н\'{а}шего. Честн\'{е}йшую Херув\'{и}м и сл\'{а}внейшую без сравн\'{е}ния Сераф\'{и}м, без истл\'{е}ния Б\'{о}га Сл\'{о}ва р\'{о}ждшую, с\'{у}щую Богор\'{о}дицу Тя велич\'{а}ем.}

\newcommand{\MolitvamiSviatyhOtecNashih}{%
Мол\'{и}твами свят\'{ы}х от\'{е}ц н\'{а}ших, Г\'{о}споди Иис\'{у}се Христ\'{е}, Б\'{о}же наш, пом\'{и}луй нас. Ам\'{и}нь.}

\newcommand{\LMolitvamiSviatyhOtecNashih}{%
\lettrine{М}{}ол\'{и}твами свят\'{ы}х от\'{е}ц н\'{а}ших, Г\'{о}споди Иис\'{у}се Христ\'{е}, Б\'{о}же наш, пом\'{и}луй нас. Ам\'{и}нь.}

\newcommand{\tolkopoblagosloveniyu}{%
{\centering\myemph{\normalfont Читаются только по благословению духовника}\par}}

\newcommand{\tolkosviashennikom}{%
{\centering\myemph{\normalfont Молитва читается священником}\par}}

\newcommand{\SymbolOfFaith}{%
В\'{е}рую во ед\'{и}наго Б\'{о}га Отц\'{а}, Вседерж\'{и}теля, Творц\'{а} н\'{е}бу и земл\'{и}, в\'{и}димым же всем и нев\'{и}димым.
И во ед\'{и}наго Г\'{о}спода Иис\'{у}са Христ\'{а}, С\'{ы}на Б\'{о}жия, Единор\'{о}днаго, \'{И}же от Отц\'{а} рожд\'{е}ннаго пр\'{е}жде всех век; Св\'{е}та от Св\'{е}та, Б\'{о}га \'{и}стинна от Б\'{о}га \'{и}стинна, рожд\'{е}нна, несотвор\'{е}нна, единос\'{у}щна Отц\'{у}, \'{И}мже вся б\'{ы}ша.
Нас р\'{а}ди челов\'{е}к и н\'{а}шего р\'{а}ди спас\'{е}ния сш\'{е}дшаго с неб\'{е}с и воплот\'{и}вшагося от Д\'{у}ха Св\'{я}та и Мар\'{и}и Д\'{е}вы и вочелов\'{е}чшася.
Расп\'{я}таго же за ны при Понт\'{и}йстем Пил\'{а}те, и страд\'{а}вша, и погреб\'{е}нна.
И воскр\'{е}сшаго в тр\'{е}тий день по Пис\'{а}нием.
И возш\'{е}дшаго на небес\'{а}, и сед\'{я}ща одесн\'{у}ю Отц\'{а}.
И п\'{а}ки гряд\'{у}щаго со сл\'{а}вою суд\'{и}ти жив\'{ы}м и м\'{е}ртвым, Ег\'{о}же Ц\'{а}рствию не б\'{у}дет конц\'{а}.
И в Д\'{у}ха Свят\'{а}го, Г\'{о}спода, Животвор\'{я}щаго, \'{И}же от Отц\'{а} исход\'{я}щаго, \'{И}же со Отц\'{е}м и С\'{ы}ном споклан\'{я}ема и ссл\'{а}вима, глаг\'{о}лавшаго прор\'{о}ки.
Во ед\'{и}ну Свят\'{у}ю, Соб\'{о}рную и Ап\'{о}стольскую Ц\'{е}рковь.
Испов\'{е}дую ед\'{и}но крещ\'{е}ние во оставл\'{е}ние грех\'{о}в.
Ч\'{а}ю воскрес\'{е}ния м\'{е}ртвых, и ж\'{и}зни б\'{у}дущаго в\'{е}ка. Ам\'{и}нь.}

\newcommand{\TroparPomilujNas}{Пом\'{и}луй нас, Г\'{о}споди, пом\'{и}луй нас; вс\'{я}каго бо отв\'{е}та недоум\'{е}юще, си\'{ю} Ти мол\'{и}тву \'{я}ко Влад\'{ы}це гр\'{е}шнии прин\'{о}сим: пом\'{и}луй нас.

\slavan

Г\'{о}споди, пом\'{и}луй нас, на Тя бо упов\'{а}хом; не прогн\'{е}вайся на ны зел\'{о}, ниж\'{е} помян\'{и} беззак\'{о}ний н\'{а}ших, но пр\'{и}зри и н\'{ы}не \'{я}ко благоутр\'{о}бен, и изб\'{а}ви ны от враг н\'{а}ших; Ты бо ес\'{и} Бог наш, и мы л\'{ю}дие Тво\'{и}, вс\'{и} дел\'{а} р\'{у}ку Тво\'{е}ю, и \'{и}мя Тво\'{е} призыв\'{а}ем.

\inynen

Милос\'{е}рдия дв\'{е}ри отв\'{е}рзи нам, благослов\'{е}нная Богор\'{о}дице, над\'{е}ющиися на Тя да не пог\'{и}бнем, но да изб\'{а}вимся Тоб\'{о}ю от бед: Ты бо ес\'{и} спас\'{е}ние р\'{о}да христи\'{а}нскаго.}



\renewcommand{\pripevc}[1]{#1}

\begin{document}



\mypart{ЕДИНОВЕРЧЕСКИЕ МОЛИТВЫ}\label{_content_edinovercheskie-molitvi}
%http://www.molitvoslov.org/content/edinovercheskie-molitvi

 

\mychapter{Молитвенное правило}
%http://www.molitvoslov.org/content/molitvennoepravilo

 

\section{Молитвы утренние и вечерние}\begin{mymulticols}
%http://www.molitvoslov.org/text1000.htm 
 


\myemph{ Утром, вставши от сна, еще на постели, перекрестись с молитвою:}


Господи, Исусе Христе, Сыне Божии, помилуй мя грешнаго.


\myemph{ Вставши с постели и умыв лице свое, а вечером, отходя ко сну, встань с благоговением пред святыми иконами и, взирая на них, мысль свою устремляя к невидимому Богу и его святым, истово, не спеша, ограждая себя крестным знамением и кланяясь, произноси с умилением молитву мытаря:}


Боже, милостив буди мне грешному \myemph{(поклон)}. Создавыи мя, Господи, и помилуй мя \myemph{(поклон)}. Без числа согреших, Господи, помилуй и прости мя, грешнаго \myemph{(поклон)}.


Достойно есть, яко воистину блажити Тя Богородице, присно блаженную и пренепорочную, и Матерь Бога нашего. Честнейшую херувим и славнейшую воистину серафим, без истления Бога Слова рождьшую, сущую Богородицу, Тя величаем \myemph{ (поклон всегда земной)}.


Слава Отцу и Сыну и Святому Духу \myemph{(поклон)}. И ныне и присно и во веки веком, аминь \myemph{(поклон)}.


Господи помилуй, Господи помилуй, Господи благослови \myemph{(поклон)}. 


Господи, Исусе Христе, Сыне Божии, молитв ради Пречистыя Твоея Матере, силою Честнаго и Животворящаго Креста, и святаго Ангела моего хранителя, и всех ради святых, помилуй и спаси мя грешнаго, яко Благ и Человеколюбец. Аминь. \myemph{ (поклон земной, без крестнаго знамения)}.


\myemph{ Эти молитвы называются «началом» или «приходными и исходными поклонами» потому, что они совершаются в начале и после всякого молитвенного правила.}


После этого повтори молитву мытаря с поклонами:


Боже, милостив буди мне грешному \myemph{(поклон)}. Создавыи мя, Господи, и помилуй мя \myemph{(поклон)}. Без числа согреших, Господи, помилуй и прости мя грешнаго \myemph{(поклон)}. 


\myemph{ И с благоговением начинай утренние молитвы.}


За молитв святых отец наших, Господи Исусе Христе, Сыне Божии, помилуй нас. Аминь \myemph{ (поклон всегда поясной). Перекрестись и говори трижды:}


Слава Тебе, Боже наш, слава Тебе всяческих ради. 


\myemph{ Далее:} Боже, очисти мя грешнаго, яко николиже благо сотворих пред Тобою \myemph{(поклон)}, но избави мя от лукаваго, и да будет во мне воля Твоя \myemph{(поклон)}, да неосужденно отверзу уста моя недостойная и восхвалю имя Твое святое: Отца и Сына и Святаго Духа, ныне и присно и во веки веком. Аминь \myemph{(поклон)}.

\myemph{ Обведенные чертой молитвы вечером не читаются.} 


\myemph{ Далее перекрестись и читай: }


Царю небесныи, Утешителю, Душе истинныи, Иже везде сыи и вся исполняя, Сокровище благих и жизни Подателю, прииди и вселися в ны, и очисти ны от всякия скверны, и спаси, Блаже, душа наша.


Святыи Боже, Святыи Крепкии, Святыи Безсмертныи, помилуй нас \myemph{ (трижды с поклонами)}. Слава Отцу и Сыну и Святому Духу, и ныне и присно и во веки веком. Аминь. Пресвятая Троице, помилуй нас. Господи, очисти грехи наша. Владыко, прости беззакония наша. Святыи, посети и исцели немощи наша, Имене Твоего ради. Господи помилуй \myemph{(трижды)}. Слава Отцу и Сыну и Святому Духу, и ныне и присно и во веки веком. Аминь. Отче наш, Иже еси на небесех, да святится имя Твое, да приидет царствие Твое, да будет воля Твоя, яко на небеси и на земли. Хлеб наш насущныи даждь нам днесь. И остави нам долги наша, яко же и мы оставляем должником нашим. И не введи нас во искушение, но избави нас от лукаваго. Господи Исусе Христе Сыне Божии, помилуй нас. Аминь. Господи помилуй \myemph{ (12 раз)}. 


\myemph{ Если утро "--- читается: }


От сна востав, благодарю Тя, Всесвятая Троице, яко многия ради благости и долготерпения не прогневася на мя грешнаго и лениваго раба Твоего, и не погубил еси мене со беззаконии моими, но человеколюбствова. И в нечаянии лежаща, воздвиже мя утреневати и славословити державу Твою непобедимую. И ныне, Владыко, Боже Пресвятыи, просвети очи сердца моего и отверзи ми устне поучатися словесем Твоим, и разумети заповеди Твоя, и творити волю Твою, и пети Тя во исповедании сердечнем. Воспевати же и славити пречестное и великолепое имя Твое: Отца и Сына и Святаго Духа, ныне и присно и во веки веком. Аминь.


\myemph{ Если вечер:}


Слава Отцу и Сыну и Святому Духу. И ныне, и присно, и во веки веком. Аминь.


\myemph{ Далее, как утром, так и вечером, читай:}


Приидите, поклонимся Цареви нашему Богу \myemph{(поклон)}. Приидите, поклонимся Христу, Цареви и Богу нашему \myemph{(поклон)}. Приидите, поклонимся и припадем к самому Господу Исусу Христу, Цареви и Богу нашему \myemph{(поклон)}. 


\mysubtitle{ПСАЛОМ 50 (покаянный)}


Помилуй мя, Боже, по велицей милости Твоей. И по множеству щедрот Твоих очисти беззаконие мое. Наипаче омыи мя от беззакония моего и от греха моего очисти мя. Яко беззаконие мое аз знаю и грех мой предо мною есть выну. Тебе единому согреших и лукавое пред Тобою сотворих. Яко да оправдишися в словесех Своих и победиши внегда судити. Се бо в беззакониих зачат есмь и во гресех роди мя мати моя. Се бо истину возлюбил еси; безвестная и тайная премудрости Твоея явил ми еси. Окропиши мя иссопом и очищуся. Омыеши мя и паче снега убелюся. Слуху моему даси радость и веселие: возрадуются кости смиренныя. Отврати лице Твое от грех моих и вся беззакония моя очисти. Сердце чисто созижди во мне, Боже, и дух прав обнови во утробе моей. Не отверзи мене от Лица Твоего, и Духа Твоего Святаго не отыми от мене. Воздаждь ми радость спасения Твоего и духом владычным утверди мя. Научу беззаконныя путем Твоим и нечестивии к Тебе обратятся. Избави мя от кровий, Боже, Боже спасения моего; возрадуется язык мой правде Твоей. Господи, устне мои отверзеши, и уста моя возвестят хвалу Твою. Яко аще бы восхотел жертвы, дал бых убо; всесожжения не благоволиши. Жертва Богу "--- дух сокрушен: сердце сокрушенно и смиренно Бог не уничижит. Ублажи, Господи, благоволением Твоим Сиона; и да созиждутся стены Иеросалимския. Тогда благоволиши жертву правды, возношение и всесожигаемая. Тогда возложат на олтарь Твой тельца. 

\myemph{ Ограждая себя благоговейно крестным знамением, произносим СИМВОЛ ВЕРЫ "--- слова святых отец Первого и Второго Вселенского собора (крестное знамение без поклона): }


Верую во единаго Бога Отца Вседержителя, Творца небу и земли, видимым же всем и невидимым. И во единаго Господа, Исуса Христа, Сына Божия, Единороднаго, Иже от Отца рожденнаго прежде всех век. Света от Света, Бога истинна от Бога истинна, рождена, а не сотворена, единосущна Отцу, Им же вся быша. Нас ради человек, и нашего ради спасения сшедшаго с небес, и воплотившагося от Духа Свята и Марии Девы вочеловечшася. Распятаго за ны при Понтийстем Пилате, страдавша и погребенна. И воскресшаго в третии день по писаниих. И возшедшаго на небеса, и седяща одесную Отца. И паки грядущаго со славою судити живым и мертвым, Его же царствию несть конца. И в Духа Святаго, Господа истиннаго и Животворящаго, Иже от Отца исходящаго, иже со Отцем и Сыном споклоняема и сславима, глаголавшаго пророки. И во едину святую соборную и апостольскую Церковь. Исповедую едино Крещение во оставление грехов. Чаю воскресения мертвым. И жизни будущаго века. Аминь.


\myemph{ Далее: }


Богородице Дево, радуися, обрадованная Марие, Господь с Тобою, благословенна Ты в женах и благословен Плод чрева Твоего, яко родила еси Христа Спаса, Избавителя душам нашим \myemph{ (с поклонами трижды)}.


О! Всепетая Мати, рождьшая всех святых святейшее Слово, нынешнее приношение приемши, от всякия напасти избави всех, и грядущия изми муки, вопиющия Ти: Аллилуия \myemph{ (трижды, с поклонами земными)}.


Непобедимая и божественная сила честнаго и животворящаго Креста Господня, не остави мя грешнаго, уповающаго на Тя \myemph{(поклон)}. Пресвятая Владычице моя Богородице, помилуй мя, и спаси мя, и помози ми ныне, в жизни сей, и во исход души моея, и в будущем веце \myemph{(поклон)}. Вся небесныя силы, святии ангели и архангели, херувими и серафими, помилуйте мя, и помолитеся о мне грешнем ко Господу Богу, и помозите ми ныне, в жизни сей, и во исход души моея, и в будущем веце \myemph{(поклон)}. Ангеле Христов, хранителю мой святыи, помилуй мя и помолися о мне грешнем ко Господу Богу, и помози ми ныне, в жизни сей, и во исход души моея, и в будущем веце \myemph{(поклон)}. Святыи великии Иоанне, пророче и Предотече Господень, помилуй мя, и помолися о мне грешнем ко Господу Богу, и помози ми ныне, в жизни сей, и во исход души моея, и в будущем веце \myemph{(поклон)}. Святии славнии апостоли, пророцы и мученицы, святителие, преподобнии и праведнии и вси святии, помилуйте мя, и помолитеся о мне грешнем ко Господу Богу, и помозите ми ныне, в жизни сей, и во исход души моея, и в будущем веце \myemph{(поклон)}. 


\myemph{ После сего молись по три раза с поклонами следующие молитвы.}


Пресвятая Троице Боже наш, слава Тебе.Господи Исусе Христе, Сыне Божии, помилуй мя грешнаго. Слава, Господи, Кресту Твоему Честному. Пресвятая Госпоже Богородице, спаси мя грешнаго раба Твоего. Ангеле Христов, хранителю мой святыи, спаси мя грешнаго раба Твоего. Святии архангели и ангели, молите Бога о мне грешнем. Святыи великии Иоанне, пророче и Предотече, Крестителю Господень, моли Бога о мне грешнем. Святыи славныи пророче Илие, моли Бога о мне грешнем. Святии праотцы, молите Бога о мне грешнем. Святии пророцы, молите Бога о мне грешнем. Святии апостоли, молите Бога о мне грешнем. Святии славнии апостоли и евангелисти: Матфее, Марко, Луко и Иоанне Богослове, молите Бога о мне грешнем. Святии славнии верховнии апостоли Петре и Павле, молите Бога о мне грешнем. Святии велиции трие святителие: Василие Великии, Григорие Богослове и Иоанне Златоусте, молите Бога о мне грешнем. Святителю Христов Николае, моли Бога о мне грешнем. Преподобныи отче Сергие, моли Бога о мне грешнем. Святыи священномучениче и исповедниче Аввакуме, моли Бога о мне грешнем. Святителю Христов и исповедниче Амвросие, моли Бога о мне грешнем. Преподобнии и богоноснии отцы наши, пастырие и учителие вселенней, молите Бога о мне грешнем. Святии вси, молите Бога о мне грешнем.


\myemph{ После этого помолись святому, имя которого носишь, и святому, празднуемому этого числа, также и другим святым, кому хочешь. Не забудь помолиться и епитимию, какие поклоны имеешь от духовного отца.}


Потом молись о здравии правящего епископа, духовного отца, родителей, родственников и близких, по три раза произнося с поклонами о здравии и о спасении: 


Милостиве Господи, спаси и помилуй раб Своих \myemph{ \myemph{(поклон)} (называй имена, за кого молишься)}. Избави их от всякия скорби, гнева и нужды \myemph{(поклон)}. От всякия болезни душевныя и телесныя \myemph{(поклон)}. И прости им всякое согрешение, вольное и невольное \myemph{(поклон)}. И душам нашим полезная сотвори \myemph{(поклон)}. 


\myemph{ Потом помолись за упокой духовных отцов, родителей и близких, и за кого имеешь усердие, по три раза произнося с поклонами: }


Покой, Господи, душа усопших раб Твоих \myemph{ \myemph{(поклон)} (называй имена, за кого молишься)}. И елика в житии сем яко человецы согрешиша, Ты же, яко Человеколюбец Бог, прости их и помилуй \myemph{(поклон)}. Вечныя муки избави \myemph{(поклон)}. Небесному царствию причастники учини \myemph{(поклон)}. И душам нашим полезная сотвори \myemph{(поклон)}.


\myemph{ Заканчивая молитвы, говори: }


Господи, или словом, или делом, или помышлением согреших во всей жизни моей, помилуй мя и прости мя, милости Твоея ради \myemph{ (поклон земной)}. Все упование мое к Тебе возлагаю, Мати Божия, сохрани мя во Своем си крове \myemph{ (поклон земной)}. Упование ми Бог, и прибежище мое Христос, и покровитель ми есть Дух Святыи \myemph{ (поклон земной)}.


Достойно есть, яко воистину блажити Тя Богородице, присно блаженную и пренепорочную, и Матерь Бога нашего. Честнейшую херувим и славнейшую воистину серафим, без истления Бога Слова рождьшую, сущую Богородицу, Тя величаем \myemph{ (поклон всегда земной)}.


Слава Отцу и Сыну и Святому Духу \myemph{(поклон)}.И ныне и присно и во веки веком, аминь \myemph{(поклон)}. Господи помилуй, Господи помилуй, Господи благослови \myemph{(поклон)}. 


\mysubtitle{И отпуст:}


Господи, Исусе Христе, Сыне Божии, молитв ради Пречистыя Твоея Матере и преподобных и богоносных отец наших и всех святых, помилуй и спаси мя грешнаго, яко Благ и Человеколюбец. Аминь.


\myemph{ И, наклонившись до земли, без осенения себя крестным знамением, читай прощение: }


Ослаби, остави, отпусти, Боже, согрешения моя, вольная и невольная, яже в слове и в деле, и яже в ведении и не в ведении, яже во уме и в помышлении, яже во дни и в нощи, вся ми прости, яко Благ и Человеколюбец. Аминь. 


\myemph{ Вставши, читай молитву сию с поклонами:}


Ненавидящих и обидящих нас прости, Господи Человеколюбче. Благотворящим благо сотвори, братиям и всем сродником нашим, иже и уединившимся, даруй им вся, яже ко спасению прошения и живот вечныи \myemph{(поклон)}. В болезнех сущия посети и исцели, в темницах сущих свободи, по водам плавающим Правитель буди и иже в путех шествующим, исправи и поспеши \myemph{(поклон)}.Помяни, Господи, и плененныя братию нашу, единоверных православныя веры, и избави их всякаго злаго обстояния \myemph{(поклон)}. Помилуй, Господи, давших нам милостыню и заповедавших нам, недостойным, молитися о них, прости их и помилуй \myemph{(поклон)}. Помилуй, Господи, труждающихся и служащих нам, милующих и питающих нас, и даруй им вся, яже ко спасению, прошения и живот вечныи \myemph{(поклон)}. Помяни, Господи, прежде отшедшия отцы и братию нашу и всели их, идеже присещает свет лица Твоего \myemph{(поклон)}. Помяни, Господи, и нашу худость и убожество, и просвети наш ум светом разума святаго Евангелия Твоего, и настави нас на стезю заповедей Твоих, молитвами Пречистыя Твоея Матере и всех святых Твоих, аминь \myemph{(поклон)}.


\myemph{ Оканчиваются сии молитвы обычным семипоклонным началом (смотри «приходные и исходные поклоны» в начале). }


По окончании молитв, как утром, так и вечером, ограждаясь своим нательным крестом, говори: Господи, Исусе Христе, Сыне Божии, благослови и освяти, и сохрани мя силою Живоноснаго Креста Твоего. 


\myemph{ После этого целуй крест.}


\myemph{ И читай молитву Кресту, перекрестясь:}


Да воскреснет Бог, и разыдутся врази Его, и да бежат от лица Его ненавидящии Его, яко исчезает дым, да исчезнут. Яко тает воск от лица огня, тако да погибнут беси от лица любящих Бога, и знаменающихся крестным знамением, и да возвеселимся рекуще: радуися, Кресте Господень, прогоняя бесы силою на Тебе пропятаго Господа нашего Исуса Христа, во ад сшедшаго, и поправшаго силу диаволю, и давшаго нам Крест Свой Честныи на прогнание всякаго супостата.


О! Пречестныи и Животворящии Кресте Господень, помогай ми, с Пресвятою Госпожею Богородицею, и со всеми святыми небесными силами, всегда и ныне и присно и во веки веком. Аминь.


\end{mymulticols}

\section{Чин чтения 12-ти псалмов}\begin{mymulticols}
%http://www.molitvoslov.org/text1002.htm 
 


Чин, како подобает пети дванадесять псалмов особь, ихже пояху преподобнии отцы пустыннии во дни и в нощи, о нихже вспоминается в книгах отеческих и в житиях и мучениих святых многих.


Сей же чин принесе от святыя горы преподобный Досифей архимандрит Киево-Печерский.


\myemph{ По обычнем начале глаголем: }


За молитв святых отец наших, Господи, Исусе Христе, Сыне Божии, помилуй нас, аминь. \myemph{(поклон)}


\myemph{ Ограждая лице свое, читаем:}


Царю небесныи, Утешителю, Душе истинныи, Иже везде сыи и вся исполняя, сокровище благих и жизни Подателю, прииди и вселися в ны, и очисти ны от всякия скверны, и спаси Блаже, душа наша.

Святыи Боже, Святыи Крепкии, Святыи Безсмертныи, помилуй нас. \myemph{(трижды с поклонами)}.

Слава Отцу и Сыну и Святому Духу, и ныне и присно и во веки веком, аминь.

Пресвятая Троице, помилуй нас. Господи, очисти грехи наша; Владыко, прости беззакония наша; Святыи, посети и исцели немощи наша; имене Твоего ради.

Господи помилуй. \myemph{(трижды)}

Слава Отцу и Сыну и Святому Духу, и ныне и присно и во веки веком, аминь.

Отче наш, Иже еси на небесех; да святится имя Твое; да приидет царствие Твое; да будет воля Твоя, яко на небеси и на земли; хлеб наш насущныи даждь нам днесь; и остави нам долги наша, яко же и мы оставляем должником нашим; и не введи нас во искушение; но избави нас от лукаваго.

Господи Исусе Христе, Сыне Божии, помилуй нас.

Аминь.

Господи, помилуй, 12.

Слава Отцу и Сыну и Святому Духу, и ныне и присно и во веки веком, аминь.

Приидите, поклонимся Цареви нашему Богу \myemph{(поклон)}.

Приидите, поклонимся Христу, Цареви и Богу нашему \myemph{(поклон)}.

Приидите, поклонимся и припадем к Самому Господу Исусу Христу, Цареви и Богу нашему \myemph{(поклон)}.


\mysubtitle{Псалом 26}


Господь просвещение мое и Спаситель мой, кого ся убою? Господь Защититель животу моему, от кого ся устрашу? Внегда приближатися на мя злобующе, снести плотии моих, оскорбляющии мя, и врази мои, тии изнемогоша и падоша. Аще ополчится на мя полк, не убоится сердце мое, аще востанет на мя брань, на Нань аз уповаю. Едино просих от Господа, то взыщу: еже жити ми в дому Господни вся дни живота моего, зрети ми красоту Господню и посещати церковь святую Его. Яко скры мя в крове Своем в день зла моего, покры мя в тайне крова Своего. На камень вознесе мя, и ныне се вознесе главу мою на враги моя, обыдох и пожрох в крове Его жертву хваления и воскликновения, пою и воспою Господеви. Услыши, Господи, глас мой, им же воззвах, помилуй мя и услыши мя. Тебе рече сердце мое: Господа взыщу. Взыска Тебе лице мое, лица Твоего, Господи, взыщу. Не отврати лица Твоего от мене и не уклонися гневом от раба Твоего, помощник ми буди, не отрини мене, и не остави мене, Боже Спасителю мой. Яко отец мой и мати моя остависта мя, Господь же восприят мя. Законоположи ми, Господи, в пути Твоем и настави мя на путь правыи,  враг моих ради. Не предаждь мене в душах стужающим ми, яко восташа на мя свидетели неправеднии, и солга неправда себе. Верую видети благая Господня на земли живых. Потерпи Господа, мужайся, и да крепится сердце твое, и потерпи Господа.

\mysubtitle{Псалом 31}


Блажени, им же отпустишася беззакония, и им же прикрышася греси. Блажен муж, ему же не вменит Господь греха, ниже есть во устех его лесть. Яко умолчах, обетшаша кости моя, зовущу ми весь день. Яко день и нощь отяготе на мне рука Твоя, возвратихся на страсть, егда унзе ми терн. Беззаконие мое познах и греха моего не покрых, рех: исповем на мя беззаконие мое Господеви, и Ты отпустил еси нечестие сердца моего. За то помолится к Тебе всяк преподобныи во время благопотребно, обаче в потопе вод мног, к нему не приближатся. Ты еси прибежище мое, от скорби одержащия мя, радосте моя, избави мя от обышедших мя. Вразумлю тя и наставлю тя на путь сей, вонь же пойдеши, утвержу на тя очи Мои. Не будите яко конь имск, имже несть разума, броздами и уздою челюсти их востягнеши, неприближающимся к тебе. Многи раны грешному, уповающаго же на Господа милость обыдет. Веселитеся о Господе, и радуйтеся, праведнии, и хвалитеся, вси правии сердцем.


\mysubtitle{Псалом 56}


Помилуй мя, Боже, помилуй мя: яко на Тя упова душа моя. И на сень крилу Твоею надеюся, дондеже прейдет беззаконие. Воззову к Богу Вышнему, Богу, благодеявшему мне. Посла с Небесе и спасе мя, даде в поношение попирающыя мя. Посла Бог милость Свою и истину Свою, и изъят душу мою от среды скимен, поспах смущен. Сынове человечестии, зубы их оружия и стрелы, и язык их меч остр. Вознесися на Небеса, Боже, и по всей земли слава Твоя. Сеть уготоваша ногам моим, и смириша душу мою. Ископаша пред лицем моим яму и впадошася в ню. Готово сердце мое, Боже, готово сердце мое, пою и воспою во славе моей. Востани слава моя, востани псалтырю и гусли, востану рано. Исповемся Тебе в людех, Господи, пою Тебе во языцех. Яко возвеличися до Небес милость Твоя, и даже до облак истина Твоя. Вознесися на Небеса, Боже, и по всей земли слава Твоя.


\myemph{ Таже, Трисвятое, и по Отче наш:}


\mysubtitle{Тропари, глас 1}


Объятия Отча отверсти ми потщися, блудно изжих мое житие, но на богатство неизживущее взираю щедрот Твоих, Спасе: ныне обнищавшее мое сердце не презри. Тебе бо, Господи, умилением зову: согреших, Отче, на небо и пред Тобою.


Слава: Егда приидеши, Боже, судити на землю со славою, и вострепещут всяческая, река же огненная пред судищем потечет, книги разгнутся, и тайная обличатся, тогда избави мя огня негасимаго, и сподоби мя одесную Тебе стати, Судие праведныи.


И ныне: Матерь Тя Божию молим вси Дево; в щедроты Твои прибегаем любовию и в благодать Твою; Тебе бо имамы грешнии спасение, и Тебе стяжахом в напастех Едину всенепорочную.


Посем: Господи, помилуй, 30.

Слава Отцу и Сыну и Святому Духу, и ныне и присно и во веки веком, аминь.

Честнейшую херувим, и славнейшую воистину серафим, без истления Бога, Слова рождьшую, сущую Богородицу Тя величаем.

И поклон.

Именем Господним благослови, отче.

За молитв святых отец наших, Господи, Исусе Христе, Сыне Божии, помилуй нас. Аминь.


\myemph{ В Великий пост (кроме суббот и воскресений), а также и в другие посты (и на масленнице в среду и пятницу):}


Господи помилуй, 40, Слава, и ныне. (И 17 поклонов с молитвою св. Ефрема Сирина).

\begin{enumerate}

\item[1.] Честнейшую херувим… (поклон земной великий). Именем Господним… За молитв святых отец наших…


\myemph{ И молитва св. Ефрема:}


\item[2.] Господи и Владыко животу моему, дух уныния, небрежения, сребролюбия и празднословия отжени от мене (земной великий поклон).


\item[3.] Дух же целомудрия, смирения, терпения и любве даруй ми, рабу Твоему (земной великий поклон).


\item[4.] Ей, Господи Царю, даждь ми зрети моя согрешения и еже не осуждати брата моего, яко благословен еси во веки, аминь (земной великий поклон).


\item[5--6.] Господи Исусе Христе, Сыне Божии, помилуй мя грешнаго (дважды с поклонами).

\item[7.] Боже, милостив буди мне грешному \myemph{(поклон)}.


\item[8.] Боже, очисти грехи моя и помилуй мя \myemph{(поклон)}.


\item[9.] Создавыи мя Господи, помилуй. \myemph{(поклон)}.


\item[10.] Без числа согреших Господи, прости мя \myemph{(поклон)}.


\item[11--16.] Повторяем молитвы 5--10 с земными поклонами.


\item[17.] Затем еще раз читаем всю молитву св. Ефрема Сирина и земной великий поклон.

\end{enumerate}

Приидите, поклонимся, трижды.


\mysubtitle{Псалом 33}


Благословлю Господа на всяко время, выну хвала Его во устех моих. О Господе похвалится душа моя, да услышат кротцыи и возвеселятся. Возвеличите Господа со мною, и вознесем имя Его вкупе. Взысках Господа и услыша мя, и от всех скорбий моих избави мя. Приступите к Нему и просветитеся, и лица ваша не постыдятся. Сей нищий возва, и Господь услыша и, и от всех скорбий его спасе и. Ополчится Ангел Господень окрест боящихся Его, и избавит их. Вкусите и видите, яко благ Господь, блажен муж, иже уповает Нань. Бойтеся Господа, вси святии Его, яко несть лишения боящимся Его. Богатии обнищаша и взалкаша, взыскающии же Господа не лишатся всякаго блага. Приидете, чада, послушайте мене, страху Господню научу вас. Кто есть человек хотяи живот, любяи дни видети благи? Удержи язык свой от зла, и устне свои, еже не глаголати льсти. Уклонися от зла и сотвори благо, взыщи мир, и пожени и. Очи Господни на праведныя, и уши Его в молитву их. Лице же Господне, на творящия злая, еже потребити от земли память их. Возваша праведнии, и Господь услыша их, и от всех печалей их избави их. Близ Господь сокрушенных сердцем, и смиренныя духом спасет. Многи скорби праведным, и от всех их избавит я Господь. Хранит Господь вся кости их, ни едина от них сокрушится. Смерть грешников люта, и ненавидящии праведнаго прегрешат. Избавит Господь душа раб Своих, и не прегрешат вси, уповающии Нань.


\mysubtitle{Псалом 38}


Рех: сохраню пути моя, еже не согрешати языком моим: положих устом моим хранило, внегда востати грешному предо мною. Онемех и смирихся, и умолчах от благ, и болезнь моя обновися. Согреяся сердце мое во мне, и в поучении моем разгорится огнь. Глаголах языком моим: скажи ми, Господи, кончину мою и число дний моих, кое есть, да разумею, чесо лишаюся аз? Се пядию измерены положил еси дни моя, и состав мой яко нивочто же пред Тобою; обаче всяческая суета всяк человек живыи. Иде убо образом ходит человек, обаче всуе мятется: сокровищует, и невесть, кому собирает я. И ныне кто терпение мое, не Господь ли? и состав мой от Тебе есть. От всех беззаконии моих избави мя, поношение безумному дал мя еси. Онемех и не отверзох уст моих, яко Ты сотвори. Отстави от мене раны Твоя, от крепости бо руки Твоея аз исчезох. Во обличении о беззаконии показал еси человека, и истаял еси, яко паучину, душу его; обаче всуе всяк человек. Услыши молитву мою Господи, и моление мое внуши; слез моих не премолчиши, яко пресельник аз от Тебе и пришлец, якоже вси отцы мои. Ослаби ми, да почию, прежде даже не отъиду, и к тому не буду.


\mysubtitle{Псалом 40}


Блажен разумеваяи на нища и убога, в день лют избавит и Господь. Господь сохранит и, и живит и, и ублажит и на земли, и не предаст его в руки врагов его. Господь поможет ему на одре болезни его; все ложе его обратил еси в болезнь его. Аз рех: Господи, помилуй мя, исцели душу мою, яко согреших Ти. Врази мои реша мне злая: когда умрет и погибнет имя его? И вхождаше видети; всуе глаголаше: сердце его собра беззаконие ему; исхождаше вон, и глаголаше вкупе. На мя шептаху вси врази мои, на мя помышляху злая мне. Слово законопреступное возложиша на мя, еда спяи не приложит воскреснути. Ибо человек мира моего, нань же уповах, ядыи хлебы моя, возвеличи на мя пяту. Ты же, Господи, помилуй мя и воздвигни мя и воздам им. О сем познах, яко восхоте мя; яко не возрадуется враг мой о мне. Мене же за незлобие прият, и утвердил мя еси пред Собою в век. Благословен Господь Бог Израилев от века и до века, будет, будет.

Таже Трисвятое, и по Отче наш:


\mysubtitle{Тропари, глас 4:}


Смиренную душу мою посети Господи, иже во гресех  житие все изжившу. Но яко  блудницу приими мене и спаси мя.

Слава: Все житие мое блудно изжих Господи, со блудницами окаянныи; яко же блудныи умилением зову: Отче Небесныи, согреших, очисти, и приими мя, и не отрини мене от Себе, от Тебе удалившагося, и неплодными делы ныне обнищавша.

И ныне: К Богородице прилежно ныне притецем грешнии, со смирением припадающе и покаянием, вопиюще из глубины душевныя: Владычице, помози, мсилосердовавши на ны, потщися, яко изгибаем от множества грехов, не отврати раб Своих тощь, Тебе бо едину помощницу имамы.


Посем: Господи, помилуй, 30. И далее, якоже пренаписахом.


Таже, Приидите, поклонимся, трижды.


\mysubtitle{Псалом 69}


Боже, в помощь мою воньми, Господи, помощи ми потщися. Да постыдятся и посрамятся ищущии душу мою. Да возвратятся вспять и постыдятся, хотящии ми злая. Да возвратятся абие стыдящеся глаголющии ми: благо же, благо же. Да возрадуются и возвеселятся о Тебе, вси ищущии Тебе Боже, и да глаголют выну: да возвеличится Господь, любящии спасение Твое. Аз же нищ есмь и убог, Боже, помози ми. Помощник мой и Избавитель мой еси Ты, Господи, не закосни.


\mysubtitle{Псалом 70}


На Тя, Господи, уповах, да не постыжуся в век. Правдою Твоею изми мя и избави мя. Приклони ко мне ухо Твое и спаси мя. Буди ми в Бог Защититель и в место крепко спасти мя, яко утверждение мое и прибежище мое еси Ты. Боже мой, изми мя из руки грешнаго, из руки законопреступнаго и обидящаго. Яко Ты еси терпение мое, Господи, Господи упование мое от юности моея. В Тебе утвердихся от утробы, от чрева матере моея Ты еси мой Покровитель. О Тебе пение мое выну, яко чудо бых многим, и Ты помощник мой крепок. Да исполнятся уста моя похвалы, яко да воспою славу Твою, весь день великолепие Твое. Не отверзи мене во время старости, внегда исчезати крепости моей, не остави мене. Яко реша врази мои мне, и стрегущии душу мою совещаша вкупе. Глаголюще: Бог оставил есть его, поженете и имете его, яко несть избавляяи. Боже мой, не удалися от мене, Боже мой, в помощь мою вонми. Да постыдятся и исчезнут оклеветающии душу мою, да облекутся в студ и срам ищущии злая мне. Аз же всегда уповаю на Тя и приложу на всяку похвалу Твою. Уста моя возвестят правду Твою, весь день спасение Твое. Яко не познах книжна, вниду в силе Господни, Господи, помяну правду Твою Единаго. Боже мой, яже научил мя еси от юности моея, и до ныне возвещу чудеса Твоя. И даже до старости и маторства, Боже мой, не остави мене, дондеже возвещу мышцу Твою роду всему грядущему. Силу Твою и правду Твою, Боже, даже до вышних, яже сотворил ми еси величия. Боже, кто подобен Тебе? Елики явил ми еси скорби многи и злы, и обращь оживил мя еси, и от бездн земли возведе мя. Умножил еси на мне величествие Твое, и обращь утешил мя еси, и от бездн земли паки возведе мя. Ибо аз исповемся Тебе в людех, Господи, в сосудех псаломских истину Твою, Боже, пою Тебе в гуслех, Святый Израилев. Возрадуетеся устне мои, егда пою Тебе, и душа моя, юже еси избавил. Еще же и язык мой весь день поучится правде Твоей, егда постыдятся и посрамятся ищущии злая мне.


\mysubtitle{Псалом 76}


Гласом моим ко Господу воззвах, гласом моим к Богу, и внят ми. В день печали моея Бога взысках, руками моима нощию пред Ним, и не прельщен бых. Отвержеся утешитися душа моя, помянух Бога и возвеселихся, возскорбех и пренеможе дух мой. Предваристе стражбы очи мои, смутихся и не глаголах. Помыслих дни первыя, и лета вечная помянух, и поучихся. Нощию сердцем моим глумляхся, и тужаше дух мой. Еда во веки отринет Господь, и не приложит благоволити паки? Или до конца милость Свою отсечет, сконча глагол от рода в род? Еда забудет ущедрити Бог? Или удержит во гневе Своем щедроты Своя? И рех: ныне начах, си измена десница Вышняго. Помянух дела Господня, яко помяну от зачала чудес Твоих. И поучуся во всех делех Твоих, и в зачинаниих Твоих поглумлюся. Боже, во святем путь Твой; кто Бог велии яко Бог наш? Ты еси Бог творяи чудеса. Познал еси в людех силу Твою, избавил еси мышцею Твоею люди Твоя, сыны Ияковля и Иосифовы. Видеша Тя воды, Боже, видеша Тя воды, и убояшася; смутишася бездны множеством шума вод. Глас даша облацы, ибо стрелы Твоя преходят, глас грома Твоего в колеси. Осветиша молния Твоя вселенную, подвижася и трепетна бысть земля. В мори путие Твои, и стези Твоя в водах многих, и стопы Твоя не познаются. Наставил еси яко овца люди Твоя, рукою Моисеовою и Ааронею.

Таже, Трисвятое, и по Отче наш.

Молитва Исусова.


\mysubtitle{Тропари, глас 6:}


Помышляю день страшныи и плачуся деянии моих лукавых; како отвещаю Безсмертному Царю? Коим ли дерзновением воззрю на Судию блудныи аз? Благоутробныи Отче, Сыне Единородныи, Душе Святыи, помилуй нас.

Слава. Во юдоли плачевне на месте идеже положи, егда сядеши Милостиве, сотворити праведныи суд, не обличи моя сокровенная, ниже посрами мене пред ангелы, но пощади мя, Боже, и помилуй мя.

И ныне: Милосердия двери отверзи нам, Благословенная Богородице Дево, надеющиися на Тя  не  погибнем, но да избавимся Тобою от бед, Ты бо еси спасение роду християнскому.

Посем: Господи, помилуй, 30, и далее якоже преднаписася.

Таже, Приидите, поклонимся, трижды.


\mysubtitle{Псалом 101}


Господи, услыши молитву мою, и вопль мой к Тебе да приидет. Не отврати лица Твоего от мене, в онь же день аще скорблю. Приклони ко мне ухо Твое, в онь же  день аще призову Тя, скоро услыши мя. Яко исчезоша яко дым дние мои, и кости моя яко сушило сосхошася. Уязвен бых яко трава, и изсше сердце мое, яко забых снести хлеб мой. От гласа воздыхания моего прильпе кость моя плоти моей. Уподобихся неясыти пустынному, бых яко нощныи вран на нырищи. Забдех и бых яко птица особящаяся на зде. Весь день поношаху ми врази мои, и хвалящии мя мною кленяхуся. Зане пепел яко хлеб ядях, и питие мое с плачем растворях. От лица гнева Твоего и ярости Твоея: яко вознес низверже мя. Дние мои яко сень уклонишася, и аз яко сено изсхох. Ты же Господи, во веки пребываеши, и память Твоя в род и род. Ты воскрес ущедриши Сиона, яко время ущедрити его, яко прииде время. Яко благоволиша раби Твои камение его, и персть его ущедрят. И убоятся языцы имени Господня, и вси царие земстии славы Твоея. Яко созиждет Господь Сиона, и явится в славе Своей. Призре на молитву убогих, и не уничижи моления их. Да напишется сие в род ин, и людие зиждемии восхвалят Господа. Яко призре с высоты святыя Своея, Господь с Небесе на землю призре. Услышати воздыхания окованных, разрешити сыны умерщвленных. Возвестити в Сионе имя Господне, и хвалу Его во Иеросалиме. Внегда собратися людем вкупе, и царие  работати Господеви. Отвеща ему на пути крепости его, умаление дний моих возвести ми. Не возведи мене в преполовение дний моих, в род и род лета Твоя. В началех Ты Господи, землю основа, и дела рук Твоих суть небеса. Та погибнут, Ты же пребываеши; и вся, яко риза обетшают, и яко одежду  свиеши их, и изменятся. Ты же тожде еси, и лета твоя не оскудеют. Сынове раб Твоих вселятся, и семя их во век исправится.


\mysubtitle{Молитва Манассии, царя Иудейска}


Господи Вседержителю, Боже отец наших, Авраамов и Исааков и Ияковль, и семене их праведнаго. Сотворивыи небо и землю со всею лепотою их, и сопныи море словом повеления Твоего. Затворивыи бездну, и запечатлев ю страшным и славным именем Твоим, Его же вся боятся, и трепещут от лица славы Твоея. Яко непостоянна велелепота славы Твоея, и не стерпим гнев еже на грешники прещения Твоего. Безчислена же и неизследованна милость обещания Твоего. Ты бо еси Господь вышнии, милосерд, долготерпелив и многомилостив, и каяся о злобах человеческих. Но Ты Господи, по множеству благости Твоея, обеща покаяние и оставление согрешшим к Тебе, и множеством щедрот Твоих нарече покаяние грешником во спасение. Ты убо Господи Боже праведных, неси положил покаяние праведным Твоим, Аврааму и Исааку и Иякову, не согрешившим пред Тобою, но положил еси покаяние мне грешному, зане согреших Ти паче числа песка морскаго. Умножишася беззакония моя Господи, умножишася, и несмь достоин воззрети и видети высоту небесную, от множества неправд моих связан есмь многими юзами железными. Яко не возвести ми главы моея, и несть ми восклонения. Зане прогневах ярость Твою, и лукавое пред Тобою сотворих. И не сотворих воли Твоея, не сохраних повелении Твоих. И ныне поклоняю колена сердца моего, и молю яже от Тебе благость. Согреших Господи, согреших, и беззакония моя аз свем. Но прошу и молюся Тебе: отради ми Господи, отради ми, и не погуби мене со беззаконьми моими. Ниже в век враждовав соблюдеши зол моих, и не осуди мене в преисподних земли. Зане Ты еси Боже, Бог кающихся, да и на мне явиши всю благость Твою, яко недостойна суща спасеши мя, по мнозей милости Твоей. И восхвалю Тя всегда, во вся дни живота моего. Яко Тебе поют вся силы небесныя, и Твоя есть слава во веки, аминь.

Слава в вышних Богу, и на земли мир, в человецех благоволение. Хвалим Тя, благословим Тя \myemph{(поклон)}, кланяемтися, славословим Тя \myemph{(поклон)}, благодарим Тя, великия ради славы Твоея \myemph{(поклон)}. Господи, Царю небесныи. Боже Отче Вседержителю, и Господи Сыне Единородныи, Исусе Христе, и Святый Душе. Господи Боже Агньче Божии, Сыне Отечь, вземляи грех мiра, помилуй нас, Вземляи грехи мiра, приими молитвы наша; седяи одесную Отца, помилуй нас. Яко Ты еси Един Свят, Ты еси Един Господь, Исус Христос, в славу Богу Отцу, аминь.

На всяк день благословим Тя, и восхвалим  имя Твое во веки и в век века.

Господи, прибежище бысть нам в род и род. Аз рех: Господи, помилуй мя, и исцели душу мою, яко согреших Тебе. Господи, к Тебе прибегох, научи мя творити волю Твою, яко Ты еси Бог мой. Яко от Тебе есть источник живота, во свете Твоем узрим свет. Пробави милость Твою ведущим Тя.

Сподоби, Господи, в день сей без греха сохранитися нам. Благословен еси, Господи, Боже отец наших, и хвально, и прославлено имя Твое во веки, аминь.

Буди, Господи, милость Твоя на нас, якоже уповахом на Тя. Благословен еси, Господи, научи нас оправданием Твоим. Благословен еси, Владыко, вразуми нас оправданием Твоим. Благословен еси, Святыи, просвети нас оправданием Твоими.

Господи, милость Твоя во веки, и дела  руку Твоею не презри. Тебе подобает хвала, Тебе подобает пение, Тебе слава подобает, Отцу, и Сыну, и Святому Духу, ныне и присно и во веки веком, аминь.


\mysubtitle{Молитва великаго Евстратия}

Величая величаю Тя Господи, яко призрел еси на смирение мое, и неси мене затворил в руках вражиих, но спасл еси от бед душу мою. И ныне Владыко, да покрыет мя рука Твоя, и да приидет на мя милость Твоя, яко смятеся душа моя, и болезненна есть во исхождении своем, от окаяннаго ми и смраднаго телесе сего. Да некогда лукавыи супостата совет срящет, и препнет ю во тьме, за неведомыя и ведомыя в житии сем бывшия ми грехи. Милостив ми буди, Владыко, и да не узрит душа моя темнаго взора лукавых бесов, но да приимут ю ангели Твои светлии  и пресветлии. Даждь славу имени Твоему святому, и Твоею силою возведи мя на божественныи Твой суд. Внегда судитимися, да не приимет мене рука князя мiра сего, исторгнути мя грешника во глубину адову. Но предстани ми, и буди мой Спаситель и Заступник. Телесныя бо сия муки, веселие суть рабом Твоим. Помилуй, Господи, осквернившуюся страстьми жттия сего душу мою, и чисту ю ради покаяния и исповедания приими. Яко благословен еси во веки веком, аминь.


Таже, Трисвятое, и по Отче наш.

Молитва Исусова.


\mysubtitle{Тропари, глас 8}


Оком милосердым Господи, виждь мое смирение. Яко по малу жизнь моя скончевается, и от дел несть спасения. Сего ради молюся: оком благоутробным Ти Господи, виждь смирение мое и спаси мя.

Слава: Век мой скончавается, и страшный Твой престол готовится, житие мое мимо ходит, суд мене ждет, претя мне огненною мукою, и пламенем негасимым. Слезам тучу подаждь ми и угаси его силу, хотяи спастися всем человеком.

И ныне: Иже нас ради рождеися от Девы, и распятие претерпев, Благии, испровергии смертию смерть, и воскресение явлеи, яко Бог. Не презри их же созда рукою Своею. Яви человеколюбие Свое, Милостиве. Приими Рождьшую Тя Богородицу молящуюся за ны, и спаси Спасе наш, люди согрешшыя.

Таже, Господи помилуй, 30. 

Честнейшую херувим, и славнейшую воистину серафим, без истления Бога, Слова рождьшую, сущую Богородицу Тя величаем. \myemph{(поклон)}.

Слава Отцу и Сыну и Святому Духу \myemph{(поклон)}, и ныне и присно и во веки веком, аминь \myemph{(поклон)}.

Господи помилуй, Господи помилуй, Господи благослови \myemph{(поклон)}.

И отпуст: Господи Исусе Христе, Сыне Божии, молитв ради Пречистыя Твоея Матере, и преподобных и богоносных отец наших, и всех ради святых, помилуй и спаси мя грешнаго, яко Благ и Человеколюбец. Аминь.

Затем "--- прощение.

Аще ли пост, творим поклоны по преждеписанному.

По 17 поклонех Слава, и ныне. Господи помилуй (дважды), Господи благослови (без поклонов). И отпуст псалтырный: Господи Исусе Христе, Сыне Божии, молитв ради Пречистыя Твоея Матере, Силою Честнаго и Животворящаго Креста, и святых небесных сил безплотных, и преподобных и богоносных отец наших, и святаго пророка Давыда (и святаго, его же есть день) и всех святых, помилуй и спаси мя грешнаго, яко Благ и Человеколюбец. Аминь.

И прощение. Наклонившись до земли, не крестясь, читай:

Ослаби, остави, отпусти Боже, согрешения моя, вольная и невольная, яже в слове и в деле, и яже в ведении и не в ведении, яже во уме и в помышлении, яже во дни и в нощи, вся ми прости, яко Благ и Человеколюбец, аминь.

Вставши, читай молитву сию с поклонами:

Ненавидящих и обидящих нас прости, Господи Человеколюбче. Благотворящим благо сотвори, братиям и всем сродником нашим, иже и уединившимся, даруй им вся, яже ко спасению прошения и живот вечныи \myemph{(поклон)}.

В болезнех сущия посети и исцели, в темницах сущих свободи, по водам плавающим Правитель буди и иже в путех шествующим исправи и поспеши \myemph{(поклон)}.

Помяни Господи, и плененныя братию нашу, единоверных православныя веры, и избави их всякаго злаго обстояния \myemph{(поклон)}.

Помилуй Господи, давших нам милостыню и заповедавших нам, недостойным, молитися о них, прости их и помилуй \myemph{(поклон)}.

Помилуй Господи, труждающихся и служащих нам, милующих и питающих нас, и даруй им вся, яже ко спасению, прошения и живот вечныи \myemph{(поклон)}.

Помяни Господи, прежде отшедшия отцы и братию нашу и всели их, идеже присещает свет лица Твоего \myemph{(поклон)}.

Помяни Господи, и нашу худость и убожество, и просвети наш ум светом разума святаго Евангелия Твоего, и настави нас на стезю заповедей Твоих, молитвами Пречистыя Твоея Матере и всех святых Твоих, аминь \myemph{(поклон)}.

Господи помилуй \myemph{(трижды)}. 

И обычные исходные поклоны.


\end{mymulticols}

\section{Келейное правило}\begin{mymulticols}
%http://www.molitvoslov.org/text1001.htm 
 


Начало правилу келейному

Подобает ведати, како начати правило свое в келии. Став убо на обычном своем месте, сотвори со смирением сердца и сокрушенным помыслом поклоны до земли, глаголя сице:

\begin{enumerate}

\item Боже, очисти мя, грешнаго, и помилуй мя \myemph{(поклон)}.

\item Создавыи мя, Господи, помилуй мя \myemph{(поклон)}.

\item Без числа согреших, Господи, прости мя \myemph{(поклон)}.

\item Боже, милости буди мне, грешному \myemph{(поклон)}.

\item Боже, прости беззакония моя и согрешения \myemph{(поклон)}.

\item Пресвятая Владычице моя Богородице, помилуй мя и спаси мя, и помози ми ныне в жизни сей, и во исход души моея, и в будущем веце \myemph{(поклон)}.

\item Непобедимая и божественная сила, честнаго и животворящаго Креста Господня, не остави мя, грешнаго, уповающаго на тя \myemph{(поклон)}.

\item Вся небесныя силы, святии ангели и архангели, херувими и серафими, помилуйте мя и помолитеся о мне, грешнем, ко Господу Богу, и помозите ми ныне, в жизни сей, и во исход души моея, и в будущем веце \myemph{(поклон)}.

\item Ангеле Христов, хранителю мой святыи, помилуй мя и помолися о мне грешнем ко Господу Богу, и помози ми ныне, в жизни сей, и во исход души моея, и в будущем веце \myemph{(поклон)}.

\item Святыи великии Иоанне, пророче и Предотече Господень, помилуй мя и помолися о мне грешнем ко Господу Богу, и помози ми ныне, в жизни сей, и во исход души моея, и в будущем веце \myemph{(поклон)}.

\item Святии славнии апостоли, пророцы и мученицы, святителие, лреподобнии и праведнии и вси святии, помилуйте мя и помолитеся о мне грешнем ко Господу Богу, и помозите ми ныне, в жизни сей, и во исход души моея, и в будущем веце \myemph{(поклон)}.

\item Господи, или словом, или делом, или помышлением, согреших во всей жизни моей, помилуй мя и прости мя милости Твоея ради \myemph{(поклон)}.

\end{enumerate}

\mysubtitle{Посем, молитва святаго Макария, 1-я}

Боже, очисти мя, грешнаго, яко николиже благо сотворих пред Тобою, но избави мя от лукаваго, и да будет во мне воля Твоя, да неосужденно отверзу уста моя недостойная, и восхвалю имя Твое святое, Отца и Сына и Святаго Духа, ныне и присно и во веки веком. Аминь.


\mysubtitle{Молитва 2-я, святаго Исаака Сириянина}

Господи Исусе Христе Боже мой, посещаяи тварь Свою, Тебе явлены страсти моя, и немощь естества моего, и крепость супостата моего. Ты Сам, Владыко, покрый мя от злобы его, занеже сила его крепка и естество мое страстно, и сила немощна. Ты убо Благии, ведыи немощь мою, иже нося неудобьство безсилия моего, сохрани мя от помысла смущена, и потопа страстей, и достойна мя сотвори, сея службы святыя, да не како в страстех моих растлю сладость ея, и обрящуся безстуден пред Тобою и дерзостив. Но милостию Си помилуй мя, яко благословен еси во веки. Аминь.


\mysubtitle{Молитва 3-я, святаго Стефана Фивейскаго}

Владыко Господи Исусе Христе, Боже мой, Ты Помощник ми буди, в руку Твоею аз, да не оставиши мене согрешати, яко прельщен есмь, да не оставиши мене последовати хотению моему лукавому, да не оставиши мене погибнути во гресех моих, ущедри создание Свое, не отверзи мене от лица Твоего грех ради моих, яко к Тебе прибегох, исцели душу мою, яко согреших Ти, пред Тобою суть вси оскорбляющие мя, и ищущие дущу мою изъяти ю, и несть ми прибежища иного, токмо к Тебе Господи, Господи, спаси мя милости Твоея ради, яко Ты еси Господи, сильныи во всяческих. Яко Твое есть Царство и сила и слава, Отца и Сына и Святаго Духа, ныне и присно и во веки веком. Аминь.


\mysubtitle{Молитва 4-я, святаго Иоанна Златоустаго}

Господи, аще хощу, аще не хощу, спаси мя, понеже бо аз яко кал любовещныи, греховную скверну желаю, но Ты яко Благ и Всесилен, можеши ми возбранити. Аще бо праведнаго помилуеши, ничтоже велие, ащи чистаго спасеши, ничтоже дивно, достойни бо суть милости Твоея. Но на мне паче Владыко, окаяннем и грешнем и сквернем удиви милость Свою, покажи благоутробие Свое, Тебе бо оставлен есмь нищии, обнищав всеми благими делы. Господи, спаси мя милости Твоея ради, яко благословен еси во веки. Аминь.


\mysubtitle{Молитва 5-я}

Владыко, помилуй мя благости Твоея ради, и не остави мене заблудити от Твоея воли. И не отверзи убогия молитвы моея от Твоего лица, но услыши Господи, глас молитвы моея, егда молюся Тебе во дни и в нощи, и приими яко избранно кадило, и не возбрани грехов моих ради благости Твоея, но спаси мя имене Твоего ради святаго. Твое бо есть единаго, еже миловати и спасати нас, и Тебе славу возсылаем, Отцу и Сыну и Святому Духу, ныне и присно, и во веки веком. Аминь.


Таже, помолився, и предочистив себя, целуй крест свой, иже носиши на персех своих, и знаменаяся им, глаголи:

Господи Исусе Христе, Сыне Божии, благослови и освяти и сохрани мя, силою Креста живоноснаго Ти.

\mysubtitle{Посем:}

Слава Тебе, Боже наш, слава Тебе всяческих ради. \myemph{(трижды)}.

За молитв святых отец наших, Господи Исусе Христе, Сыне Божии, помилуй нас. Аминь. Таже:

Царю Небесныи.

Трисвятое. По Отче наш:

Господи помилуй, 12. Слава, и ныне.

Приидите поклонимся, трижды.

Псалом 50: Помилуй мя, Боже:

Символ Веры: Верую во единаго Бога:

Таже, или каноны глаголи, или молитвы Исусовы, или поклоны.

По завершении канонов глаголи:

Достойно есть. Трисвятое. По Отче наш: тропари канонов.

Господи помилуй, 40.

Таже Честнейшую херувим.

Слава, и ныне, Господи помилуй, дважды, Господи, благослови (с поклонами).

И отпуст, по обычаю.

Неции же глаголют последи, Господи помилуй, 40. Молитвы канонов. И исповедание.

Таже: Честнейшую. И отпуст.
\longpage[2]


\end{mymulticols}

\mychapterending


\mychapter{Суточный круг}
%http://www.molitvoslov.org/content/sutochniykrug

 

\section{Шестый час}\begin{mymulticols}
%http://www.molitvoslov.org/text1013.htm 
 


Приидите, поклонимся (трижды, и поклоны три).


\mysubtitle{Псалом 53}


Боже, во имя Твое спаси мя, и в силе Твоей суди ми. Боже, услыши молитву мою, внуши глаголы уст моих. Яко чуждии восташа на мя, и крепцыи взыскаша душу мою, и не предложиша Бога пред собою. Се бо Бог помогает ми, и Господь заступник души моей. Отвратит злая врагом моим, истиною Твоею потреби их. Волею пожру Тебе, исповемся имени Твоему, Господи, яко благо. Яко от всякия печали избави мя, и на враги моя воззре око мое.


\mysubtitle{Псалом 54}


Внуши Боже, молитву мою, и не презри моления моего. Вонми ми, и услыши мя. Возскорбех печалию моею, и смутихся от гласа вражия, и от стужения грешнича. Яко уклониша на мя беззаконие, и во гневе враждоваху ми. Сердце мое смутися во мне, и страх смерти нападе на мя. Боязнь и трепет прииде на мя, и покры мя тма. И рех: кто даст ми криле, яко голуби, и полещу и почию. Се удалихся бегая, и водворихся в пустыни. Чаях Бога спасающаго мя, от малодушия и бури. Потопи Господи, и раздели языки их. Яко видех беззаконие и пререкание во граде. День и нощь обыдет и по стенам его, беззаконие и труд посреде его и неправда. И не оскуде от пути его лихва и лесть. Яко аще бы враг поносил ми, претерпел бых убо. И аще бы ненавидяи мя на мя велеречевал, укрылбыхся от него. Ты же человече равнодушне, владыко мой и знаемый мой, иже купно насладил мя еси брашна, во храме Божии ходихове единомышлением. Да приидет же смерть на ня, и снидут во ад живи, яко лукавство в жилищих их, посреде их. Аз к Богу возвах, и Господь услыша мя. Вечер и заутра и полудне, повем и возвещу, и услышит глас мой. Избавит миром душу мою от приближающихся мне, яко во мнозе бяху со мною. Услышит Бог и смирит их, Сыи прежде век. Несть бо им изменения, яко не убояшася Бога. Прострет руку свою на воздаяние, оскверниша завет его. Разделишася от гнева лица его, и приближишася сердца их, умякнуша словеса их паче елея, и та суть стрелы. Возверзи на Господа печаль твою, и Той тя препитает, не даст в век молвы праведнику. Ты же Боже, низведеши их в студенец истления. Мужие крове и льсти не преполовят дней своих, аз же Господи, уповаю на Тя.


\mysubtitle{Псалом 90. }


Живыи в помощи Вышняго, в крове Бога небеснаго водворится. Речет Господеви: Заступник мой еси, и прибежище мое, Бог мой и уповаю Нань. Яко той избавит тя от сети ловчи, и от словесе мятежна. Плещьма Своима осенит тя, и под крыле Его надеешися. Оружие обыдет тя истина Его, не убоишися от страха нощнаго, от стрелы летящия во дне. От вещи во тме преходящия, от сряща и беса полуденнаго. Падет от страны твоея тысяща, и тма одесную тебе, к тебе же не приближится. Обаче очима своима смотриши, и воздаяние грешником узриши. Яко Ты Господи, упование мое; Вышняго положил еси прибежище твое. Не приидет к тебе зло, и рана не приближится к телеси твоему. Яко ангелом Своим заповесть о тебе, сохранити тя во всех путех твоих. На руках возмут тя, да некогда преткнеши о камень ноги твоея. На аспида и василиска наступиши, и попереши льва и змия. Яко на Мя упова, и избавлю и, покрыю и, яко позна имя мое. Воззовет ко Мне, и услышу и, с ним есмь в скорби, изму и, и прославлю его. Долготу дней исполню и, и явлю ему спасение Мое. 

Слава, и ныне. Аллилуия \myemph{(трижды)}.


\myemph{ Тропарь по уставу.} 

Аще ли же пост, глаголем сии тропарь трижды: Иже в шестыи день же и час, на Кресте пригвождеи, Иже в раи дерзновенныи от Адама грех, и согрешении наших рукописание раздери, Христе Боже, и спаси нас.

\myemph{ Стих: }Аз к Богу возвах, и Господь услыша мя.

Стих: Вечер и заутра и полудне, повем и возвещу, и услышит глас мой.

Слава, и ныне, богородичен: Яко не имамы дерзновения, за премногия грехи наша, но Ты, иже от Тебе рождьшагося, моли Богородице Дево, много бо может молитва Матерня, на умоление Владыки. Не презри грешных мольбы Всечистая, яко милостив есть, и спасти могии, Иже страдати нас ради изволивыи. 

Аще ли пост, чтется паремия, и в лествице. Таже, Скоро да предварят ны щедроты Твоя Господи, яко обнищахом зело, помози нам, Боже Спасителю наш; славы ради имене Твоего Господи, избави нас, очисти грехи наша, имене Твоего ради. Трсвятое, и по Отче наш, кондак, по уставу. 

\myemph{ Аще ли же пост, глаголем тропари сия, глас 2: }Спасение содея посреде земли, Христе Боже, на Кресте пречистеи руце Свои простер, собирая вся языки вопиющия: Господи, слава Тебе. 

Слава. Пречистому Ти образу покланяемся Благии, просяще прощения прегрешением нашим, Христе Боже, волею бо благоизволил еси взыти на Крест, да избавиши их же созда от работы вражия. Тем благодаряще вопием Ти: радости вся исполнивыи, Спасе наш, пришедыи спасти мир (поклон великий). 

И ныне. Милосердия сущи источник, милости сподоби нас, Богородице, призри на люди согрешившия, яви яко присно силу Твою, на Тя бо уповающе: радуися, вопием Ти, якоже иногда Гавриил, безплотным архистратиг. 

Господи помилуй (40). Иже на всяко время… Господи помилуй \myemph{(трижды)}. Слава, и ныне. Честнейшую херувим… Именем Господним благослови, отче. За молитв святых отец наших… И сотворяем три великия поклоны, глаголюще молитву: Господи и Владыко животу моему… И поклон един. Таже, Владыко Боже, Отче Вседержителю… 

\myemph{ Аще ли кроме поста, глаголем молитву сию:} Боже и Господи силам, и всея твари Содетелю, Иже за милосердие неизреченныя милости Твоея, единороднаго Сына Твоего, Господа нашего Исуса Христа, низпослав на спасение роду нашему, и честным Его Крестом, рукописание грех наших растерзав, и обличив тем начала и власти тмы. Сам, Владыко Человеколюбче, приими от нас грешных, благодарныя сия молитвы, и избави нас от всякаго всегубительнаго и мрачнаго грехопадения, и всех иже озлобити нас ищущих, видимых и невидимых враг. Пригвозди страсе Твоем плоти наша, и не уклони сердец наших в словеса и помышления лукавствия, но любовию Твоею уязви душа наша, да к Тебе всегда взирающе, и еже от Тебе светом наставляеми, Тебе неприкосновеннаго и присносущнаго зряще Света, непрестанное Тебе исповедание, и благодарение возсылаем, безначальному Отцу, со единородным Ти Сыном, и с Пресвятым и Благим и Животворящим Ти Духом, ныне и присно и во веки веком, аминь.


\end{mymulticols}

\section{Третий час}\begin{mymulticols}
%http://www.molitvoslov.org/text1012.htm 
 


Приходные поклоны 

За молитв святых отец наших, Господи Исусе Христе Сыне Божии, помилуй нас. Аминь 

Царю небесныи. Трисвятое. И по Отче наш. Молитва Исусова. Аминь. Господи помилуй (12). Слава, и ныне. 

Приидите поклонимся Цареви нашему Богу \myemph{(поклон)}.

Приидите поклонимся Христу, Цареви и Богу нашему \myemph{(поклон)}.

Приидите поклонимся и припадем к Самому Господу Исусу Христу, Цареви и Богу нашему \myemph{(поклон)}.


\mysubtitle{Псалом 16}


Услыши Господи правду мою, вонми моление мое. Внуши молитву мою, не во устнах льстивых. От лица Твоего судьба моя изыдет, очи мои да видите правоты. Искусил еси сердце мое, посетил еси нощию, искусил мя еси, и не обретеся во мне неправда. Яко да не возглаголют уста моя дел человеческих за словеса устен Твоих. Аз сохраних пути жестоки. Соверши стопы моя во стезях Твоих, да не подвижутся стопы моя. Аз возвах, яко услыша мя Боже, приклони ухо Твое ко мне, и услыши глаголы моя. Удиви милости Твоя, спасаяи уповающих на Тя. От противящихся деснице Твоей, сохрани мя Господи, яко зеницу ока. В крове крилу Твоею покрыеши мя, от лица нечестивых острастьших мя. Врази мои душу мою удержаша. Тук свой затвориша, уста их глаголаша гордыню. Изгонящии мя ныне обыдоша мя, очи свои возложиша уклонити на землю. Объяша мя яко лев готов на лов, и яко скимен обитая в сокровищих. Воскресни Господи, предвари я, и запни им; изми душу мою от нечестивых оружия Твоя, от рук враг Твоих. Господи, отмаль от земли раздели я в животе их, и сокровенных Твоих исполнися утроба их. Насытишася сынов, и оставиша останки младенцем своим. Аз же праводю явлюся лицу Твоему, насыщуся внегда явитмися слава Твоя.


\mysubtitle{Псалом 24}


К Тебе Господи, воздвигох душу мою. Боже мой, на Тя уповах, да не постыжуся во веки. Ни посмеютмися врази мои, ибо вси терпящии Тя не постыдятся. Да постыдятся беззаконнующии всуе. Пути Твоя Господи, скажи ми, и стезям Твоим научи мя. Настави мя на истину Твою, и научи мя, яко Ты еси Бог Спас мой, и Тебе терпех весь день. Помяни щедроты Твоя Господи, и милости Твоя, яко от века суть. Грех юности моея и невежествия моего не помяни. По милости Твоей помяни мя Ты, ради благости Твоея, Господи. Благ и прав Господь, сего ради закон положит согрешающим на пути. Наставит кроткия на суд, и научит кроткия путем своим. Вси путие Господни, милость и истина, взыскающим завета Его и свидения Его. Ради имене Твоего Господи, очисти грех мой, мног бо есть. Кто есть человек бояися Господа? Закон положит ему на пути, иже изволи. Душа его во благих водворится, и семя его наследит землю. Держава Господь боящимся Его, и завет Его явит им. Очи мои выну ко Господу, яко Той исторгнет от сети нозе мои. Призри на мя и помилуй мя, яко единород и нищь есмь аз. Скорби сердца моего умножишася, и от бед моих изведи мя. Виждь смирение мое, и труд мой, и отпусти вся грехи моя. Виждь враги моя, яко умножишася, и ненавидением неправедным возненавидеша мя. Сохрани душу мою, и избави мя, да не постыжуся яко уповах на Тя. Незлобивии и правии прилепляхуся мне, яко потерпех Тя, Господи. Избави Боже, Израиля от всех скорбей его. 


\mysubtitle{Псалом 50}


Помилуй мя Боже, по велицей милости Твоей. И по множеству щедрот Твоих очисти беззаконие мое. Наипаче омый мя от беззакония моего и от греха моего очисти мя. Яко беззаконие мое аз знаю и грех мой предо мною есть выну. Тебе единому согреших и лукавое пред Тобою сотворих. Яко да оправдишися в словесех Своих и победиши внегда судити. Себо в беззакониих зачат есмь и во гресех роди мя мати моя. Се бо истину возлюбил еси, безвестная и тайная премудрости Твоея явил ми еси. Окропиши мя иссопом и очищуся. Омыеши мя и паче снега убелюся. Слуху моему даси радость и веселие: возрадуются кости смиренныя. Отврати лице Твое от грех моих и вся беззакония моя очисти. Сердце чисто созижди во мне Боже, и дух прав обнови во утробе моей. Не отверзи мене от лица Твоего, и Духа Твоего Святаго не отыми от мене. Воздаждь ми радость спасения Твоего и духом владычным утверди мя. Научу беззаконныя путем Твоим и нечестивии к Тебе обратятся. Избави мя от кровий Боже, Боже спасения моего; возрадуется язык мой правде Твоей. Господи, устне мои отверзеши, и уста моя возвестят хвалу Твою. Яко аще бы восхотел жертвы, дал бых убо; всесожжения не благоволиши. Жертва Богу "--- дух сокрушен: сердце сокрушенно и смиренно Бог не уничижит. Ублажи Господи, благоволением Твоим Сиона; и да созиждутся стены Иеросалимския. Тогда благоволиши жертву правде, возношение и всесожигаемая. Тогда возложат на олтарь Твой тельца.

Слава, ныне. Аллилуия \myemph{(трижды)}. Господи помилуй \myemph{(трижды)}: Слава, тропарь дню, или святому. И ныне: Богородице, Ты еси лоза истинная…


\myemph{ Аще ли есть пост, глаголи сие, на глас 6, трижды:} Господи Иже Пресвятыи Свой Дух, в третии час апостолом Своим пославыи, того Благии не отими от нас, но обнови нас молящихтися.

Стих: Сердце чисто созижди во мне Боже, и Дух прав обнови во утробе моей.


\myemph{ Стих:} Не отверзи мене от лица Твоего, и Духа Твоего Святаго не отыми от мене. И по коемждо стисе: Господи, Иже Пресвятый Свой Дух… весь до конца. И поклоны три. Посем, Слава, и ныне. 

Богородице, Ты еси лоза истинная, возрастившая нам плод животный, Тебе молимся молися Владычице, со святыми апостолы, помиловатися душам нашим.


\myemph{ Аще ли велкий пост, бывает чтение в лествице.} По скончании же чтения, псаломщик глаголет: аминь. Господь Бог благословен, благословен Господь день дне. Поспешит нам Бог спасении наших, Бог наш, Бог спасати. Трисвятое. И по Отче наш, кондак дню, или святому, или празднику.


\myemph{ Аще ли же пост, глаголем тропари сия, глас 8.} Благословен еси, Христе Боже наш, Иже премудры ловцы явлеи, низпослав им Дух Святыи, и теми уловлеи вселенную. Человеколюбче, слава Тебе.

Слава. Егда сшед языки размеси, разделяше племена Вышнии. Егда же огненыя языки раздая, в совокупление вся призва, единогласно славити Всесвятыи Дух.

И ныне. Надеже и заступнице, и прибежище християном, и необоримая стено, изнемогающим и труждающимся в пучине, пристанище тихое Ты еси, Пречистая Богородице. Но яко мир избавляеши непрестанною Си молитвою, помяни нас, Дево Всепетая. Господи помилуй (40).


\myemph{ Посем молитва святаго Великаго Василия:} Иже на всяко время, и на всяк час, на небеси и на земли, покланяемыи и славимыи, Боже Благии, долготерпеливе и многомилостиве; Иже праведники любя и грешных милуя; Иже всех зовыи ко спасению, обещания ради будущих благ; сам Господи, приими молитвы наша в час сий, и в благости исправи живот наш, к заповедем Твоим. Душа наша освяти, и телеса очисти, помышления исправи, и мысли очисти, разум уцеломудри и истрезви, и избави нас от всякия скорби, зол и болезней, и душевных страстей, и огради нас святыми ангелы Твоими, яко да ополчением их соблюдаеми и наставляеми, достигнем во единство веры, и в разум неприкосновенныя Ти славы, яко благословен еси во веки, аминь. 

Таже, Господи помилуй \myemph{(трижды)}. Слава, и ныне. Честнейшую херувим, и славнейшую воистину серафим, без истления Бога Слова рождьшую, сущую Богородицу Тя величаем \myemph{(поклон)}. Именем Господним благослови, отче. За молитв святых отец наших, Господи Исусе Христе, Сыне Божии, помилуй нас. Аминь.


Аще есть пост, или кроме поста, егда поем Аллилуия, и творим три поклоны великия, глаголюще молитву святаго Ефрема: Господи и Владыко животу моему…. И прочих поклонов, 12.


\myemph{ И паки:}Господи и Владыко животу моему…, с поклоном. 

И по сем молитва святаго Великаго Василия. Владыко Боже, Отче Вседержителю, и Господи Сыне Единородныи Исусе Христе, и Святыи Душе. Едино Божество и едина сила, помилуй мя грешнаго, и ими же веси судьбами, спаси мя недостойнаго раба Твоего, яко благословен еси во веки, аминь.


\end{mymulticols}

\section{Первый час}\begin{mymulticols}
%http://www.molitvoslov.org/text1010.htm 
 


\myemph{ Начало Перваго часа }


\myemph{ Глаголет псаломщик:} Аминь.

Приидите, поклонимся \myemph{ (трижды, и поклоны три)}.


\mysubtitle{Псалом 5}


Глаголы моя внуши Господи, разумей звание мое. Воньми глас молитвы моея. Царю мой и Боже мой. Яко к Тебе помолюся Господи, заутра услыши глас мой. Заутра предстану Ти, и узриши мя, яко Бог не хотя беззаконие Ты еси. Не приидет к Тебе лукавнуяи, не пребудут же законопреступницы пред очима Твоима. Возненавидел еси вся творящия беззаконие, погубиши вся глаголющия лжу. Мужа крови и льстива гнушается Господь. Аз же множеством милости Твоея, вниду в дом Твой. Поклонюся к церкви святей Твоей, во страсе Твоем. Господи, настави мя правдою Твоею, враг моих ради, исправи пред Собою путь мой. Яко несть во устех их истины, сердце их суетно. Гроб отверст гортань их, языки своими льщаху: суди им Боже. Да отпадут от мыслей своих, по множеству нечестия их изрини я, яко прогневаша Тя, Господи. И возвеселятся вси уповающии на Тя, во веки возрадуются и вселишися в них. И похвалятся о Тебе любящии имя Твое, яко Ты благословиши праведника. Господи, яко оружием благоволения венчал еси нас.


\mysubtitle{Псалом 89}


Господи, прибежище бысть нам в род и род. Прежде даже горам не быти, и создатися земли и вселенней, от века и до века Ты еси. Не возврати человека во смирение и рекл еси: обратитеся сынове человечестии. Яко тысяща лет пред очима Твоима Господи, яко день вчерашнии, иже мимо иде, и стража нощная. Уничижения их еще будут. Утро яко трава мимо идет, утро процветет и прейдет, на вечер отпадет, ожестеет и исхнет. Яко исчезохом гневом Твоим, и яростию Твоею смутихомся. Положил еси беззакония наша пред Собою, век наш в просвещение лица Твоего. Яко вси дние наши оскудеша, и гневом Твоим исчезохом. Лета наша яко паучина паучахуся, дние лет наших, в них же седмьдесят лет. Аще ли же в силах, осмьдесят лет; и множае их труд и болезнь. Яко прииде кротость на ны, и накажемся. Кто свесть державу гнева Твоего, и от страха Твоего ярость Твою исчести; Десницу Твою тако скажи ми, и окованныя сердцем в мудрости. Обрати Господи, доколе; и умолен буди на рабы Своя. Исполнихомся заутра милости Твоея Господи, и возрадовахомся и возвеселихомся. Во вся дни наша возвеселихомся, задни в няже смирил ны еси; и лета в няже видехом злая. И призри на рабы Своя, и на дела Твоя и настави сыны их. И буди светлость Господа Бога нашего на нас, и дела рук наших исправи на нас, и дело рук наших исправи.


\mysubtitle{Псалом 100}


Милость и суд воспою Тебе, Господи. Пою и разумею в пути непорочне, когда приидеши ко мне? Прехождах в незлобии сердца моего, посреде дому моего. Не предлагах пред очима моима вещь законопреступну, творящия преступление возненавидех. Не прильпе мне сердце лукаво, уклоняющагося от мене лукаваго не познах. Оклеветающаго тай искреняго своего, сего изгонях. Гордым оком и несытым сердцем, с сим не ядях. Очи мои на верныя земли, посаждати их с собою. Ходяи по пути непорочну, сей ми служаше; не живяше посреде дому моего, творяи гордыню. Глаголяи неправедная, не исправляше пред очима моима. Во утрия избивах вся грешныя земли, потребити от града Господня вся делающия беззаконие. 

Слава, и ныне. Аллилуия \myemph{(трижды)}. Господи помилуй \myemph{(трижды)}. Слава, тропарь дню или святому. И ныне, богородичен: Что Тя наречем…


\myemph{ Аще ли же пост, поем на глас 6:} Заутра услыши глас мой, Царю мой и Боже мой \myemph{(трижды)}, и поклоны три.

Стих. яко к Тебе помолюся, Господи.

Стих. Глаголы моя внуши Господи, разумей звание мое.

Слава, и ныне, богородичен: Что Тя наречем, О! Обрадованная? Небо, яко восияла еси Солнце правде. Рай, яко возрастила еси Цвет нетления. Деву, яко пребыла еси нетленна, Чистая; Матерь, яко имела еси на руку Сына всех Бога. Того моли спастися душам нашим.


\myemph{ Аще ли пост, сия два тропари поем по трижды:} Стопы моя направи по словеси Твоему, и да не одолеет ми всяко беззаконие, избави мя от клеветы человеческия, и сохраню заповеди Твоя. Лице Твое просвети на раба Твоего, и научи мя оправданием Твоим.

Да исполнятся уста моя похвалы, яко да воспою славу Твою, весь день велелепоту Твою. 

Трисвятое. И по Отче наш, кондак святому, или дневи. Господи помилуй (40).


\myemph{ Таже, молитва Великаго Василия:} Иже на всяко время, и на всяк час, на небеси и на земли, покланяемыи и славимыи, Боже Благии, долготерпеливе и многомилостиве; Иже праведники любя и грешных милуя; Иже всех зовыи ко спасению, обещания ради будущих благ; сам Господи, приими молитвы наша в час сий, и в благости исправи живот наш, к заповедем Твоим. Душа наша освяти, и телеса очисти, помышления исправи, и мысли очисти, разум уцеломудри и истрезви, и избави нас от всякия скорби, зол и болезней, и душевных страстей, и огради нас святыми ангелы Твоими, яко да ополчением их соблюдаеми и наставляеми, достигнем во единство веры, и в разум неприкосновенныя Ти славы, яко благословен еси во веки, аминь. 

Таже, Господи помилуй \myemph{(трижды)}. Слава, и ныне. Честнейшую херувим, и славнейшую воистину серафим, без истления Бога Слова рождьшую, сущую Богородицу Тя величаем \myemph{(поклон)}. Именем Господним благослови, отче. За молитв святых отец наших, Господи Исусе Христе, Сыне Божии, помилуй нас. Аминь.


\myemph{ Иерей глаголет:} За молитв святых отец наших… И мы глаголем: Аминь. Таже молитва: Христе, Свете истинный… И конечный отпуст.


\myemph{ Аще ли пост, по Отче наш, глаголются тропари сия, глас 4.} Скоро предвари, даже не поработимся, врази бо хулят Тя, и претят нам, Христе Боже. Но порази Крестом Своим борющихся с нами, да разумеют, колико может православных вера, молитвами Богородицы, едине Человеколюбче.

Слава. Премудрости наставниче, и смыслу давче, немудрым наказателю, и нищим защитителю, утверди и вразуми сердце мое, Владыко. Ты даждь ми слово, еже отчее единородное Слово: се бо устнама моима не возбраню, еже звати Ти: Милостиве, помилуй мя падшаго.

И ныне. Преславную Христову Матерь, и святых ангел святейшу, немолчно воспоем сердцем и устнами, Богородицу Тя исповедающе, яко воистину рождьшую нам Бога воплощена, и молящуюся непрестанно о душах наших. 

Господи помилуй (40). Иже на всяко время… Господи помилуй \myemph{(трижды)}. Слава, и ныне. Честнейшую херувим… (поклон великий). Именем Господним благослови, отче. За молитв святых отец наших… Аминь. И творим великия поклоны, с молитвою святаго Ефрема. И прочих 12 метаний. И паки вышеписаную молитву: Господи и Владыко животу моему… всю до конца. И поклон един. Посем, Трисвятое. И по Отче наш, Господи помилуй (12).


\myemph{ И молитва:} Христе Свете истинныи, иже просвещая и освещая всякаго человека грядущаго в мир, да знаменается на нас свет лица Твоего Господи, яко да в нем ходяще, узрим свет неприкосновенныя Ти славы, исправи стопы наша, к деланию заповедей Твоих, молитвами Пречистыя Ти Матере, и всех святых Твоих, яко благословен еси во веки, аминь. Таже, обычный отпуст дня того, и настоящаго святаго.


\end{mymulticols}

\section{Чин обедницы}\begin{mymulticols}
%http://www.molitvoslov.org/text1015.htm 
 


\myemph{ По начале Божественныя службы, и по первом возгласе, псаломщик глаголет псалом:}


Благослови душе моя Господа. 

\myemph{ Егда же во уставе укажет на литургии антифоны празднику, и тогда часы глаголются вси вкупе.}


\mysubtitle{Псалом 102}


Благослови, душе моя, Господа, и вся внутреняя моя имя святое Его. Благослови, душе моя, Господа, и не забывай всех воздаянии Его. Очищающаго вся беззакония твоя, исцеляющаго вся недуги твоя. Избавляющаго от истления живот твой, венчающаго тя милостию и щедротами. Исполняющаго во благих желание твое, обновится яко орлу юность твоя. Творяи милостыню Господь, и судьбу всем обидимым. Сказа пути Своя Моисеови, сыновом израилевым хотения Своя. Щедр и милостив Господь, долготерпелив, и многомилостив. Не до конца прогневается, ни в век враждует. Не по беззаконием нашим сотворил есть нам, ни по грехом нашим воздал есть нам. Яко по высоте небесней от земли, утвердил есть Господь милость Свою на боящихся Его. Елико отстоят востоцы от запад, удалил есть от нас беззакония наша. Яко же щедрит отец сыны, ущедрит Господь боящихся Его. Яко Той позна создание наше, помяну яко персть есмы. Человек, яко трава, дние его, яко цвет сельныи, тако оцветет. Яко дух пройдет в нем и не будет, и не познает к тому места своего. Милость же Господня от века и до века, на боящихся Его. И правда Его на сынех сынов, хранящих завет Его, и помнящих заповеди Его творити я. Господь на небеси уготова престол Свой, и царство Его всеми обладает. Благословите Господа вси ангели Его, сильнии крепостию творящии слово Его, услышати глас словес Его. Благословите Господа вся силы Его, слуги Его, творящии волю Его. Благословите Господа вся дела Его, на всяком месте владычества Его, благослови, душе моя Господа.


\mysubtitle{Аще ли Литургия, глаголем по 2-м возгласе, псалом 145}


Хвали, душе моя, Господа, восхвалю Господа в животе моем, пою Богу моему дондеже есмь. Не надейтеся на князи на сыны человеческия, в них же несть спасения. Изыдет дух его, и возвратится в землю свою, в той день погибнут вся помышления его. Блажен ему же Бог Ияковль помощник его, упование его на Господа Бога своего. Сотворшаго небо и землю, море и вся яже в них. Хранящаго истину в век, творящаго суд обидимым, дающаго пищу алчущим. Господь решит окованных, Господь умудряет слепца, Господь возводит низверженныя, Господь любит праведники. Господь хранит пришельца, сира и вдову приемлет, и путь грешных погубит. Воцарится Господь во веки. Бог Твой Сионе в род и род. 

Слава, и ныне.


\myemph{ Клирицы поют:} Единородный Сыне, и Слово Божие, безсмертне сыи, изволивыи спасения нашего ради воплотитися от святыя Богородицы и присно Девы Марии непреложно вочеловечивыися; распятся, Христе Боже, и смертию на смерть наступивыи, Едине сыи Святыя Троицы, спрославляемыи Отцу и Святому Духу, спаси нас.


\myemph{ По скончании пения глаголет канархист блаженны, по уставу.}


Во царствии Си егда приидеши, помяни нас, Господи. Блажени нищии духом, яко тех есть царство небесное. Блажени плачущии, яко тии утешатся. Блажени кротцыи, яко тии наследят землю. Блажени алчущии и жаждущии правды, яко тии насытятся. Блажени милостивии, яко тии помиловани будут. Блажени чистии сердцем, яко тии Бога узрят. Блажени миротворцы, яко тии сынове Божии нарекутся. Блажени изгнани правды ради, яко тех есть царство небесное. Блажени есте егда поносят вам, и изженут вы, и рекут всяк зол глагол на вы лжуще Мене ради. Радуйтеся и веселитеся, яко мзда ваша многа на небесех. Слава, и ныне.


\myemph{ Аще ли кроме Божественныя литургии, посем чтем:} Апостол, и Евангелие.


\myemph{ Таже, глаголем:} Помяни нас Господи, егда приидеши во царствии Си \myemph{(поклон)}.

Помяни нас Владыко, егда приидеши во царствии Си \myemph{(поклон)}.

Помяни нас Святыи, егда приидеши во царствии Си \myemph{(поклон)}.


\myemph{ Лик небесныи поет Тя и глаголет:} свят, свят, свят Господь Саваоф, исполнь небо и землю славы Твоея. Приступите к Нему и просветитеся, и лица ваша не постыдятся. Лик небесныи поет Тя и глаголет: свят, свят, свят Господь Саваоф, исполнь небо и землю славы Твоея. 

Слава. Собор святых ангел и архангел, со всеми небесными силами поет Тя и глаголет: свят, свят, свят Господь Саваоф, исполнь небо и землю славы Твоея. 

И ныне. Исповедание православныя веры, Перваго собора.

Верую во единаго Бога Отца, Вседержителя, Творца небу и земли, видимым же всем и невидимым.

И во единаго Господа Исуса Христа Сына Божия, Единороднаго, иже от Отца рожденнаго прежде всех век. Света от Света, Бога истинна от Бога истинна, рождена, а не сотворена, единосущна Отцу, Им же вся быша.

Нас ради человек, и нашего ради спасения сшедшаго с небес и воплотившагося от Духа Свята и Марии Девы вочеловечшася.

Распятаго за ны при понтийстем Пилате, страдавша и погребена.

И воскресшаго в третии день по писаниих.

И возшедшаго на небеса, и седяща одесную Отца.

И паки грядушаго со славою судити живым и мертвым, Его же царствию несть конца.

Втораго Собора:

И в Духа Святаго, Господа истиннаго и Животворящаго, иже от Отца исходящаго, иже со Отцем и Сыном спокланяема и сславима, глаголавшаго пророки.

И во едину святую соборную и апостольскую Церковь.

Исповедую едино Крещение во оставление грехов.

Чаю воскресение мертвым.

И жизни будущаго века. Аминь.


\myemph{ В литургии сие не глаголем:} Ослаби, остави, отпусти Боже, согрешения наша, вольная и невольная, яже в слове и в деле, и яже в ведении и не в ведении, яже во уме и в помышлении, яже во дни в нощи, вся ми прости, яко Благ и Человеколюбец.


\myemph{ Таже, Отче наш. И по возгласе или молитве Исусовой глаголются кондаки, преже храму, и потом дню: таже, святому.} 

Слава. Со святыми покой, Христе, душа раб Своих, идеже несть болезни, ни печали, ни воздыхания, но жизнь вечная.

И ныне, богородичен: Заступнице християном непостыдная, Ходатаице ко Творцу непреложная, не презри грешных моления гласы, но предвари яко Блага на помощь нашу, верно вопиющих Ти: ускори на молитву, и потщися на умоление, заступающи присно Богродице чтущих Тя. 

Господи помилуй (12). Таже Всесвятая Троице.


\myemph{ Аще ли пост великии, глаголется:} Господи помилуй (40). Слава, и ныне. Честнейшую херувим (поклон великий). Именем Господним благослови, отче. За молитв святых отец наших… Аминь. И сотворяем три великия поклоны с молитвою: Господи и Владыко животу моему. И прочих метаний (12). Посем, Трисвятое.


\myemph{ И по Отче наш:} Господи помилуй (12).


\myemph{ Таже, сию молитву:} Всесвятая Троице единосущная, неодержимая держава, и неразделимое царство, яже всем благим виновна, благоволи в настоящии сей час о мне грешнем, и вся ми омый скверны, и просвети ми смысл, яко да всегда воспеваю Тя, и славословлю и глаголю: един Свят, един Господь Исус Христос, в славу Богу Отцу, аминь. 

Буди имя Господне благословено от ныне и до века (трижды с поклонами). 

Слава, и ныне.


\mysubtitle{Таже, псалом 33}


Благословлю Господа на всяко время, выну хвала Его во устех моих. О Господе похвалится душа моя, да услышат кротцыи и возвеселятся. Возвеличите Господа со мною, и вознесем имя Его вкупе. Взысках Господа и услыша мя, и от всех скорбий моих избави мя. Приступите к Нему и просветитеся и лица ваша не постыдятся. Се нищии возва, и Господь услыша и, и от всех скорбий его спасе и. Ополчится ангел Господень окрест боящихся Его, и избавит их. Вкусите и видите, яко благ Господь, блажен муж, иже уповает Нань. Бойтеся Господа вси святии Его, яко несть лишения боящимся Его. Богатии обнищаша и взалкаша, взыскающии же Господа не лишатся от всякаго блага. Приидите чада послушайте мене, страху Господню научу вас. Кто есть человек хотяи живот, любяи дни видети благи; Удержи язык свой от зла, и устне свои еже не глаголати льсти. Уклонися от зла, и сотвори благо, взыщи мир, и пожени и. Очи Господни на праведныя, и уши Его в молитву их. Лице же Господне на творящия злая, еже потребити от земли память их. Возваша праведнии, и Господь услыша их, и от всех печалей их избави их. Близ Господь сокрушенных сердцем, и смиренныя духом спасет. Многи скорби праведным, и от всех их избавит я Господь. Хранит Господь вся кости их, и не едина от них сокрушится. Смерть грешником люта, и ненавидящии праведнаго прегрешат. Избавит Господь душа раб Своих, и не прегрешат вси уповающии Нань. 

Достойно есть, яко воистину блажити Тя, Богородице, присноблаженную и пренепорочную, и Матерь Бога нашего. Честнейшую херувим, и славнейшую воистину серафим, без истления Бога Слова рождьшую, сущую Богородицу, Тя величаем. И поклон до земли вси равно. Слава, и ныне. Господи помилуй (дважды). Господи благослови. И три поклона. 


\myemph{ Таже, отпуст почину. По отпусте же клирицы поют:} Аминь. Многолетны соблюди Господи, и помилуй державу Российскую, спаси Господи, и помилуй господина нашего Преосвященнейшего Митрополита (имя рек) и отец наших духовных, и вся христианы Господи, спаси. 

Господи помилуй \myemph{(трижды)}.


\myemph{ И обычные исходные поклоны.}


\end{mymulticols}

\section{Утреня}\begin{mymulticols}
%http://www.molitvoslov.org/text1009.htm 
 


\myemph{ По обычнем начале глаголем:}


За молитв святых отец наших, Господи Исусе Христе, Сыне Божии, помилуй нас. Аминь. Трисвятое. И по Отче наш. Молитва Исусова. Аминь. Господи помилуй, (12). Слава и ныне. Приидите, поклонимся (трижды с поклонами).


\mysubtitle{Псалом 19}


Услышит тя Господь в день печали, защитит тя имя Бога Ияковля. Послет ти помощь от Святаго, и от Сиона заступит тя. Помянет всяку жертву твою, и вся сожжения твоя тучна буди. Даст ти Господь по сердцу твоему, и весь совет твой исполнит. Возрадуемся о спасении твоем, и во имя Господа Бога нашего возвеличимся. Исполнит Господь вся прошения твоя, ныне познах, яко спасл есть Господь христа Своего. Услышит его с Небесе святаго Своего, в силах спасение десница Его. Сии на колесницах и сии на конех, мы же во имя Господа Бога нашего призовем. Тии спяти быша и падоша, мы же востахом и исправихомся. Господи, спаси царя и услыши ны, вонь же день аще призовем Тя. 


\mysubtitle{Псалом 20}


Господи, силою Твоею возвеселится царь, и о спасении Твоем возрадуется зело. Желание сердца его дал еси ему, и хотения устну его неси лишил его. Яко предварил еси его благословением благостынным, положил еси на главе его венец от камене честна. Живота просил есть у Тебе, и дал еси ему долготу дний, в век века. Велия слава его спасением Твоим, славу и велелепие возложиши нань. Яко даси ему благословение в век века, возвеселиши его радостию с лицем Твоим. Яко царь уповает на Господа, и милостию Вышняго не подвижится. Обрящется рука Твоя всем врагом Твоим, десница Твоя обрящет вся ненавидящия Тебе. Яко положиши их яко пещь огненую во время лица Твоего, Господь гневом Своим смутит я, и снесть их огнь. Плод их от земли погубиши, и семя их от сынов человеческих. Яко уклониша на Тя злая, помыслиша советы, их же не возмогоша составити. Яко положиши я хребет во избытцех Своих, уготоваеши лице их. Вознесися, Господи, силою Твоею, воспоем и поем силы Твоя. 

Слава, и ныне. Трисвятое. И по Отче наш. Тропари сия, глас 1. Спаси Господи, люди Своя и благослови достояние Свое, победы державе Российстей на сопротивныя даруй, и Своя сохраняя Крестом люди.

Слава. Вознесыися на Крест волею, тезоименитому ныне граду Твоему, щедроты Твоя даруй, Христе Боже. Возвесели силою Своею державу Российскую, победы дая ей на сопостаты, пособие имущу Твое оружие, миру непобедимую победу.

И ныне. Предстательнице страшная и непостыдная, не презри Благая молитв наших. Всепетая Богородице, утверди православных житие и спаси державу Российскую, подай же ей небесную победу Богом, Его же родила еси, едина Благословенная. Господи помилуй (12). Слава Отцу и Сыну и Святому Духу, и ныне, и присно, и во веки веком. Аминь.

Слава в вышних Богу, и на земли мир, в человецех благоволение (трижды без поклонов).


\myemph{ Таже:} Господи, устне мои отверзеши и уста моя возвестят хвалу Твою (дважды). 

И глаголет екса псалмы, легко, с тихостию и со всяким вниманием и страхом Божиим, яко к самому Богу беседуя невидимо, и о своих гресех умоляя Его. Братия же стоят руце имуше согбени к персем, главы же мало преклонены.


\mysubtitle{Псалом 3}


Господи, чтося умножиша стужающии ми, мнози востают на мя. Мнози глаголют души моей: несть спасения ему о Бозе его. Ты же Господи, заступник мой еси, слава моя, и вознося главу мою. Гласом моим ко Господу возвах, и услыша мя от горы святыя Своея, Аз уснух и спах; востах, яко Господь заступит мя. Не убоюся от тем людей, окрест нападающих на мя. Воскресни Господи, спаси мя, Боже мой. Яко Ты порази вся враждующия ми всуе, зубы грешником сокрушил еси. Господне есть спасение, и на людех Твоих благословение Твое. Аз уснух и спах; востах, яко Господь заступит мя.


\mysubtitle{Псалом 37}


Господи, не яростию Твоею обличи мене, ни гневом Твоим покажи мене. Яко стрелы Твоя унзоша ми, и утвердил еси на мне руку Твою. Несть исцеления в плоти моей, от лица гнева Твоего; несть мира в костех моих, от лица грех моих. Яко беззакония моя превзыдоша главу мою, яко бремя тяжко отяготеша на мне, возсмердеша и согнишася раны моя, от лица безумия моего. Пострадах и смирихся до конца, весь день сетуя хождах. Яко лядвия моя наполнишася поругании, и несть исцеления в плоти моей. Озлоблен бых и смирихся до зела, рыках от воздыхания сердца моего. Господи, пред Тобою все желание мое, и воздыхание мое от Тебе не утаися. Сердце мое смутися, остави мя сила моя, и свет очию моею, и той несть со мною. Друзи мои, и искрении мои, прямо мне приближишася, и сташа. И ближнии мои отдалече мене сташа, и нуждахуся, ищущии душу мою, и ищущии злая мне: глаголаху суетная и льстивная, весь день поучахуся. Аз же яко глух: не слышах; и яко нем: не отверзая уст своих. И бых яко человек не слыша, и не имыи во устех своих обличения. Яко на Тя Господи, уповах, Ты услышиши, Господи Боже мой. Яко рех: да некогда порадуютмися врази мои, и внегда подвижатися ногам моим, на мя велеречеваша. Яко аз на раны готов, и болезнь моя предо мною есть выну. Яко беззаконие мое аз возвещу, и попекуся о гресе моем. Врази же мои живут, и укрепишася паче мене, и умножишася ненавидящии мя без правды. Воздающии ми зла возблагая, оболгаху мя: зане гонях благостыню. Не остави мене Господи Боже мой, не отступи от мене. Воньми в помощь мою, Господи спасения моего. Не остави мене, Господи Боже мои, не отступи от мене. Воньми в помощь мою, Господи спасения моего.

Не остави мене, Господи Боже мои, не отступи от мене. Воньми в помощь мою, Господи спасения моего.


\mysubtitle{Псалом 62}


Боже, Боже мой, к Тебе утренюю, возжада Тебе душа моя, коль множицею Тебе плоть моя, в земли пусте и непроходне и безводне. Тако во святем явихся Тебе, видети силу Твою и славу Твою. Яко лучши милость Твоя паче живот, устне мои похвалите Тя. Тако благословлю Тя в животе моем, о имени Твоем воздежу руце мои. Яко от тука и масти исполнися душа моя, и устне радости восхвалят Тя уста моя. аще поминах Тя на постели моей, на утрених поучахся в Тя. Яко бысть Помощник мой, и в крове крилу Твоею возрадуюся. Прильпе душа моя по Тебе, мене же прият десница Твоя. Тии же всуе искаша душу мою, внидут в преисподняя земли. Предадятся в руки оружию, части лисовом будут. Царь же возвеселится о Бозе, похвалится всяк кленыися Им. Яко заградишася уста глаголющих неправду.

На утрених поучахся в Тя. Яко бысть Помощник мой, и в крове крилу Твоею возрадуюся. Прильпе душа моя по Тебе, мене же прият десница Твоя. 

Слава, и ныне. Аллилуия, Аллилуия, слава Тебе, Боже (трижды, без поклонов). Господи помилуй \myemph{(трижды)}. Слава, и ныне.


\mysubtitle{Псалом 87}


Господи Боже спасения моего, во дни возвах и в нощи пред Тобою. Да внидет пред Тя молитва моя, приклони ухо Твое к молению моему. Яко исполнися зол душа моя, и живот мой аду приближися. Привменен бых с низходящими в ров, бых яко человек без помощи в мертвых свободь. Яко язвении спящии во гробех, их же не помянул еси к тому и тии от руки Твоея отриновени быша. Положиша мя в рове преисподнем, в темных и сени смертней. На мне утвердися ярость Твоя, и вся волны Твоя наведе на мя. Удалил еси знаемых моих от мене, положиша мя мерзость себе, предан бых и не исхождах. Очи мои изнемогосте от нищеты. Возвах к Тебе Господи, весь день, воздех к Тебе руце мои. Еда мертвыми твориши чудеса? ли врачеве воскресят и исповедятся Тебе? Еда повесть кто во гробе милость Твою, и истину Твою в погибели? Еда познана будут во тме чудеса Твоя, и правда Твоя в земли забвенней? Аз к Тебе Господи, возвах и утро молитва моя предварит Тя. Вскую Господи, отрееши душу мою? отвращаеши лице Свое от мене. Нищь есмь аз, и в трудех от юности моея, вознес же смирихся и изнемогох. На мне преидоша гневи Твои, устрашения Твоя возмутиша мя. Обыдоша мя яко вода, весь день одержаша мя вкупе. Удалил еси от мене друга и искреняго, и знаемых моих от страстей.

Господи Боже спасения моего, во дни возвах и в нощи пред Тобою. Да внидет пред Тя молитва моя, приклони ухо Твое к молению моему.


\mysubtitle{Псалом 102}


Благослови, душе моя, Господа, и вся внутреняя моя имя святое Его. Благослови, душе моя, Господа, и не забывай всех воздаянии Его. Очищающаго вся беззакония твоя, исцеляющаго вся недуги твоя. Избавляющаго от истления живот твой, венчающаго тя милостию и щедротами. Исполняющаго во благих желание твое, обновится яко орлу юность твоя. Творяи милостыню Господь, и судьбу всем обидимым. Сказа пути Своя Моисеови, сыновом израилевым хотения Своя. Щедр и милостив Господь, долготерпелив, и многомилостив. Не до конца прогневается, ни в век враждует. Не по беззаконием нашим сотворил есть нам, ни по грехом нашим воздал есть нам. Яко по высоте небесней от земли, утвердил есть Господь милость Свою на боящихся Его. Елико отстоят востоцы от запад, удалил есть от нас беззакония наша. Яко же щедрит отец сыны, ущедрит Господь боящихся Его. Яко Той позна создание наше, помяну яко персть есмы. Человек, яко трава, дние его, яко цвет сельныи, тако оцветет. Яко дух пройдет в нем и не будет, и не познает к тому места своего. Милость же Господня от века и до века, на боящихся Его. И правда Его на сынех сынов, хранящих завет Его, и помнящих заповеди Его творити я. Господь на небеси уготова престол Свой, и царство Его всеми обладает. Благословите Господа вси ангели Его, сильнии крепостию творящии слово Его, услышати глас словес Его. Благословите Господа вся силы Его, слуги Его, творящии волю Его. Благословите Господа вся дела Его, на всяком месте владычества Его, благослови, душе моя Господа.

На всяком месте владычествия Его, благослови душе моя Господа.


\mysubtitle{Псалом 142}


Господи, услыши молитву мою, внуши моление мое воистине Твоей, услыши мя в правде Твоей. И не вниди в суд с рабом Твоим, яко не оправдится пред Тобою всяк живыи. Яко погна враг душу мою, и смирил есть в земли живот мой. Посадил мя есть в темных, яко мертвыя веку и уны во мне дух мой, во Мне смутися сердце мое. Помянух дни древния, поучихся во всех делех Твоих, и в делех руку Твоею поучахся, воздех к Тебе руце мои, душа моя, яко земля безводная Тебе. Скоро услыши мя Господи, исчезе дух мой. Не отврати лица Твоего от мене, и уподоблюся низходящим в ров. Слышану сотвори мне заутра милость Твою, яко на Тя уповах. Скажи мне Господи, путь, воньже пойду, яко к Тебе взях душу мою. Изми мя от враг моих, Господи, к Тебе прибегох. Научи мя творити волю Твою, яко Ты еси Бог мой. Дух Твой благии наставит мя на землю праву. Имене Твоего ради Господи, живиши мя, правдою Твоею изведеши от печали душу мою. И милостию Твоею потребиши враги моя, и погубиши вся стужающия души моей, яко аз раб Твой есмь.

Услыши мя в правде Твоей Господи, услыши, и не вниди в суд с рабом Твоим. Дух Твой Благии наставит мя на землю праву. 

Слава, и ныне. Аллилуия, Аллилуия, слава Тебе, Боже (трижды, с поклонами). Господи помилуй (12). Слава, и ныне. 

Таже, Бог Господь и явися нам, благословен грядыи во имя Господне.

Стих 1. Исповедайтеся Господеви, яко благ, яко в век милость Его.


\myemph{ Стих 2.}


Обышедше обыдоша мя, и именем Господним противляхся им.


\myemph{ Стих 3.}


Не умру, но жив буду, и повем дела Господня.


\myemph{ Стих 4.}


Камень, Его же небрегоша зиждущии, Сей бысть во главу углу, от Господа бысть Сей, и есть дивна во очию нашею. 

Таже тропарь дню, или святому, или празднику.


\myemph{ Аще ли пост, поется Аллилуия, на глас октая. Стихи же глаголем:}


\myemph{ Стих 1.}


От нощи утренюет дух мой к Тебе, Боже.


\myemph{ Стих 2.}


Правде научитеся живущии на земли.


\myemph{ Стих 3.}


Господи Боже наш, мир даждь нам, вся бо воздал еси нам.


\myemph{ Стих 4.}


Тако быхом возлюбленному Твоему.


Таже троичны во октаи, и обычная стихология, и седальны, и прочая по уставу.


\mysubtitle{Таже, псалом 50}


Помилуй мя Боже…


\myemph{ Посем песни пророческия и каноны. На 9 песни канона поем песнь Пресвятей Богородице (от Луки, зачало 4):} Величит душа моя Господа, и возрадовася дух мой о Бозе Спасе моем (правый лик). И по коемждо стисе припеваем: Честнейшую херувим, и славнейшую воистину серафим, без истления Бога Слова рождьшую, сущую Богородицу Тя величаем. И поклон до земли. Аще ли неделя, или празднество, поклоны творим малыя, токмо последнии поклон земныи.

Яко призре на смирение Рабы Своея, се бо от ныне блажат Мя вси роди (левый лик). Яко сотвори мне величие Сильныи, и свято имя Его, и милость Его в род и род на боящихся Его (правый лик). 

Сотвори державу мышцею Своею, расточи гордыя мысли сердца их (левый лик). Низложи сильныя со престол и вознесе смиренныя, алчущия исполни благ, и богатящияся отпусти тща (правый лик). 

Восприят Израиля отрока Своего, помянути милости, якоже глагола ко отцем нашим, Аврааму и семени его до века (левыи лик).


\myemph{ Таже, песнь 9-я канонов. По 9-й песни оба лика совокупльшеся поем:} Достойно есть, и творим поклон вси вкупе до земли. Таже, ектения малая. Аще ли неделя; поем на глас октая: Свят Господь Бог наш \myemph{(трижды)}, таже светилен.


\mysubtitle{Посем, псалом 148}


Хвалите Господа с небес, хвалите Его в вышних. Хвалите Его вси ангели Его. Хвалите Его вся силы Его. Хвалите Его солнце и луна, хвалите Его вся звезды и свет. Хвалите Его небеса небес и вода яже превыше небес, да восхвалят имя Господне. Яко Той рече, и быша; Той повеле, и создашася. Постави я в век и в век века. Повеление положи, и не мимо идет. Хвалите Господа от земли, змиеве и вся бездны. Огнь, град, снег, голоть, дух бурен, творящая слово Его. Горы и вси холми, древа плодоносна и вси кедри. Зверие и вси скоти, гади и птицы пернаты. Царие земстии и вси людие, князи и вси судии земстии. Юноши и девы, старцы с юнотами, да восхвалят имя Господне. Яко вознесеся имя Того единаго. Исповедание Его на земли и на небеси, и вознесет рог людей Своих. Песнь всем преподобным Его, сыновом израилевом, людем приближающимся Ему.


\mysubtitle{Псалом 149}


Воспойте Господеви песнь нову, хваления Его в церкви преподобных. Да возвеселится Израиль о сотворшем Его, и сынове Сиони возрадуются о Цари своем. Да восхвалят имя Его в лице, в тимпане и псалтыри да поют Ему. Яко благоволит Господь в людех Своих, и вознесет кроткия во спасение. Восхвалятся преподобнии во славе, и возрадуются на ложах своих. Возношения Божия в гортани их, и мечи обоюду остры в руках их. Сотворити отмщение во языцех, обличение в людех. Связати царя их путы, и славныя их ручными оковы железными. Сотворити в них суд написан, слава си есть всем преподобным Его.


\mysubtitle{Псалом 150}


Хвалите Бога во святых Его, хвалите Его во утвержении сил Его. Хвалите Его на силах Его, хвалите Его по премногому величествию Его. Хвалите Его во гласе трубнем, хвалите Его во псалтыри и гуслех. Хвалите Его в тимпане и лице, хвалите Его во струнах и органех. Хвалите Его в кимвалех доброгласных, хвалите Его в кимвалех восклицания. Всяко дыхание да хвалит Господа.


\myemph{ Таже, стихеры хвалитныя, и Слава, и И ныне по уставу. Посем глаголем:} Слава показавшему нам свет, и прочая. Аще ли на хвалитех стихер нет и мы глаголем: Тебе слава подобает, Господи Боже наш, и Тебе славу возсылаем, Отцу и Сыну и Святому Духу, и ныне и присно и во веки веком, аминь. Слава показавшему нам свет. 

Слава в вышних Богу, и на земли мир, в человецех благоволение. Хвалим Тя, благословим Тя \myemph{(поклон)}, кланяемтися, славословим Тя \myemph{(поклон)}, благодарим Тя великия ради славы Твоея \myemph{(поклон)}. Господи Царю небесныи, Боже Отче Вседержителю, и Господи Сыне Единородныи Исусе Христе, и Святыи Душе: Господи Боже, Агньче Божии, Сыне Отечь, вземляи грехи мира, помилуй нас. Вземляи грехи мира, приими молитвы наша. Седяи одесную Отца, помилуй нас. Яко Ты един свят, Ты един Господь Исус Христос, в славу Богу Отцу, аминь. На всяку нощь благословлю Тя, и восхвалю имя Твое во веки, и в век века. Господи, Прибежище бысть нам, в род и род. Аз рех: Господи, помилуи мя и исцели душу мою, яко согреших Тебе. Господи, к Тебе прибегох, научи мя творити волю Твою, яко Ты еси Бог мой, яко от Тебе есть источник живота, во свете Твоем узрим свет. Пробави милость Твою ведущим Тя. Сподоби Господи, в нощь сию без греха сохранитися нам. Благословен еси, Господи Боже отец наших, и хвально и прославлено имя Твое во веки, аминь. Буди Господи, милость Твоя на нас, якоже уповахом на Тя. Благословен еси Господи, научи нас оправданием Твоим. Благословен еси Владыко, вразуми нас оправданием Твоим. Благословен еси Святыи, просвети нас оправданием Твоим. Господи, милость Твоя во веки, и дела руку Твоею не презри. Тебе подобает хвала, Тебе подобает пение. Тебе слава подобает: Отцу и Сыну и Святому Духу, ныне и присно и во веки веком, аминь.


\myemph{ Посем ектения.} Исполним заутрения. Таже, на стиховне стихеры. Посем: Слава, и ныне, Богородичен.


\myemph{ Таже:} Благо есть исповедатися Господеви, и пети имени Твоему, Вышнии, возвещати заутра милость Твою, и истину Твою на всяку нощь. Аще есть пост, глаголем дважды. Таже, Трисвятое. И по Отче наш, тропарь дню, или святому по уставу, и ектения и по возгласе глаголет псаломщик: Аминь. Приидите, поклонимся, трижды. И псалмы Перваго часа. И прочее. И отпуст.


\myemph{ Аще ли пост, то глаголем по Отче наш тропари сии:} В церкви стояще славы Твоея, на небеси стояти мнимся, Богородице дверь небесная, отверзи нам двери милости Твоея. Господи помилуй (40). Господи благослови. Господи Исусе Христе Сыне Божии, помилуй нас. Глаголем: Аминь. Небесныи Царю, державу нашу укрепи, веру утверди, языки укроти, мир умири, и святыи храм сей добре сохрани, и прежде отшедшия отцы и братию нашу в кровех с праведными учини и нас в православней вере и в покаянии, Господи, приими и помилуй, яко Благ и Человеколюбец. Таже, Господи помилуй \myemph{(трижды)}. Слава, и ныне. Честнейшую херувим, и славнейшую воистину серафим, без истления Бога Слова рождьшую, сущую Богородицу Тя величаем (поклон великий). Именем Господним благослови, отче.


\myemph{ И творим великия поклоны, с молитвою святаго Ефрема, якоже предуказася на вседневней полунощнице. }

\myemph{ И абие начинаем первый час.}


\end{mymulticols}

\section{Девятый час}\begin{mymulticols}
%http://www.molitvoslov.org/text1014.htm 
 


Приидите, поклонимся (трижды, и поклоны три)


\mysubtitle{Псалом 83}


Коль возлюблена села Твоя, Господи сил, желает и скончевается душа моя во дворы Господня. Сердце мое и плоть моя возрадовастася о Бозе живе. Ибо птица обрете себе храмину, и горлица гнездо себе, идеже положи птенцы своя. Олтари Твоя, Господи Боже сил, Царю мой и Боже мой. Блажени живущии в дому Твоем, в век века восхвалят Тя. Блажен муж ему же есть заступление его от Тебе. Восхождения в сердцы своем положи, во юдоли плачевне, в месте идеже положи. Ибо благословение даст закон даяи. Пойдут от силы в силу, явится Бог богом в Сионе. Господи Боже сил, услыши молитву мою, внуши Боже Ияковль. Защитниче наш, виждь Боже, и призри на лице Христа Твоего. Яко лучше день един во дворех Твоих, паче тысящ. Изволих приметатися в дому Бога моего паче, неже жити ми в селех грешничих. Яко милость и истину любит Господь. Бог благодать и славу даст. Господь не лишит блага, ходящих незлобием. Господи Боже сил, блажен человек уповаяи на Тя.


\mysubtitle{Псалом 84}


Благоволил еси Господи, землю Твою, возвратил еси плен Ияковль Оставил еси беззакония людем Твоим, покрыл еси вся грехи их. Укротил еси весь гнев Твой, возвратил еси от гнева ярости Твоея. Возврати нас, Боже спасении наших, и возврати ярость Твою от нас. Еда во веки прогневаешися на ны, или простреши гнев свой от рода в род. Боже, Ты обращ живиши ны, и людие Твои возвеселятся о Тебе. Яви нам Господи, милость Твою, и спасение Твое даждь нам. Услышу, что речет о мне Господь Бог; яко речет мир на люди Своя и на преподобныя Своя, и на обращающия сердца к Нему. Обаче близ боящихся Его, спасение Его, вселити славу в землю нашу. Милость и истина сретостеся, правда и мир облобызастася. Истина от земли восия, и правда с небесе приниче. Ибо Господь даст благость, и земля наша даст плод свой. Правда пред Ним предъидет, и положит в путь стопы своя.


\mysubtitle{Псалом 85}


Приклони Господи, ухо Твое и услыши мя, яко нищ и убог есмь аз. Сохрани душу мою, яко преподобен есмь, спаси раба Твоего, Боже мой, уповающаго на Тя. Помилуй мя Господи, яко к Тебе воззову весь день. Возвесели душу раба Твоего, яко к Тебе взях душу мою. Яко Ты Господи, благ, и долготерпелив, и многомилостив всем призывающим Тя. Внуши Господи, молитву мою, и воньми гласу моления моего. В день печали моея возвах к Тебе, яко услыша мя. Несть подобен Тебе в бозех Господи, и несть по делом Твоим. Вси языцы елико сотвори приидут, и поклонятся пред Тобою Господи, и прославят имя Твое. Яко велии еси Ты, и творяи чудеса, Ты еси Бог един. Настави мя Господи, на путь Твой, и пойду во истине Твоей: да возвеселится сердце мое, боятися имене Твоего. Исповемся Тебе, Господи Боже мой, всем сердцем моим, и прославлю имя Твое в век. Яко милость Твоя велия на мне, и избави душу мою от ада преисподняго. Боже, законопреступницы восташа на мя, и сонм державных взыскаша душу мою, и не предложиша Тебе пред собою. И Ты Господи Боже мой, щедр и милостив, долготерпелив и многомилостив и истинен. Призри на мя и помилуй мя, даждь державу Твою отроку Твоему, и спаси сына рабы Твоея. Сотвори со мною знамение во благо, и да видят ненавидящии мя и постыдятся; яко Ты Господи, поможе ми, и утешил мя еси.

Слава, и ныне. Аллилуия \myemph{(трижды)}. Господи помилуй \myemph{(трижды)}. Тропарь по уставу.


\myemph{ Аще ли же пост, глаголем сий тропарь, глас 8:} Иже в девятыи час, нас ради плотию смерть вкусивыи, умертви плоти нашея мудрование, Христе Боже, спаси нас.


\myemph{ Стих 1:}


Да приближится молитва моя пред Тя Господи, по словеси Твоему вразуми мя.


\myemph{ Стих 2:}


Да внидет прошение мое пред Тя Господи, по словеси Твоему избави мя.


Слава, и ныне. Иже нас ради рождеися от Девы, и распятие претерпев Благии, испровергии смертию смерть, и воскресение явль яко Бог, не презри их же созда рукою Своею, яви человеколюбие Свое Милостиве, приими рождьшую Тя Богородицу, молящуся заны, и спаси, Спасе наш, люди согрешьшия.


\myemph{ Аще ли пост, бывает чтение. Таже:} Не предаждь нас до конца, имене Твоего ради, и не разори завета Твоего, и не отстави милости Твоея от нас, Авраама ради возлюбленнаго от Тебе, и за Исаака раба Твоего, и Израиля святаго Твоего. Таже, Трисвятое. И по Отче наш, кондак дню, или святому.


\myemph{ Аще ли же пост, глаголем тропари сия, глас 8:} Видя разбойник Начальника жизни на Кресте висяща, и глаголаше: аще небы Бог был воплощься, Иже с нами распныися, небы солнце лучи свои потаило, ниже земля трепещущи тряслася, но Иже за вся терпяи, помяни мя Спасе, егда приидеши во царствии Си. 

Слава. Посреди двою разбойнику мерило праведное обретеся Крест Твой: овому убо низходящу во ад тяготою хуления, другому же легчащуся от грех к разуму богословия. Христе Боже, слава Тебе. 

И ныне. Агньца и пастыря и Спаса миру, на Кресте зрящи рождьшая Тя, глаголаше слезящии: мир убо радуется, приемля Тобою избавление, утроба же Моя горит, зрящи Твое распятие, еже за всех терпиши, Сыне и Боже мой. 


\myemph{ Таже,} Господи помилуй (40). Иже на всяко время… Господи помилуй \myemph{(трижды)}. Слава, и ныне. Честнейшую херувим… Именем Господним благослови, отче. За молитв святых отец наших… Аминь. И сотворяем три великия поклоны, со обычною молитвою: Господи и Владыко животу моему… И прочих метаний (12). Посем молитва. Владыко Боже, Отче Вседержителю…


\myemph{ Аще ли великий пост, и мы по поклонех молитвы:} Владыко Боже, Отче… не глаголем, но поем на оба лика: Во царствии Си…


\myemph{ Аще ли кроме поста, глаголем молитву сию:} Владыко, Господи Исусе Христе, Боже наш, долготерпеливе о наших согрешениих, даже и до нынешняго часа приведыи нас, воньже на животворящем древе вися, благоразумному разбойнику, Иже в рай путь сотворил еси вход, и смертию смерть разрушил еси, очисти нас грешных, и недостойных раб Твоих. Согрешихом бо и беззаконновахом, и несмы достойни возвести очию нашею, и возрети на высоту небесную, зане же оставихом путь правды Твоея, и ходихом в волях сердец наших, но молим Твою благость: пощади нас Господи, по множеству милости Твоея, и спаси нас имене Твоего ради святаго, яко исчезоша в суете дние наши. Изми нас из руки сопротивнаго, и остави нам прегрешения наша, и умертви плотьское наше мудрование, да ветхаго отложивше человека, и в новаго облецемся, и Тебе поживем, нашему Владыце и Благодателю, и тако Твоим повелением последующе, в вечныи покой достигнем, идеже есть всем веселящимся жилище. Ты бо еси воистину истинное веселие, и радость любящим Тя, Христе Боже наш, и Тебе славу возсылаем, со безначальным Ти Отцем, и с Пресвятым и Благим и Животворящим Ти Духом, ныне и присно и во веки веком, аминь.


\end{mymulticols}

\section{Павечерница великая}\begin{mymulticols}
%http://www.molitvoslov.org/text1006.htm 
 


Начало Великия павечерницы. 

Певаема бывает в Великии пост, и в пост святых Апостол, и в пост Святыя Богородицы, и в пост Рожества Христова, егда есть Аллилуия, с поклоны.

По обычнем начале глаголем:

За молитв святых отец наших, Господи Исусе Христе, Сыне Божии, помилуй нас.

Аминь.

Царю Небесныи. Трисвятое. И по Отче наш. 

Молитва Исусова. 

Господи помилуй (12).

Слава, и ныне. 

Приидите, поклонимся \myemph{(трижды)}.


\mysubtitle{Аще ли есть первая неделя Великаго поста, глаголем, псалом 69:}


Боже, в помощь мою воньми. Господи, помощи ми потщися. Да постыдятся и посрамятся ищущии душу мою. Да возвратятся вспять и постыдятся, мыслящии ми злая. Да возвратятся абие стыдящеся, глаголющии ми: благо же, благо же. Да возрадуются и возвеселятся о Тебе, вси ищущии Тебе Боже, и глаголют выну: да возвеличится Господь, любящии спасение Твое. Аз же нищь есмь и убог, Боже помози ми. Помощник мой и Избавитель мой еси Ты. Господи, не закосни.

Посем канон святаго Андрея Критскаго.

По каноне же глаголем псалом, 4-й Внегда возвах.

Аще ли кроме первыя недели святаго поста по изглаголании: Приидите, поклонимся, глаголем сия псалмы.

Внегда возвах, услыша мя Боже правды моея, в скорби распространил мя еси. Ущедри мя, и услыши молитву мою. Сынове человечестии, доколе тяжкосердии; вскую любите суетная, и ищете лжи; И уведите, яко удиви Господь преподобнаго Своего, Господь услышит мя внегда взову к Нему. Гневайтеся и не согрешайте, яже глаголете в сердцах ваших, на ложах ваших умилитеся. Пожрите жертву правды, и уповайте на Господа, мнози глаголют: кто явит нам благая? 3наменася на нас свет лица Твоего Господи, дал еси веселие в сердцы моем. От плода пшеницы, вина и елея своего умножишася. В мире вкупе усну и почию, яко Ты Господи, единаго на упование вселил мя еси.


\mysubtitle{Псалом 6}


Господи, да не яростию Твоею обличиши мене, ниже гневом Твоим покажеши мене. Помилуй мя Господи, яко немощен есмь, исцели мя Господи, яко смутишася кости моя. И душа моя смутися зело, и Ты Господи, доколе? Обрати Господи, изми душу мою, спаси мя ради милости Твоея. Яко несть в смерти поминаяи Тебе, во аде же кто исповестьтися? Утрудихся воздыханием моим, измыю на всяку нощь ложе мое, слезами моими постелю мою омочу. Смутися от ярости око мое, обетшах во всех вразех моих. Отступите от мене вси творящии беззаконие, яко услыша Господь глас плача моего. Услыша Господь моление мое, Господь молитву мою прият. Да постыдятся и смятутся вси врази мои, возвратятся и устыдятся зело вскоре.


\mysubtitle{Псалом 12}


Доколе Господи, забудеши мя до конца? Доколе отвращаеши лице Свое от мене? Доколе положу советы в души моей, болезни в сердцы моем день и нощь? Доколе вознесется враг мой на мя? Призри и услыши мя, Господи Боже мой. Просвети очи мои да некогда усну в смерть, да некогда речет враг мой: укрепихся нань. Стужающии ми возрадуются, аще ся подвижу, аз же на милость Твою уповах. Возрадуется сердце мое о спасении Твоем, воспою Господеви благодеявшему мне, и пою имени Господню Вышнему.

Слава, и ныне. 

Аллилуия, Аллилуия, слава Тебе, Боже (трижды и поклоны три). 

Господи помилуй \myemph{(трижды)} . Слава, и ныне.


\mysubtitle{Псалом 24}


К Тебе Господи, воздвигох душу мою, Боже мой, на Тя уповах, да не постыжуся во веки. Ни посмеютмися врази мои, ибо вси терпящии Тя не постыдятся. Да постыдятся беззаконнующии всуе. Пути Твоя Господи, скажи ми, и стезям Твоим научи мя. Настави мя на истину Твою, и научи мя, яко Ты еси Бог Спас мой, и Тебе терпех весь день. Помяни щедроты Твоя Господи, и милости Твоя, яко от века суть. Грех юности моея, и невежествия моего не помяни, по милости Твоей помяни мя Ты, ради благости Твоея, Господи. Благ и прав Господь, сего ради закон положит согрешающим на пути. Наставит кроткия на суд, и научит кроткия путем Своим. Вси путие Господни милость и истина, взыскающим завета Его и свидения Его. Ради имене Твоего Господи, и очисти грех мой, мног бо есть. Кто есть человек бояися Господа; закон положит ему на пути иже изволи. Душа его во благих водворится, и семя его наследит землю. Держава Господь боящимся Его, и завет Его явит им. Очи мои выну ко Господу, яко Той исторгнет от сети нозе мои. Призри на мя и помилуй мя, яко единород и нищь есмь аз. Скорби сердца моего умножишася, и от бед моих изведи мя. Виждь смирение мое и труд мой, и отпусти вся грехи моя. Виждь враги моя, яко умножишася, и ненавидением неправедным возненавидеша мя. Сохрани душу мою и избави мя, да не постыжуся, яко уповах на Тя. Незлобивии и правии прилепляхуся мне, яко потерпех Тя, Господи. Избави Боже, Израиля от всех скорбий его.


\mysubtitle{Псалом 30}


На Тя Господи, уповах, да не постыжуся во веки, правдою Твоею избави мя, и изми мя. Приклони ко мне ухо Твое, ускори изъяти мя. Буди ми в Бог защититель, и в дом прибежища спасти мя. Яко держава моя и прибежище мое еси Ты, и имене Твоего ради наставиши мя, и препитаеши мя. Исторгнеши мя от сети сея, юже скрыша ми, яко Ты еси Защититель мой, Господи. В руце Твои предаю дух мой, избавил мя еси, Господи Боже истине. Возненавидел еси хранящия суетная затщее, аз же на Господа уповах. Возрадуюся и возвеселюся о милости Твоей, яко призрел еси на смирение мое, спасл еси от бед душу мою. И неси мене затворил в руках вражииих, поставил еси на пространне нозе мои. Помилуй мя Господи, яко скорблю, смутися от ярости око мое, душа моя и утроба моя. Яко исчезе в болезни живот мой, и лета моя в воздыханиих, изнеможе нищетою крепость моя, и кости моя смутишася. От всех враг моих бых поношение, и соседом моим зело, и страх знаемым моим. Видящии мя вон бежаша от мене, забвен бых яко мертв от сердца. Бых яко сосуд погублен, яко слышах гаждение многих живущих окрест, внегда собратися им вкупе на мя, прияти душу мою совещаша. Аз же на Тя Господи, уповах, рек: Ты еси Бог мой, в руку Твоею жребии мои. Избави мя из руки враг моих, и от гонящих мя. Просвети лице Твое на раба Твоего, спаси мя милостию Твоею Господи, да не постыжуся, яко призвах Тя. Да постыдятся нечестивии, и снидут во ад. Немы да будут устны льстивыя, глаголющия на праведнаго беззаконие, гордынею и уничижением. Коль много множество благости Твоея Господи, юже скрыл еси боящимся Тебе, соделал еси уповающим на Тя, пред сыны человеческими. Скрыеши их в тайне лица Твоего, от мятежа человеческа, покрыеши их в крове от пререкания язык. Благословен Господь, яко удиви милость Свою во граде обстояния. Аз же рех во изступлении моем: отвержен есмь от лица очию Твоею; сего ради услыша глас молитвы моея, внегда возвах к Тебе. Возлюбите Господа вси преподобнии Его, яко истины взыскает Господь, и воздает излиха творящим гордыню. Мужайтеся и да крепится сердце ваше, вси уповающии на Господа.


\mysubtitle{Псалом 90}


Живыи в помощи Вышняго, в крове Бога небеснаго водворится. Речет Господеви: Заступник мой еси, и прибежище мое, Бог мой, и уповаю Нань. Яко Той избавит тя от сети ловчи, и от словесе мятежна. Плещьма Своима осенит тя, и под криле Его надеешися. Оружие обыдет тя истина Его, не убоишися от страха нощнаго, от стрелы летящия во дни. От вещи во тме преходящия, от сряща и беса полуденнаго. Падет от страны твоея тысяща, и тма одесную тебе, к тебе же не приближится. Обаче очима своима смотриши, и воздаяние грешииком узриши. Яко Ты Господи, упование  мое; Вышняго положил еси прибежище твое. Не приидет к тебе зло, и рана не приближится телеси твоему. Яко ангелом Своим заповесть о тебе, сохранити тя во всех путех твоих. На руках возмут тя, да некогда преткнеши о камень ноги твоея. На аспида и василиска наступиши, и попереши льва и змия. Яко на Мя упова, и избавлю и, покрыю и, яко позна имя Мое. Воззовет ко Мне, и услышу и, с ним есмь в скорби, изму и, и прославлю его. Долготу дний исполню и, и явлю ему спасение Мое.


Посем: Слава, и ныне. 

Аллилуия, Аллилуия, слава Тебе, Боже (трижды и поклоны три).

Господи помилуй \myemph{(трижды)}. Слава, и ныне.

Таже, песнь святаго пророка Исаии.

С нами Бог, разумейте языцы и покаряйтеся, яко с нами Бог.

Услышите и до последних земли, яко с нами Бог. Могущии покаряйтеся, яко с нами Бог. Аще бо паки возможете, и паки побеждени будете, яко с нами Бог. И иже аще совет совещаваете, и разорит и Господь, яко с нами Бог. И слово еже аще возглаголете, не имать пребывати в вас, яко с нами Бог. Страха же вашего не имам убоятися, ниже соблазнитися, яко с нами Бог. Господа же Бога нашего святите, и Той будет вам в боязнь, яко с нами Бог. И иже аще Нань надеешися, и Той будет ти во освящение, яко с нами Бог. И уповающе будем Нань, и спасемся Его ради, яко с нами Бог. Се аз и дети, яже ми дал Бог, яко с нами Бог. Людие ходящии во тме, видеша свет велии, яко с нами Бог. Живущии во стране и сени смертней, свет восияет на вы, яко с нами Бог. Яко Отроча родися нам, Сын и дастся нам, яко с нами Бог. Ему же власть бысть на раме Его, яко с нами Бог. И миру Его несть предела, яко с нами Бог. И нарицается имя Его велика совета Ангел, яко с нами Бог. Чуден Советник, яко с нами Бог. Бог крепок, Владыка, Князь смирения, яко с нами Бог. Отец будущаго века, яко с нами Бог.

Слава: Яко с нами Бог.

И ныне: Яко с нами Бог. 

Таже, сей стих глаголем трижды: С нами Бог, разумейте языцы и покаряйтеся, яко с нами Бог. 

Таже: День пребыв благодарю Тя  Господи, вечер молюся с нощию, безнаветие подаждь ми Спасе, и спаси мя.

Слава. День изшед славословлю Тя Владыко, вечер молюся с нощию, безсоблазньство подаждь ми Спасе, и спаси мя.

И ныне. День прешед песнословя Тя Святыи, вечер молюся с нощию, без греха сохрани мя Спасе, и спаси мя.

Безплотное естество херувимское, немолчными песньми Тя славословят, шестокрыльная и животная серафим, непрестанными гласы Тя превозносят, ангельская же вся воинства, трисвятыми песньми Тя восхваляют. Прежде бо всех еси сыи Отец, и собезначальна имаши Своего Сына, и равночестен нося Дух жизни, Троицу являеши нераздельну.

Пресвятая Дево, Мати Божия, и иже Слову самовидцы и слуги, пророк же и мученик, и вси ликове, яко безсмертну имуще жизнь, о всех молитеся прилежно, яко вси есмы в бедах, да льсти избавльшеся лукаваго, ангельскую вопием песнь: Святыи, Святыи, Святыи, Трисвятыи Господи, помилуй и спаси нас. 

Таже, исповедание православныя веры, перваго Собора:

Верую во единаго Бога Отца, Вседержителя, Творца небу и земли, видимым же всем и невидимым. И во единаго Господа Исуса Христа, Сына Божия, Единороднаго, Иже от Отца рожденнаго прежде всех век. Света от света, Бога истинна от Бога истинна, рожденна, а не сотворена, единосущна Отцу, Им же вся быша. Нас ради человек, и нашего ради спасения, сшедшаго с небес, и воплотившагося от Духа Свята, и Марии Девы вочеловечьшася. Распятаго за ни при понтийстем Пилате, страдавша и погребена, и воскресшаго в третий день по писаниих. И возшедшаго на небеса, и седяща одесную Отца. И паки грядущаго со славою, судити живым и мертвым, Его же царствию несть конца.

Втораго Собора: И в Духа Святаго, Господа истиннаго и животворящаго, Иже от Отца исходящаго, Иже со Отцем и Сыном спокланяема и сславима, глаголавшаго пророки. И во едину святую соборную и апостольскую Церковь. Исповедую едино Крещение, во оставление грехов. Чаю воскресения мертвым. И жизни будущаго века, аминь.

Посем глаголет, начинаяи лик:

О! Пресвятая Госпоже Владычице Богородице Всемилостивая, моли Бога о нас грешных. 

И творим метание на кииждо стих. И другии лик той же стих глаголет. Наченыи же глаголет: 

Вся небесныя силы, святии ангели и архангели, молите Бога о нас грешных.

Святыи великии Иоанне пророче и Предотече, Крестителю Господень, моли Бога о нас грешных.

Святии славнии апостоли, пророцы и мученицы, и вси святии, молите Бога о нас грешных.

Преподобнии и богоноснии отцы наши, пастырие и учителие вселенней, молите Бога о нас грешных.

Зде рцы святаго, его же есть храм.

Непобедимая и божественная сила, честнаго и животворящаго Креста Господня, не остави нас грешных, уповающих на Тя. 

Боже, милостив буди нам грешным. 

Егда же достигнем сего стиха: Боже, очисти грехи наша, и помилуй нас, глаголем его трижды.

Посем, Трисвятое, и поклоны обычныя (3). 

По Отче наш. Молитва Исусова.


\myemph{ В понедельник и среду вечер, тропари сия глас 2.}

Просвети очи мои Христе Боже, да некогда усну в смерть, да некогда речет враг мой: укрепихся нань.

Стих. Призри и услыши мя, Господи Боже мой. 

Паки той же тропарь: Просвети очи мои Христе Боже, да некогда усну в смерть, да некогда речет враг мой: укрепихся нань.

Слава. Заступник души моей буди Боже, яко посреде хожду сетей многих. Избави мя от них, и спаси мя Блаже, яко Человеколюбец.

И ныне, Богородичен. Яко не имамы дерзновения, за премногия грехи наша, но Ты, Иже от Тебе рождьшагося, моли Богородице Дево. Много бо может молитва Матерня, на умоление Владыки. Не презри грешных мольбы Всечистая, яко милостив есть, и спасти могии, Иже страдати нас ради изволивыи.


\myemph{ Ины тропари, во вторник и в четверток вечер, глас 8.}

Сия обои тропари пременяти чрез день.

Невидимых враг моих неусыпание свеси Господи, и окаянныя плоти моея неможение веси, Создавыи мя. Тем же в руце Твои предаю дух мой, покрый мя крилома Твоея благости, да некогда усну в смерть. И умнии мои очи просвети, наслаждением божественных словес Твоих, и воздвигни мя во время  благопотребно, к Твоему славословию, яко един Благ и Человеколюбец.

Стих. Призри и услыши мя, Господи Боже мой. 

Яко страшен суд Твой Господи, ангелом предстоящим, человеком вводимым, книгам разгибаемым, делом испытаемым, помыслом истязуемым. Кии суд будет мне, зачатому во гресех? кто ми пламень угасит? кто ми тму просветит? аще не Ты Господи, помилуеши мя Блаже, яко Человеколюбец.

Слава. Слезы ми даждь Боже, якоже древле жене грешнице. И сподоби мя омочити нозе Твои, яже мя от пути льстиваго свободившая и мvро благоухания Тебе приносити, житие чисто, покаянием ми стяжанно. Да услышу и аз благаго Твоего гласа: вера твоя спасе тя, иди в мир.

И ныне, Богородичен. Непостыдную, Богородице, надежду Твою имея спасуся, и заступление Твое стяжав Пречистая, не убоюся, но пожену враги моя, и побежду я, во едино оболкся, яко во броня, в кров Твой, и всемогущую Твою помощь. И моляся вопию Ти: Владычице, спаси мя мольбами Си, и воздвигни мя от мрачнаго сна, ко Твоему славословию, силою из Тебе воплотившагося Христа Бога.

Господи помилуй (40). 

Слава, и ныне.

Честнейшую херувим (поклон един).

Именем Господним благослови, отче.

За молитв святых отец наших. 

Аминь.


\myemph{ Таже, молитва Великаго Василия.}

Господи, Господи, избави нас от всякия стрелы летящия во дни. И избави нас от всякия вещи во тме приходящия, приими жертву вечернюю рук наших воздеяние. Сподоби же нас и нощное поприще без порока преити, неискусныи от злых, и избави нас от всякаго смущения и страха, иже от диявола нам прибывающаго. Даруй душам нашим умиление, и помыслом нашим попечение, еже на страшнем и праведнем Твоем суде испытания. Пригвозди страсе Твоем плоти наша, и умертви уды наша сущия на земли. Яко да и сонным безмолвием, просветимся зрением судеб Твоих. Отими от нас всяко мечтание неподобно, и похоть вредну. Востави же нас во время молитвы, утвержены в вере, и преспевающих в заповедех Твоих, благоволением и благостию, единороднаго Сына Твоего, с Ним же благословен еси, и с пресвятым и животворящим Ти Духом, ныне и присно и во веки веком, аминь.

Приидите, поклонимся (трижды, и поклоны три).


\mysubtitle{Псалом 50}


Помилуй мя Боже, по велицей милости Твоей. И по множеству щедрот Твоих очисти беззаконие мое. Наипаче омый мя от беззакония моего и от греха моего очисти мя. Яко беззаконие мое аз знаю и грех мой предо мною есть выну. Тебе единому согреших и лукавое пред Тобою сотворих. Яко да оправдишися в словесех Своих и победиши внегда судити. Се бо в беззакониих зачат есмь и во гресех роди мя мати моя. Се бо истину возлюбил еси, безвестная и тайная премудрости Твоея явил ми еси. Окропиши мя иссопом и очищуся. Омыеши мя и паче снега убелюся. Слуху моему даси радость и веселие: возрадуются кости смиренныя. Отврати лице Твое от грех моих и вся  беззакония моя очисти. Сердце чисто созижди во мне Боже, и дух прав обнови во утробе моей. Не отверзи мене oт лица Твоего, и Духа Твоего Святаго не отыми от мене. Воздаждь ми радость спасения Твоего и Духом Владычным утверди мя. Научу беззаконныя путем Твоим и нечестивии к Тебе обратятся. Избави мя от кровий Боже, Боже спасения моего; возрадуется язык мой правде Твоей. Господи, устне мои отверзеши, и уста моя возвестят хвалу Твою. Яко аще бы восхотел жертвы, дал бых убо; всесожжения не благоволиши. Жертва Богу "--- дух сокрушен: сердце сокрушенно и смиренно Бог не уничижит. Ублажи Господи, благоволением Твоим Сиона; и да созиждутся стены Иеросалимския. Тогда благоволиши жертву правде, возношение и всесожигаемая. Тогда возложат на олтарь Твой тельца.


\mysubtitle{Псалом 101}


Господи, услыши молитву мою, и вопль мой к Тебе да приидет. Не отврати лица Твоего от мене, вонь же день аще скорблю, приклони ко мне ухо Твое, вонь же день, аще призову Тя, скоро услыши мя. Яко исчезоша, яко дым дние мои, и кости моя, яко сушило сосхошася. Уязвен бых, яко трава, и изсше сердце мое, яко забых снести хлеб мой. От гласа воздыхания моего, прильпе кость моя плоти моей. Уподобихся неясыти пустынному, бых яко нощныи вран на нырищи. Забдех, и бых яко птица особящаяся назде. Весь день поношаху ми врази мои, и хвалящии мя мною кленяхуся. Зане пепел, яко хлеб ядях, и питие мое с плачем растворях. От лица гнева Твоего, и ярости Твоея, яко вознес низверже мя. Дние мои, яко сень уклонишася, и аз яко сено изсхох. Ты же Господи, во веки пребываеши, и память Твоя в род и род. Ты воскрес ущедриши Сиона, яко время ушедрити его, яко прииде время. Яко благоволиша раби Твои камение его, и персть его ущедрят. И убоятся языцы имене Господня, и вси царие земстии славы Твоея. Яко созиждет Господь Сиона, и явится в славе Своей. Призре на молитву убогих, и не уничижи моления их. Да напишутся сии в род их, и людие зиждемии восхвалят Господа. Яко призре с высоты святыя Своея, Господь с небесе на землю призре. Услышати воздыхание окованных, разрешити сыны умерщвенных. Возвестити в Сионе имя Господне, и хвалу Его во Иеросалиме. Внегда собратися людем вкупе, и царие работати Господеви. Отвеща ему на пути крепости его, умаление днии моих возвести ми. Не возведи мене в преполовение днии моих, в род и род лета Твоя. В началех Ты Господи, землю основа, и дела рук Твоих суть небеса. Та погибнут, Ты же пребываеши, и вся яко риза обетшают, и яко одежду свиеши их, и изменятся. Ты же тожде еси, и лета Твоя не оскудеют. Сынове раб Твоих вселятся, и семя их в век исправится


\myemph{ Таже молитва Манасии, царя иудейска:} Господи Вседержителю, Боже отец наших, Авраамов и Исааков и Ияковль, и семене их праведнаго. Сотворивыи небо и землю со всею лепотою их, и сопныи море словом повеления Твоего. Затворивыи бездну, и запечатлев ю, страшным и славным именем Твоим. Его же вся боятся, и трепещут от лица славы Твоея. Яко непостоянна велелепота славы Твоея, и не стерпим гнев еже на грешники прещения Твоего. Безчислена же и неизследованна милость обещания Твоего. Ты бо еси Господь вышнии, милосерд, долготерпелив и многомилостив, и каяся о злобах человеческих. Но Ты Господи, по множеству благости Твоея, обеща покаяние и оставление согрешшим к Тебе, и множеством щедрот Твоих нарече покаяние грешником во спасение. Ты убо Господи Боже праведных, неси положил покаяние праведным Твоим, Аврааму и Исааку и Иякову, несогрешившим пред Тобою, но положил еси покаяние мне грешному, зане согреших Ти паче числа песка морскаго. Умножишася беззакония моя Господи, умножишася, и несмь достоин воззрети и видети высоту небесную, от множества неправд моих связан есмь многими юзами железными. Яко не возвести ми главы моея, и несть ми восклонения. Зане прогневах ярость Твою, и лукавое пред Тобою сотворих. И не сотворих воли Твоея, ни сохраних повелении Твоих. И ныне поклоняю колена сердца моего, и молю яже от Тебе благость. Согреших Господи, согреших, и беззакония моя аз свем. Но прошу и молюся Тебе: отради ми Господи, отради ми, и не погуби мене со беззаконьми моими. Ниже в век враждовав соблюдеши зол моих, и не осуди мене в преисподних земли. Зане Ты еси Боже, Бог кающихся, да и на мне явиши всю благость Твою, яко недостойна суща спасеши мя, по мнозей милости Твоей. И восхвалю Тя всегда, во вся дни живота моего. Яко Тебе поют вся силы небесныя, и Твоя есть слава во веки, аминь.


\myemph{ Таже: Трисвятое. }

\myemph{ По Отче наш. Молитва Исусова.}

Аминь.


\mysubtitle{Глаголем тропари сия, глас 6:}


Помилуй нас Господи, помилуй нас, всякаго бо ответа недоумеюще, сию Ти молитву, яко Владыце, грешнии раби Твои приносим. Помилуй нас Господи, помилуй нас.

Слава. Господи, помилуй нас, на Тя бо уповахом. Не прогневайся на ны зело, ниже помяни беззаконии наших. Но призри на ны, яко Милосерд, и избави нас от враг наших. Ты бо еси Бог наш, и мы людие Твои, и вси дела руку Твоею, и имя Твое призываем.

И ныне, Богородичен. Милосердия двери отверзи нам, благословеная Богородице Дево, надеющиися на Тя не погибнем, но да избавимся Тобою от бед. Ты бо еси спасение роду християнскому.

Господи помилуй (40), Слава, и ныне.

Честнейшую херувим (поклон великий).

Именем Господним благослови, отче.

За молитв святых отец наших.

Аминь.


\myemph{ Таже молитва Великаго Василия.}


Владыко Боже Отче Вседержителю, и Господи Сыне Единородныи Исусе Христе, и Святыи Душе, едино Божество, и едина сила, помилуй мя грешнаго и ими же веси судьбами, спаси мя недостойнаго раба Твоего, яко благословен еси во веки, аминь.

Приидите, поклонимся (трижды, и поклоны три).


\mysubtitle{Аще первая неделя Великаго поста, глаголем псалом 142:}


Господи, услыши молитву мою.


\mysubtitle{В прочая же недели глаголем псалом 69}


Боже, в помощь мою вонми. Господи, помощи ми потщися. Да постыдятся и посрамятся ищущии душу мою. Да возвратятся вспять и постыдятся мыслящии ми злая. Да возвратятся абие стыдящеся, глаголющии ми: благо же, благо же. Да возрадуются и возвеселятся о Тебе, вси ищущии Тебе Боже, и глаголют выну: да возвеличится Господь, любящии спасение Твое. Аз же нищь есмь и убог, Боже помози ми. Помощник мой и Избавитель мой еси Ты, Господи, не закосни.


\mysubtitle{Псалом 142}


Господи, услыши молитву мою, внуши моление мое во истине Твоей, услыши мя в правде Твоей. И не вниди в суд с рабом Твоим, яко не оправдится пред Тобою всяк живыи. Яко погна враг душу мою, и смирил есть в земли живот мой. Посадил мя есть в темных, яко мертвыя веку, и уны во мне дух мой, во мне смутися сердце мое. Помянух дни древния, поучихся во всех делех Твоих, и в делех руку Твоею поучахся. Воздех к Тебе руце мои, душа моя яко земля безводная Тебе. Скоро услыши мя Господи, исчезе дух мой. Не отврати лица Твоего от мене, и уподоблюся низходящим в ров. Слышану сотвори мне заутра милость Твою, яко на Тя уповах. Скажи мне Господи, путь, воньже поиду, яко к Тебе взях душу мою. Изми мя от враг моих Господи, к Тебе прибегох. Научи мя творити волю Твою, яко Ты еси Бог мой. Дух Твой Благии наставит мя на землю праву. Имене Твоего ради Господи, живиши мя, правдою Твоею изведеши от печали душу мою. И милостию Твоею потребиши враги моя и погубиши вся стужающия души моей, яко аз раб Твой есмь.

Слава в вышних Богу, и на земли мир, в человецех благоволение. Хвалим Тя, благословим Тя \myemph{(поклон)}, кланяемтися, славословим Тя \myemph{(поклон)}, благодарим Тя великия ради славы Твоея \myemph{(поклон)}. Господи Царю Небесныи, Боже Отче Вседержителю, и Господи Сыне Единородныи Исусе Христе, и Святыи Душе. Господи Боже Агньче Божии, Сыне Отечь. Вземляи грехи мiра, помилуй нас. Вземляи грехи мiра, приими молитвы наша. Седяи одесную Отца, помилуй нас. Яко Ты един свят, Ты един Господь Исус Христос, в славу Богу Отцу, аминь.

На всяку нощь благословим Тя, и восхвалим имя Твое во веки, и в век века. Господи, прибежище бысть нам, в род и род. Аз рех: Господи, помилуй мя, и исцели душу мою, яко согреших Тебе. Господи, к Тебе прибегох, научи мя творити волю Твою, яко Ты еси Бог мой. Яко от Тебе есть источник живота. Во свете Твоем узрим свет, пробави милость Твою ведущим Тя. Сподоби Господи, в нощь сию, без греха сохранитися нам. Благословен еси Господи Боже отец наших, и хвально и прославлено имя  Твое во веки, аминь. Буди  Господи, милость Твоя на нас, яко же уповахом на Тя. Благословен еси Господи, научи нас оправданием Твоим. Благословен еси Владыко, вразуми нас оправданием Твоим. Благословен еси Святыи, просвети нас оправданием Твоим. Господи, милость Твоя во веки, и дела руку Твоею не презри. Тебе подобает хвала. Тебе подобает пение. Тебе слава подобает, Отцу и Сыну и Святому Духу, ныне и присно и во веки веком, аминь.

Аще ли первая неделя святаго поста глаголем:

Трисвятое. 

По Отче наш. Молитва Исусова.

Аминь.

И поем тропари: Господи сил с нами.

В прочая же недели и посты, чтем канон Пресвятей Богородице. Таже, поем: Достойно есть. И поклон до земли. Трисвятое, и по Отче наш.


\myemph{ Аще есть пост Великии, начинает первый лик тропарь, глас 6:}

Господи сил с нами буди, иного бо разве Тебе помощника в скорбех не имамы, Господи сил помилуй нас.

И другии лик поет той же стих: Господи сил с нами.


\myemph{ Посем первый лик поет:} Хвалите Бога во святых Его, хвалите Его во утвержении сил Его.


\myemph{ И абие вторыи лик поет тропарь:} Господи сил с нами.


\myemph{ Тако и прочая стихи правыи лик поет. А тропарь поет левыи лик.}


\myemph{ Стих 2.}


Хвалите Его на силах Его, хвалите Его по премногому величествию Его. Господи сил.


\myemph{ Стих 3.}


Хвалите Его во гласе трубнем, хвалите Его во псалтыри и гуслех. Господи сил.


\myemph{ Стих 4.}


Хвалите Его в тимпане и лице, хвалите Его во струнах и органех. Господи сил.


\myemph{ Стих 5.}


Хвалите Его в кимвалех доброгласных, хвалите Его в кимвалех восклицания. Всяко дыхание да хвалит Господа. Господи сил.


\myemph{ Таже, оба лика совокупльшеся вкупе, сий стих поют:} Хвалите Бога во святых Его, хвалите Его во утвержении сил Его. Господи сил.


\myemph{ Слава, глас 6.}


Господи, аще не быхом святыя Твоя имели молитвеники, и благостыню Твою милующую нас, како смели быхом Спасе, пети Тя, Его же славословят безпрестани ангели. Сердцеведче, пощади душа наша.

И ныне. Многая премножества моих Богородице, согрешении, к Тебе прибегох Пречистая, спасения требуя. Посети болезнующую мою душу, и моли Сына Своего и Бога нашего, подати ми оставление, их же содеях зле, едина Благословеная.

Всесвятая Богородице, во время живота моего не остави мене, и человечестей помощи не ввери мене, но Сама мя заступи и помилуй мя.

Все упование мое к Тебе возлагаю, Мати Божия, сохрани мя во Своем си крове.

Господи помилуй (40). Посем молитва Великаго Василия:

Иже на всяко время, и на всяк час, на небеси и на земли, покланяемыи и славимыи, Боже Благии, долготерпеливе и многомилостиве. Иже праведники любя и грешных милуя, Иже всех зовыи ко спасению, обещания ради будущих благ; сам Господи, приими молитвы наша в час сей, и в благости исправи живот наш к заповедем Твоим. Душа наша освяти и телеса очисти, помышления исправи и мысли очисти. Разум уцеломудри и истрезви, и избави нас от всякия скорби, зол и болезней, и душевных страстей. И огради нас святыми ангелы Твоими, яко да ополчением их соблюдаеми и наставляеми, достигнем во единство веры, и в разум неприкосиовенныя Ти славы, яко благословен еси во веки, аминь.

Господи помилуй \myemph{(трижды)}.

Слава, и ныне.

Честнейшую херувим, и славнейшую воистину серафим, без истления Бога Слова рождьшую, сущую Богородицу, Тя величаем (поклон великий).

Именем Господним благослови, отче.

За молитв святых отец наших. Аминь.

И творим великия поклоны с молитвою святаго Ефрема, якоже явлено в Вечерни.


\myemph{ По сих, молитва Пресвятей Богородице, Великаго Василия:}


Нескверная, Неблазная, Нетленная, Пречистая Богоневесто Владычице, яже Бога Слова человеком преславным Си Рожеством соединивши, и падшее естество рода нашего Небесным совокупльши, яже ненадежным едина Надежда, и побеждаемым Помощница; готовое Заступление к Тебе прибегающим, и всем християном Прибежище; не гнушайся мене грешнаго и сквернаго, иже скверными помыслы и словесы, и делы и деяньми, всего себе непотребна сотворша, и сластем уныния, и лености нравом раба бывша. Но яко человеколюбиваго Бога Мати, человеколюбне умилосердися о мне грешнем и блуднем, и приими мое еже от скверных устен приносимое Тебе моление, и Своего Сына и нашего Владыку и Господа, Матерьним Си дерзновением обьемши умоли, яко да отверзет и мне человеколюбное милосердие Своея благости, и презрев моя безчисленая прегрешения, обратит мя на покаяние, и Своим заповедем делателя искусна явит мя; предстани ми присно милостивая и милосердая и человеколюбивая в настоящей сей жизни, теплая Предстательнице и Помощнице, сопротивных полки отгоняющи и ко спасению наставляющи мя, и во время исхода моего окаянную ми душу соблюдающи, и темныя образы лукавых бесов далече от нея отревающи. В страшныи же день праведнаго суда вечныя мя избавляющи муки, и неизреченней славе Твоего Сына и Бога нашего наследника мя показающи. Ея же да улучу, Владычице моя Пресвятая Богородице, Твоего ради ходатайства и заступления, благодатию и человеколюбием, единороднаго Сына Твоего, Господа Бога и Спаса нашего Исуса Христа, Ему же подобает всяка слава, честь и покланяние, со безначальным Его Отцем, и с Пресвятым и Благим и Животворящим Духом, ныне и присно и во веки веком, аминь.


\myemph{ Аще ли есть Великии пост, глаголем:} Трисвятое. По Отче наш, возгласа несть. Господи помилуй (12).


\myemph{ Таже, молитва святаго Антиоха мниха:} Даждь нам Владыко, на сон грядущим, покой души и телу, и сохрани нас от мрачнаго сна греховнаго, и всякаго темнаго и нощнаго сладострастия. Утиши стремление страстем и угаси телесное разжжение и стрелы лукаваго, яже на ны лукавно движимыя. И востания плоти нашея утоли и всяко земное и вещное мудрование успи, и даруй нам Боже, бодр ум и целомудр помысл, сердце трезвящеся, сон легок, и всякаго сатанина мечтания неискусен. Востави же нас во время молитвы утвержены в заповедех Твоих, и память судеб Твоих в себе выну имуща,  и всенощное славословие нам даруй, пети и благословити и славити пречестное и великолепое имя Твое, Отца и Сына н Святаго Духа, и ныне и присно и во веки веком, аминь.

Преславная и присно Дево Богородице Марие, Мати Христа Бога нашего, приими молитвы наша, и принеси я Сыну Своему и Богу нашему, да спасет и просветит Тебе ради душа наша.

Упование нам Бог и прибежище наше Христос, и Покровитель нам есть Дух Святыи.

Слава, и ныне. 

Господи помилуй, Господи помилуй, Господи благослови.


\myemph{ И отпуст седмичный, иже в субботу.}


\myemph{ Таже, прощение: Ослаби, остави. }

\myemph{ И молитва:} Ненавидящих и обидящих нас прости, Господи Человеколюбче.

Господи помилуй \myemph{(трижды)}. 

\myemph{ И исходные поклоны.}


\end{mymulticols}

\section{Павечерница малая}\begin{mymulticols}
%http://www.molitvoslov.org/text1005.htm 
 


\myemph{ По обычнем начале глаголем:}


За молитв святых отец наших, Господи Исусе Христе, Сыне Божии, помилуй нас. Аминь.

Царю небесныи, Утешителю, Душе истинныи. Иже везде сыи, и вся исполняя, сокровище благих, и жизни Подателю, прииди и вселися в ны, и очисти ны от всякия скверны, и спаси Блаже, душа наша.

Святыи Боже. Святыи Крепкии, Святыи безсмертныи, помилуй нас (трижды, и поклоны три).

Слава Отцу и Сыну и Святому Духу, ныне и присно и во веки веком, аминь.

Пресвятая Троице, помилуй нас. Господи, очисти грехи наша: Владыко, прости беззакония наша: Святыи, посети и исцели немощи наша, имене твоего ради. 

Господи помилуй \myemph{(трижды)}. Слава, и ныне.

Отче наш, Иже еси на небесах. да святится имя Твое. да приидет царствие Твое. да будет воля Твоя, яко на небеси и на земли. хлеб наш насущныи даждь нам днесь. и остави нам долги наша, якоже и мы оставляем должником нашим. и не введи нас во искушение. но избави нас от лукаваго.

Господи Исусе Христе, Сыне Божии, помилуй нас.

Аминь.

Господи помилуй (12). Слава, и ныне.

Приидите, поклонимся Цареви нашему Богу \myemph{(поклон)}.

Приидите, поклонимся Христу, Цареви и Богу нашему \myemph{(поклон)}.

Приидите, поклонимся и припадем к Самому Господу Исусу Христу, Цареви и Богу нашему \myemph{(поклон)}.


\mysubtitle{Псалом 50}


Помилуй мя Боже, по велицей милости Твоей, И по множеству щедрот Твоих, очисти беззаконие мое. Наипаче омый мя от беззакония моего, и от греха моего очисти мя. Яко беззаконие мое аз знаю, и грех мой предо мною есть выну. Тебе единому согреших, и лукавое пред Тобою сотворих. Яко да оправдишися в словесех Своих, и победиши внегда судити. Се бо в беззакониих зачат есмь, и во гресех роди мя мати моя. Се бо истину возлюбил еси, безвестная и тайная премудрости Твоея явил ми еси. Окропиши мя иссопом и очищуся, омыеши мя и паче снега убелюся. Слуху моему даси радость и веселие, возрадуются кости смиренныя. Отврати лице Твое от грех моих, и вся беззакония моя очисти. Сердце чисто созижди во мне Боже, и дух прав обнови во утробе моей. Не отверзи мене от лица Твоего и Духа Твоего Святаго не отыми от мене. Воздаждь ми радость спасения Твоего, и Духом Владычным утверди мя. Научу беззаконныя путем Твоим, и нечестивии к Тебе обратятся. Избави мя от кровий Боже, Боже спасения моего; возрадуется язык мой правде Твоей. Господи, устне мои отверзеши, и уста моя возвестят хвалу Твою. Яко аще бы восхотел жертвы, дал бых убо: всесожжения не благоволиши. Жертва Богу дух сокрушен: сердце сокрушено и смирено Бог не уничижит. Ублажи, Господи, благоволением Твоим Сиона, и да созиждутся стены иеросалимския. Тогда благоволиши жертву правды, возношение и всесожегаемая. Тогда возложат на олтарь Твой тельца.


\mysubtitle{Псалом 69}


Боже, в помощь мою воньми. Господи, помощи ми потщися. Да постыдятся и посрамятся ищущии душу мою. Да возвратятся вспять и постыдятся мыслящии ми злая. Да возвратятся абие стыдящеся глаголющии ми: благо же, благо же. Да возрадуются и возвеселятся о Тебе, вси ищущии Тебе, Боже, и глаголют выну: да возвеличится Господь, любящии спасение Твое. Аз же нищь есмь и убог, Боже помози ми. Помощник мой, и Избавитель мой еси Ты  Господи, не закосни.


\mysubtitle{Псалом 142}


Господи, услышы молитву мою, внуши моление мое воистине Твоей, услыши мя в правде Твоей. И не вниди в суд с рабом Твоим, яко не оправдится пред Тобою всяк живыи. Яко погна враг душу мою, и смирил есть в земли живот мой. Посадил мя есть в темных, яко мертвыя веку и уны во мне дух мой, во Мне смутися сердце мое. Помянух дни древния, поучихся во всех делех Твоих. и в делех руку Твоею поучахся, воздех к Тебе руце мои, душа моя, яко земля безводная Тебе. Скоро услыши мя Господи, исчезе дух мой. Не отврати лица Твоего от мене, и уподоблюся низходящим в ров. Слышану сотвори мне заутра милость Твою, яко на Тя уповах. Скажи мне Господи, путь, воньже пойду, яко к Тебе взях душу мою. Изми мя от враг моих, Господи, к Тебе прибегох. Научи мя творити волю Твою, яко Ты еси Бог мой. Дух Твой Благии наставит мя на землю праву. Имене Твоего ради Господи, живиши мя, правдою Твоею изведеши от печали душу мою. И милостию Твоею потребиши враги моя, и погубиши вся стужающия души моей, яко аз раб Твой есмь.


Слава в вышних Богу, и на земли мир, в человецех благоволение. Хвалим Тя, благословим Тя \myemph{(поклон)}, кланяемтися, славословим Тя \myemph{(поклон)}, благодарим Тя великия ради славы Твоея \myemph{(поклон)}. Господи Царю небесныи, Боже Отче Вседержителю, и Господи Сыне Единородныи Исусе Христе, и Святыи Душе: Господи Боже, Агньче Божии, Сыне Отечь, вземляи грехи мiра, помилуй нас. Вземляи грехи мiра. приими молитвы наша. Седяи одесную Отца, помилуй нас. Яко Ты един свят, Ты един Господь Исус Христос, в славу Богу Отцу, аминь. На всяку нощь благословлю Тя, и восхвалю имя Твое во веки, и в век века. Господи, Прибежище бысть нам, в род и род. Аз рех: Господи, помилуи мя и исцели душу мою, яко согреших Тебе. Господи, к Тебе прибегох, научи мя творити волю Твою, яко Ты еси Бог мой, яко от Тебе есть источник живота, во свете Твоем узрим свет. Пробави милость Твою ведущим Тя. Сподоби Господи, в нощь сию без греха сохранитися нам. Благословен еси, Господи Боже отец наших, и хвально и прославлено имя Твое во веки,  аминь. Буди Господи, милость Твоя на нас, якоже уповахом на Тя. Благословен еси Господи, научи нас оправданием Твоим. Благословен еси Владыко, вразуми нас оправданием Твоим. Благословен еси Святыи, просвети нас оправданием Твоим. Господи, милость Твоя во веки, и дела руку Твоею не презри. Тебе подобает хвала, Тебе подобает пение. Тебе слава подобает: Отцу и Сыну и Святому Духу, ныне и присно и во веки веком, аминь.

Исповедание православныя веры, перваго Собора (перекрестись без поклона): 

Верую во единаго Бога Отца, Вседержителя, Творца небу и земли, видимым же всем и невидимым.

И во единаго Господа Исуса Христа, Сына  Божия, Единороднаго, Иже от Отца рожденнаго прежде всех век. Света от Света, Бога истинна от Бога истинна, рождена, а не сотворена, единосущна Отцу, Им же вся быша.

Нас ради человек, и нашего ради спасения сшедшаго с небес, и воплотившагося от Духа Свята и Марии Девы вочеловечьшася.

Распятаго за ны при понтийстем Пилате, страдавша и погребена.

И воскресшаго в третии день по писаниих.

И возшедшаго на небеса, и седяща одесную Отца.

И паки грядущаго со славою судити живым и мертвым, Его же царствию несть конца.

Втораго собора:

И в Духа Святаго, Господа истиннаго и Животворящаго, Иже от Отца исходящаго, Иже со Отцем и Сыном спокланяема  и сславима, глаголавшаго пророки.

И во едину святую соборную и апостольскую Церковь.

Исповедую едино Крещение во оставление грехов.

Чаю воскресения мертвым.

И жизни будущаго века. Аминь.


\myemph{ Посем, канон Пресвятей Богородице, во октаи гласу или ин, по уставу. В  келиях глаголем канон Пресвятей Богородице, творение Феофаново, глас8. }


\mysubtitle{Песнь 1, ирмос}


Воду прошед, яко по суху, из египетска зла избежав. Израильтянин вопияше: Избавителю Богу нашему поем.


\myemph{ Запев:} Пресвятая Госпоже Богородице, спаси нас \myemph{(поклон)}.

Многими содержим напастьми, к Тебе прибегаю, спасении искии, О! Мати Слову и Дево, от лютых и нужных мя спаси.


\myemph{ Запев.}


Страстей мя смущают прилози, и многия печали исполняют ми душу; умири Отроковице, тишиною Сына и Бога Твоего, Всенепорочная.


Слава: Спаса рождшую Тя и Бога, молю Девице избавити мя от лют. К Тебе бо ныне прибегаю, простирая душу и помышление.

И ныне. В скорби мя суща и в тузе, посещения божественаго промышления, еже от Тебе сподоби, едина Богомати, яко блага, Благому же Родительнице.


\mysubtitle{Песнь 3, ирмос}


Небесному кругу, Верхотворче Господи, и церкви Зиждителю, Ты мене утверди в любви Своей, желанием сыи край, и верным утвержение,  едине Человеколюбче.


\myemph{ Запев.}


Заступление и покров жизни моей положих Тя, Богородительнице Дево, Ты мене окорми  к пристанищу Своему, благим виновна, и верным утвержение, едина всепетая.


\myemph{ Запев.}


Молю Девице, душевное смущение, и печали бурю разорити ми. Ты бо, Богоневестная, Начальника тишине Христа родила еси, едина Пречистая.


Слава.


Благодетеля рождьши добрым виновнаго, благодеяния богатство всем источи. Вся бо можеши яко сильна в крепости, Христа рождьши, Богоблаженная.

И ныне. Лютыми недуги, и болезненными страстьми истязаему, Дево, Ты ми помози, исцелением бо неоскудно Тя знаю сокровище, некрадомо и неиздаемо.

Господи помилуй \myemph{(трижды)}. 

Слава, и ныне:


\mysubtitle{Седален, глас 8}


Яко Всенепорочная Невесто Творцу, яко Неискусомужная Мати Избавителя, приятелище яко сущи Утешителя всепетая, беззаконию мя суща скверно жилище и бесом игралище в разуме бывша, потщися сих злодействия мя избавити, светло жилище добродетельми соделавши, светодательная и нетленная, разжени


\mysubtitle{Песнь 4, ирмос}


Услышах Господи, смотрения Твоего таиньство, разумех дела Твоя, и прославих Твое Божество.

Запев. Страстей моих смущение, яже кормчию  рождьши Господа, бурю утиши моих прегрешении, Богоневестная.

Запев. Милосердия бездну призывающу Твою подаждь ми, яже благосердаго рождьши, и Спаса всем поющим Тя.

Слава. Надежу и утверждение, и спасения стену недвижиму, стяжавше Тя Всенепорочная, злолютых всех избавляемся.

И ныне. На одре болезни моея греховне слежаща мя, яко человеколюбива, помози ми Богородице, Мати присно Дево.


\mysubtitle{Песнь 5, ирмос}


Просвети нас повелением Си Господи, и мышцею высокою Твоего мира подаждь нам, Едине Человеколюбче.

Запев. Исполни Чистая, веселия живота нашего, Твою нетленную дающи радость, веселия рождьши сущаго Вину.

Запев. Избави нас от бед, Богородице Чистая, вечное рождьши избавление, и мира всяк ум преимущаго.

Слава. Разреши мглу прегрешении моих, Богоневесто, светом Твоея светлости, яже Свет рождьшая Божественыи и Превечныи.

И ныне. Исцели Чистая, души моея неможение посещением Си убеждьшися, и здравие мольбами Си подаждь ми.


\mysubtitle{Песнь 6 , ирмос}


Молитву пролию ко Господу, и Тому возвещу печаль мою, яко зол душа моя наполнися, и живот мой аду приближися. Но молюся яко Иона: от тли, Боже мой, возведи мя.

Запев. Смерти и тли яко спасл есть, Сам ся издав на смерть, тлением и смертию мое естество ято бывшее, Дево, моли Господа и Сына Своего, врагов злодействия мя избавити.

Запев. Заступницу Тя животу свем, и хранительницу тверду Девице, напастей решащу молвы, и налоги бесом отгонящу. И молюся всегда, страстей мя избавити, Всенепорочная.

Слава. Яко стену и прибежище стяжахом, и душам Тя совершеное спасение, и пространство в скорбех Отроковице, и светом Ти присно радуемся, О! Владычице, ныне нас от страстей и бед спаси.

И ныне. На одре ныне немощию низлежу, и несть исцеления плоти моей. Но яже Бога и Спаса мiру, и Избавителя недугом рождьшая, Тебе молюся благой: от истлення болезней востави мя.

Господи помилуй \myemph{(трижды)}. 

Слава, и ныне.


\mysubtitle{Кондак, глас 6}


Заступнице християном непостыдная, Ходатаице ко Творцу непреложная, не презри грешных моления гласы. Но предвари яко блага на помощь нашу, верно вопиющих Ти: Ускори на молитву, и потщися на умоление, заступающи присно, Богородице, чтущих Тя.

Икос. Простри длани Свои, имаже всех Владыку, яко Младенца прият, за множество благости, не остави нас всегда надеющихся на Тя. Бодреною Си молитвою, и неисчетною простынею ущедри нас, и подаждь душам нашим милосердие Свое во веки источающи. Тебе бо имамы грешнии заступницу, от находящих на ны бед и зол. Но яко имущи милосердия щедроты, ускори на молитву, и потщися на умоление, заступающи присно Богородице чтущих Тя.


\mysubtitle{Песнь 7, ирмос}


Иже от Июдеи дошедше отроцы в Вавилон древле, верою утвержении, пламень пещныи попраша глаголюще: отец наших Боже,   благословен еси.

Запев. Наше спасение якоже восхотел, Спасе, устроити, во чрево Девыя вселися, юже мiру Заступницу показа. Отец наших Боже, благословен еси.

3апев. Волителя милости, Егоже роди Мати, ныне умоли, избавити мя от прегрешении и душевныя скверны, Сыну Твоему зовуща: отец наших Боже, благословен еси.

Слава. Сокровище спасения, и источник нетления, Тебе рождьшую, столп утвержения, и дверь покаяния, вопиющим показа: отец наших Боже, благословен еси.

И ныне. Телесныя недуги, и душевныя грехи, Богородительнице, любовию приступающим к крову Твоему Девице, испелити сподоби, яже Спаса Христа нам рождьшая.


\mysubtitle{Песнь 8, ирмос}


Царя Небесного, Его же поют вои ангельстии, хвалите, и превозносите Его во веки.

Запев. Помощи яже от Тебе требующия, не презри Дево, поющия и превозносящия Христа во веки.

Запев. Неможение души моея исцеляеши и телесныя болезни Дево, яко да Тя поем, и превозносим Чистую во веки.

Слава. Исцелением богатство изливаеши, верно поющим Тя Девице, и превозносящим Христа во веки.

И ныне. Напастныя Ты прилоги отгоняеши, и страстныя находы Дево, яко да Тя поем и превозносим Чистую во веки.


\mysubtitle{Песнь 9, ирмос}


Воистину Богородицу Тя исповедающе, спасении Тобою, Девице Чистая, с небесными вои Тя величаем.

Запев \myemph{(поклон)} Тока слез моих не отвратися, яже от всякаго лица вся слезы отъемшаго, Девице, Христа  рождьшая.

Запев. Радости мое сердце исполни Дево, яже радости приемши исполнение, греховную печаль потребляющи.

Слава. Света Твоего зарями просвети Дево, мрак неведения отгонящи, благоверно Богородицу Тя исповедающих.

И ныне. В место озлобления, немощи низлежащаго, Дево исцели, из нездравия во здравие претворяющи.

По каноне же, Достойно есть. Трисвятое. И по Отче наш. Молитва Исусова. Тропарь дню и святому, его же есть храм. Аще ли храм Христов, или Богородицы, глаголи преже храму.


\myemph{ Таже:} Боже отец наших, творяи присно с нами по Твоему смотрению, не отстави милости Твоея от нас, но тех молитвами, во смирении устрои живот наш.


\myemph{ Ин тропарь, глас той же:} Иже во всем мiре мученик Твоих Господи, яко багряницею и виссом кровьми их, Церкви Твоя украсившися, теми вопиет Ти, Христе Боже: людем Твоим щедроты Твоя низпосли, и мир граду Твоему даруй, и душам нашим велию милость.

Слава. Со святыми покой, Христе, душа раб Своих, идеже несть болезни, ни печали, ни воздыхания, но жизнь вечная.

И ныне. Молитвами Господи, всех святых и Богородицы, Твой мир даждь нам, и помилуй нас, яко един Щедр.


\myemph{ Аще в пяток вечер, глаголются тропари сия:} Апостоли, пророцы и мученицы, святителие, преподобнии и праведнии, иже добре подвиг скончавше, и веру соблюдше, дерзновение имуще ко Спасу, молим, вы, о нас Того молите яко Блага, спасти душа наша.

Таже тропарь, его же церковь.

Слава. Со святыми покой.


\myemph{ И ныне, кондак, глас 8.} Яко начатки естества Содетелю твари, вселенная приносит Ти Господи, богоносныя мученики. Тех молитвами, во смирении глубоце церковь Свою, и град Свой, Богородицы ради соблюди, едине Многомилостиве.


\myemph{ Аще ли господьскии праздник, или нарочитаго святаго, глаголем кондак его.} Господи помилуй, 40.


\myemph{ Таже, молитва Великаго Василия:} Иже на всяко время, и на всяк час, на небеси и на земли, покланяемыи и славимыи, Боже Благии, долготерпеливе и многомилостиве, Иже праведники любя, и грешных милуя, Иже всех зовыи ко спасению, обещания ради будущих благ; сам Господи, приими молитвы наша в час сей, и в благости исправи живот наш,  к заповедем Твоим. Душа наша освяти, и телеса очисти,  помышления исправи, и мысли очисти, разум уцеломудри, и истрезви, и избави нас от всякия скорби, зол и болезней, и душевных страстей. И огради нас святыми ангелы Твоими, яко да ополчением их соблюдаеми и наставляеми, достигнем во единство веры, и в разум неприкосновенныя Ти славы, яко благословен еси во веки, аминь.

Господи помилуй \myemph{(трижды)}. Слава, и ныне. Честнейшую херувим.

Именем Господним благослови, отче. 

За молитв святых отец наших, Господи Исусе Христе, Сыне Божии, помилуй нас. Аминь.


\myemph{ Молитва Пресвятей Богородице, святаго Великаго Василия:} Нескверная, Неблазная, Нетленная, Пречистая, Богоневесто Владычице, яже Бога Слова человеком, преславным Си Рожеством соединивши. И падшее естество рода нашего, небесным совокупльши. Яже ненадежным едина надежда, и побеждаемым Помощница. Готовое заступление к Тебе прибегающим, и всем християном прибежище, не гнушайся мене грешнаго и сквернаго. Иже скверными помыслы и словесы, и делы и деяньми, всего себе непотребна сотворша, и сластем уныния, и лености нравом раба бывша. Но яко человеколюбиваго Бога Мати, человеколюбне умилосердися о мне грешнем и блуднем, и приими  мое еже от скверных устен, приносимое Тебе моление. И Своего Сына и нашего Владыку и Господа, матерним Си дерзновением объемши умоли, яко да отверзет и мне человеколюбное милосердие Своея благости. И презрев, моя безчисленая прегрешения, обратит мя на покаяние. И Своим заповедем делателя искусна явит мя. Предстани ми присно милостивая, и милосердая и человеколюбивая в настоящей сей жизни, теплая Предстательнице и Помощнице, сопротивных полки отгоняющи, и ко спасению наставляющи мя, и во время исхода моего, окаянную ми душу соблюдающи, и темныя образы лукавых бесов, далече от нея отревающи. В страшныи же день праведнаго суда, вечныя мя избавляющи муки. И неизреченней славе Твоего Сына и Бога нашего, наследника мя показающи. Ея же да улучу, Владычице моя Пресвятая Богородице, Твоего ради ходатайства и заступления, благодатию и человеколюбием, единороднаго Сына Твоего, Господа Бога и Спаса нашего Исуса Христа. Ему же подобает всяка слава, честь и покланяние, со Безначальным Его Отцем, и с Пресвятым и Благим, и Животворящим Духом, ныне и присно и во веки веком, аминь.

Молитва святаго Антиоха мниха: Даждь нам Владыко, на сон грядущим, покой души и телу, и сохрани нас от мрачнаго сна греховнаго, и всякаго темнаго и нощнаго сладострастия. Утиши стремление страстем, и угаси телесное разжжение, и стрелы лукаваго, яже на ны лукавно движимыя, и востания плоти нашея утоли, и всяко земное и вещное мудрование успи. И даруй нам Боже, бодр ум и целомудр помысл, сердце трезвящеся, сон легок, и всякаго сатанина мечтания неискусен. Востави же нас во время молитвы, утвержены в заповедех Твоих. И память судеб Твоих в себе выну имуща. И всенощное славословие нам даруй, пети и благословити и славити, пречестное и великолепое имя Твое, Отца и Сына и Святаго Духа, ныне и присно и во веки веком, аминь.

Преславная и присно Дево Богородице Марие, Мати Христа Бога нашего, приими молитвы наша, и принеси я Сыну Своему и Богу нашему, да спасет и просветит Тебе ради душа наша.

Упование нам Бог, и прибежище наше Христос, и покровитель нам есть Дух Святыи.

Слава, и ныне. Господи помилуй (дважды), Господи благослови.


\myemph{ И отпуст:} Господи Исусе Христе, Сыне Божии, молитв ради Пречистыя Ти Матере, и преподобных и богоносных отец наших, и всех святых, помилуй и спаси мя грешнаго, яко Благ и Человеколюбец. Аминь. 

Таже прошение: Ослаби, остави, отпусти Боже, согрешения моя, вольная и невольная, яже в слове и в деле, и яже в ведении и не в ведении, яже во уме и в помышлении, яже во дни и в нощи, вся ми прости, яко Благ и Человеколюбец, аминь.

Посем вместо ектении глаголем молитву сию: Ненавидящих и обидящих нас прости, Господи Человеколюбче. Благотворящим благо сотвори, братиям и всем сродником нашим, иже и уединившимся, даруй им вся, яже ко спасению прошения и живот вечныи \myemph{(поклон)}.

В болезнех сущия посети и исцели, в темницах сущих свободи, по водам плавающим Правитель буди и иже в путех шествующим исправи и поспеши \myemph{(поклон)}.

Помяни Господи, и плененныя братию нашу, единоверных православныя веры, и избави их всякаго злаго обстояния \myemph{(поклон)}.

Помилуй Господи, давших нам милостыню и заповедавших нам, недостойным, молитися о них, прости их и помилуй \myemph{(поклон)}.

Помилуй Господи, труждающихся и служащих нам, милующих и питающих нас, и даруй им вся, яже ко спасению, прошения и живот вечныи \myemph{(поклон)}.

Помяни Господи, прежде отшедшия отцы и братию нашу и всели их, идеже присещает свет лица Твоего \myemph{(поклон)}.

Помяни Господи, и нашу худость и убожество, и просвети наш ум светом разума святаго Евангелия Твоего, и настави нас на стезю заповедей Твоих, молитвами Пречистыя Твоея Матере и всех святых Твоих, аминь \myemph{(поклон)}.

Господи помилуй \myemph{(трижды)}. И обычные исходные поклоны.


\end{mymulticols}

\section{Вечерня}\begin{mymulticols}
%http://www.molitvoslov.org/text1004.htm 
 


\myemph{ По обычнем начале глаголем:}


За молитв святых отец наших, Господи Исусе Христе, Сыне Божии, помилуй нас.

Псаломщик глаголет кротким и равным гласом с тихостию, и со всяким вниманием и со страхом Божиим во услышание всем: Аминь.

Царю Небесныи, Утешителю, Душе Истинныи, Иже везде сыи, и вся исполняя, сокровище благих, и жизни Подателю, прииди и вселися в ны, и очисти ны от всякия скверны, и спаси Блаже, душа наша.

Святыи Боже, Святыи Крепкии, Святыи Безсмертныи, помилуй нас (трижды, и поклоны три). 

Слава Отцу и Сыну и Святому Духу, и ныне и присно и во веки веком, аминь.

Пресвятая Троице, помилуй нас. Господи, очисти грехи наша. Владыко, прости беззакония наша. Святыи, посети и исцели немощи наша, имене Твоего ради.

Господи,помилуй \myemph{(трижды)}.

Слава Отцу и Сыну и Святому Духу, ныне и присно и во веки веком, аминь.

Отче наш, Иже еси на небесех, да святится имя Твоеда приидет царствие Твое, да будет воля Твоя, яко на небеси и на земли, хлеб наш насущныи даждь Нам днесь, и остави нам долги наша, яко же и мы оставляем должником нашим, и не введи нас во искушение, но избави нас от лукаваго.

Господи Исусе Христе, Сыне Божии, помилуй нас. Аминь.

Господи помилуй, 12. Слава, и ныне.

Приидите, поклонимся Цареви нашему Богу \myemph{(поклон)}.

Приидите, поклонимся Христу, Цареви и Богу нашему \myemph{(поклон)}.

Приидите, поклонимся и припадем к самому Господу Исусу Христу, Цареви и Богу нашему \myemph{(поклон)}.


\mysubtitle{Таже, псалом 103}


Благослови душе моя Господа. Господи Боже мой возвеличился еси зело. Во исповедание и в велелепотуся облече. Одеяися светом яко ризою, пропинаяи небо яко кожу. Покрываяи водами превыспреняя Своя, полагаяи облаки на восхождение Свое. Ходяи на крылу ветреню. Творяи ангелы Своя духи, и слуги Своя огнь палящь. Основаяи землю на тверди Своей, не преклонится в век века. Бездна яко риза одеяние ея. На горах станут воды. От запрещения Твоего побегнут, от гласа грома Твоего устрашатся. Восходят горы, и низходят поля, в место еже основал еси им. Предел положи, его же не прейдут, ниже обратятся покрыти землю. Посылая источники в дебрех, посреде гор пройдут воды. Напаяют вся звери сельныя. Ждут онагри  в жажду свою. На ты птицы небесныя привитают, от  среды камения дадят глас. Напаяя горы от превыспрених Своих, от плода дел Твоих насытится земля. Прозябая пажити скотом, и траву на службу человеком. Извести хлеб от земли, и вино веселит сердце человеку, умастити лице елеом. И хлеб сердце человеку укрепит. Насытятся древа польская, кедри ливаньстии их же еси насадил. Ту птицы вогнездятся, еродиево жилище обладает ими. Горы высокия еленем, камень прибежище заяцем. Сотворил еси луну во времена, солнце позна запад свой. Положи тму и бысть нощь, в ней же пройдут вси зверие дубравнии. Скимни рыкающе восхитити, и испросити от Бога пищу себе. Восия солнце, и собрашася, и в ложах своих лягут. Изыдет человек на дело свое, и на делание свое до вечера. Яко возвеличишася дела Твоя Господи, вся премудростию сотворил еси. Исполнися земля твари Твоея. Се море велико и пространно, ту гади, им же несть числа, животна малая с великими. Ту корабли преплавают, змий сей, его же созда ругатися ему. Вся к Тебе чают, дати пищу им во благо время. Давшу Тебе им, соберут. Отверзшу Тебе руку, всяческая исполнятся благости. Отвращьшу же Тебе лице, возмятутся. Отимеши духи их, и исчезнут, и в персть свою возвратятся. Послеши дух Свой и созиждутся, и обновиши лице земли. Буди слава Господня во веки, возвеселится Господь о делех Своих. Призираяи на землю, и творя ю трястися. Прикасаяся горах и воздымятся. Воспою Господеви в животе моем, пою Богу моему дондеже есмь. Да насладится ему беседа моя, аз же возвеселася о Господе. Скончаются грешницы от земли, и беззаконницы, яко не быти им. Благослови душе моя Господа. 

По конце псалма рцы стих: Яко возвеличишася дела Твоя Господи, вся премудростию сотворил еси. 

Таже, Слава, и ныне.

Аллилуия, Аллилуия, слава Тебе, Боже (три и поклоны три).

Господи помилуй (12), Слава, и ныне. Кафисма рядовая.

Егда же бывает малая вечерня, и по псалме глаголем: Слава, и ныне. Аллилуия, Аллилуия, слава Тебе, Боже \myemph{(трижды)}. Господи помилуй \myemph{(трижды)}. Слава, и ныне.


\mysubtitle{Таже псалом 140}


Господи возвах к Тебе, услыши мя. Воньми глас молитвы моея, егда воззову к Тебе. Да ся исправит молитва моя, яко кадило пред Тобою, воздеяние руку моею: жертва вечерняя. Положи Господи, хранение устом моим, и дверь ограждения о устнах моих. Не уклони сердце мое в словеса лукавствия, непщевати вины о гресех. С человеки творящими беззаконие, и не счетаюся со избранными их. Покажет мя праведник милостию, и обличит мя. Елей же грешнаго, да не намастит главы моея. Яко еще и молитва моя во благоволениих их. Пожерты быша при камени судии их. Услышатся глаголи мои, яко возмогоша, яко толща земли проседеся на земли, разсыпашася кости их при аде. Яко к Тебе Господи, Господи очи мои, на Тя уповах, не отъими душу мою. Сохрани мя от сети юже составиша ми, и от соблазн делающих беззаконие. Впадутся во мрежу свою грешницы, един есмь аз дондеже преиду.


\mysubtitle{Псалом 141}


Гласом моим ко Господу возвах, гласом моим ко Господу помолюся. Пролию пред Ним молитву мою, и печаль мою пред ним возвещу. Внегда исчезает дух мой, и Ты позна стези мои. На пути сем, по нему же хождах, скрыша сеть мне. Смотрях одесную и возглядах, и небе знаяи мене. Погибе бегство от мене, и несть взыскаяи душу мою. Возвах к Тебе Господи, рек: Ты еси упование мое, часть моя еси на земли живых. Воньми молитву мою, яко смирихся зело. Избави мя от гонящих мя, яко укрепишася паче мене. Изведи из темницы душу мою, исповедатися имени Твоему. Мене ждут праведницы, дондеже воздаси мне.


\mysubtitle{Псалом 129}


Из глубины возвах к Тебе Господи, Господи, услыши глас мой. Будете уши Твои, внемлюще глас молитвы моея. Аще беззакония назриши Господи, Господи, кто постоит, яко от Тебе очищение есть. Имене ради Твоего потерпех Тя Господи, потерпе душа моя в слово Твое, упова душа моя на Господа. От стражи утрения до нощи, от стражи утрения, да уповает Израиль на Господа. Яко от Господа милость, и много от Него избавление: и Той избавит Израиля от всех беззаконий его.


\mysubtitle{Псалом 116}


Хвалите Господа вси языцы, похвалите Его вси людие. Яко утвердися милость Его на нас, и истина Господня пребывает во веки.

Таже, стихеры настоянного дня по уставу. Посем, сей стих, творение святаго священномученика Анфиногена.

Свете тихии, святыя славы, безсмертнаго Отца небеснаго, святаго блаженнаго, Исуса Христа, Сына Божия, пришедшаго на запад солнцу, видевше свет вечернии, поем Отца и Сына и Святаго Духа, Бога. Достоин еси во вся времена, пет быти гласы преподобными, Сыне Божии, живот даяи всему мiру: его же ради весь мiр славит Тя.

Таже, прокимны дневныя. В субботу вечер прокимен: Господь воцарися, в лепотуся облече. 

Стих первыи. Облечеся Господь в силу, и препоясася.

Стих вторыи. Ибо утверди вселенную, яже не подвижится.

Стих третии. Дому Твоему подобает, святыни Господи в долготу днии.


\myemph{ В неделю вечер:} Се ныне благословите Господа вси раби Господни.

Стих. Стоящии в храме Господни, во дворех дому Бога нашего.


\myemph{ В понедельник вечер:} Господь услышит мя, егда взову к Нему.

Стих. Внегда возвах, услыша мя Боже правды моея.


Во вторник вечер: Милость Твоя Господи, поженет мя, во вся дни живота моего.

Стих. Господь пасет мя, и ничто же мя лишит.


\myemph{ В среду вечер:} Боже, во имя Твое спаси мя, и в силе Твоей суди ми.

Стих. Боже, услыши молитву мою, внуши глаголы уст моих.


\myemph{ В четверток вечер:} Помощь моя от Господа, сотворшаго небо и землю.

Стих. Возведох очи мои в горы, отнюду же приидет помощь моя.


\myemph{ В пяток вечер:} Боже, заступник мой еси Ты, и милость Твоя предварит мя.

Стих. Изми мя от враг моих Боже, и от востающих на мя избави мя.


\mysubtitle{ Егда же несть тропаря святому, тогда поем Аллилуию, на глас 6. Стихи же глаголем сия:}


\myemph{ В понедельник вечер:} Господи, не яростию Твоею обличи мене, ни гневом Твоим покажеши мене.


\myemph{ Во вторник и в четверток вечер:} Возносите Господа Бога нашего, и кланяйтеся подножию ногу Его, яко свято есть.


\myemph{ В среду вечер:} Во всю землю изыде вещание их, и в концы вселенныя глаголы их.


\myemph{ Аще ли суббота за упокой, поем:} Аллилуия, Аллилуия, Аллилуия, на глас 8.

Стих. Блажени яже избра и прият я Господь.

Стих. Память их от рода в род.

Стих. Душа их во благих водворятся.

И еще стих. Аллилуия, Аллилуия, Аллилуия.


\myemph{ По прокимне, или по ектении, или по аллилуии псаломшик глаголет:} Сподоби Господи, в вечер сей, без греха сохранитися нам \myemph{(поклон)}. Благословен еси Господи Боже отец наших \myemph{(поклон)}, и хвально и прославлено имя Твое во веки, аминь \myemph{(поклон)}. Буди Господи, милость Твоя на нас, якоже уповахом на Тя. Благословен еси  Господи, научи нас оправданием Твоим; благословен еси Владыко, вразуми нас оправданием Твоим; благословен еси Святыи, просвети нас оправданием Твоим. Господи, милость Твои во веки, и дела руку Твоею не презри. Тебе подобает хвала, Тебе подобает пение. Тебе слава подобает, Отцу и Сыну и Святому Духу, ныне и присно и во веки веком, аминь.


\myemph{ Таже ектения. И на стиховне стихеры, по уставу.}

\myemph{ Таже, молитва святаго Симеона Богоприимца:} Ныне отпущаеши раба Твоего Владыко, по глаголу Твоему с миром, яко видеста очи мои спасение Твое, еже еси уготовал пред лицем всех людей, свет во откровение языком, и славу людий Твоих Израиля.

\myemph{ Таже. Трисвятое, и по Отче наш.}

Исусова молитва.

Аминь.

\myemph{ Тропарь по уставу. Посем ектения. И отпуст.}

\myemph{ Аще ли есть пост, или егда поем Аллилуиа, глаголем тропари сия, глас 4.}

Богородице Дево, радуися, обрадованная Марие, Господь с Тобою, благословена Ты в женах, и благословен плод чрева Твоего, яко родила еси Христа Спаса, Избавителя душам нашим (поклон земной).

Слава. Крестителю Христов, тебе молимся: всех нас помяни, да избавимся от беззаконии наших. Тебе бо дана бысть благодать, молитися за мы (поклон земной).

И ныне. Молите за ны святии апостоли, пророцы, и мученицы, и вси святии, да избавимся от бед и скорбей, вас бо теплыя предстатели ко Спасу вси стяжахом (поклон земной).

\myemph{ Таже, Богородичен.} Под Твою милость прибегаем, Богородице Дево, молитв наших не презри в скорбех. Но от бед избави нас, едина Чистая и Благословеная (без поклона).

Господи помилуй (40), кротким и тихим гласом. Господи благослови.

\myemph{ Исусова молитва.} Аминь. 

Небесныи Царю, державу нашу укрепи, веру утверди, языки укроти, мiр умири, и святыи храм сей добре сохрани, и прежде отшедшия, отцы и братию нашу, в кровех с праведными учини. И нас в православней вере и в покаянии Господи, приими и помилуй, яко Благ и Человеколюбец. 

Таже, Господи помилуй, трижды.

Слава, и ныне.

Честнейшую херувим, и славнейшую воистину серафим, без истления Бога Слова рождьшую, сущую Богородицу Тя величаем (поклон земной).

Именем Господним благослови, отче.

За молитв святых отец наших, Господи Исусе Христе, Сыне Божии, помилуй нас. Аминь. 

И глаголем молитву святаго Ефрема Сирина, творя поклоны земные.

Господи и Владыко животу моему, дух уныния, небрежения, сребролюбия и празднословия, отжени от мене (поклон великий).

Дух же целомудрия, смирения, терпения и любве, даруй ми рабу Твоему (поклон великий).

Ей, Господи Царю, даждь ми зрети моя согрешения, и еже не осуждати брата моего, яко благословен еси во веки, аминь (поклон великий).

И прочих поклонов (12), глаголюще в себе: Господи Исусе Христе, Сыне Божии, помилуй мя грешнаго (дважды с поклоны); Боже, милостив буди мне грешному \myemph{(поклон)}. Боже, очисти грехи моя и помилуй мя \myemph{(поклон)}. Создавыи мя Господи, помилуй \myemph{(поклон)}. Без числа согреших Господи, прости мя \myemph{(поклон)}. И паки, скончав поклоны, глаголем молитву всю выше писаную: Господи и Владыко животу моему (поклон великий един).

И по сем, Трисвятое, и по Отче наш. Господи помилуй (12). Слава, и ныне. Господи помилуй (дважды), Господи благослови. \myemph{ И отпуст седмичный.}

Господи помилуй \myemph{(трижды)}. исходные поклоны.


\end{mymulticols}

\section{Полунощница повседневная}\begin{mymulticols}
%http://www.molitvoslov.org/text1007.htm 
 


\myemph{ По обычнем начале глаголем со умилением и сокрушением сердца:} За молитв святых отец наших, Господи Исусе Христе Сыне Божии, помилуй нас. \myemph{(поклон)}. Аминь. 


Слава Тебе Боже наш, слава Тебе Всяческих ради. \myemph{(трижды)}. Боже очисти мя грешнаго, яко николиже благо сотворих пред Тобою \myemph{(поклон)}. Но избави мя от лукаваго, и да будет во мне воля Твоя \myemph{(поклон)}. Да не осужденно отверзу уста моя недостойная, и восхвалю имя Твое святое, Отца и Сына, и Святаго Духа, ныне и присно, и во веки веком, аминь \myemph{(поклон)}. 

Царю небесныи, Утешителю, Душе истинныи, иже везде сыи, и вся исполняя, сокровище благих, и жизни Подателю, прииди и вселися в ны, и очисти ны от всякия скверны, и спаси Блаже душа наша. 

Святыи Боже, Святыи Крепкии, Святыи Безсмертныи, помилуй нас (трижды, и поклоны три). Слава Отцу и Сыну и Святому Духу, и ныне и присно, и во веки веком, аминь. 

Пресвятая Троице помилуй нас. Господи очисти грехи наша. Владыко прости беззакония наша. Святыи посети, и исцели немощи наша имене Твоего ради. 

Господи помилуй \myemph{(трижды)}. 

Слава Отцу и Сыну и Святому Духу, и ныне и присно, и во веки веком, аминь. 

Отче наш, иже еси на небесех. Да святится имя Твое. Да приидет царствие Твое. Да будет воля твоя, яко на небеси и на земли. Хлеб наш насущныи, даждь нам днесь. И остави нам долги наша, якоже и мы оставляем должником нашим. И не введи нас во искушение. Но избави нас от лукаваго. 

Господи Исусе Христе, Сыне Божии, помилуй нас. Аминь. 

Господи помилуй (12). 

\myemph{ Посих, молитва:} От сна восстав, благодарю тя Всесвятая Троице, яко многия ради благости и долготерпения, не прогневася на мя грешнаго, и лениваго раба твоего, и не погубил еси мене со беззаконии моими, но человеколюбствова. И в нечаянии лежаща воздвиже мя, утреневати и славословити державу Твою непобедимую. И ныне Владыко Боже пресвятыи, просвети очи сердца моего. И отверзи ми устне поучатися словесем Твоим, и разумети заповеди Твоя, и творити волю Твою, и пети Тя во исповедании сердечнем. Воспевати же и славити пречестное и великолепое имя Твое, Отца и Сына и Святаго Духа, ныне и присно и во веки веком, аминь. 

Приидите поклонимся Цареви нашему Богу \myemph{(поклон)}.

Приидите поклонимся Христу, Цареви и Богу нашему \myemph{(поклон)}.

Приидите поклонимся и припадем к Самому Господу Исусу Христу, Цареви и Богу нашему \myemph{(поклон)}. \myemph{ Таже, Псалом 50}. Помилуй мя Боже. \myemph{ Таже, кафисма 17}.


Аллилуия, \myemph{псалом 118}.


Блажени непорочнии в путь, ходящии в законе Господни. Блажени испытающии свидения Его, всем сердцем взыщут Его. Неделающии бо беззакония в путех Его ходиша. Ты заповеда заповеди Твоя сохранити зело. Еда исправилися быша путие мои, сохранити оправдания Твоя. Тогда не постыжуся, егда призрю на вся заповеди Твоя. Исповемся Тебе в правости сердца, внегда научитимися судбам правды Твоея. Оправдания Твоя сохраню, не остави мене до зела. В чесом исправит юныи путь свой, внегда сохранити словеса Твоя. Всем сердцем моим взысках Тебе, не отрини мене от заповедей Твоих. В сердцы моем скрых словеса Твоя, да не согрешу Тебе. Благословен еси Господи, научи мя оправданием Твоим. Устнама моима возвестих вся судьбы уст Твоих. На пути свидении Твоих насладихся, яко о всяком богатьстве. В заповедех Твоих поглумлюся, и разумею пути Твоя. Во оправданиих Твоих поучуся, не забуду словес Твоих. Воздаждь рабу Твоему, живи мя, и сохраню словеса Твоя. Открый очи мои, и разумею чудеса от закона Твоего. Пришлец аз есмь на земли, не скрый от мене заповеди Твоя. Возлюби душа моя вожделети судьбы Твоя на всяко время. Запретил еси гордым, прокляти уклоняющиися от заповедей Твоих. Отъими от мене понос и уничижение, яко свидении Твоих взысках. Ибо седоша князи, и на мя клеветаху, раб же Твой глумляшеся во оправданиих Твоих. Ибо свидения Твоя поучение мое есть, и совети мои оправдания Твоя. Прильпе земли душа моя, живи мя по словеси Твоему. Пути моя исповедах, и услыша мя, научи мя оправданием Твоим. Пути оправдании Твоих вразуми мя, и поглумлюся в чудесех Твоих. Воздрема душа моя от уныния, утверди мя в словесех Твоих. Путь неправды отстави от мене, и законом Твоим помилуй мя. Путь истинныи изволих, и судьбы Твоя не забых. Прилепихся свидении Твоих, Господи не посрами мене. Путь заповедей Твоих текох, егда разширил еси сердце мое. Законоположи мне Господи, путь оправдании Твоих, и взыщу и выну. Вразуми мя, и испытаю закон Твой, и сохраню и всем сердцем моим. Настави мя на путь заповедей Твоих, яко той восхотех. Приклони сердце мое во свидения Твоя, а не в лихоимство. Отврати очи мои не видети суеты, в пути Твоем живи мя. Постави рабу Твоему слово Твое в страх Твой. Отъими поношение мое, еже непщевах, яко будьбы Твоя благи. Се вожделех заповеди Твоя, в правде Твоей живи мя. И да приидет на мя милость Твоя Господи, спасение Твое по словеси Твоему. И отвещаю поношающим ми слово, яко уповах на словеса Твоя. И не отъими от уст моих словесе истинна до зела, яко на судьбы Твоя уповах. И сохраню закон Твой всегда: в век и в век века. И хождах в широте, яко заповеди Твоя взысках. И глаголах о свидениих Твоих пред цари и не стыдяхся. И поучахся в заповедех Твоих, их же возлюбих зело. И воздвигох руце мои к заповедем Твоим их же возлюбих, и глумляхся во оправданиих Твоих. Помяни словес Твоих рабу Твоему, их же упование дал ми еси. То мя утеши во смирении моем, яко слово Твое живи мя. Гордии законопреступноваху зело, от закона же Твоего не уклонихся. Помянух судьбы Твоя от века Господи, и утешихся. Печаль прият мя от грешник, оставляющих закон Твой. Пета бяху мне оправдания Твоя, на месте пришельствия моего. Помянух в нощи имя Твое Господи, и сохраних закон Твой. Си бысть мне, яко оправдании Твоих взысках. Часть моя еси Господи, рех: сохранити закон Твой. Помолихся лицу Твоему всем сердцем моим, помилуй мя по словеси Твоему. Помыслих пути Твоя, и возвратих нозе мои, во свидения Твоя. Уготовихся и не смутихся сохранити заповеди Твоя. Южа грешник обязашася мне, и закона Твоего не забых. Полунощи востах, исповедатися Тебе на судьбы правды Твоея. Прчастник аз есмь всем боящимся Тебе, и хранящим заповеди Твоя. Милости Твоея Господи, исполнь земля, оправданием Твоим научи мя. Благость сотворил еси с рабом Твоим Господи, по словеси Твоему. Благости и наказанию и разуму научи мя, яко заповедем Твоим веровах. Прежде даже не смирихся, аз прегреших, сего ради слово Твое сохраних. Благ еси Ты Господи, и благостию Твоею научи мя оправданием Твоим. Умножися на мя неправда гордых, аз же всем сердцем моим испытаю заповеди Твоя. Усырися яко млеко сердце их, аз же закону Твоему поучихся. Благо мне, яко смирил мя еси, да научуся оправданием Твоим. Благ мне закон уст Твоих, паче тысящ злата и сребра.


Слава, и ныне. Аллилуия, Аллилуия, слава Тебе Боже \myemph{(трижды)}. Господи помилуй \myemph{(трижды)}. Слава, и ныне.


Руце Твои сотвористе мя, и создасте мя, вразуми мя, и испытаю заповеди Твоя. Боящиися Тебе узрят мя и возвеселятся, яко на словеса Твоя уповах. Разумех Господи, яко правда повеления Твоя, и воистину смирил мя еси. Буди же милость Твоя, да утешит мя по словеси Твоему раба Твоего. Да приидут мне щедроты Твоя и жив буду, яко закон Твой поучение мое есть. Да постыдятся гордии, яко без правды беззаконноваша на мя, аз же поучуся в заповедех Твоих. Да обратят мя боящиися Тебе, и ведящии свидения Твоя. Буди сердце мое непорочно во оправданиих Твоих, яко да ся не постыжу. Исчезает во спасение Твое душа моя, в слово Твое уповах. Исчезоша очи мои в слово Твое глаголюще, когда утешиши мя; Яко бых яко мех на слане, оправдании Твоих не забых. Колико есть дний раба Твоего, когда сотвориши от гонящих мя суд; Поведаша мне законопреступницы глумления, но не яко закон Твой, Господи. Вся заповеди Твоя истина, без правды погнаша мя помози ми. Мала нескончаша мене на земли, аз же не оставих заповеди Твоя. По милости Твоей живи мя, и сохраню свидения уст Твоих. Во веки Господи, слово Твое пребывает на небеси. В род и род истина Твоя. Основал еси землю и пребывает. Учинением Твоим пребывает день, яко всяческая работна Тебе. Яко аще не закон Твой поучение мое есть, тогда убо погибл бых во смирении моем. Во веки не забуду оправдании Твоих, яко в них живил мя еси.

Твой есмь аз спаси мя, яко оправдании Твоих взысках. Мене ждаша грешницы, погубити мя. Свидения Твоя разумех. Всякия кончины видех конец, широка заповедь Твоя зело. Коль возлюбих закон Твой Господи, весь день поучение мое есть. Паче враг моих умудрил мя еси заповедию Твоею, яко в век моя есть. Паче всех учащих мя разумех, яко свидения Твоя поучение мое есть. Паче старец разумех, яко заповеди Твоя взысках. От всякаго пути лукава, возбраних ногам моим, яко да сохраню словеса Твоя. От судеб Твоих не удалихся, яко Ты законоположил ми еси. Коль сладка гортани моему словеса Твоя; паче меда устом моим. От заповедей Твоих разумех, сего ради возненавидех всяк путь неправды. Светильник ногам моим закон Твой, и свет стезям моим. Кляхся и поставих, сохранити судьбы правды Твоея. Смирихся до зела Господи, живи мя по словеси Твоему. Вольная уст моих благоволи же Господи, и судьбам Твоим научи мя. Душа моя в руку Твоею выну, и закона Твоего не забых. Положиша грешницы сеть мне, и от заповедей Твоих не заблудих. Наследовах свидения Твоя в век, яко радование сердца моего суть. Обратих сердце мое сотворити оправдания Твоя, в век за воздаяние. Законопреступныя возненавидех, закон же Твой возлюбих. Помощник мой и заступник мой еси Ты, на словеса Твоя уповах. Уклонитеся от мене лукавнующии, и испытаю заповеди Бога моего. Заступи мя по словеси Твоему, и жив буду, и не посрами мене от чаяния моего. Помози ми и спасуся, и поучуся во оправданиих Твоих выну. Уничижил еси вся отступающия от оправдании Твоих, яко неправедно помышление их. Преступающия непщевах вся грешныя земли, сего ради возлюбих свидения Твоя. Пригвозди от страха Твоего плоти моя, от судеб бо Твоих убояхся. Сотворих суд и правду, не предаждь мене обидящим мя. Восприими раба Твоего во благо, да не оклеветают мене гордии. Очи мои исчезосте во спасение Твое, и в слово правды Твоея. Сотвори с рабом Твоим по милости Твоей, и судьбам Твоим научи мя. Раб Твой есмь аз, вразуми мя, и научуся свидением Твоим. Время сотворити Господеви, разориша закон Твой. Сего ради возлюбих заповеди Твоя, паче злата и топазия. Сего ради ко всем заповедем Твоим направляхся, всяк путь неправды возненавидех. Дивна свидения Твоя, сего ради испытает я душа моя. Явление словес Твоих просвещает и вразумляет младенца. Уста моя отверзох и превлекох дух, яко заповеди Твоя желах. 

Слава, и ныне. Аллилуия, Аллилуия, слава Тебе Боже \myemph{(трижды)}. Господи помилуй \myemph{(трижды)}. Слава, и ныне.


Призри на мя и помилуй мя, по суду любящих имя Твое. Стопы моя направи по словеси Твоему, и да не одолеет ми всяко беззаконие. Избави мя от клеветы человеческия, и сохраню заповеди Твоя. Лице Твое просвети на раба Твоего, и научи мя оправданием Твоим. Исходища водная изведосте очи мои, понеже не сохраних закон Твой. Праведен еси Господи, и прави суди Твои. Заповеда правду свидения Твоя, и истину зело. Истаяла мя есть жалость Твоя, яко забыша словеса Твоя врази мои. Разжжено слово Твое зело, и раб Твой возлюби е. Юноша аз есмь и уничижен, оправдании Твоих не забых. Правда Твоя правда в век, и закон Твой истина. Скорби и нужды обретоша мя, заповеди Твоя поучение мое. Правда свидения Твоя в век, вразуми мя и жив буду. Возвах всем сердцем моим, услыши мя Господи, оправдании Твоих взыщу. Возвах Ти, спаси мя и сохраню свидения Твоя. Предварих в безгодии и возвах, на словеса Твоя уповах. Предваристе очи мои ко утру, поучитися словесем Твоим. Глас мой услыши Господи, по милости Твоей и по судбе Твоей живи мя. Приближиша гонящии мя беззаконие, от закона же Твоего удалишася. Близ еси Ты Господи, и вси путие Твои истина. Исперва познах от свидении Твоих, яко в век основал я еси. Виждь смирение мое, и изми мя, яко закона Твоего не забых. Суди суд мой, и избави мя, по словеси Твоему живи мя. Далече от грешник спасение, яко оправдании Твоих не взыскаша. Щедроты Твоя многи Господи, по судбе Твоей живи мя. Мнози изгонящии мя и стужающии ми, от свидении Твоих не уклонихся. Видех неразумевающия, и истаях, яко словес Твоих не сохраниша. Виждь, яко заповеди Твоя возлюбих Господи, по милости Твоей живи мя. Начало словес Твоих истина, и в век вся судбы правды Твоея. Князи погнаша мя туне, и от словес Твоих устрашися сердце мое. Возрадуюся аз о словесех Твоих, яко обретаяи корысть многу. Неправду возненавидех, и омерзе ми, закон же Твой возлюбих. Седмицею днем хвалих Тя, о судьбах правды Твоея. Мир мног любящим закон Твой, и несть им соблазны. Чаях спасение Твое Господи, и заповеди Твоя возлюбих. Сохрани душа моя свидения Твоя, и возлюби я зело. Сохраних заповеди Твоя, и свидения Твоя, яко вси путие мои пред Тобою, Господи. Да приближится молитва моя пред Тя Господи, по словеси Твоему вразуми мя. Да внидет прошение мое пред Тя Господи, по словеси Твоему избави мя. Отрыгнут устне мои пение, егда научиши мя оправданием Твоим. Провещает язык мой словеса Твоя, яко вся заповеди Твоя правда. Буди рука Твоя спасти мя, яко заповеди Твоя изволих. Вожделех спасение Твое Господи, и закон Твой поучение мое есть. Жива будет душа моя и восхвалит Тя, и судьбы Твоя помогут мне. Заблудих, яко овча погибшее, взыщи раба Твоего, яко заповеди Твоя не забых. 


Слава, и ныне, без Аллилуии. 

\myemph{ Таже, исповедание православныя веры.} Верую во единаго Бога. Трисвятое. И по Отче наш. 


\mysubtitle{Тропари сия, глас 8:}


Се Жених грядет в полунощи, и блажен раб, его же обрящет бдяща. Недостоин же паки, его же обрящет ленящася. Блюди убо душе моя, да не сном отяготишися и да не смерти предана будеши, и царствия вне затворишися, но воспряни зовущи: свят, свят, свят еси, Боже, Богородицы ради, помилуй нас. 

Слава. День он страшныи помышляющи, душе моя побди, вжигающи свещу свою, и маслом просвещающи, не веси бо когда приидет к тебе глас глаголющии: се Жених. Блюди убо, душе моя, да не воздремлеши, и пребудеши вне толкущи, яко пять дев; но бодрено побди, яко да усрящеши Христа милостива, и даст ти чертог божественыя славы Своея. 

И ныне. Тебе необоримую стену имуще вернии, и спасения утвержение, Богородице Дево молим: сопротивных советы разори и людей Своих печаль на радость преложи. Мир Свой умири и православных утверди, и о смирении мира молися, яко Ты еси Богородице, упование наше. Господи помилуй (40). 


\myemph{ Таже молитва Великаго Василия:} Иже на всяко время, и на всяк час, на небеси и на земли, покланяемыи и славимыи, Боже Благии, долготерпеливе и многомилостиве; Иже праведники любя и грешных милуя; Иже всех зовыи ко спасению, обещания ради будущих благ; сам Господи, приими молитвы наша в час сий, и в благости исправи живот наш, к заповедем Твоим. Душа наша освяти, и телеса очисти, помышления исправи, и мысли очисти, разум уцеломудри и истрезви, и избави нас от всякия скорби, зол и болезней, и душевных страстей, и огради нас святыми ангелы Твоими, яко да ополчением их соблюдаеми и наставляеми, достигнем во единство веры, и в разум неприкосновенныя Ти славы, яко благословен еси во веки, аминь. 

\myemph{ Таже, Господи помилуй} \myemph{(трижды)}. Слава, и ныне. 

Честнейшую херувим, и славнейшую воистину серафим, без истления Бога Слова рождьшую, сущую Богородицу Тя величаем \myemph{(поклон)}. Именем Господним благослови, отче. За молитв святых отец наших, Господи Исусе Христе, Сыне Божии, помилуй нас. Аминь. 

\myemph{ Аще есть пост, или кроме поста, егда поем Аллилуия, и творим три поклоны великия, глаголюще молитву святаго Ефрема}. Господи и Владыко животу моему, дух уныния, небрежения, сребролюбия и празднословия, отжени от мене \myemph{(поклон)}.Дух же целомудрия, смирения, терпения и любве, даруй ми рабу Твоему \myemph{(поклон)}.Ей, Господи Царю, даждь ми зрети моя согрешения, и еже не осуждати брата моего, яко благословен еси во веки, аминь \myemph{(поклон)}.

\myemph{ И прочих поклонов 12, глаголюще в себе:} Господи Исусе Христе, Сыне Божии, помилуй мя грешнаго (дважды с поклонами).

Боже, милостив буди мне грешному \myemph{(поклон)}.

Боже, очисти грехи моя и помилуй мя \myemph{(поклон)}.

Создавыи мя Господи, помилуй \myemph{(поклон)}.

Без числа согреших Господи, прости мя \myemph{(поклон)}.

\myemph{ И паки скончав поклоны, глаголем молитву всю выше писаную:} Господи, и Владыко животу моему… И поклон един. 


И по изглаголании молитвы, или по Честнейшую херувим, глаголем сию молитву Великаго Василия: Владыко Боже, Отче Вседержителю, и Господи Сыне Единородныи Исусе Христе, и Святыи Душе, едино Божество, и едина сила, помилуй мя грешнаго, и ими же веси судьбами, спаси мя недостойнаго раба Твоего, яко благословен еси во веки, аминь. 


Таже, Приидите, поклонимся \myemph{(трижды)}. И поклоны три. 


\mysubtitle{Псалом 120}


Возведох очи мои в горы, отнюдуже приидет помощь моя. Помощь моя от Господа, сотворшаго небо и землю. Не даждь во смятение ноги твоея, ниже воздремлет храняи тя. Се не воздремлет, ни уснет храняи Израиля. Господь сохранит тя, Господь покров Твой на руку десную твою. Во дни солнце не ожжет тебе, ни луна нощию, Господь сохранит тя от всякаго зла, сохранит душу твою Господь. Господь сохранит вхождение твое, и исхождение твое от ныне и до века.


\mysubtitle{Псалом 133}


Се ныне благословите Господа вси раби Господни, стоящии в храме Господни, во дворех дому Бога нашего. В нощех воздежите руки ваша во святая, и благословите Господа. Благословит тя Господь от Сиона, сотворивыи небо и землю. 

Слава, и ныне. Трисвятое, и поклоны три. По Отче наш, тропари сия: Помяни Господи, яко благ, рабы Своя, и елика в житии сем согрешиша, прости. Никтоже бо без греха, токмо Ты могии преставленым дати покой. Иже глубинами мудрости, человеколюбне вся строя, и еже на пользу всем подавая, едине Содетелю, покой Господи, душа усопших раб Своих, на Тя бо упование возложиша, Творца и Зиждителя и Бога нашего. 

Слава. Со святыми покой, Христе, душа раб Своих, идеже несть болезни, ни печали, ни воздыхания, но жизнь вечная. 

И ныне, богородичен: Блажим, Тя вси роди, Богородице Дево, в Тя бо невместимыи Христос Бог наш, вместитися изволи, блажени есмы и мы помощницу Тя имуще, день бо и нощь молишися о нас, и державы царствия Твоими молитвами утвержаются. Тем благодаряще вопием Ти: радуися, Обрадованная, Господь с Тобою. 

Господи помилуй (12).


\myemph{ Таже, молитву сию:} Помяни Господи, иже в надежи воскресения, и жизни вечныя, усопшия отцы и братию нашу, и всех иже во благочестивей вере скончавшихся, и прости им всяко согрешение, вольное же и невольное, словом и делом, и помышлением согрешенное ими и всели их в места светла, в места прохладна, в места покойна, отнюдуже отбеже всяка болезнь, и печаль и воздыхание, идеже присещает свет лица Твоего, и веселит вся иже от века святыя Твоя. И даруй им царствие Твое, и причастие неизреченных и вечных Твоих благ, и Твоея безконечныя, и блаженныя жизни наслаждение. Ты бо еси воскресение и живот и покой усопшим рабом Твоим, Христе Боже наш, и Тебе славу возсылаем, со безначальным Ти Отцем, и с Пресвятым и Благим и Животворящим Ти Духом, ныне и присно и во веки веком, аминь. 

\myemph{ Таже, Слава, и ныне. Господи помилуй (дважды). Господи благослови. }

\myemph{ И отпуст:} Господи Исусе Христе, Сыне Божии, молитв ради Пречистыя Твоея Матере, преподобных и богоносных отец наших и всех ради святых, помилуй и спаси нас, яко Благ и Человеколюбец, аминь. 

\myemph{ И прощение:} Ослаби, остави, отпусти Боже, согрешения наша, вольная и невольная, яже в слове и в деле, и яже в ведении и не в ведении, яже во уме и в помышлении, яже во дни в нощи, вся ми прости, яко Благ и Человеколюбец, аминь.


\myemph{ И молитва с поклонами:} Ненавидящих и обидящих нас прости, Господи Человеколюбче. Благотоворящим благо сотвори, братиям и всем сродником нашим, иже и уединившимся, даруй им вся, яже ко спасению прошения и живот вечныи \myemph{(поклон)}. 

В болезнех сущия посети и исцели, в темницах сущия свободи, по водам плавающим Правитель буди и иже в путех шествующим исправи и поспеши \myemph{(поклон)}. 

Помяни Господи, и плененныя братию нашу, единоверных православныя веры, и избави их всякаго злаго обстояния \myemph{(поклон)}. 

Помилуй Господи, давших нам милостыню и заповедавших нам, недостойным, молитися о них, прости их и помилуй \myemph{(поклон)}. 

Помяни Господи, прежде отшедшия отцы и братию нашу и всели их, идеже присещает свет лица Твоего \myemph{(поклон)}. 

Помяни Господи, и нашу худость и убожество, и просвети наш ум светом разума святаго Евангелия Твоего, и настави нас на стезю заповедей Твоих, молитвами Пречистыя Твоея Матере и всех святых Твоих, аминь \myemph{(поклон)}. 

Господи помилуй \myemph{(трижды)},

и обычные исходные поклоны.

\end{mymulticols}

\mychapterending


%\newcommand{\pominalnayaslava}{

\myparsep

{\small Слава Отцу и Сыну и Святому Духу, и ныне и присно и во веки веков. Аминь. 

Аллилуиа, аллилуиа, аллилуиа, слава Тебе, Боже. \myemph{(Трижды)}

Господи, помилуй. \myemph{(Трижды)}

Помяни, Господи, Боже наш, в вере и надежди живота вечнаго преставльшагося раба Твоего, брата нашего \myemph{(имярек)}, и яко Благ и Человеколюбец отпущаяй грехи и потребляяй неправды, ослаби, остави и прости вся вольная его согрешения и невольная, избави его вечная муки и огня геенскаго и даруй ему причастие и наслаждение вечных Твоих благих, уготованных любящим Тя: аще бо и согреши, но не отступи от Тебе, и несумненно во Отца и Сына и Святаго Духа, Бога Тя в Троице славимаго верова, и Ед\'{и}ницу в Троице и Троицу во Единстве, православно даже до последняго своего издыхания исповеда. Темже милостив тому буди, и веру, яже в Тя вместо дел вмени, и со святыми Твоими, яко Щедр, упокой: несть бо человека, иже поживет и не согрешит. Но Ты Един еси кроме всякаго греха, и правда Твоя, правда во веки, и Ты еси Един Бог милостей и щедрот, и человеколюбия, и Тебе славу возсылаем, Отцу и Сыну и Святому Духу, ныне и присно и во веки веков. Аминь. 

Слава Отцу и Сыну и Святому Духу, и ныне и присно и во веки веков. Аминь.}

\myparsep

}

\mypart{НАПУТСТВИЕ ХРИСТИАНИНА ПЕРЕД СМЕРТЬЮ И ЗАУПОКОЙНЫЕ МОЛИТВЫ}\label{_content_molitvi-pered-smertyu}
%http://www.molitvoslov.org/content/molitvi-pered-smertyu

 

\mychapter{Кончина человека}\begin{mymulticols}
%http://www.molitvoslov.org/text182.htm 
 

Если положение больного безнадежно, то при явных признаках приближающейся смерти священник читает отходную молитву "--- «Канон молебный при разлучении души от тела» или более полно он называется «Канон молебный ко Господу нашему Иисусу Христу и Пречистой Богородице Матери Господни при разлучении души от тела всякаго правовернаго». Родственники сами могут прочитать этот канон, если невозможно пригласить священника, кроме чтения «молитвы, от иерея глаголемой на исход души», которая находится в конце канона. Этот канон читается «от лица человека с душею разлучающагося и не могущаго глаголати» и имеется в православных молитвословах. Чтение канона мирскими людьми начинается возгласом: «Молитвами святых отец наших Господи Иисусе Христе Боже наш, помилуй нас», затем следуют предначинательные молитвы: «Трисвятое», «Пресвятая Троице», «Отче наш» и далее по молитвослову. 

При чтении канона возжигается свеча и лампадка перед домашней святой иконой. Если дома иконы нет, то нужно обязательно приобрести в храме иконы Спасителя и Божией Матери. Для умирающих младенцев (детей до семи лет) из-за отсутствия грехов, перечисляемых в каноне, которые несвойственны им по малолетству, канон не читается. 

Кроме канона при разлучении души от тела еще существует «Чин, бываемый на разлучение души от тела, когда человек долго страждет». Этот чин читается над человеком, который испытывает тяжкие предсмертные мучения и никак не может умереть (как правило, читается священником). 

После смерти человека над ним немедленно читается Последование по исходе души от тела. 

\end{mymulticols}

\mychapterending

\mychapter{Приготовление усопшего к погребению}\begin{mymulticols}
%http://www.molitvoslov.org/text237.htm 
 

Умерших мы именуем усопшими, то есть уснувшими. Называем их мы так по нашей христианской вере, что души после смерти не уничтожаются, не исчезают в небытии, а отлучаются от тела и переходят из этой жизни в другую "--- загробную. Там пребывают они после частного суда по делам земным в приличествующем им месте до Страшного суда Божия, когда по слову Господа души всех умерших людей воссоединяться с телами и воскреснут. И тогда окончательно определится участь каждого: праведные наследуют Царство Небесное, блаженную вечность с Богом, а грешные "--- вечное наказание. 

Историческое обоснование погребения усопших дано в образе погребения Иисуса Христа. По примеру благочестивой древности и в настоящее время погребение предваряется совершением различных многозначительных символических действий. 

Тело усопшего омывают теплой водой, чтобы предстал он перед Богом по воскресении в чистоте и непорочности. При омовении читают «Трисвятое»: «Святый Боже, Святый Крепкий, Святый Бессмертный, помилуй нас» или «Господи, помилуй». Возжигается лампадка или свеча и горит до тех пор, пока в доме находится покойник. После омовения тело христианина одевают в чистые и, по возможности, новые одежды "--- по званию и служению его, на усопшем обязательно должен быть нательный крестик. Омовение обычно совершают люди старшего возраста, а если таковых нет, то омыть тело усопшего может любой из родственников, за исключением женщин, находящихся в данный момент в естественной нечистоте. Обычай указывает, что в омовении тела женщины участвуют только женщины. Если известно, что почивший был монахом (монахиней) или священнослужителем, то надо о его кончине сообщить в храм. Тело покойного полагают на стол, покрывают белым покрывалом "--- саваном. Затем умершего покрывают особым освященным покровом (погребальным покрывалом), на котором изображены крест, лики святых и молитвенные надписи. Все это означает, что умерший оставался верным Богу и теперь остается под покровом Божиим. Глаза должны быть закрыты, уста сомкнуты, руки сложены крестообразно, правая поверх левой. Руки и ноги усопшего связывают, чтобы развязать перед последним прощанием. В руки покойного полагают погребальный крест, на грудь кладут святую икону, для мужчин "--- образ Спасителя, для женщин "--- образ Божией Матери. На лоб умершего возлагают венчик "--- полоску бумаги с изображением Спасителя, Божией Матери и Иоанна Предтечи. Изображения эти обрамлены надписью «Трисвятого». Венчик, символизирующий соблюдение веры усопшим христианином и совершение им христианского жизненного подвига, возлагается в надежде, что почивший в вере получит по воскресении от Бога небесную награду и нетленный венец. Как правило, венчик отпечатан на одном листке с разрешительной молитвой. После приобретения молитвы-венчика в храме венчик отрезается ножницами (после отпевания листок с молитвой будет вложен в руку почившего). Перед положением усопшего во гроб тело его и гроб окропляют святой водой, причем гроб окропляют извне и изнутри. В гроб покойного полагают лицом вверх, под голову кладут подушку, набитую соломой или опилками. Гроб обыкновенно ставят посреди комнаты перед домашними иконами, головой к образам. Вокруг гроба возжигают четыре свечи: у головы, в ногах и по обеим сторонам на уровне скрещенных рук. Возженные свечи вместе изображают крест и символизируют переход умершего в Царство Истинного Света. 

К православным традициям не относятся бытующие во многих семьях различные суеверия, связанные с покойником, такие, как занавешивание зеркал, отказ пользоваться вилками во время поминальной трапезы, обычай оставлять часть блюд или стакан с водой (а еще хуже с водкой) перед портретом усопшего и т. п. 

Все эти суеверия никакого отношения к православию не имеют. Занавешивание зеркал в доме, где лежит тело покойного, оправдано лишь в том случае, когда мы, думая об усопшем, уходим от внешней суеты и отдаём свой последний молитвенный вздох об упокоении преставльшейся души.

\end{mymulticols}

\mychapterending

\mychapter{17-я кафизма (поминальная), чтомая в дни особого поминовения усопших}
%http://www.molitvoslov.org/text227.htm 
 
{\centering \myemph{(Читается ежедневно в течение 40 дней по смерти)}\par}

\begin{mymulticols}

\MolitvamiSviatyhOtecNashih 

Царю Небесный, Утешителю, Душе истины, Иже везде сый и вся исполняяй, Сокровище благих и жизни Подателю, прииди и вселися в ны, и очисти ны от всякия скверны, и спаси, Блаже, души наша.

\TrisviatoePoOtcheNash

\mysubtitle{Тропарь:}

Помилуй нас, Господи, помилуй нас; всякаго бо ответа недоумеюще, сию Ти молитву яко Владыце грешнии приносим: помилуй нас. 

\slavan

Честное пророка Твоего, Господи, торжество, небо Церковь показа, с человеки ликуют ангели: того молитвами, Христе Боже, в мире управи живот наш, да поем Ти: Аллилуиа. 

\inynen

Многая множества моих, Богородице, прегрешений, к Тебе прибегох, Чистая, спасения требуя: посети немощствующую мою душу, и моли Сына Твоего и Бога нашего дать ми оставление, яже содеях лютых, Едина Благословенная. 

Господи, помилуй. \myemph{(Сорок раз)}

\myemph{И поклоны по силе.}

\priiditepoklonimsia

\myemph{И псалмы:}

\mysubtitle{Псалом 118}

Блажени непорочнии в путь, ходящий в законе Господни. Блажени испытающии свидения Его, всем сердцем взыщут Его, не делающий бо беззакония, в путех Его ходиша. Ты заповедал еси заповеди Твоя сохранити зело. Дабы исправилися путие мои, сохранити оправдания Твоя. Тогда не постыжуся, внегда призрети ми на вся заповеди Твоя. Исповемся Тебе в правости сердца, внегда научитимися судьбам правды Твоея. Оправдания Твоя сохраню: не остави мене до зела. В чесом исправит юнейший путь свой? Внегда сохранити словеса Твоя. Всем сердцем моим взысках Тебе, не отрини мене от заповедей Твоих. В сердце моем скрых словеса Твоя, яко да не согрешу Тебе. Благословен еси, Господи, научи мя оправданием Твоим. Устнама моима возвестих вся судьбы уст Твоих. На пути свидений Твоих насладихся, яко о всяком богатстве. В заповедех Твоих поглумлюся, и уразумею пути Твоя. Во оправданиих Твоих поучуся, не забуду словес Твоих. Воздаждь рабу Твоему: живи мя, и сохраню словеса Твоя. Открый очи мои, и уразумею чудеса от закона Твоего. Пришлец аз есмь на земли: не скрый от мене заповеди Твоя. Возлюби душа моя возжелати судьбы Твоя на всякое время. Запретил еси гордым: прокляти уклоняющиися от заповедей Твоих. Отыми от мене понос и уничижение, яко свидений Твоих взысках. Ибо седоша князи, и на мя клеветаху, раб же Твой глумляшеся во оправданиих Твоих. Ибо свидения Твоя поучение мое есть, и совети мои оправдания Твоя. Прильпе земли душа моя, живи мя по словеси Твоему. Пути моя возвестих, и услышал мя еси: научи мя оправданием Твоим. Путь оправданий Твоих вразуми ми, и поглумлюся в чудесех Твоих. Воздрема душа моя от уныния, утверди мя в словесех Твоих. Путь неправды отстави от мене, и законом Твоим помилуй мя. Путь истины изволих, и судьбы Твоя не забых. Прилепихся свидением Твоим, Господи, не посрами мене. Путь заповедей Твоих текох, егда разширил еси сердце мое. Законоположи мне, Господи, путь оправданий Твоих, и взыщу и выну. Вразуми мя, и испытаю закон Твой, и сохраню и всем сердцем моим. Настави мя на стезю заповедей Твоих, яко тую восхотех. Приклони сердце мое во свидения Твоя, а не в лихоимство. Отврати очи мои, еже не видети суеты, в пути Твоем живи мя. Постави рабу Твоему слово Твое в страх Твой. Отыми поношение мое, еже непщевах, яко судьбы Твоя благи. Се, возжелах заповеди Твоя, в правде Твоей живи мя. И да приидет на мя милость Твоя, Господи, спасение Твое по словеси Твоему, и отвещаю поношающым ми слово, яко уповах на словеса Твоя. И не отыми от уст моих словесе истинна до зела, яко на судьбы Твоя уповах. И сохраню закон Твой выну, в век и в век века. И хождах в широте, яко заповеди Твоя взысках, и глаголах о свидениих Твоих пред цари, и не стыдяхся. И поучахся в заповедех Твоих, яже возлюбих зело: и воздвигох руце мои к заповедем Твоим, яже возлюбих, и глумляхся во оправданиих Твоих. Помяни словеса Твоя рабу Твоему, ихже упование дал ми еси. То мя утеши во смирении моем, яко слово Твое живи мя. Гордии законопреступоваху до зела, от закона же Твоего не уклонихся. Помянух судьбы Твоя от века, Господи, и утешихся. Печаль прият мя от грешник, оставляющих закон Твой. Пета бяху мне оправдания Твоя на месте пришельствия моего. Помянух в нощи имя Твое, Господи, и сохраних закон Твой. Сей бысть мне, яко оправданий Твоих взысках. Часть моя еси, Господи: рех сохранити закон Твой. Помолихся лицу Твоему всем сердцем моим: помилуй мя по словеси Твоему. Помыслих пути Твоя, и возвратих нозе мои во свидения Твоя. Уготовихся и не смутихся сохранити заповеди Твоя. Ужя грешник обязашася мне, и закона Твоего не забых. Полунощи востах исповедатися Тебе о судьбах правды Твоея. Причастник аз есмь всем боящымся Тебе и хранящым заповеди Твоя. Милости Твоея, Господи, исполнь земля, оправданием Твоим научи мя. Благость сотворил еси с рабом Твоим, Господи, по словеси Твоему. Благости и наказанию и разуму научи мя, яко заповедем Твоим веровах. Прежде даже не смиритимися, аз прегреших, сего ради слово Твое сохраних. Благ еси Ты, Господи, и благостию Твоею научи мя оправданием Твоим. Умножися на мя неправда гордых, аз же всем сердцем моим испытаю заповеди Твоя. Усырися яко млеко сердце их, аз же закону Твоему поучихся. Благо мне, яко смирил мя еси, яко да научуся оправданием Твоим. Благ мне закон уст Твоих, паче тысящ злата и сребра. 

\pominalnayaslava

Руце Твои сотвористе мя и создасте мя, вразуми мя, и научуся заповедем Твоим. Боящиися Тебе узрят мя и возвеселятся, яко на словеса Твоя уповах. Разумех, Господи, яко правда судьбы Твоя, и войстинну смирил мя еси. Буди же милость Твоя, да утешит мя по словеси Твоему рабу Твоему. Да приидут мне щедроты Твоя, и жив буду, яко закон Твой Поучение мое есть. Да постыдятся гордии, яко неправедно беззаконноваша на мя, аз же поглумлюся в заповедех Твоих. 

Да обратят мя боящиися Тебе и водящии свидения Твоя. Буди сердце мое непорочно во оправданиих Твоих, яко да не постыжуся. Исчезает во спасение Твое душа моя, на словеса Твоя уповах. Исчезоша очи мои в слово Твое, глаголюще: когда утешиши мя? Зане бых яко мех на слане, оправданий Твоих не забых. Колико есть дний раба Твоего? Когда сотвориши ми от гонящих мя суд? Поведаша мне законопреступницы глумления, но не яко закон Твой, Господи. Вся заповеди Твоя истина; неправедно погнаша мя, помози ми. Вмале не скончаша мене на земли, аз же не оставих заповедей Твоих. По милости Твоей живи мя, и сохраню сведения уст Твоих. Во век, Господи, слово Твое пребывает на Небеси. В род и род истина Твоя. Основал еси землю, и пребывает. Учинением Твоим пребывает день: яко всяческая работна Тебе. Яко аще бы не закон Твой поучение мое был, тогда убо погибл бых во смирении моем. Во век не забуду оправданий Твоих, яко в них оживил мя еси. 

\pominalnayaslava

Твой есмь аз, спаси мя: яко оправданий Твоих взысках. Мене ждаша грешницы погубити мя, свидения Твоя разумех. Всякия кончины видех конец, широка заповедь Твоя зело. Коль возлюбих закон Твой, Господи, весь день поучение мое есть. Паче враг моих умудрил мя еси заповедию Твоею, яко в век моя есть. Паче всех учащих мя разумех, яко свидения Твоя поучение мое есть. Паче старец разумех, яко заповеди Твоя взысках. От всякаго пути лукава возбраних ногам моим, яко да сохраню словеса Твоя. От судеб Твоих не уклонихся, яко Ты законоположил ми еси. Коль сладка гортани моему словеса Твоя: паче меда устом моим. От заповедей Твоих разумех: сего ради возненавидех всяк путь неправды. Светильник ногама моима закон Твой, я свет стезям моим. Кляхся, и поставих сохранити судьбы правды Твоей. Смирихся до зела, Господи, живи мя по словеси Твоему. Вольная уст моих благоволи же, Господи, и судьбам Твоим научи мя. Душа моя в руку Твоею выну, и закона Твоего не забых. Положиша грешницы сеть мне, и от заповедей Твоих не заблудих. Наследовах свидения Твоя во век, яко радование сердца моего суть. Приклоних сердце мое, сотворити оправдания Твоя в век за воздаяние. Законопреступныя возненавидех, закон же Твой возлюбих. Помощник мой и Заступник мой еси Ты, на словеса Твоя уповах. Уклонитеся от мене лукавнующии, и испытаю заповеди Бога моего. Заступи мя по словеси Твоему, и жив буду, и не посрами мене от чаяния моего. Помози ми, и спасуся, и поучуся во оправданиих Твоих выну. Уничижил еси вся отступающыя от оправданий Твоих, яко неправедно помышление их. Преступающыя непщевах вся грешныя земли, сего ради возлюбих свидения Твоя. Пригвозди страху Твоему плоти моя: от судеб бо Твоих убояхся. Сотворих суд и правду, не предаждь мене обидящым мя. Восприими раба Твоего во благо, да не оклеветают мене гордии. Очи мои исчезосте во спасение Твое и в слово правды Твоея: сотвори с рабом Твоим по милости Твоей, и оправаданием Твоим научи мя. Раб Твой есмь аз: вразуми мя, и увем свидения Твоя. Время сотворити Господеви: разориша закон Твой. Сего ради возлюбих заповеди Твоя паче злата и топазия. Сего ради ко всем заповедем Твоим направляхся, всяк путь неправды возненавидех. Дивна свидения Твоя: сего ради испыта я душа моя. Явление словес Твоих просвещает и вразумляет младенцы. Уста моя отверзох, и привлекох дух, яко заповедай Твоих желах. 

\pominalnayaslava

Призри на мя и помилуй мя, по суду любящих имя Твое. Стопы моя направи по словеси Твоему, и да не обладает мною всякое беззаконие. Избави мя от клеветы человеческия, и сохраню заповеди Твоя. Лице Твое просвети на раба Твоего, и научи мя оправданием Твоим. Исходища водная изведосте очи мои: понеже не сохраних закона Твоего.. Праведен еси, Господа, и прави суди Твои. Заповедал еси правду свидения Твоя, и истину зело. Истаяла мя есть ревность Твоя: яко забыша словеса Твоя врази мои. Разжжено слово Твое зело, и раб Твой возлюби е. Юнейший аз есмь и уничижен: оправданий Твоих не забых. Правда Твоя правда во век, и закон Твой истина. Скорби и нужды обретоща мя: заповеди Твоя поучение мое. Правда свидения Твоя в век: вразуми мя, и жив буду. Воззвах всем сердцем моим, услыши мя, Господи, оправдания Твоя взыщу. Воззвах Ти, спаси мя, и сохраню свидения Твоя. Предварих в безгодии и воззвах, на словеса Твоя уповах. Предваристе очи мои ко утру, поучитися словесем Твоим. Глас мой услыши, Господи, по милости Твоей: по судьбе Твоей живи мя. Приближишася гонящий мя беззаконием: от закона же Твоего удалишася. Близ еси Ты, Господи, и вси путие Твои истина. Исперва познах от свидений Твоих, яко в век основал я еси. Виждь смирение мое, и изми мя, яко закона Твоего не забых. Суди суд мой, и избави мя: словесе ради Твоего живи мя. Далече от грешник спасение, яко оправданий Твоих не взыскали. Щедроты Твоя многи, Господи, по судьбе Твоей живи мя. Мнози изгонящии мя и стужающии ми: от свидений Твоих не уклонихся. Видех неразумевающыя и истаях: яко словес Твоих не сохранила. Виждь, яко заповеди Твоя возлюбих: Господи, по милости Твоей живи мя. Начало словес Твоих истина, и во век вся судьбы правды Твоея. Князи погнаша мя туне: и от словес Твоих убояся сердце мое. Возрадуюся аз о словесех Твоих, яко обретаяй корысть многу. Неправду возненавидех и омерзих: закон же Твой возлюбих. Седмерицею днем хвалих Тя о судьбах правды Твоея. Мир мног любящым закон Твой, и несть им соблазна. Чаях спасения Твоего, Господи, и заповеди Твоя возлюбих. Сохрани душа моя свидения Твоя, и возлюби я зело. Сохраних заповеди Твоя и свидения Твоя, яко вси путие мои пред Тобою, Господи. Да приближится моление мое пред Тя, Господи: по словеси Твоему вразуми мя. Да внидет прошение мое пред Тя, Господи, по словеси Твоему избави мя. Отрыгнут устне мои пение, егда научиши мя оправданием Твоим. Провещает язык мой словеса Твоя, яко вся заповеди Твоя правда. Да будет рука Твоя еже спасти мя, яко заповеди Твоя изволих. Возжелах спасение Твое, Господи, и закон Твой поучение мое есть. Жива будет душа моя, и восхвалит Тя: и судьбы Твоя помогут мне. Заблудих, яко овча погибшее: взыщи раба Твоего, яко заповедей Твоих не забых. 

Слава Отцу и Сыну и Святому Духу, и ныне и присно и во веки веков. Аминь. 

Аллилуиа, аллйлуиа, аллилуиа, слава Тебе Боже. \myemph{(Трижды)} 

Господи помилуй. \myemph{(Трижды)}

\mysubtitle{По 17-й кафизме}

\TrisviatoePoOtcheNash

\mysubtitle{И тропари, глас 2-й}

Согреших к Тебе, Спасе, яко блудный сын: приими мя, Отче, кающагося, и помилуй мя, Боже. 

\slavan

Зову к Тебе, Христе Спасе, мытаревым гласом: очисти мя якоже онаго, и помилуй мя, Боже. 

\inynen

 Богородице, не презри мя требующа заступления Твоего: на Тя бо упова душа моя, и помилуй мя. 

Господи, помилуй. \myemph{(40 раз)}

\mysubtitle{Молитва}

Владыко Господи Вседержителю и Творче всех, щедрот Отец, и милости Бог, от земли создавый человека, и показавый его по образу Твоему и по подобию, да и тем прославится великолепое имя Твое на земли, и исторгнена убо преступлением Твоих заповедей, паки на лучшее возсоздавый его во Христе Твоем и возведый на Небеса: благодарю Тя, яко умножил еси на мне величие Твое, и не предал мя еси врагом моим в конец, исторгнути мя ищущым в пропасть адову, ниже оставил мя еси погибнути со беззаконии моими. Ныне убо, Многомилостиве и Любоблаже Господи, не хотяй смерти грешнаго, но обращения ожидаяй и приемляй: Иже низверженныя исправляяй, сокрушенныя исцеляяй, обрати и мене к покаянию, и низвержаннаго исправи, и сокрушеннаго исцели: помяни Твоя щедроты, и яже от века Твою непостижимую благость и моя безмерная забуди беззакония, яже делом, и словом, и мыслию соверших: разреши ослепление сердца моего, и даждь ми слезы умиления на очищение скверны мысли моея. Усльши, Господи, вонми, Человеколюбче, очисти, Благоутробне, и от мучительства во мне царствующих страстей окаянную мою душу свободи. И нектому да содержит мя грех: ниже да возможет на мя боритель демон, ниже к своему хотению да ведет мя, но державною Твоею рукою, его владычества исхитивый мя, Ты царствуй во мне, Благий и Человеколюбивый Господи, и всего Твоего быти, и жити мне прочее по Твоей благоволи воли. И подаждь ми неизреченною благостию сердца очищение, уст хранение, правоту деяний, мудрование смиренное, мир помыслов, тишину душевных моих сил, радость духовную, любовь истинную, долготерпение, благость, кротость, веру нелицемерну, воздержание обдержательное, и всех мя благих плодов исполни дарованием Святаго Твоего Духа. И не возведи мене в преполовение дний моих, ниже неисправлену и неготову душу мою восхитиши: но соверши мя Твоим совершенством, и тако мя настоящаго жития изведи, яко да невозбранно прошед начала и власти тьмы, Твоею благодатию узрю и аз неприступныя Твоея славы доброту неизреченную, со всеми святыми Твоими, в нихже освятися и прославися всечестное и великолепое имя Твое, Отца и Сына и Святаго Духа, ныне и присно и во веки веков. Аминь.

\end{mymulticols}

\mychapterending

\mychapter{Значение 17-й кафизмы}\begin{mymulticols}
%http://www.molitvoslov.org/text229.htm 
 

В течение всех сорока дней по смерти человека его родные и близкие должны читать Псалтирь. Сколько кафизм в день "--- зависит от времени и сил читающих, но непременно чтение должно быть ежедневным. Когда прочитана вся Псалтирь, она читается сначала. Только не следует забывать, что после каждой «Славы…» надо читать молитвенное прошение о поминовении усопшего (из «Последования по исходе души от тела»). Многие родные и близкие умершего, ссылаясь на то, что нет времени или не имеют Псалтири, или не умеют читать по-церковнославянски, доверяют это чтение другим (чтецам) за плату или иное вознаграждение. Но молитва будет сильнее, искреннее, чище, если родной или близкий умершему человек будет сам просить Бога о помиловании усопшего. 

В третий, девятый, сороковой дни следует читать по усопшему 17-ю кафизму. 

В этой кафизме изображается блаженство ходивших в законе Господнем, т.~е. блаженство праведных людей, старавшихся жить по заповедям Божиим. 

Смысл и значение 118 псалма раскрываются в 19-м стихе: «Пришлец (странник) аз есмь на земли: не скрый от мене заповеди Твоя». Толковая Библия под ред. А.П. Лопухина дает этому стиху следующее объяснение: «Жизнь на земле есть странствование, путешествие, совершаемое человеком для достижения своего отечества и постоянного, вечного местопребывания. Очевидно, последнее "--- не на земле, а за гробом. Если же так, то земная жизнь должна быть подготовленном к загробной и к ней может привести только безошибочно избранный на земле путь. Как и где найти последний? Этот путь указан в заповедях Закона. Кто не следует им, тот заблуждается и не достигнет загробной обители, т.~е. загробного упокоения, как награды за понесенные труды к его достижению. Здесь "--- довольно ясное учение о цели земного существования, бессмертии человеческой души и загробном мздовоздаянии». 

\end{mymulticols}

\mychapterending

\mychapter{Последование панихиды}\begin{mymulticols}
%http://www.molitvoslov.org/text231.htm 
 

\myemph{ После обычного начала и 90-го псалма:} «Живый в помощи Вышняго…» \myemph{ Великая ектения} «Миром Господу помолимся» \myemph{ с заупокойными прошениями, затем} «Аллилуиа» \myemph{ со стихами, потом:}

\mysubtitle{Тропарь, глас 8-й:}

Глубиною мудрости человеколюбно вся строяй, и полезное всем подаваяй, Едине Содетелю, упокой, Господи, души раб Твоих: на Тя бо упование возложиша, Творца и Зиждителя и Бога нашего. 

Слава:, \inynen

 Тебе и стену и пристанище имамы, и Молитвенницу благоприятну к Богу, Егоже родила еси, Богородице Безневестная, верных спасение.

\mysubtitle{Тропари за упокой, глас 5-й}

\pripev[Запев:]{Благословен еси, Господи, научи мя оправданием Твоим. }

Святых лик обрете источник жизни и дверь райскую, да обрящу и аз путь покаянием, погибшее овча аз есмь, воззови мя, Спасе, и спаси мя. 

\pripev[Запев:]{Благословен еси, Господи, научи мя оправданием Твоим. }

Агнца Божия проповедавше и заклани бывше якоже агнцы, и к жизни нестареемей, святии; и присносущней преставльшеся: Того прилежно, мученицы, молите долгов разрешение нам даровати. 

\pripev[Запев:]{Благословен еси, Господи, научи мя оправданием Твоим. }

В путь узкий хождшии прискорбный, вси в житии крест яко ярем вземшии и Мне последовавший верою, приидите насладитеся, ихже уготовах вам почестей и венцев небесных. 

\pripev[Запев:]{Благословен еси, Господи, научи мя оправданием Твоим. }

Образ есмь неизреченныя Твоея славы, аще и язвы ношу прегрешений, ущедри Твое создание, Владыко, и очисти Твоим благоутробием, и возжеленное Отечество подаждь ми, рая паки жителя мя сотворяя. 

\pripev[Запев:]{Благословен еси, Господи, научи мя оправданием Твоим. }

Древле убо от не сущих создавый мя и образом Твоим Божественным почтый, преступлением же заповеди, паки мя возвративый в землю, от неяже взят бых, но еже по подобию возведи древнею добротою возобразитися. 

\pripev[Запев:]{Благословен еси, Господи, научи мя оправданием Твоим. }

Упокой, Боже, рабы Твоя и учини я в раи, идеже лицы святых, Господи, и праведницы сияют яко светила; усопшия рабы Твоя упокой, презирая их вся согрешения. 

Слава Отцу и Сыну и Святому Духу. 

Трисиятельное Единаго Божества, благочестно поем вопиюще: Свят еси, Отче Безначальный, Собезначальный Сыне, и Божественный Душе; просвети нас, верою Тебе служащих, и вечнаго огня исхити. 

И ныне и присно и во веки веков. Аминь. 

Радуйся, Чистая, Бога плотию рождшая во спасение всех, Еюже род человеческий обрете спасение, Тобою да обрящем рай, Богородице Чистая, Благословенная! 

Аллилуиа, аллилуиа, аллилуиа, слава Тебе, Боже. (\myemph{ Трижды}) 

\myemph{ Посем ектения:} Паки и паки… 

Еще молимся о упокоении душ усопших рабов Божиих (\myemph{ имярек}), и о еже простится им всякому прегрешению, вольному же и невольному. 

\myemph{ Хор:} Яко да Господь Бог учинит души их, идеже праведнии упокояются. 

\myemph{ Хор:} Господи, помилуй. 

Милости Божия, Царства Небеснаго, и оставления грехов их, у Христа Безсмертнаго Царя и Бога нашего просим. 

\myemph{ Хор:} Подай, Господи.

\mysubtitle{Молитва}

Боже духов и всякия плоти, смерть поправый и диавола упразднивый, и живот миру Твоему даровавый! Сам Господи, упокой души усопших рабов Твоих (\myemph{ имярек}), в месте светле, в месте злачне, в месте покойне, отнюдуже отбеже болезнь, печаль и воздыхание. Всякое согрешение, содеянное ими, словом, или делом, или помышлением, яко Благий Человеколюбец Бог прости, яко несть человек, иже жив будет и не согрешит, Ты бо един, кроме греха, правда Твоя правда во веки, и слово Твое истина. 

\myemph{ Возглас:} Яко Ты еси Воскресение, и Живот, и покой усопших рабов Твоих…

\mysubtitle{Седален, глас 5-й:}

Покой, Спасе наш, с праведными рабы Твоя, и сия всели во дворы Твоя, якоже есть писано, презирая яко Благ, прегрешения их, вольная и невольная, и вся яже в ведении и не в ведении, Человеколюбче. 

\slavainynen

 От Девы возсиявый миру, Христе Боже, сыны света Тою показавый, помилуй нас.

\mysubtitle{Ирмосы канона, глас 8:}

\mysubtitle{ПЕСНЬ 1:} Яко по суху пешешествовав Израиль, по бездне стопами, гонителя фараона видя потопляема, Богу победную песнь поим вопияше. 

\myemph{Катавасия:} Воду прошед яко сушу, и египетскаго зла избежав, израильтянин вопияше: 

Избавителю и Богу нашему поим. 

Запевы: Покой, Господи, души усопших раб Твоих. 

Слава Отцу и Сыну и Святому Духу. 

И ныне и присно и во веки веков. Аминь.

\mysubtitle{ПЕСНЬ 3:} Несть свят, якоже Ты, Господи Боже мой, вознесый рог верных Твоих, Блаже, и утвердивый нас на камени исповедания Твоего. 

\myemph{Катавасия:} Небеснаго круга Верхотворче, Господи, и Церкве Зиждителю, Ты мене утверди в любви Твоей, желаний краю, верных утверждение, Едине Человеколюбче. 

\mysubtitle{ПЕСНЬ 6:} Житейское море, воздвизаемое зря напастей бурею, к тихому пристанищу Твоему притек, вопию Ти: возведи от тли живот мой, Многомилостиве. 

\myemph{Катавасия:} Молитву пролию ко Господу, и Тому возвещу печали моя, яко зол душа моя исполнися, и живот мой аду приближися, и молюся яко Иона: от тли, Боже, возведи мя. 

\myemph{Кондак, глас 8-й:} Со святыми упокой, Христе, душы раб Твоих, идеже несть болезнь, ни печаль, ни воздыхание, но жизнь безконечная. 

\myemph{Икос:} Сам Един еси Безсмертный, сотворивый и создавый человека: земнии убо от земли создахомся, и в землю туюжде пойдем, якоже повелел еси, Создавый мя и рекий ми: яко земля еси и в землю отыдеши, аможе вси человецы пойдем, надгробное рыдание творяще песнь: аллилуиа, аллилуиа, аллилуиа. 

\myemph{ Иерей:} Богородицу и Матерь Света в песнех возвеличим. 

\myemph{ Хор:} Дуси и души праведных восхвалят Тя, Господи. 

\mysubtitle{ПЕСНЬ 9:} Бога человеком невозможно видети, на Него же не смеют чини ангельстии взирати: Тобою бо, Всечистая, явися человеком Слово воплощенно, Егоже величающе, с небесными вой Тя ублажаем. 

\myemph{Катавасия:} Ужасеся о сем небо, и земли удивишася концы, яко Бог явися человеком плотски, и чрево Твое бысть пространнейшее небес. Тем Тя, Богородицу, ангелов и человек чиноначалия величают. 

\TrisviatoePoOtcheNash

\myemph{Тропари, глас 4-й:} Со духи праведных скончавшихся души раб Твоих, Спасе, упокой, сохраняя их во блаженней жизни, яже у Тебе, Человеколюбче. 

В покоищи Твоем, Господи, идеже вси святии Твои упокоеваются, упокой и души раб Твоих, яко Един еси Человеколюбец. 

\slavan

Ты еси Бог, сошедый во ад и узы окованных разрешивый, Сам и души раб Твоих упокой. 

\inynen

 Едина Чистая и Непорочная Дево, Бога без семене рождшая, моли спастися душам их. 

\myemph{Посем ектения:} Помилуй нас, Боже… 

\myemph{Отпуст и возглашение иерея или диакона:}

Во блаженном успении вечный покой подаждь, Господи, усопшим рабом Твоим \myemph{(имярек)} и сотвори им вечную память. 

\myemph{ Хор:} Вечная память. (\myemph{ Трижды})

\end{mymulticols}

\mychapterending

\mychapter{Чин литии, совершаемой мирянином дома и на кладбище}\begin{mymulticols}
%http://www.molitvoslov.org/text233.htm 
 

\MolitvamiSviatyhOtecNashih 

Слава Тебе, Боже наш, слава Тебе. 

Царю Небесный, Утешителю, Душе истины, Иже везде сый и вся исполняяй, Сокровище благих и жизни Подателю, прииди и вселися в ны, и очисти ны от всякия скверны, и спаси, Блаже, души наша. 

\TrisviatoePoOtcheNash

Господи, помилуй. \myemph{(12 раз)}

\slavainynen

\priiditepoklonimsia

\mysubtitle{Псалом 90}

Живый в помощи Вышняго, в крове Бога Небеснаго водворится. Речет Господеви: Заступник мой еси, и Прибежище мое, Бог мой, и уповаю на Него. Яко Той избавит тя от сети ловчи и от словесе мятежна, плещма Своима осенит тя, и под криле Его надеешися: оружием обыдет тя истина Его. Не убоишися от страха нощнаго, от стрелы летящия во дни, от вещи во тьме преходящия, от сряща и беса полуденнаго. Падет от страны твоея тысяща, и тма одесную тебе, к тебе же не приближится: обаче очима твоима смотриши, и воздаяние грешников узриши. Яко Ты, Господи, упование мое, Вышняго положил еси прибежище твое. Не приидет к тебе зло, и рана не приближится телеси твоему. Яко ангелом Своим заповесть о тебе, сохранити тя во всех путех твоих. На руках возмут тя, да не когда преткнеши о камень ногу твою. На аспида и василиска наступиши, и попереши льва и змия. Яко на Мя упова, и избавлю и; покрыю и, яко позна имя Мое. Воззовет ко Мне, и услышу его; с ним есмь в скорби, изму его, и прославлю его; долготою дний исполню его, и явлю ему спасение Мое. 

Слава Отцу и Сыну и Святому Духу, и ныне и присно и во веки веков. Аминь. 

Аллилуиа, аллилуиа, аллилуиа, слава Тебе, Боже. \myemph{(Трижды)}

\mysubtitle{Тропарь, глас 4-й:}

Со духи праведных скончавшихся, душу раба Твоего, Спасе, упокой, сохраняя ю во блаженной жизни, яже у Тебе, Человеколюбче. 

В покоищи Твоем, Господи, идеже вси святии Твои упокоеваются, упокой и душу раба Твоего, яко Един еси Человеколюбец. 

\slavan

Ты еси Бог, сошедый во ад, и узы окованных разрешимый, Сам и душу раба Твоего упокой. 

\inynen

 Едина Чистая и непорочная Дево, Бога без семени рождшая, моли спастися душе его.

\mysubtitle{Седален, глас 5-й:}

Покой, Спасе наш, с праведными раба Твоего, и сего всели во дворы Твоя, якоже есть писано, презирая, яко Благ, прегрешения его вольная и невольная, и вся яже в ведении и не в ведении, Человеколюбче. 

\mysubtitle{Кондак, глас 8-й:}

Со святыми упокой, Христе, душу раба Твоего, идеже несть болезнь, ни печаль, ни воздыхание, но жизнь безконечная.

\mysubtitle{Икос:}

Сам Един еси Безсмертный, сотворивый и создавый человека, земнии убо от земли создахомся, и в землю туюжде пойдем, якоже повелел еси, Создавый мя и рекий ми: яко земля еси, и в землю отъидеши, аможе вси человецы пойдем, надгробное рыдание творяще песнь: аллилуиа, аллилуиа, аллилуиа. 

\Chestneyshuyu 

Слава Отцу и Сыну и Святому Духу, и ныне и присно и во веки веков. Аминь. 

Господи, помилуй \myemph{(Трижды)}, благослови. 

\MolitvamiSviatyhOtecNashih

Во блаженном успении вечный покой подаждь, Господи, усопшему рабу Твоему \myemph{(имя)}, и сотвори ему вечную память. 

Вечная память. \myemph{(Трижды)}

Душа его во благих водворится, и память его в род и род. 

\end{mymulticols}

\mychapterending

\mychapter{Акафист за единоумершего}
%http://www.molitvoslov.org/text230.htm 
 

{\centering\myemph{ Читается ежедневно в течение 40 дней по смерти и столько же перед годовщиной смерти.}\par} 

\begin{mymulticols}

\newcommand{\akafistZaEdinKondakOne}{%
\mysubtitle{Кондак 1}

Избранный Ходатаю и Первосвященниче, о спасении мира грешнаго душу Свою положивый, давый нам власть чадами Божиими быти и обитати в невечернем дни Царствия Твоего! Даруй прощение и вечную радость усопшему, о немже с мольбою взываем Ти: Иисусе, Судие Всемилостивый, рая сладости сподоби раба Твоего.
}

\newcommand{\akafistZaEdinIkosOne}{%
\mysubtitle{Икос 1}

Господом данный Ангеле хранителю святый, прииди помолиться о рабе твоем, егоже на всех путех жизни сопровождал, хранил и наставлял еcи, воззови с нами ко Спасу Всещедрому. 

Иисусе, истреби рукописание грехов раба Твоего (\myemph{ имярек}). 

Иисусе, уврачуй его язвы душевные. 

Иисусе, да не будет на земли горьких воспоминаний о нем. 

Иисусе, сего ради помилуй огорчавших его и обиженных им. 

Иисусе покрый его несовершенства светоносною ризою Твоего искупления. 

Иисусе, возвесели его милосердием Твоим. 

Иисусе неизреченный, великий и чудный, явися ему Сам. 

Иисусе, Судие Всемилостивый, рая сладости сподоби раба Твоего.
}

\akafistZaEdinKondakOne

\akafistZaEdinIkosOne

\mysubtitle{Кондак 2}

Яко безутешная горлица носится душа над юдолию земною, созерцая с высоты божественного разумения грехи и соблазны минувшаго пути, горько скорбя о каждом невозвратном дне, ушедшем без пользы, но помилуй раба Твоего, Владыко, да внидет он в покой Твой, взывая: Аллилуиа.

\mysubtitle{Икос 2}

Если Ты страдал о всем мире, если Ты проливал слезы и кровавый пот о живых и мертвых, то кто нас удержит от молитвы за усопшаго. Подражая Тебе, сошедшему даже до ада, молимся о спасении раба Твоего. 

Иисусе, жизни Подателю, озари его светом Твоим. 

Иисусе, да будет он едино с Тобою и Отцем. 

Иисусе, всех призываяй в виноградник Твой, не забуди озарити его светом Твоим. 

Иисусе, щедрый Раздаятелю вечных наград, яви его светом чертога Твоего. 

Иисусе, верни душе его благодатныя силы первозданныя чистоты. 

Иисусе, да умножатся во имя его добрыя дела. 

Иисусе, согрей осиротевших Твоею таинственною отрадою. 

Иисусе, Судие Всемилостивый, рая сладости сподоби раба Твоего.

\mysubtitle{Кондак 3}

Связанный узами плоти раб Твой падал греховно, но дух его томился по Твоей вечной правде и святости, ныне же, когда немощь плотская скована могильным тлением, да вознесется душа его превыше солнца к Тебе Всесвятому и воспоет песнь избавления: Аллилуиа.

\mysubtitle{Икос З}

Верховный апостол Твой в холодную нощь у костра трижды отрекся от Тебе и Ты спасл еси его. Единый ведый человеческого естества немощь, прости и рабу Твоему (\myemph{ имярек}) многовидныя отпадения от воли Твоея. 

Иисусе, всели его там, где нет заблуждений. 

Иисусе, избави его от тягостных терзаний совести. 

Иисусе, да сгинет навеки память грехов. 

Иисусе, соблазнов юности его не помяни. 

Иисусе, от тайных беззаконий очисти его. 

Иисусе, осени его тихим светом спасения. 

Иисусе, Судие Всемилостивый, рая сладости сподоби раба Твоего.

\mysubtitle{Кондак 4}

Бури жизни миновали, страдания земные окончены, безсильны враги с их злобою, но сильна любовь, избавляющая от вечного мрака и спасающая всех, о ком возносится Тебе дерзновенная песнь: Аллилуиа.

\mysubtitle{Икос 4}

Ты без числа милосерд к нам. Ты единый Избавитель, что мы прибавим к подвигу спасающей любви Твоей, но Симон Киринейский помогал нести Крест Тебе Всесильному, так и ныне благости Твоей угодно спасение близких совершать с участием нашим. 

Иисусе, Ты заповедал друг друга тяготы носити. 

Иисусе, союз любви положивый между мертвыми и живыми. 

Иисусе, да послужат подвиги любящих во спасение рабу Твоему (\myemph{ имярек}). 

Иисусе, услыши его вопль сердечный, возносимый нашими устами. 

Иисусе, в наших слезах приими его покаяние. 

Иисусе, Судие Всемилостивый, рая сладости сподоби раба Твоего.

\mysubtitle{Кондак 5}

Боже, да будет принят Тобою его предсмертный вздох сокрушения, якоже мольба благоразумнаго разбойника. Он угас на жизненном кресте, дай наследовать ему Твоя обетования, якоже тому: «Аминь, глаголю тебе, со Мною будеши в раю», где сонмы раскаявшихся грешников радостно поют: Аллилуиа.

\mysubtitle{Икос 5}

За нас Распятый, за нас измученный, простри руку с Твоего Креста, каплями крови Твоея безследно смой излитыя согрешения его, благообразною Твоею наготою согрей обнаженную осиротевшую душу. 

Иисусе, Ты знал его жизнь до рождения и возлюбил его. 

Иисусе, Ты видел его из далече с высоты Креста Твоего. 

Иисусе, Ты простирал ему в даль грядущему изъязвленные объятия Твои. 

Иисусе, Ты Взывал о прощении его на кровавой Голгофе. 

Иисусе, Ты кротко умирал за него в тяжких муках. 

Иисусе, претерпевый во гробе положение, освяти его могильный покой. 

Иисусе воскресший, вознеси ко Отцу озлобленную миром и Тобою спасенную душу. 

Иисусе, Судие Всемилостивый, рая сладости сподоби раба Твоего.

\mysubtitle{Кондак 6}

Спит он вечным сном могилы, но душа его не дремлет, чает Тебя Господи, жаждет Тебя вечного Жениха. Да исполнятся на умершем слова Твои: «Ядый Мою Плоть и пияй Мою Кровь, имать живот вечный». Дай ему ясти от манны сокровенныя и пети у Престола Твоего: Аллилуиа.

\mysubtitle{Икос 6}

Смерть разлучила со всеми ближними, стала дальше душа, знаемии сокрушаются, и токмо Ты Един остался близок. Разрушились преграды плоти и Ты открылся в неприступном величии Божества с ожиданием ответа. 

Иисусе, Любовь превыше всякого разумения, помилуй раба Твоего. 

Иисусе, удаляясь от Тебя, он тяжко страдал. 

Иисусе, прости неверность его сердца. Иисусе, обманутые надежды рождали тоску по Тебе. 

Иисусе, вспомни те часы, когда душа его трепетала восторгом Твоим. 

Иисусе, дай скончавшемуся неземную радость и покой. 

Иисусе, единый верный, неизменный, приими его. 

Иисусе, Судие Всемилостивый, рая сладости сподоби раба Твоего.

\mysubtitle{Кондак 7}

Веруем, что недолгой будет разлука наша. Мы хороним тебя, как на ниве зерно, ты произрастешь в иной стране. Да погибнут в могиле плевелы грехов твоих, а дела добрые там просияют, где семена добра приносят нетленные плоды, где души святые поют: Аллилуиа.

\mysubtitle{Икос 7}

Когда уделом умершего станет забвение, когда образ его поблекнет в сердцах и время изгладит место с могилою и ревность молитвы о нем, тогда Ты не остави его, дай отраду одинокой душе. 

Иисусе, Твоя любовь не охладевает. 

Иисусе, неистощимо Твое благоволение. 

Иисусе, в неумолкаемых мольбах Церкви да омоются грехи его приношением Жертвы Безкровной. 

Иисусе, предстательством всех святых даруй ему благодать молитвы о живых. 

Иисусе, во дни испытаний наших приими его ходатайство о нас. 

Иисусе, Судие Всемилостивый, рая сладости сподоби раба Твоего.

\mysubtitle{Кондак 8}

Будем молиться со слезами, когда память об усопшем мучительно свежа, будем поминать имя его в нощи и во дни, раздавая милостыню, питая гладных, из глубины души взывая: Аллилуиа.

\mysubtitle{Икос 8}

Тайнозритель Иоанн Богослов созерцал у Престола Агнца Божия великое множество людей, облеченных в белые ризы: это те, кто пришел от великой скорби. Они радостно служат Богу день и ночь, и Бог обитает с ними, и не коснется их мука. 

Иисусе, причти к ним и раба Твоего (\myemph{ имярек}). 

Иисусе, он много страдал и томился. 

Иисусе, известны Тебе все горькие часы и тягостные минуты его. 

Иисусе, на земли он имел печали и скорби, дай на небе отраду. 

Иисусе, услади его от источников живыя воды. 

Иисусе, отыми всяку слезу от очей его. 

Иисусе, всели его, где не опаляет, но живит солнце правды Твоея. 

Иисусе, Судие Всемилостивый, рая сладости сподоби раба Твоего.

\mysubtitle{Кондак 9}

Кончено странствование земное, какой благодатный переход в мир Духа, какое созерцание новых неведомых миру земному вещей и небесных красот, душа возвращается в отечество свое, где светлое солнце, правда Божия просвещает поющих: Аллилуиа.

\mysubtitle{Икос 9}

Если отблеск и след Твой полагает сияние на лице смертных, то каков же Ты Сам. Если плоды Твоих рук так прекрасны, и земля, отражающая только тень Твою, полна невыразимым величием, то каков же невидимый лик Твой. Открой славу Твою усопшему рабу Твоему (\myemph{ имярек}). 

Иисусе, обостри его слух к восприятию Твоего Божества. 

Иисусе, обостри его слух к разумению небесных. 

Иисусе, да будет его радость преисполненная. 

Иисусе, подкрепи его надеждою встречи в обителях блаженных. 

Иисусе, дай нам почувствовать благодатную силу заупокойной молитвы. 

Иисусе, Судие Всемилостивый, рая сладости сподоби раба Твоего.

\mysubtitle{Кондак 10}

Отче наш, скончавшагося в Царствие Твое приими, где нет греха и зла, где нерушима Святая воля, где в сонмах чистейших душ и непорочных ангелов святится Твое благодатное имя и благоухает хвала: Аллилуиа.

\mysubtitle{Икос 10}

В тот день ангелы поставят Престол Твой, Судие, и Ты возсияеши во славе Отца Твоего, неся воздаяние всякому человеку. О, воззри тогда милостиво на смиреннаго раба Твоего (\myemph{ имярек}), рцы ему: «Прииди одесную Мене». 

Иисусе, яко Бог власть имаши оставляти грехи. 

Иисусе, прости его согрешения забытыя или стыдом утаенныя. 

Иисусе, отпусти беззакония немощи и неведения. 

Иисусе, избави его от несвятимых глубин адскаго отчаяния. 

Иисусе, да унаследует он Твои животворящия обетования. 

Иисусе, сопричти его благословенным Отца Твоего. 

Иисусе, дай ему во веки нескончаемое блаженство. 

Иисусе, Судие Всемилостивый, рая сладости сподоби раба Твоего.

\mysubtitle{Кондак 11}

Владыко Всеблагий, да отверзятся усопшему солнцевидные врата райские, да встретят его с ликованием соборы праведных и святых, сонмы близких и любящих его, да возрадуются о нем светоносные ангелы Твои, да узрит он и Присноблаженную Матерь Твою там, где победно звучит: Аллилуиа.

\mysubtitle{Икос 11}

Под дыханием Твоим оживают цветы, воскресает природа, пробуждаются сонмы мельчайших тварей, Твой взор светлее весенних небес, Твоя любовь, Иисусе, теплее лучей солнечных. Ты из праха земного воскресил бренную плоть человеческую к расцвету вечной нетленной жизни весны, тогда озари и раба Твоего (\myemph{ имярек}) светом милостей Твоих. 

Иисусе, в Твоей деснице благоволение и жизнь. 

Иисусе, во взоре Твоем свет и любовь. 

Иисусе, избави усопшего от вечныя смерти духовныя. 

Иисусе, он уснул с надеждою, подобно Нилу-реке пред холодной зимой. 

Иисусе, пробуди его, когда терния земли облекутся цветом вечности. 

Иисусе, да не омрачит ничто земное его последнего сна. 

Иисусе, Счастье Неизменное и цель нашего бытия. 

Иисусе, Судие Всемилостивый, рая сладости сподоби раба Твоего.

\mysubtitle{Кондак 12}

Христе! Ты "--- Царство Небесное, Ты "--- земля кротких, Ты "--- обитель многих, Ты "--- питие совершенно новое, Ты "--- одеяние и венец преподобных, Ты "--- ложе упокоения святых, Ты "--- Сладчайший Иисусе! Тебе подобает хвала: Аллилуиа.

\mysubtitle{Икос 12}

Под образом тихих садов неземной красоты, и светлых как солнце обителей и в великолепии небесных песнопений Ты открыл нам блаженство любящих Тя. 

Иисусе, да внидет раб Твой в радость Твою. 

Иисусе, облеки его сиянием славы Отца. 

Иисусе, просвяти его озарением Духа Святаго. 

Иисусе, да будет он восходить от славы в славу. 

Иисусе, да узрит он Тебя лицем к лицу. 

Иисусе, Судие Всемилостивый, рая сладости сподоби раба Твоего.

\mysubtitle{Кондак 13}

О, Женише Безсмертный, в полночь греха и неверия грядый с небес со ангелами судити миру всему. Отверзи двери славного чертога Твоего рабу Твоему (\myemph{ имярек}), да в безчисленных сонмах святых во веки поет: Аллилуиа, Аллилуиа, Аллилуиа. 

\myemph{ Этот кондак читается трижды, затем 1-й икос «Господом данный Ангеле…» и 1-й кондак «Избранный Ходатаю».}

\akafistZaEdinIkosOne

\akafistZaEdinKondakOne

\end{mymulticols}

\mychapterending

\mychapter{Молитва родителей за детей}\begin{mymulticols}
%http://www.molitvoslov.org/text225.htm 
 

Господи Иисусе Христе, Боже наш, Владыко живота и смерти. Утешителю скорбящих! С сокрушенным и умиленным сердцем прибегаю к Тебе и молюся Ти: помяни, Господи, во Царствии Твоем усопшего раба Твоего (-ую рабу Твою), чадо мое (\myemph{ имярек}), и сотвори ему (ей) вечную память. Ты, Владыко живота и смерти, даровал еси мне чадо сие. Твоей же благой и премудрой воле изволися и отьяти е у мене. Буди благословенно Имя Твое, Господи. Молю Тя, Судие неба и земли, безконечною любовию Твоею к нам, грешным, прости усопшему чаду моему вся согрешения его, вольная и невольная, яже словом, яже делом, яже ведением и неведением. Прости, Милостиве, и наша родительская согрешения, да не пребудут они на чадах наших: вем, яко множицею согрешихом пред Тобою, множицею не соблюдохом, не сотворихом, якоже заповедал еси нам. Аще же усопшее чадо наше, нашея или своея ради вины, бяше в житии сем, работая миру и плоти своея и не паче Тебе, Господу и Богу своему: аще возлюби прелести мира сего, а не паче Слово Твое и Заповеди Твоя, аще предавшеся сластем житейским, а не паче сокрушению о гресех своих, и в невоздержании бдение, пост и молитву забвению предавше,"--- молю Тя усердно, прости преблагий Отче, чаду моему вся таковыя прегрешения его, прости и ослаби, аще и ино злое сотвори в житии сем, Христе Иисусе! Ты воскресил еси дщерь Иаира по вере и молитве отца ея. Ты исцелил еси дщерь жены-хананеянки по вере и прошению матери ея: услыши убо и молитву мою, не презри и моления моего о чаде моем. Прости, Господи, прости, вся согрешения его и, простив и очистив душу его, изми муки вечныя и всели ю со всеми святыми Твоими, от века благоугодившими Тебе, идеже несть болезнь, ни печаль, ни воздыхание, но жизнь безконечная: яко несть человек, иже жив будет и не согрешит. Ты Един кроме греха: да егда имаши судити мирови, услышит чадо мое превожделенный глас Твой: приидите, благословении Отца Моего, и наследуйте уготованное вам Царствие от сложения мира. Яко Ты еси Отец милостей и щедрот, Ты живот и воскрешение наше, и Тебе славу возсылаем со Отцем и Святым Духом, ныне и присно и во веки веков. Аминь. 

\end{mymulticols}

\mychapterending

\mychapter{Молитва вдовца за супругу}\begin{mymulticols}
%http://www.molitvoslov.org/text223.htm 
 

Христе Иисусе, Господи и Вседержителю! В сокрушении и умилении сердца моего молюся Тебе: упокой, Господи, душу усопшия рабы Твоея (\myemph{ имярек}), в Небесном Царствии Твоем. Владыко Вседержителю! Ты благоволил еси супружеский союз мужа и жены, егда рекл еси: не добро быти человеку единому, сотворим ему помощника по нему. Ты освятил еси союз сей во образ духовнаго союза Христа с Церковью. Верую, Господи, и исповедую, яко Ты благословил еси сочетати и мене сим святым союзом с единою из рабынь Твоих. Твоей же благой и премудрой воле изволися отъяти у мене сию рабу Твою, юже дал еси мне, яко помощницу и сопутницу жизни моея. Преклоняюся пред сею Твоею волею, и молюся Ти от всего сердца моего, приими моление мое сие о рабе Твоей (\myemph{ имярек}), и прости ей, аще согреши словом, делом, помышлением, ведением и неведением; аще земное возлюби паче небеснаго; аще о одежде и украшении тела своего печеся паче, неже о просвещении одеяния души своея; или аще небреже о чадех своих; аще преогорчи кого словом или делом; аще проропта в сердце своем на ближняго своего или осуди кого или ино что от таковых злых содела. Вся сия прости ей, яко благий и человеколюбивый: яко несть человек, иже жив будет и не согрешит. Не вниди в суд с рабою Твоею, яко созданием Твоим, не осуди ю по грехом ея на вечныя муки, но пощади и помилуй по велицей милости Твоей. Молю и прошу Тя, Господи Сил, даруй ми по вся дни жизни моея не преставати молитися о усопшей рабе Твоей, и даже до кончины живота моего просити ей у Тебе, Судии всего мира, оставления согрешений ея. Да якоже Ты, Боже, положил еси на главу ея венец от камене честна, венчая ю зде на земли; тако увенчай ю вечною Твоею славою в Небесном Царствии Твоем, со всеми святыми, тамо ликующими, да вкупе с ними вечно воспевает Всесвятое Имя Твое с Отцем и Святым Духом. Аминь. 

\end{mymulticols}

\mychapterending

\mychapter{Молитва вдовицы за супруга}\begin{mymulticols}
%http://www.molitvoslov.org/text226.htm 
 

Христе Иисусе, Господи и Вседержителю! Ты "--- плачущих утешение, сирых и вдовиц заступление. Ты рекл еси: призови Мя в день скорби твоея, и изму тя. Во дни скорби своея прибегаю к Тебе аз и молюся Ти: не отврати лица Твоего от мене и услыши моление мое, приносимое Тебе со слезами. Ты, Господи Владыко всяческих, благословил еси сочетати мя со единем из рабов Твоих, во еже быти нам едино тело и един дух; Ты дал еси мне сего раба, яко сожителя и защитника. Твоей же благой и премудрой воле изволися отьяти от мене сего раба Твоего и оставити мя едину. Преклоняйся пред сею Твоею волею и к Тебе прибегаю во дни скорби моея: утоли печаль мою о разлучении с рабом Твоим, другом моим. Аще отъял еси его от меня, не отыми от меня Твоея милости. Якоже некогда приял вдовицы две лепты, тако приими и сие моление мое. Помяни, Господи, душу усопшаго раба Твоего (\myemph{ имярек}), прости ему вся согрешения его, вольная и невольная, аще словом, аще делом, аще ведением и неведением, не погуби его со беззаконьми его и не предаждь вечной муки, но по велицей милости Твоей и по множеству щедрот Твоих ослаби и прости вся согрешения его и вчини его со святыми Твоими, идеже несть болезнь, ни печаль, ни воздыхание, но жизнь безконечная. Молю и прошу Тя, Господи, даруй ми во вся дни жизни моея не преставати молитися о усопшем рабе Твоем, и даже до исхода моего просити у Тебе, Судии всего мира, оставления всех согрешений его и вселения его в небесныя обители, яже еси уготовал любящим Тя. Яко аще бо и согреши, но не отступи от Тебе, и несумненно Отца и Сына и Святаго Духа православно даже до последняго своего издыхания исповеда; тем же веру его, яже в Тя, вместо дел ему вмени: яко несть человек, иже жив будет и не согрешит, Ты Един кроме греха, и правда Твоя "--- правда во веки. Верую Господи, и исповедую, яко Ты услышиши моление мое и не отвратиши лица Твоего от мене. Видя вдовицу, зельне плачущу, умилосердився, сына ея, на погребение несома, воскресил еси: тако, умилосердився, утиши и скорбь мою. Яко же отверзл еси рабу Твоему Феофилу, отшедшему к Тебе, двери милосердия Твоего и простил еси ему прегрешения его по молитвам святыя Церкви Твоея внемля молитвам и милостыням супруги его: сице и аз молю Тя, приими и мое моление о рабе Твоем, и введи его в жизнь вечную. Яко Ты еси упование наше, Ты еси Бог, еже миловати и спасати, и Тебе славу возсылаем со Отцем и Святым Духом, ныне и присно и во веки веков. Аминь!

\end{mymulticols}

\mychapterending{}

\mychapter{О мертворожденных младенцах}
%http://www.molitvoslov.org/text222.htm 
 

Господи, помилуй чад моих, умерших во утробе моей. За веру и слезы мои, ради милосердия Твоего, Господи, не лиши их света Твоего Божественного! 

\mychapterending

\mychapter{Молитва детей за родителей}\begin{mymulticols}
%http://www.molitvoslov.org/text224.htm 
 

Господи, Иисусе Христе, Боже наш! Ты сирых хранитель, скорбящих прибежище и плачущих утешитель. Прибегаю к Тебе аз, сирый, стеня и плача, и молюся Тебе: услыши моление мое и не отврати лица Твоего от воздыханий сердца моего и от слез очей моих. Молюся Тебе, милосердный Господи, утоли скорбь мою о разлучении с родителем (материю) моим (-ею) (\myemph{ имярек}), душу же его (ея), яко отшедшею (-ия) к Тебе с истинною верою в Тя и твердою надеждою на Твое человеколюбие и милость, приими в Царство Твое Небесное. Преклоняюсь пред Твоею святою волею, еюже отъят (-а) бысть у мене, и прошу Тя, не отыми точию от него (нея или них) милости и благосердия Твоего. Вем, Господи, яко Ты, Судия мира сего, грехи и нечестия отцев наказуеши в детях, внуках и правнуках даже до третьяго и четвертаго рода: но и милуеши отцев за молитвы и добродетели чад их, внуков и правнуков. С сокрушением и умилением сердца молю Тя, милостивый Судие, не наказуй вечным наказанием усопшаго (-ую) незабвеннаго (-ую) для мене раба (-у) Твоего (ю), родителя (матерь) моего (-ю) (имярек), но отпусти ему (ей) вся согрешения его (ея) вольная и невольная, словом и делом, ведением и неведением сотворенная им (ею) в житии его (ея) зде на земли, и по милосердию и человеколюбию Твоему, молитв ради Пречистая Богородицы и всех святых, помилуй его (ю) и вечныя муки избави. Ты, милосердный Отче отцев и чад! Даруй мне, во вся дни жизни моея, до последняго издыхания моего, не преставати памятовати о усопшем родителе (матери) моем (ей) в молитвах своих, и умоляти Тя, праведнаго Судию, да вчиниши его (ю) в месте светле, в месте прохладно и в месте покойне, со всеми святыми, отнюдуже отбеже всяка болезнь, печаль и воздыхание. Милостиве Господи! Приими днесь о рабе Твоем (Твоей) (имярек) теплую молитву мою сию и воздай ему (ей) воздаянием Твоим за труды и попечения воспитания моего в вере и христианском благочестии, яко научившему (-ей) мя первее всего ведети Тя, своего Господа, в благоговении молитися Тебе, на Тебе Единого уповати в бедах, скорбех и болезнех и хранити Заповеди Твоя; за благопопечение его-(ея) о моем духовном преуспеянии, за тепле приносимыя им (ею) о мне моления пред Тобою и за все дары, им (ею) испрошенные мне от Тебе, воздай ему (ей) Своею милостию, Своими небесными благами и радостями в вечном Царствии Твоем. Ты бо еси Бог милостей и щедрот и человеколюбия. Ты покой и радость верных рабов Твоих, и Тебе славу возсылаем со Отцем и Святым Духом, и ныне и присно и во веки веков. Аминь. 

\end{mymulticols}

\mychapterending{}

\mychapter{Молитва за усопших внезапною (скоропостижною) смертию}\begin{mymulticols}
%http://www.molitvoslov.org/text221.htm 
 

Неисповедимы судьбы Твои, Господи! Неизследимы пути Твои! Даяй дыхание всякой твари и вся от не сущих в бытие приведый, Ты овому посылаеши Ангела смерти в День, егоже не весть, и в час, егоже не чает; оваго же исхищаеши из руки смерти, даруеши живот при последнем издыхании; овому долготерпиши и даеши время на покаяние; оваго же, яко злак, сечеши мечем смерти во едином часе, во мгновении ока; оваго поражаеши громом и молниею, оваго сожигаеши пламенем, оваго же предаеши в снедь зверем дубравниим; овых повелеваеши поглотити волнам и безднам морским и пропастям земным; овых похищаеши язвою губительною, идеже смерть яко жнец пожинает и разлучает отца или матерь от чад их, брата от брата, супруга от супруги, младенца отторгает от лона матерняго, бездыханными повергает сильных земли, богатых и убогих. Что убо сие есть? Дивное и недоумеваемое нами смотрение Твое, Боже! Но Господи, Господи! Ты токмо Един, ведый вся, веси, чесо ради бывает сия и чесо ради быть, яко раб Твоей (раба Твоя) (\myemph{ имя}) во едино мгновение ока пожерт(-та) бысть зиянием смерти. Аще наказуеши за многая, тяжкия прегрешения его(ея), молим Тя, Многомилостиве и Всемилостиве Господи, да не яростию Твоею обличиши его (ю) и накажеши вконец, но, по благости Твоей и по безприкладному Твоему милосердию, яви ему (ей) велию милость Твою в прощении и оставлении грехов. Аще ли же усопший (-шая) раб Твой (раба Твоя), в житии сем помышляя День Судный, позна свое окаянство и возжела принести Тебе плоды достойные покаяния, но не достиже сего, зане позван(-на) бысть Тобою в день, егоже не веда, и в час, егоже не жда, того убо ради паче молим Тя, Преблагий и Премилосердый Господи, несоделанное покаяние, еже видесте очи Твои, и неоконченное дело спасения его(ея) исправи, устрой, соверши Твоею неизреченною благостию и человеколюбием; едину бо надежду имамы на Твое безконечное милосердие: у Тебе бо суд и наказание, у Тебе правда и неистощимая милость; Ты наказуеши, вкупе же и милуеши; биеши, вкупе же и приемлеши; прилежно молим Тя, Господи Боже наш, не накажи внезапу позванного(-ую) к Тебе Страшным Судом Твоим, но пощади, пощади его (ю) и не отрини от лица Твоего. О, страшно внезапу впасти в руце Твои, Господи, и предстати на суд Твой нелицеприятный! Страшно отъити к Тебе без благодатнаго напутствия, без Покаяния и причащения Святых Твоих страшных и животворящих Тайн, Господи! Аще внезапу усопший(-шая) приснопоминаемый(-мая) нами раб Твой (раба Твоя) толико многогрешен(-на), толико повинен(-на) есть осуждению на Твоем праведном суде, молим Тя, умилосердися над ним (нею), не осуди его (ю) на вечное мучение, на вечную смерть; потерпи нам еще, даждь нам долготу дний наших, еже молитися Ти по вся дни о усопшем(-шей) рабе Твоем(ей), дондеже услышиши нас и приимеши милостию Твоею внезапу отшедшаго(-ую) к Тебе; и даждь нам, Владыко, омыти грехи его (ея) слезами сокрушения и воздыханиями нашими пред Тобою, да не низведен(-на) будет по грехом своим на место мучения раб Твой (раба Твоя) (имя), но да вселится в место упокоения. Ты Сам, Господи, повелеваеши ударяти в двери милосердия Твоего, молим Тя убо, Прещедрый Царю, и не престанем умоляти Твое милосердие и взывати с кающимся Давидом: помилуй, помилуй раба Твоего (рабу Твою), Боже, по велицей Твоей милости. Аще же недовольно Ти словес наших, сего малаго моления нашего, умоляем Тя, Господи, верою в спасительныя заслуги Твои, упованием на искупительную и чудодейственную силу Твоея жертвы, принесенныя Тобою за грехи всего мира; молим Тя, о Сладчайший Иисусе! Ты еси Агнец Божий, вземляй грехи мира, распныйся ради нашего спасения! Молим Тя, яко Спасителя и Искупителя нашего, спаси и помилуй и вечныя муки избави душу приснопоминаемаго(-ыя) нами внезапу усопшаго(-шия) раба Твоего (рабы Твоея) (имя), не остави его (ю) погибнути во веки, но сподоби достигнути тихаго пристанища Твоего и упокой тамо, идеже вси святии Твои упокоеваются. Вкупе же молим Тя, господи Иисусе Христе Боже наш, приими милостию Твоею и всех внезапу преставльшихся к Тебе рабов Твоих (\myemph{ имена}), ихже вода покры, трус объят, убийцы убиша, огнь попали, град, снег, мраз, голоть и дух бурен умертвиша, гром и молния попали, губительная язва порази, или иною коею виною умреша, по Твоему изволению и попущению, молим Тя, приими их под Твое благоутробие и воскреси их в жизнь вечную, святую и блаженную. Аминь.

\end{mymulticols}

\mychapterending

\mychapter{Об ослаблении вечных мук умерших некрещеными}
%http://www.molitvoslov.org/text236.htm 
 

\section{Мученику Уару}\begin{mymulticols}

\mysubtitle{Тропарь, глас 4-й:}

Воинством святых страстотерпец страждущие законне, зря онех, показал еси мужески крепость свою. И устремився на страсть волею, и умрети вожделе за Христа, Иже приял еси почесть победы твоего страдания, Уаре, моли спастися душам нашим.

\mysubtitle{Кондак, глас 4:}

Христу последуя, мучениче Уаре, Того испив чашу, и венцем мучения увязеся, и со Ангелы ликовствуеши, моли непрестанно за души наша

\mysubtitle{Молитва:}

О, святым мучениче Уаре досточудный, ревностию по Владыце Христу разжигаемь, Небеснаго Царя пред мучителем исповедал еси, и о Нем усердно пострадал еси, и ныне предстоиши Ему со ангелы, и в вышних ликуеши, и зриши ясно Святую Троицу, и светом Безначальнаго Сияния наслаждаешися, воспомяни и наших сродников томление, умерших в нечестии, приими наше прошение, и якоже Клеопатрин род неверный молитвами твоими от вечных мук свободил еси, тако воспомяни елицы противобожне погребенныя, умершия некрещенными, потщися испросити оным от вечныя тьмы избавление, да вси едиными усты и единем сердцем восхвалим Премилосерднаго Творца во веки веков. Аминь

\end{mymulticols}

\mychapterending

\mychapter{Об ослаблении вечных мук умерших без покаяния}
%http://www.molitvoslov.org/text235.htm 
 

\section{Преподобному Паисию Великому}\begin{mymulticols}

\mysubtitle{Тропарь, глас 2-й}

Божественною любовию от юности распалаемь, преподобне, вся красная, яже в мире, возненавиде, Христа единаго возлюбил еси, сего ради в пустыню вселися, идеже сподобися Божественнаго посещения, Егоже неудобь зрети и ангельскима очима, паде, поклонися. Великий же Дародатель, яко Человеколюбец рече к тебе: не ужасайся, возлюбленный Мой, дела твоя угодна Мне. Се даю тебе дар: о коем-либо грешнице помолишися, отпустятся ему греси. Ты же в чистоте сердца твоего возгореся, прием воду и прикоснуся Неприкосновенному, умы нозе Его и, воду пив, даром чудес обогатися, болящия исцеляти, бесы от человек отгоняти и грешники от муки молитвою своею избавляти. О преподобне отче Паисие, молю тя, да умолиши и о мне, якоже тебе Бог обетова, ибо от сих грешник первый аз есмь, да даст ми Господь время покаяния и простит мое согрешение, яко Благий и Человеколюбец, да со всеми и аз воспою Ему: аллилуия.

\mysubtitle{Кондак, глас 2-й:}

Житейских молв оставль, безмолвное житие возлюбил еси, Крестителю подобяся всеми образы, с нимже тя почитаем, отче отцев Паисие.

\mysubtitle{Молитва:}

Страстей победителя, душам помощника, о всех молебника, всем спасения ходатая и наставника, из глубины сердца воздыхая, усердно и пламенно молим тя, Паисие Преподобне! Внемли и помози нам, не отринь и не презри нас, но абие услыши в смирении сердца притекающих к тебе. Ты, преподобне, к спасению ближних прилежно стремился и многих грешников к свету спасения привел еси. Подвиги чрезмерные успокоением считал по себе, пречудне, и, любовью ко Господу всегда горя, явления Христа Спаса сподобился еси, и Ему, за людей Умершему, любовию подражая, и об отрекшихся от Христа молился еси. Услыши нас, Паисие прехвальне, ибо недостойны есмы молитися о даровании нам великой милости Господней, понеже грешны есмы, и уста оскверненныя и сердца отягощенныя имеем, и под бременем прегрешений страждем, и не достигает молитва наша до Господа. Сего ради помолися за нас мольбою твоею крепкою и богоприятною, святый Паисие, да избавлены будут скончавшиеся без покаяния сродники, ближние и знаемые наши от муки вечныя, и молитву твою во благоволении приимет Спас наш и милосердие Свое вместо добрых дел их даст им, свободит их, веруем, от страданий и вселит в селениях праведных, и нас в покаянии скончатися удостоит, да прославим вкупе Всесвятое и великолепое имя Отца и Сына и Святаго Духа, во веки веков. Аминь.

\end{mymulticols}

\mychapterending

\mychapter{КАНОН о самовольне живот свой скончавших}
%http://www.molitvoslov.org/text916.htm 
 

{\centering\myemph{ (для келейного чтения)\symbolfootnote[1]{Печатается по изданию Свято-Успенского Псково-Печерского монастыря, 1993.}}

}

\begin{mymulticols}

В 1925 году на Сергиевском Подворье в Париже "--- Рю де Крим, 93 "--- покончил с собой выстрелом один солдат. Он, несомненно, был нездоров умом…

У нас его отпевали…

Но, смущенный смертью его на нашей земле, я обратился через одно лицо к преподобному о. Нектарию Оптинскому с вопросом: что нам делать? Он ответил, чтобы я нашел еще двух человек, которые бы согласились в течение 40 дней читать за него заупокойный канон, а затем предать его в волю Божию.

Разумеется, канон читать нужно было дома, а не за богослужением в храме. Так мы трое тогда и сделали…

После ко мне обращались и другие по этому вопросу; а недавно просили даже отпевать самоубийцу женщину. Я отказывался,"--- согласно канонам "--- правило Тимофея Александрийского "--- и по внутреннему противлению души моей; но помня наставление о.Нектария, а также и указание митрополита Новгородского и Петроградского Григория, что можно молиться о самоубийцах в домашней молитве, но и тогда "--- с предварительною молитвою о помиловании нас молящихся, да не прогневаем Господа,"--- я обратился в каноник "--- «можно и во всякую субботу в Oктоихе найти заупокойные каноны»; но написанный там канон совершенно непригоден был для молитвы о самоубийцах, ибо там говорилось о скончавшихся «благочестно», «в вере», «в упокоении» и прочее,"--- все это не подходило к вольным, «не в сумасшествии», самоубийцам. Поэтому нужно было целиком почти переделать все тропари.

Может, кому-либо пригодится…

А меня да помилует Господь за дерзновение такое…

Многогрешный митрополит Вениамин (Федченков)

\mysubtitle{КАНОН, ГЛАС 2-й} {\centering \myemph{(читается келейно и с благословения духовника)}

}

\mysubtitle{ПЕСНЬ 1}

\irmos{Во глубине постла иногда фараонитское всевоинство преоруженная сила; воплощшееся же Слово всезлобный грех потребило есть: препрославленный Господь, славно бо прославися.}

\pripev{Упокой, Господи, душу усопшаго раба Твоего.}

Всемилостивый Господи, помилуй мя по милости Твоей, пети Тебе начинающа и молитися о убившемся дерзко, да не в суд ми сие будет, ни в осуждение, Судие Праведнейший.

\pripev{Упокой, Господи, душу усопшаго раба Твоего.}

Да не будет молитва моя во грех, яко кроткий Давид глаголет: ниже во отягощение души самовольней смерти предавшемуся; но мученик мольбами оставление прегрешений наших даруй нам.

\slava

Богатую милость присно всем источая, Всеблагий Христе Боже, и ныне рабу Твоему \myemph{ (имя)}, нами Тебе молящемуся, помилование даруй и с преставльшимися во благочестии причти, аще и нечестиво скончался в неразумии.

\inyne

Утвердивши колеблемый ум мой, Мати Божия, укрепи мя Божественными повеленьми Рождшагося из Твоего освященного чрева и упразднившего, Владычице, адово мрачное царство, со дерзновением неосужденно молитися об усопшем \myemph{ (имя)}. 

\mysubtitle{ПЕСНЬ 3}

\irmos{На камени мя веры утвердив, разширил еси уста моя на враги моя: возвесели бо ся дух мой, внегда пети: несть свят, якоже Бог наш, и несть праведен паче Тебе, Господи.}

\pripev{Упокой, Господи, душу усопшаго раба Твоего.}

Твоея красоты сиянием озаритися усопшему благоволи, в вере с почившими, богатый в милости; Ты бо еси Бог наш Единый безгрешный и многомилостивый.

\pripev{Упокой, Господи, душу усопшаго раба Твоего.}

На месте прохлаждения в вечнем покоище сподоби прияти и раба Твоего, малодушие жизнь свою пресекшаго: Ты бо еси Един Сильный и немощи наша на Ся восприявый, Господи Боже наш.

\slava

В чертозе небеснем водворитися, Владыко, приими раба Твоего \myemph{ (имя рек)}, да не услышит он праведнаго Твоего гласа: «не вем тя», аще и воистину повинен есть сему.

\inyne

Умерщвленных нас воздвигла еси, Богородице, жизни нетленней Жизнодавца рождши; и раба Твоего \myemph{ (имя)} возведи из ада преисподнейшаго, Сына Твоего о милости умоляющи. 

\mysubtitle{ПЕСНЬ 4}

\irmos{Пою Тя, слухом бо. Господа услышах и ужасохся: до мене бо идеши, мене ища заблуждшаго. Тем многое Твое снизхождение, еже на мя, прославляю, Многомилостиве.}

\pripev{Упокой, Господи, душу усопшаго раба Твоего.}

Во уповании на Тя, Боже, изнемогшего, и насилием вражиим одоленнаго, раба Твоего пощади, Спасе презирая, яко Благ, прегрешения его вольная и невольная.

\pripev{Упокой, Господи, душу усопшаго раба Твоего.}

Милосердия имый неизследимую бездну, Господи, бездну греховную раба Твоего покрый, Человеколюбче,"--- омыв Кровию Своею грех неведения и маловерия его.

\slava

Живыми господствуя и мертвыми владычествуя, помилуй и раба Твоего \myemph{ (имя)}, власть бо имаши всякую на небеси и на земли, Христе Боже, Отцу соседый и Духом вся очищаяй.

\inyne

Исцелила еси, Богородительнице, Евино поползновение окаянное; Зиждителя бо родила еси, исправити могущаго тварь погибающую: Того моли и сего раба \myemph{ (имя)} неключимаго, заповедь Божию преступившего, спасти по единой милости. 

\mysubtitle{ПЕСНЬ 5}

\irmos{Света подателю и веков Творче, Господи, во свете Твоих повелений настави нас: разве бо Тебе иного Бога не знаем.}

\pripev{Упокой, Господи, душу усопшаго раба Твоего.}

Умертвишася и во дно адово поползнувшагося возведи, Боже, преступление того заглаждая; на Кресте бо грехи наша пришел еси омыти честною Кровию Своею, Христе, Искупителю всяческих.

\pripev{Упокой, Господи, душу усопшаго раба Твоего.}

Спасти нас пришедый, Христе Боже, приими ныне и наша недостойная моления о рабе Твоем \myemph{ (имя)}, да муки вечныя избавиши, помилован будет с нами, Спасителю Преблагий.

\slava

Молим Тя, Слове Божий, аще и дерзновенно сие, прогнева бо Тя раб Твой \myemph{ (имя)}, обаче благодатию Твоею умоляем, не возгнушайся молитв наших и помилуй нас, вкупе вси бо повинны есмы пред Тобою, Всесвятый Владыко.

\inyne

Под кровом Твоим, Мати Божия, спасаются вси надеющийся на Тя: Твоим заступлением защити и нас, и сподоби со дерзновением неосужденно молити Сына Твоего о живот свой скончавших безнадежно \myemph{ (имя)}. 

\mysubtitle{ПЕСНЬ 6}

\irmos{В бездне греховней валяяся, неизследную милосердия Твоего призываю бездну: от тли, Боже, мя возведи.}

\pripev{Упокой, Господи, душу усопшаго раба Твоего.}

Из земли создавши человека, Боже, и в землю ныне возвратишася, пощади, Господи, в неразумии сокрушившаго скудельный сосуд свой: можеши бо, Владыко Творче, и паки возставити его от истления.

\pripev{Упокой, Господи, душу усопшаго раба Твоего.}

Неизреченный и непостижимый, Господи, щедротами несказаннаго Человеколюбия Твоего помилуй живот свой скончавшаго, презирая согрешения его, яже в ведении и неведении содеянная.

\slava

Озаритися милостию Твоею сподоби, Владыко, вольне оставльшаго \myemph{ (имя)} житие свое и к свету Твоему возведи, из тьмы несветимыя изымая, милосердия ради.

\inyne

Избавительнице, усердно призывающим Тя, Пречистая Владычице, Сына Твоего умоли скончавшагося \myemph{ (имя)} ущедрити: живота и смерти Владыку родила еси. 

\mysubtitle{ПЕСНЬ 7}

\irmos{Богопротивное веление беззаконнующаго мучителя высок пламень вознесло есть; Христос же простре благочестивым отроком росу духовную: Сый благословен и препрославлен.}

\pripev{Упокой, Господи, душу усопшаго раба Твоего.}

Во ад с душею, яко Бог, сошедый, и в рай разбойника с Собою введый, помяни и ныне раба Твоего \myemph{ (имя)}, преступлением многим Тя прогневавша, ныне же Тебе нами молящася бедне.

\pripev{Упокой, Господи, душу усопшаго раба Твоего.}

Един безгрешный, в мертвых волею вменился еси, да умерщвленных силою человекоубийцы воскресивши с Собою, от державы того исхити и раба Твоего, Всесильне.

\slava

Радость присносущную даровати нам Воскресением Своим изволивый, от мрачнаго горькаго царства изми раба Твоего \myemph{ (имя)}, да от муки вечныя свободившися, узрит себе одесную Тебе, Судие Многомилостивый.

\inyne

Сияние славы Отчей, плотию от Тебе явившегося умоли, Пресвятая, из адовы мглы свободити непотребнаго раба Твоего \myemph{ (имя)}, той бо яко Матерь, Тебе нами просит. 

\vspace{-\baselineskip}\mysubtitle{ПЕСНЬ 8}

\irmos{Пещь иногда огненная в Вавилоне действа разделяше, Божиим велением халдеи опаляющая, верныя же орошающая, поющия: Благословите, вся дела Господня, Господа.}

\pripev{Упокой, Господи, душу усопшаго раба Твоего.}

Смертное царство Своею смертию разоривый, Божественною же силою вход в вечный живот отверзый, не затвори благоутробия Твоего и рабу Твоему, аще и прегрешившему, обаче из ада ныне горце взывающа.

\pripev{Упокой, Господи, душу усопшаго раба Твоего.}

Красоты Твоея насладитися сподоби, студ греха безобразнаго очищая умершаго смертию поносной: Един бо еси, Владыко, греха свободен, крестною смертию за нас умертвился еси.

\slava

Исправити хотя падших в персть смертную, Сам смерти причастен явился еси, Безсмертный, от смертных удолий силою Твоею возведи раба Твоего \myemph{ (имя)}, лютою смертию поползнувшагося.

\inyne

Велия воистину есть Рождества Твоего тайна, Богомати: Бога воплотившагося родила еси; Его же моли, яко Сына Твоего, спасти бедне скончавшагося \myemph{ (имя)}: имаши бо дерзновение к Рождшемуся от Тебя Благоутробному Спасу мира. 

\vspace{-\baselineskip}\mysubtitle{ПЕСНЬ 9}

\irmos{Яже прежде солнца Светильника Бога возсиявшаго, плотски к нам пришедшаго, из боку девичу неизреченно воплотившая, Благословенная, Всечистая, Тя, Богородице, величаем.}

\pripev{Упокой, Господи, душу усопшаго раба Твоего.}

Имеяй власть живых и мертвых, пощади, Спасе, раба Твоего \myemph{ (имя)}, смертию горькою скончавшагося: Тебе бо от Отца Твоего дадеся живити и миловати, ихже хощеши, Сыне Божий.

\pripev{Упокой, Господи, душу усопшаго раба Твоего.}

Преславное множество святых Твоих, Христе, непрестанно молит Тя за грешныя рабы Твоя, да помилуеши ихже волиши милосердием Твоим и геенны вечныя избавиши.
%\end{mymulticols}

%\newpage

%\begin{mymulticols}

\slava

Благости воистину пучина неизследимая, благодатию Твоею помилуй прогневавшаго Тя раба Твоего \myemph{ (имя)}, и нас молящихся.

\inyne

Достойно Тя, Богородице, ублажающе, к Тебе, яко Матери нашей, притекаем надежно; сами себе и брата нашего \myemph{ (имя)}, малодушие скончавшагося, Тебе Заступнице и с Тобою Христу Богу предаем смиренно.

\end{mymulticols}

\mychapterending

\mychapter{Молитва за благодетелей, особенно руководивших к добродетели}\begin{mymulticols}
%http://www.molitvoslov.org/text197.htm 
 

Господи Иисусе Христе Боже наш! Ты заповедал еси нам воздавати всем должное, егда рекл еси: «Якоже хощете да творят вам человецы, и вы творите им такожде». Что убо сотворю, Господи, что воздам ныне, во мнозе благотворившему(-шей) мне, но уже отшедшему (-шей) от земли живых рабу Твоему (рабе Твоей)! У Тебе, Господи, у Тебе, богатый в милостех Боже наш, всякое воздаяние есть. Сего ради молю Тя, Подателю благих, приими от мене сие малое моление мое, вкупе же и благодарение за все благодеяния Твои, дарованныя мне в приснопоминаемом(-мой) усопшем(-шей) рабе Твоем (Твоей). Молю Тя, Боже, Спасителю мой призри с высоты Святаго жилища Твоего на молитву мою и помяни милостию Твоею отшедшаго(-ую) к Тебе раба Твоего (рабу Твою) (имя), толико благодеявшаго(-ую) ми; яви ему (ей) силу благости Твоей и неизреченное Твое милосердие в прощении и оставлении грехов и покрый его (ю) оправданием Твоим; подаждь ему, Господи, Царство Небесное и вечный покой, яко помогавшему(-шей) ми и ближним моим в нуждех и лишениях, в скорбех и болезнех и подававшему (-шей) ми благая от избытков даров Твоих. Молюся убо Тебе, Благость вечная, воздай ему (ей) за благотворения его (ея) сторицею "--- за земныя блага небесными, за временныя вечными, и подаждь ему (ей) вечный покой; и яко научавшему(-шей) мя делам милосердия, кротости, смиренномудрию, терпению, воздержанию, вселявшему (-шей) в мя дух мира и любви христианской, подававшему (-шей) ми благие советы и благие образы христианскаго жития, даруй, Премилосердый Господи, достойно узрети невечерний свет Твоего Царствия и насыти его (ю) от небесныя трапезы Твоея, яко во имя Твое препитавшаго(-ую) жаждущую и алчущую мою душу. Помяни, Господи, и всех отшедших к Тебе, в вере и добродетельном житии подвизавшихся во славу Пресвятаго Имени Твоего, добро творивших сирым, нищим и убогим; упокой в лоне Авраама, Исаака, Иакова, посещавших недугующих, страждущих и в темнице заключенных; приими в чертог Твой принимавших в дому своем странна и пришельца. Помяни, Господи, во Царствии Твоем усопших рабов Твоих, выну предстоявших и молившихся Тебе теплою молитвою о нашем душевном спасении и благоденствии, о сохранении и благостоянии святых Твоих церквей, о мире и спасении всего мира; упокой со святыми Твоими рабов Твоих, молившихся о обращении к Тебе заблудших и отпадших от веры и о спасении погибающих во грехах. Молю Тя, Христе Боже наш, молю Тя, яко Судию, имущаго судити мирови, воздай им по человеколюбию Твоему, и сподоби их услышати вожделенный оный глас: «Приидите, благословеннии Отца Моего, и наследуйте уготованное вам Царствие от сложения мира», и да внидут в онь, хваля и прославляя Твое великолепое имя в безконечные веки. Аминь. 

\end{mymulticols}

\mychapterending

\mychapter{Молитва за скончавшихся вне своего Отечества, за безродных и убогих}\begin{mymulticols}
%http://www.molitvoslov.org/text217.htm 
 

Господи Иисусе Христе Боже наш, сирых Защитниче, странных Наставниче, плавающих Кормчий, обуреваемых пристанище, с путешествующими путешествующий, утешение плачущих, нищих кормителю! Молим Тя, Многомилостиве Господи, помяни во Царствии Твоем души усопших раб Твоих (\myemph{ имена}), в вере преставльшихся к Тебе, но отшедших в земли чуждей, без матерняго Святыя Твоея Церкви напутствия и благословения. Приими, Владыко, под Твое благоутробие рабов преставльшихся тамо в скорби и печали о своих присных и ближних, имже никтоже бе утешаяй в страшный час смерти, идеже никтоже бе возносяй Тебе теплое моление о благотишном прохождении их от сея привременныя жизни в мир горний. Темже убо молим Тя, Преблагий, молим Тя, вездесущий и присносущный Господи, воззри на сих рабов Твоих благосердием Твоим, прости им всякое согрешение вольное и невольное, словом и делом содеянное и вся, яже в ведении и неведении. Умоляем Тя молитвами и предстательством Матере Твоея, Пресвятая Владычицы нашея Богородицы, и всех святых, не погуби отшедших к Тебе со беззаконием их, но сотвори им велию милость Твою, яви им Твое человеколюбие и щедроты Твоя, приведи их к тихому Твоему пристанищу и водвори их в Дому Твоем. 

Еще молимся о упокоении душ усопших раб Твоих (души усопшаго (-шия) раба Твоего (рабы Твоея) (\myemph{ имена}), отшедших к Тебе безродными и никогоже от присных своих о себе молитвенников имущими, разве Святыя Твоея Церкви; сего убо ради молим Тя, Господи, приими и о сих наше моление: прости им вся согрешения и помилуй их по велицей Твоей милости; всели их во Святом жилище Твоем, да обрящут в Тебе покой, яко в Спасителе и Искупителе рода человескаго. 

Господи! Господи! Аще усопший раби Твои, иже быша в нищете и убожестве, за чашу студеныя воды, за укрух хлеба, за едину лепту, поданныя им во имя Твое, приносили быша Тебе моления о нас, прошения и благодарения; аше сирые благословляли быша благотворящую им руку, моляся о ниспослании нам благ Твоих, того убо ради молим Тя ныне, Христе Боже наш, и мы о них: помяни убо во Царствии Твоем сих, яко меньших, по реченному Тобою, братий Твоих; прости им всякое прегрешение вольное и невольное, покрый их Твоим благоутробием и всели их во Святый чертог Твой; и яко получавших зде на земли пропитание чрез призывание святаго Твоего имени, препитай тамо на небеси от тука дома Твоего и да утешатся, и да насладятся нескончаемою радостию вси истинно плакавшие к Тебе; водвори их со святыми Твоими, да вси купно у Тебе, Судие живых и мертвых, составят сонмы избранных Твоих, и тако да славят, хвалят и превозносят безпредельное Твое милосердие в безпредельные веки. Аминь. 

\end{mymulticols}

\mychapterending

\mychapter{Молитва за всякого усопшего}\begin{mymulticols}
%http://www.molitvoslov.org/text196.htm 
 

Помяни, Господи Боже наш, в вере и надежди живота вечнаго преставльшагося раба Твоего, брата нашего (\myemph{ имя}), яко Благ и Человеколюбец, отпущаяй грехи и потребляяй неправды, ослаби, остави и прости вся вольная его согрешения и невольная, избави его вечныя муки и огня геенскаго, и даруй ему причастие и наслаждение вечных Твоих благих, уготованных любящым Тя: аще бо и согреши, но не отступи от Тебе, и несумненно во Отца и Сына и Святаго Духа, Бога Тя в Троице славимаго, верова, и Единицу в Троицу и Троицу в Единстве православно даже до последняго своего издыхания исповеда. Темже милостив тому буди, и веру яже в Тя вместо дел вмени, и со святыми Твоими яко Щедр упокой: несть бо человека, иже поживет и не согрешит. Но Ты Един еси кроме всякаго греха, и правда Твоя правда во веки, и Ты еси Един Бог милостей и щедрот, и человеколюбия, и Тебе славу возсылаем, Отцу и Сыну и Святому Духу, ныне и присно и во веки веков. Аминь. 

\end{mymulticols}

\mychapterending

\mychapterz{Молитва за обидевших и ненавидевших нас}{\LargeМолитва за обидевших и ненавидевших нас}
%http://www.molitvoslov.org/text200.htm 

\vspace{-\baselineskip}

\begin{mymulticols}

Владыко Человеколюбче Господи, Иисусе Христе, Сыне Божий! Ты, по неизреченной любви Своей к нам, грешным и недостойным рабам Твоим, сияеши солнце Свое на злыя и благая, дождиши на праведныя и неправедныя; Ты, Преблагий, заповедуеши нам любити врагов наших, добро творити ненавидевшим и обидевшим нас, благославляти клянущих нас и молитися за творящих нам напасть и изгонящих нас. Ты, Спасителю наш, вися на крестном древе, и Сам прощал еси врагов Своих, хульно ругавшихся Тебе, и молился за мучителей Твоих; Ты дал еси образ нам, да последуем стопам Твоим. Ты, о дражайший Искупителю наш, научивый нас прощати врагов, повелел еси вкупе и молитися за них; молюся убо Тебе, Иисусе Прещедрый, Сыне и Агнче Божий, вземляй грехи мира, прости отшедшаго (-ую) к Тебе раба твоего (рабу Твою) (\myemph{ имя}) и приими его (ю) не яко врага моего, соделавшаго ми злая, но яко согрешившаго (-ую) пред Тобою, Молю Тя, безпредельный в милостех, Господи Боже наш, приими с миром, пресельшагося (-уюся) к Тебе от мира сего без примирения со мною; спаси и помилуй его (ю), Боже, Твоею великою и богатою милостию. Господи, Господи! Да не яростию Твоею, ниже гневом Твоим накажеши раба Твоего (рабу Твою), творившаго ми напасть, оскорбление, поношение и злохуление; молю Тя, не помяни сих прегрешений его (ея), но отпусти и прости ему (ей) вся сия по человеколюбию Твоему, и помилуй по велицей Твоей милости. Молю Тя, о Преблагий и Прещедрый Иисусе, яко уз адовых Решителя, смерти Победителя, грешных Спасителя, разреши рабу Твоему (рабе Твоей) сии согрешения, имиже, яко пленицами ада, связася усопший (-шая). Ты бо, Господи, рекл еси: «аще не отпущаете человеком согрешения их, ни Отец ваш Небесный отпустит вам согрешений ваших»; о, да не будет сего! Со умилением и сокрушением сердца умоляю Тя, Спасителю Премилостивый, разреши ему (ей) сии узы злых наваждений и козней диавольских, не погуби усопшаго (-ую) гневом Твоим, но отверзи ему (ей), Жизнодавче, двери милосердия Твоего, да внидет во святый Твой град, хваля всесвятое и великолепое Имя Твое и воспевая неизреченную любовь Духа Твоего Святаго к погибающим грешникам. И якоже Ты, Благость вечная, помянул еси на кресте благоразумнаго разбойника, с Тобою распята, путесотворив ему вход в рай, сице, молю Тя, Всещедрый, помяни во Царствии Твоем и отшедшаго (-ую) к Тебе раба Твоего (рабу Твою) (\myemph{ имя}) не затвори, но отверзи и ему (ей) двери милосердия Твоего, Твое бо есть, еже миловати и спасти ны, Боже наш, и Тебе славу возсылаем со Безначальным Твоим Отцем, Пресвятым и Благим и Животворящим Твоим Духом, ныне и присно и во веки веков. Аминь. 

\end{mymulticols}

\mychapterending

\mychapter{Молитва за умирающего}\begin{mymulticols}
%http://www.molitvoslov.org/text195.htm 
 

Господи, Иисусе Христе Сыне Божий, заступи, спаси, помилуй и сохрани Боже, Твоею благодатию душу раба Твоего (\myemph{ имя}), и грехи юности и неведения его не помяни, и даруй ему кончину христианску, непостыдну и мирну, и да не узрит душа его мрачного взора лукавых демонов, да приимут его Ангели Твои светлии и пресветлии, и на Страшном Суде Твоем милостив ему буди, ибо Твое есть единого Господа, еже миловати и спасати нас. 

\end{mymulticols}

\mychapterending

\mychapter{Моление о упокоении православных воинов, за веру и Отечество на брани убиенных}\begin{mymulticols}
%http://www.molitvoslov.org/text220.htm 
 

Непобедимый, непостижимый и крепкий во бранех Господи Боже наш! Ты, по неисповедимым судьбам Твоим, овому посылаеши Ангела смерти под кровом его, овому на селе, овому на мори, овомуже на поле брани от оружий бранных, изрыгающих страшныя и смертоносныя силы, разрушающия телеса, расторгающия члены и сокрушающия кости ратующих; веруем, яко по Твоему, Господи, премудрому смотрению, такову приемлют смерть защитники веры и Отечества. 

Молим Тя, Преблагий Господи, помяни во Царствии Твоем православных воинов, на брани убиенных, и приими их в небесный чертог Твой, яко мучеников изъязвленных, обагренных своею кровию, яко пострадавших за Святую Церковь Твою и за Отечество, еже благословил еси, яко достояние Твое. Молим Тя, приими убо отшедших к Тебе воинов в сонмы воев Небесных Сил, приими их милостию Твоею, яко павших во брани за независимость земли Русския от ига неверных, яко защищавших от врагов веру православную, защищавших Отечество в тяжкие годины от иноплеменных полчищ; помяни, Господи, и всех, добрым подвигом подвизавшихся за древнехранимое Апостольское Православие, за освященную и в язык свят избранную Тобою землю Русскую, в нюже враги Креста и Православия приношаху и огнь, и меч. Приими с миром души раб Твоих (\myemph{ имена}), воинствовавших за благоденствие наше, за мир и покой наш, и подаждь им вечное упокоение, яко спасавшим грады и веси и ограждавшим собою Отечество, и помилуй павших на брани православных воинов Твоим милосердием, прости им вся согрешения, в житии сем содеянная словом, делом, ведением и неведением. Призри благосердием Твоим, о Премилосердый Господи, на раны их, мучения, стенания и страдания, и вмени им вся сия в подвиг добрый и Тебе благоугодный; приими их милостию Твоею, зане лютыя скорби и тяготу зде прияша, в нуждех, тесноте, в трудех и бдениих быша, глад и жажду, изнурение и изнеможение претерпеша, вменяеми быша яко овцы заколения. Молим Тя, Господи, да будут раны их врачеством и елеем, возлиянным на греховныя язвы их. Призри с небесе, Боже, и виждь слезы сирых, лишившихся отцев своих, и приими умиленныя о них мольбы сынов и дщерей их; услыши молитвенныя воздыхания отцев и матерей, лишившихся чад своих; услыши, благоутробне Господи, неутешных вдовиц, лишившихся супругов своих; братий и сестер, плачущих о своих присных,"--- и помяни мужей, убиенных в крепости сил и во цвете лет, старцев, в силе духа и мужества; воззри на сердечныя скорби наша, виждь сетование наше и умилосердися, Преблагий, к молящимся Тебе, Господи! Ты отъял еси от нас присных наших, но не лиши нас Твоея милости: услыши молитву нашу и приими милостивно отшедших к Тебе приснопоминаемых нами рабов Твоих (\myemph{ имена}); воззови их в чертог Твой, яко доблих воинов, положивших живот свой за веру и Отечество на полях сражений; приими их в сонмы избранных Твоих, яко послуживших Тебе верою и правдою, и упокой их во Царствии Твоем, яко мучеников, отшедших к Тебе израненными, изъязвленными и в страшных мучениях предававшими дух свой; всели во святый Твой град всех приснопоминаемых нами рабов Твоих (\myemph{ имена}), яко воинов доблих, мужественно подвизавшихся в страшных приснопамятных нам бранех; облецы их тамо в виссон светел и чист, яко зде убеливших ризы своя в крови своей, и венцев мученических сподоби; сотвори их купно участниками в торжестве и славе победителей, ратоборствовавших под знаменем Креста Твоего с миром, плотию и диаволом; водвори их в сонме славных страстотерпцев, добропобедных мучеников, праведных и всех святых Твоих. Аминь. 

\end{mymulticols}

\mychapterending

\mychapter{Молитва за наставников и воспитателей}\begin{mymulticols}
%http://www.molitvoslov.org/text199.htm 
 

Боже разума, Боже чудес, живый, приснотекущий Источниче премудрости! Молю Тя, помяни во Царствии Твоем души усопших раб Твоих, наставников моих (\myemph{ имена}), яко просвещавших мой ум и вселявших в сердце мое дух премудрости и разума, дух совета и крепости, дух ведения и правды, истины и добродетели. Молю Тя, Владыко, милостив буди о гресех их; ослаби, прости и отпусти им вся согрешения, елика сотворили в житии сем, словом, делом, помышлением, ведением и неведением. Воздаждь им, света Подателю, милостию Твоею, щедротами Твоими за вся попечения их о нашем духовном преуспеянии; яви им милость благоутробия Твоего за благая их поучения, за благопотребные их советы и вразумления, яже посеяваху и насаждаху во умах и сердцах наших, яко благия семена; воздаждь им за вся благая небесными и вечными Твоими благами, да узрят в Тебе, Боже, лицем к лицу вечную Истину и да озаряются и наслаждаются трисиятельным Светом Ея в безконечные веки во Царствии Твоем. Аминь. 

\end{mymulticols}

\mychapterending

\mychapter{Молитва четвертая за всех в вере скончавшихся православных христиан}\begin{mymulticols}
%http://www.molitvoslov.org/text194.htm 
 

Господи, Господи! коль неизбежен, коль страшен нелицеприятный и неизменимо-вечный Суд Твой! В великом и притрепетном смущении дух наш мятется, сердце содрогается и истаевает, яко воск, от единаго возвещения Твоего непреложнаго глагола о последнем Суде; немотствует язык от единаго ожидания гласа последней трубы Архангела, имущей возбудити от мертвых и воззвати от живых на Страшный Суд Твой от востока и запада, от севера и юга. О, воистину страшен день той, в оньже приидеши, Боже, на землю во славе со Ангелами и всеми святыми! Горнии и дольнии все предстанут тогда со страхом и трепетом и полчища преисподних вострепещут пред Тобою, Судие мира сего; небо и землю призовеши на Суд, судити вселенней в правде и истине! Страшен день той, в оньже престоли поставятся, книги разгнутся, и явная и тайная наша деяния, словеса и помышления обличатся пред Тобою, пред Ангелами и человеками, вкупе же и осудятся: ничтоже бо, Господи, у Тебе сокровенно, еже тогда не открыется, и тайно, еже не уведано будет. Кто возможет, Судие мира сего, устояти тогда пред лицем Твоим и целым миром горним и дольним! Аще праведник едва спасется, нечестивый и грешный како постоит? И кто убо заступит за ны, аще не Твое милосердие, Господи? Камо убежим от праведнаго Суда Твоего? О, пощади, пощади тогда нас и отшедщаго (-ую) к Тебе приснопоминаемаго (-ую) нами раба Твоего (рабу Твою) (\myemph{ имя}) и не осуди его (ю) тогда на вечное мучение по грехом его (ея). Благоволи убо, Премилосердый Господи, прияти в небесный и мысленный Твой жертвенник наши жертвы за него (ню) "--- молитвы наши и милостыни, яко кадило благовонное, вкупе с молитвами и безкровным священнодействием пастырей Святая Твоея Церкви, да не постыдится пред Тобою, пред Ангелами и целым миром усопший (-шая) раб Твой (раба Твоя). Услыши ны, Боже наш, и не отрини до конца; виждь коленопреклоненных, повергшихся долу пред Тобою, Господи; услыши просящих Твоея милости и Царства Небеснаго усопшему (-шей) рабу Твоему (рабе Твоей). Молим Тя, о Премилосердый и Прещедрый Иисусе, егда приидеши со Ангелы Святыми Твоими судити мирови, пощади, пощади тогда. Спасе, создание Твое: образ бо есть неизреченныя Твоея славы. С сокрушенным сердцем молим, просим Тя, Спасителю и Искупителю наш, не осуди раба Твоего (рабу Твою) (\myemph{ имя}) праведным судом Твоим на вечное мучение, аще достоин (-на) есть всякаго осуждения и муки; не отлучи его (ю) от стада избранных Твоих, но, по неизреченному милосердию Твоему, долготерпению и любви к нам грешным и недостойным рабам Твоим, ради Твоих дражайших заслуг сподоби его (ю), Владыко, Царствия Твоего, еже уготовал еси любящим Тя от сложения мира, и да внидет в место упокоения, хваля всесвятое и великолепое Имя Твое, славя неизреченное милосердие Твое, величая человеколюбие и щедроты Твоя, яко Ты еси Бог милости и щедрот, и Тебе славу возсылаем, со безначальным Твоим Отцем, Пресвятым и Благим и Животворящим Твоим Духом, ныне и присно и во веки веков. Аминь. 

\end{mymulticols}

\mychapterending

\mychapter{Молитва о упокоении скончавшихся после тяжкой и продолжительной болезни}\begin{mymulticols}
%http://www.molitvoslov.org/text219.htm 
 

Господи, Господи! Праведен еси Ты, и суд Твой праведен: Ты, в предвечной Своей Премудрости, положил еси предел жизни нашея, егоже никтоже прейдет. Премудры Твои законы, неизследимы путие Твои! Ты повелеваеши ангелу смерти изъяти от тела душу у младенца и старца, у мужа и юноши, у здраваго и больнаго по несказанным и недоведомым нам судьбам Твоим; но веруем, яко на сие есть Твоя святая воля, зане, по суду правды Твоея, Ты, Преблагий Господи, яко премудрый и всемогущий и всеведущий Врач душ и телес наших, болезни и недуги, беды и злоключения посылаеши человеку, яко врачевство духовное. Ты поражаеши его и исцеляеши, умерщвляеши в нем мертвенное и оживотворяеши безсмертное, и, яко чадолюбивый Отец, наказуеши, егоже приемлеши: молим Тя, Человеколюбче Господи, приими преставльшагося к Тебе раба Твоего (рабу Твою) (\myemph{ имя}), егоже (юже) взыскал еси человеколюбием Своим, наказавый тяжкою болезнию телесною, во еже спасти душу от болезни смертныя; и аще вся сия приял(-ла) есть от Тебе со смиренномудрием, терпением и любовию к Тебе, яко вседетельному Врачу душ и телес наших, яви убо ему (ей) днесь богатую милость Твою, яко претерпевшему(-ей) вся сия грех своих ради. Вмени ему (ей), Господи, сию привременную тяжкую болезнь в некое наказание за прегрешения, содеянныя в сей юдоли плача, и исцели душу его (ея) от недугов греховных. Помилуй, Господи, помилуй взысканнаго(-ую) Тобою, и, наказаннаго(-ую) временно, молю Тя, не накажи лишением вечных Твоих небесных благ, но сподоби его (ю) наслаждатися ими во Царствии Твоем. Аще же усопший(-шая) раб Твой (раба Твоя), не разсуждая в себе, чесо ради бысть сие прикосновение целительныя и промыслительныя Десницы Твоея, стропотно что рече в себе, или, по неразумию своему, возропта в сердце своем, яко бремя сие возмни себе неудобоносимым, или, по немощи естества своего, стужаше долгою болезнию и огорчися злоключением, молим Тя, Долготерпеливе и Многомилостиве Господи, прости ему (ей) сия согрешения по безпредельной Твоей милости и безприкладному Твоему милосердию к нам грешным и недостойным рабам Твоим, прости ради любви Твоей к роду человеческому; аще ли же беззакония его (ея) превзыдоша главу его (ея), болезнь же и недуги не подвигоша его (ю) к полному и искреннему покаянию, умоляем Тя, Начальниче жизни нашея, умоляем Твоими искупительными заслугами, помилуй и спаси, Спасе, раба Твоего (рабу Твою) от вечныя смерти. Господи Боже, Спасителю наш! Ты, по вере в Тя, даровал еси прощение и отпущение грехов, даровав помилование и исцеление тридесятиосмилетнему разслабленному, егда рекл еси: «Отпущаются тебе греси твои»; с сею же верою и упованием на Твою благость, прибегаем к Твоему, о Иисусе Прещедрый, неизреченному милосердию и во умилении сердца нашего молим Тя, Господи: рцы убо и днесь тако слово помилования, слово отпущения грехов усопшему(-шей), приснопоминаемому(-мой) нами рабу Твоему (рабе Твоей) (\myemph{ имя}), да исцелеет духовне, и да вселится в месте светле, в месте покойне, идеже несть болезнь, ни печаль, ни воздыхание, и да пременятся тамо болезни и недуги его (ея), слезы страданий и скорбей в источник радости о Дусе Святе. Аминь. 

\end{mymulticols}

\mychapterending

\mychapterz{Молитва за пастырей Православной Церкви}{\LargeМолитва за пастырей Православной Церкви}
%http://www.molitvoslov.org/text198.htm 

\vspace{-\baselineskip}
 
\begin{mymulticols}

Господи Иисусе Христе Боже наш, давый нам овы апостолы, овы же пророки, овы же благовестники, овы же пастыри и учители к совершению Святых, в дело служения, в созидение Тела Твоего! Молим Тя, Владыко Многомилостивый, услыши молитвенный глас наш и помяни во Царствии Твоем преставльшихся из дольняго мира сего в мир горний рабов Твоих: православных Патриархов, митрополитов (\myemph{ имена}), архиепископов и епископов (\myemph{ имена}), архимандритов, игуменов, иеромонахов, иереев (\myemph{ имена}), иеродиаконов и диаконов (\myemph{ имена}); водвори их, Господи, в небесные кровы вечнаго покоя. Яко послуживших ко славе Пресвятаго имене Твоего и во славу Святая Соборныя и Апостольския Церкви Твоея, к назиданию и спасению чад Ея; помяни и всех в монашеском чине и причте церковном Тебе послуживших. Молим Тя, Господи Иисусе Христе Боже наш, молим яко Пастыреначальника, приими милостию Твоею души усопших раб Твоих, яко избранных слуг и строителей Тайн Твоих, предстоявших страшному престолу Твоему, совершавших безкровное священнодействие и подававших нам Святые и Животворящие Дары Твои "--- Честное Твое Тело и Честную Кровь Твою, во оставление грехов и в жизнь вечную; того ради молим Тя, Боже наш, соделай и их причастниками благодатных даров Духа Твоего Святаго во Царствии Твоем. Господи! Ты дал еси им власть во имя Твое прощати и отпущати нам грехи, вязати и решити, прилежно молим Тя: прости, отпусти и им согрешения их и не помяни грехопадений их; спаси и помилуй их по велицей Твоей милости; призри на них благосердием Твоим и приими их, Христе, под Твое благоутробие, яко учивших ны блюсти вся, елика заповедал еси нам; со святыми упокой в дому Твоем, яко домостроителей, смиренно подвизавшихся в насаждении и укреплении Православныя Веры и благочестия в сердцах наших; сподоби их неосужденно предстати пред Тобою на судищи Твоем, яко пред Пастыреначальником, и всели их в небесный Твой вертоград, яко насаждения масличная. Приведи их, Светодавче Господи, к невечернему свету Царствия Твоего, яко просвещавших ны светом Евангельским, всели их во Святый Твой град, яко вселявших в нас дух премудрости и разума, дух совета и крепости, дух ведения и благочестия, дух страха Твоего,"--- и даждь им наслаждатися от тука дома Твоего, яко питавших и наслаждавших души наша от божественныя Трапезы Твоея. Молим Тя, Царю, молитвами апостолов и пророков, святителей и учителей и всех святых, прослави их, яко служителей словесе во Святой Церкви Твоей, воинствующей на земли, прослави в Церкви, торжествующей на Небеси, и сопричти их к лику святых Твоих апостолов, святителей и учителей Церкви, послуживших и благоугодивших Тебе словом и делом. Аминь. 

\end{mymulticols}

\mychapterending

\mychapter{Молитва о скончавшихся в состоянии душевного заболевания}
%http://www.molitvoslov.org/text218.htm 
 
\begin{mymulticols}

Дивна дела Твоя, Господи, и величию Твоего разума несть конца! Коль многообразными и неизследимыми стезями преводиши Ты люди Своя в веце сем от дольняго мира к горнему! Разоряеши бо премудрость премудрых, отвергаеши разум разумных, высокия и гордыя расточаеши и смиряеши, сильныя и доброродныя низлагаеши и посрамляеши; буия же и смиренныя, уничиженныя, немощныя и худородныя избираеши, возвышаеши, укрепляеши и прославляеши, открывая премудрые и благие советы всепромыслительного смотрения Твоего младенцам и утаевая их от тех, иже о себе мнятся быти совершени умы своими. Кто доволен есть уразумети вся судьбы Твоя, Боже, имиже, по богатству благости Своея и долготерпения, вразумляеши, научаеши и спасаеши нас грешных и недостойных рабов Твоих? Овому прилагаеши смышление и даруеши здравое разсуждение, оваго же обуяеши, у оваго исторгаеши сердце каменно и влагаеши плотяно; овому ниспосылаеши крепость и бодрость телесну, вкупе же и мужество душевно; овому же болети душею и телом попускаеши; оваго радуеши, овому же печаловатися хощеши; оваго живиши, оваго же мертвиши. О, великий и неисповедимый в тайнах Своего смотрения о нас, Господи Боже наш! Вем и от всего сердца исповедую пред Тобою, яко и со усопшим (-шею) приснопоминаемым (-мою) мною рабом Твоим (рабою Твоею) (имя), не без премудраго и благаго изволения и смотрения Твоего бысть сокрушительный недуг душевный, егоже истинная вина и вся сила Тебе Единому точию ведома суть: темже благоговейно повергаяся пред Твоими, Господи, неизреченными, но выну спасительными судьбами, со умилением вопию Ти: приклони ухо Твое к молению моему и человеколюбие вонми ми, и скоро услыши молитву мою о сем рабе Твоем (Твоей). Недоумеваю бо и зело скорблю о необычайных судьбах его (ея), зане, не якоже инии человецы смысленнии, самосознательни и тверди в памяти, отьиде к Тебе от нас, но яко лишен (-на) сый ума, извращен в смысле, болезненно опечалован (-на) сердцем, не имеяй самообладательныя силы своего изволения, неистово и свирепо являющися в деяниях своих. Господи! Господи! Что сие есть неудоборазумеваемое некое в душевней жизни человечестей явление? Како созданный по образу и по подобию Твоему впадает в крайнее неразумие, безумствует, неистовствует и содевает неподобная? Что убо сие есть с дыханием жизни Твоея, Боже Вседержителю? Воистинну, по Твоим, Господи, неизследимым путям и неизреченным советам вся сия премудро и благим намерением Тобою устрояются: наводиши бо, яко привременное наказание, и попускаеши толикая нестроения и превращения в чудном составе существа нашего, да сокровенными у Тебе тако путями устроиши, возставиши и исправиши вся стропотная и развращенная, яже неразсудным, злым и лукавым обычаем учинена суть, и тако видимая злая у человек, у Тебе, Всеблагий Владыко, приснотекущий Источниче Премудрости и разума, у Тебе обращаются во благая. Того ради смиренно припадая к неистощимым щедротам Твоим, Человеколюбче Спасе наш, молю Тя, Подателю благих, приими усопшаго (-ую) раба Твоего (рабу Твою) (имя), яко тако люте болезновавшаго (-ую) и страдавшаго (-ую) в житии сем, приими милостию Твоею, да не пробавиши, Блаже, и тамо, во стране отшедших от нас братий наших, стужения и страдания его (ея), яже суть безмерно лютейша неже зде земная; но отврати от него (нея) прещение гнева Твоего и воззри на него (на ню) пресветлым оком милосердия Твоего; но не предаждь люте погибнути во веки, зане и Сам, Премилосердый Спасе, не хощеши смерти грешника, не хощеши, да кто погибнет, Сладчайший Иисусе Христе, Сыне Божий, исцеляяй неисцельно-недугующих телом, прокаженных очищаяй и бесных творяй смыслящими! Молю Тя, Господи, исцели душу усопшаго (-шия) раба Твоего (рабы Твоея) (\myemph{ имя}), очисти его (ю) от всех вольных и невольных, ведомых и неведомых прегрешений и грехопадений его (ея), Сам, Спасе, очисти, якоже не возмогшаго разумно и в здравой памяти исповедатися Тебе в них, того ради умилосердися о нем (о ней), Господи, и возьми от него (от нея) все тяжкое бремя грехов. Вземляй грехи мира! Покрый и оправдай его (ю) Своею благодатию, и, имиже веси судьбами. Спаси раба (-бу) Твоего (-ю) от избранных Твоих, но сподоби общником (-цею) быти и небесных ликостояний, яко убогаго (-ую) раба Твоего (рабу Твою), иже (яже) зде на земли отлученна (-у) и отчужденна (-у) себе имяше от сродников своих, другов и ближних, и не лиши его Небеснаго Твоего Царствия, яко лишенна (-ну) зде на земли видети благая Твоя и оставлнна (-ну) влачити дни свои исполнении горькоплачевных деяний: того ради, молю Тя, Господи, упокой его (ю) со святыми Твоими в Дому Отца Твоего, с Нимже Тебе, Избавителю и Жизнодавцу нашему, купно со Пресвятым Твоим Духом славу возсылаем всегда, ныне и присно и во веки веков. Аминь

\end{mymulticols}

\mychapterending

\mychapter{Молитва по исходе души из тела}\begin{mymulticols}
%http://www.molitvoslov.org/text185.htm 
 

Боже Духов и всякия плоти! Ты твориши ангелы Своя духи, и слуги Своя пламень огненный. Пред Тобою трепещут Херувимы и Серафимы и тьмы тем и тысяща тысящ со страхом и трепетом выну предстоят Престолу Твоему. Ты за хотящих улучити спасение посылаеши на служение святых Твоих Ангелов; Ты и нам грешным коемуждо даеши святаго Ангела Твоего, яко пестуна, еже хранили ны во всех путех наших от всякаго зла и таинственно наставляти и вразумляти ны даже до последняго издыхания нашего. Господи! Ты повеле еси изъяти душу от приснопоминаемаго(-мой) нами раба Твоего (рабы Твоея) (\myemph{ имя}), воля Твоя "--- воля святая; молим Тя, Жизнодавче Господи, не отними точию ныне от души его (ея) сего пестуна и хранителя ея, не остави ю едину, яко в путь шествующую; повели ему, яко хранителю, не удалятися помощию в сем страшном прохождении ея в мир горний невидимый; молим Тя, да будет убо ей заступником и защитником от злаго сопротивника в прохождении мытарств, дондеже преведет ю к Тебе, яко к Судии неба и земли. О, страшно прохождение сие для души, грядущей на суд Твой нелицеприятный, и имущей в прохождении сем истязатися духами злобы поднебесными! Темже убо молим Тя, Преблагий Господи, благоволи и еще послали святых Твоих Ангелов душе преставльшагося(-шейся) к Тебе раба Твоего (рабы Твоея) (\myemph{ имя}), да защитят, оградят и сохранят ю от нападения и истязания страшных и злых оных духов, яко истязателей и мытарей воздушных, служителей князя тьмы; молим Тя, свободи ю сего злаго обстояния, да не срящет ю злых демонов полчище; сподоби ю безбоязненно, благотишно и невозбранно преяти от земли страшный сей путь со Ангелы Твоими, да вознесут ю поклонитися Престолу Твоему и да приведут ю к свету Твоего милосердия.

\end{mymulticols}

\mychapterending

\mychapter{Моление о том, чтобы Бог даровал нам усердие к молитве за усопших и принял бы её}\begin{mymulticols}
%http://www.molitvoslov.org/text184.htm 
 

Многомилостиве и Всемилостиве Господи, Владыко и Судие живых и мертвых, дивный и неисповедимый в судьбах Своих Боже наш! Ты зриши печали, скорби и болезни, обдержащия ны зело, и воздыхания наша слышиши, и гласу нашему внемлеши. Како убо во многих скорбех наших обрящем утешение и покой душам нашим, аще не Твоим благопоспешеством, Всесильне Спасе наш? Ты бо еси Бог наш, и разве Тебе иного не знаем. Приими, Человеколюбче Господи, молитву нашу, юже приносим Тебе от всего сердца нашего и от всея души нашея. Веси, Господи, яко не довольны есмы достойную принести Тебе молитву; но Сам, Всеблагая, научи нас, како подобает молитися Ти. Согрей сердца наша теплотою Духа Твоего Святаго, да не тща будет молитва наша за души от нас преставльшихся раб Твоих, и даждь убо нам со дерзновением и неосужденно смети призывати Тебе, Небеснаго Бога Отца, яко Ты еси упование наше и прибежище; припадаем пред Тобою, в Троице покланяемый и прославляемый Боже наш, и Тебе единому приносим молитву нашу в скорбех и болезнех наших, в радости и печали, яко ты еси утешение наше, радование наше, сокровище наше и жизни нашея Податель: от Тебе бо исходит всяко деяние благо и всяк дар совершен. Сего ради дерзаем взывала к Тебе нашими бренными устнами: услыши нас, Боже, услыши нас грешных, молящихся Тебе; приклони ухо Твое к молению нашему и вопль наш к Тебе да приидет. Приими, Преблагий Господи, молитву нашу за души усопших раб Твоих, яко благовонное кадило. Тебе приносимое. Отжени от нас весь помысл лукавый; созижди в нас, Боже, сердце чисто и дух сокрушен; всели в нас, Боже, корень благих "--- страх Твой в сердца наши; виждь слезы наши, виждь сетование наше; зри скорби наша, зри горько плачущих о преставльшемся (-шейся) к Тебе рабе Твоем (Твоей) (\myemph{ имя}) и упокой его (ю) во Царствии Твоем, яко Ты еси Живот, Воскресение и Покой усопших раб Твоих, и Тебе славу возсылаем, яко Безначальному Отцу со Единородным Твоим Сыном и Всесвятым и Благим и Животворящим Твоим Духом всегда, ныне и присно и во веки веков. Аминь. 

\end{mymulticols}

\mychapterending

\mychapter{Молитва за всех в вере скончавшихся православных христиан}\begin{mymulticols}
%http://www.molitvoslov.org/text189.htm 
 

Владыко Многомилостиве и Всемилостиве Господи Иисусе Христе Боже наш! Милосерд ты к нам грешным и милосердию Твоему несть конца: не хощеши бо нашей вечной погибели, зане рекл еси: «Не хощу смерти грешника, но еже обратитися и живу быти ему». Слава милосердию Твоему! Слава смотрению Твоему! Слава долготерпению Твоему, Господи! От всего сердца нашего, от всея души нашея благодарим Тя, Боже наш, яко сподобил еси усопшаго(ую) раба Твоего (рабу Твою) (\myemph{ имя}) на одре болезни исповедатися Тебе во гресех своих. Ты же, Господи, Един ведый вся, веси, коль тяжки греси его (ея), но обаче молим Тя, приими раба Твоего (рабу Твою), отшедшаго(ую) к Тебе аще и грешником(-цею), но не нераскаянным(-ою), и воздаждь ему (ей) по раскаянному сердцу его (ея). И якоже Ты приемлеши и единое слово, и единый вздох, и единую каплю слезную от исповедающегося Тебе с чистым и сокрушенным сердцем, темже убо молим Тя, Сердцеведче Святый, приими ныне от нас нашу молитву за усопшаго(-ую), вонми молений наших гласу, и услыши ны; со умиленною душею и сокрушенным сердцем молим Тя, Господи, воззри всемилостивым оком Своим на покаявшагося (-шуюся) Тебе усопшаго(-ую) раба Твоего (рабу Твою) (имя) и попали невещественным огнем Духа Твоего Святаго вся терния грехов его (ея), неисповеданных пред лицем Твоим. Аще же усопший(-шая) раб Твой (раба Твоя) при покаянии забвению что предаде по немощи естества, или тяжкой болезни, или страха ради часа смертнаго не возможе в сокрушении и умилении сердца покаятися пред Тобою, Господи, ниже возможе принести Тебе плодов достойных покаяния, со умилением просим, молим Тя, Спасителю наш, восполни его (ея) покаяние молитвами Святыя Твоея Церкви и всех святых Твоих, от века благоугодивших Тебе, наипаче же восполни сию духовную скудость Твоею благостию, Твоими дражайшими искупительными и спасительными заслугами; и якоже мытаря оправдал еси, и разбойника, на кресте молившагося Ти, помиловал еси, сице молим Тя: помилуй преставльшагося(-уюся) к Тебе с верою и упованием и прости ему (ей) всякое согрешение вольное и невольное, словом, делом, ведением и неведением содеянное, и такода неосужденно явится пред лицем Твоим в день великаго и Страшнаго Суда Твоего. Молим Тя, дражайший Искупителю наш, сподоби усопшаго(ую) раба Твоего (рабу Твою) (\myemph{ имя}) внити в Царствие Твое, яко сподобившагося(-уюся), по Твоей благости, причаститися святых Твоих страшных и животворящих Тайн во оставление грехов и в жизнь вечную. Ты бо, Господи, пречистыми устами рекл еси: «Ядый Мою Плоть и пияй Мою Кровь имать живот вечный». Ты, Господи, Сам сподобил еси усопшаго(-ую) раба Твоего (рабу Твою) причаститися пречистаго Твоего Тела и пречистыя Твоея Крове, яко залога вечнаго блаженства; сего ради паки молим Тя: да будут ему (ей) святыя и животворящия Тайны Твоя, яко угль, попаляющ вся согрешения его (ея), вся беззакония и неправды его (ея), да будут во очищение, во освящение, оправдание и оставление грехов, и в залог живота и блаженства вечнаго. Аминь 

\end{mymulticols}

\mychapterending

\mychapter{Молитва вторая за всех в вере скончавшихся православных христиан}\begin{mymulticols}
%http://www.molitvoslov.org/text193.htm 
 

Боже щедрот и всякия благодати! Призри с высоты Святаго жилища Твоего на нас грешных и недостойных рабов Твоих, молящихся Тебе и просящих дати долгов разрешение усопшему (-ей) рабу Твоему (рабе Твоей). Молим Тя, Благость вечная, не вниди в суд с рабом (-ою) Твоим (-ею), яко созданием Твоей Премудрости, но обаче милости сподоби его (ю); яви ему (ей) благоутробие Твое, Всеблагий Владыко, и не отвержи его (ю) от лица Твоего, Человеколюбче Господи: спаси от вечныя смерти, яже есть смерть вторая, душу его (ея), да тако усопший (-шая) не престанет славословити Твою безконечную благость и превозносити имя Твое во веки. Тебе, Сердцеведче, ведомы дела, намерения и помышления усопшаго (-шия) раба Твоего (рабы Твоея) и вся содеянная им (ею) благая же и злая пред Тобою нага и объявлена суть. Но что убо пред Тобою, Святый, вся благая наша? Яко капля в мори, тако суть пред Тобою и пред Твоею правдою и Святостию вся благая наша, содеваемая в веце сем. Безмерны согрешения наши; велико и тяжко бремя грехов наших, но неизреченно велики и безпредельны заслуги Единороднаго Сына Твоего, Господа нашего Иисуса Христа, Егоже по безпредельной любви Твоей дал еси нам грешным, да всяк веруяй в Онь, не погибнет, но имать живот вечный. Зело грешны есмы пред Тобою и недостойны Твоея милости, но Твое человеколюбие соделывает нас достойными; велики беззакония наша, но неистощимо Твое милосердие: безмерныя заслуги Возлюбленнаго Твоего Сына, принесшаго Себя в жертву за грехи наши, побеждают всяк грех, всяко беззаконие; того ради со умилением и со дерзновением молим Тя, Источниче жизни нашея, со дерзновением упования молим, да усвоятся усопшему (-ей) рабу Твоему (рабе Твоей) (\myemph{ имя}) сии драгоценныя для нас искупительныя и спасительныя заслуги Единороднаго Сына Твоего; умоляем Твое милосердие слезами сокрушения, покрый грехи усопшаго (-шия) раба Твоего (рабы Твоея) величием сих дражайших заслуг Искупителя нашего; прости все согрешения его (ея) и не погуби его (ю) со беззаконми его (ея) в геенне огненней. Пощади, пощади, Владыко, отшедшаго (-ую) к Тебе раба Твоего (рабу Твою) (\myemph{ имя}) и вечных мучений избави, избави ради безмерных и неизглаголанных душевных скорбей и болезней Единороднаго Сына Твоего, понесенных Им пред испитием горькой чаши страданий Своих; пощади его (ю) и спаси от вечныя смерти ради поношения и биения Спасителя нашего, ради заушения и заплевания; помилуй ради дражайшия Крове Его из пречистых ребр Его пролиянныя. Молим Тя, Отче щедрот и всякаго милосердия, пречистыми и животворящими Его ранами, Его святою Кровию уврачуй смертоносныя и греховныя язвы усопшаго (-шия) раба Твоего (рабы Твоея), и тако да исцелеет духовне и да сподобится в неприступном свете Твоем, Светодавче Господи, славословит и прославляти Пресвятое имя Твое во веки. Аминь. 

\end{mymulticols}

\mychapterending

\mychapter{Молитва ко Пресвятой Богородице}\begin{mymulticols}
%http://www.molitvoslov.org/text187.htm 
 

Пресвятая Владычице Богородице! К Тебе прибегаем, Заступнице наша: Ты бо скорая еси помощница, наша ходатаица у Бога неусыпающая! Наипаче же молим Тя в час сей: помози новопреставленному(-ной) рабу Твоему (рабе Твоей) (\myemph{ имя}) прейти страшный и неведомый оный путь; молим Тя, Владычице мира, силою Твоею отжени от движимой страхом души его (ея) страшныя силы темных духов, да смятутся и посрамятся пред Тобою; свободи ю от истязания воздушных мытарей, разори советы их и низложи их, яко зломыслящих врагов. Буди ей, о Всемилостивая Госпоже Богородице, заступницею и защитницею от воздушного князя тьмы, мучителя и страшных путей стоятеля; молим Тя, Пресвятая Богородице, ризою Твоею честною Защити ю, да тако безбоязненно и невозбранно прейдет от земли на небо. Молим Тя, Заступнице наша, заступи за раба Твоего (рабу Твою) Матерним Твоим у Господа дерзновением; молим Тя, Помоще наша, помози ему (ей), имущему (-щей) судитися еще прежде онаго Страшнаго судилища, помози убо оправдатися пред Богом, яко Творцем неба и земли, и умоли Единороднаго Сына Твоего, Господа Бога и Спаса нашего Иисуса Христа, да упокоит усопшаго(-ую) в недрах Авраама с праведными и всеми святыми. Аминь.

\end{mymulticols}

\mychapterending

\mychapter{Молитва Ангелу хранителю}\begin{mymulticols}
%http://www.molitvoslov.org/text186.htm 
 

Святый Ангеле хранителю, данный усопшему (-шей) рабу Божию (рабе Божией) (имя)! Не преставай охраняти душу его (ея) от злых страшных оных бесов; буди ей пестуном и утешителем и тамо, в оном невидимом мире духов; приими ю под криле своя и преведи ю невозбранно чрез врата воздушных истязателей; предстани ходатаем и молитвенником за ню у Бога,"--- моли Его Преблагаго, да не низведена будет в место мрака, но да вчинит ю, идеже пребывает Свет невечерний. 

\end{mymulticols}

\mychapterending

\mychapter{Канон молебный от лица человека с душею разлучающегося и не могущаго глаголати, глас 6-й}\begin{mymulticols}
%http://www.molitvoslov.org/text179.htm 
 

\mysubtitle{Песнь 1}

\irmos{Яко по суху пешешествовав Израиль, по бездне стопами, гонителя фараона видя потопляема, Богу победную песнь поим, вопияше. }

\pripev{Пресвятая Богородице, спаси нас.}

Каплям подобно дождевным, злии и малии дние мои, летним обхождением оскудевающе, помалу исчезают уже, Владычице, спаси мя. 

\pripev{Пресвятая Богородице, спаси нас.}

Твоим благоутробием и многими щедротами Твоими, Владычице, преклоняема естественно, в час сей ужасный предстани ми, Помощнице Непоборимая. 

\pripev{Пресвятая Богородице, спаси нас.}

Содержит ныне душу мою страх велик, трепет неисповедим и болезнен есть, внегда изыти ей от телесе, Пречистая, юже утеши. 

\slavan

Грешным и смиренным известное прибежище, о мне извести Твою милость, Чистая, и бесовския избави руки, якоже бо пси мнози обступиша мя. 

\inynen

 Се время помощи, се время Твоего заступления, се, Владычице, время, о немже день и нощь припадах тепле и моляхся Тебе.

\mysubtitle{Песнь 3}

\irmos{Несть свят, якоже Ты, Господи Боже мой, вознесый рог верных Твоих, Блаже, и утвердивый нас на камени исповедания Твоего. }

\pripev{Пресвятая Богородице, спаси нас.}

Издалеча сего дне, Владычице, провидя, и того яко пришедша помышляяй присно, слезами теплыми моляхся не забыти мене. 

\pripev{Пресвятая Богородице, спаси нас.}

Обыдоша мя мысленнии рыкающе скимны, и ищут восхитити и растерзати мя горце, ихже зубы, Чистая, и челюсти сокруши и спаси мя. 

\pripev{Пресвятая Богородице, спаси нас.}

Угасшу убо отнюд органу словесному, и связавшуся языку, и затворившуся гласу, в сокрушении сердца молю Тя, Спасительнице моя, спаси мя. 

\slavan

Приклони ухо Твое ко мне, Христа Бога моего Мати, от высоты многия славы Твоея, Благая, и услыши стенание конечное, и руку ми подаждь. 

\inynen

 Не отврати от мене многия щедроты Твоя, не затвори утробу Твою человеколюбную, Чистая; но предстани ми ныне, и в час судный помяни мя.

\mysubtitle{Песнь 4}

\irmos{Христос моя сила, Бог и Господь, честная Церковь боголепно поет, взывающи, от смысла чиста о Господе празднующи. }

\pripev{Пресвятая Богородице, спаси нас.}

Умовение согрешением, ток слезный ныне положи, Благая, сердца моего сокрушение приемлющи: о Тебе утверждающу упование, Благая, егда како страшнаго мя избавиши огненнаго мучения, яко Сама благодати еси источник, Богородительнице. 

\pripev{Пресвятая Богородице, спаси нас.}

Непостыдное и непогрешительное всем, иже в нуждах, прибежище, Владычице Пренепорочная, Ты ми буди Заступница в час испытания. 

\pripev{Пресвятая Богородице, спаси нас.}

Простерши пречистеи Твои и всечестнеи руце, яко священнии голубине криле, под кровом и сению тех покрый мя, Владычице. 

\slavan

Воздушнаго князя насильника, мучителя, страшных путей стоятеля и напраснаго сих словоиспытателя, сподоби мя прейти невозбранно отходяща от земли. 

\inynen

 Се мя, Владычице, страх усрете, егоже и бояхся: се подвиг велик объят мя, в немже буди ми Помощница, Надежде спасения моего.

\mysubtitle{Песнь 5}

\irmos{Божиим светом Твоим, Блаже, утренюющих Ти души любовию озари, молюся, Тя ведети, Слове Божий, Истиннаго Бога, от мрака греховнаго взывающа. }

\pripev{Пресвятая Богородице, спаси нас.}

Не забуди мя, Благая, ниже отврати от мене, Твоего отрока, лице Твое, но услыши мя, яко скорблю, и вонми души моей, и сию избави. 

\pripev{Пресвятая Богородице, спаси нас.}

Иже по плоти сродницы мои, и иже по духу братие, и друзи, и обычнии знаемии, плачите, воздохните, сетуйте, се бо от вас ныне разлучаюся. 

\pripev{Пресвятая Богородице, спаси нас.}

Ныне избавляяй никако и помогаяй воистинну никтоже; Ты помози ми, Владычице, да не яко человек безпомощен, в руках враг моих затворен буду. 

\slavan

Вшедше, святии мои Ангели, предстаните судищу Христову, колене свои мысленнии преклоньше, плачевне возопийте Ему: помилуй, Творче всех, дело рук Твоих, Блаже, и не отрини его. 

\inynen

 Поклоньшеся Владычице и Пречистей Матери Бога моего помолитеся, яко да преклонит колена с вами, и преклонит Его на милость: Мати бо сущи и Питательница услышана будет.

\mysubtitle{Песнь 6}

\irmos{Житейское море, воздвизаемое зря напастей бурею, к тихому пристанищу Твоему притек, вопию Ти: возведи от тли живот мой, Многомилостиве. }

\pripev{Пресвятая Богородице, спаси нас.}

Устне мои молчат, и язык не глаголет, но сердце вещает: огнь бо сокрушения сие снедая внутрь возгарается, и гласы неизглаголанными Тебе, Дево, призывает. 

\pripev{Пресвятая Богородице, спаси нас.}

Призри на мя свыше, Мати Божия, и милостивно вонми ныне на мое посещение снити, яко да видев Тя, от телесе изыду радуяся. 

\pripev{Пресвятая Богородице, спаси нас.}

Растерзаеми соузы, раздираеми закони естественнаго сгущения, и составления всего телеснаго, нужду нестерпимую и тесноту сотворяют ми. 

\slavan

Святых Ангел священным и честным рукам преложи мя, Владычице, яко да тех крилы покрывся, не вижу безчестнаго и смраднаго и мрачнаго бесов образа. 

\inynen

 Чертоже Божий Всечестне, небесному разумному чертогу сподоби мя, мою угасшую и несияющую свещу вжегши, святым елеем милости Твоея.

\mysubtitle{Кондак, глас 6}

Душе моя, душе моя, востани, что спиши, конец приближается, и нужда ти молвити: воспряни убо, да пощадит тя Христос Бог, Иже везде сый, и вся исполняяй. 

\mysubtitle{Икос} 

Христово врачевство видя отверсто, и от сего Адаму источающе здравие, пострадав уязвися диавол, яко беды приемля рыдаше, и своим другом возопи: что сотворю Сыну Мариину? Убивает мя Вифлеемлянин, Иже везде сый и вся исполняяй.

\mysubtitle{Песнь 7}

\irmos{Росодательну убо пещь содела Ангел, преподобным отроком, халдеи же опаляющее веление Божие, мучителя увеща вопити: благословен еси, Боже отец наших. }

\pripev{Пресвятая Богородице, спаси нас.}

Нощь смертная мя постиже неготова, мрачна же и безлунна, препущающи неприготовлена к долгому оному пути страшному: да спутешествует ми Твоя милость, Владычице. 

\pripev{Пресвятая Богородице, спаси нас.}

Се вси дние мои исчезоша воистинну в суете, якоже пишется, и лета моя со тщанием, сети же смертныя воистинну и горькия предвариша мою душу, яже мя обдержат. 

\pripev{Пресвятая Богородице, спаси нас.}

Множество грехов моих да не возможет победити Твоего многаго благоутробия, Владычице; но да обыдет мя Твоя милость и вся да покрыет беззакония моя. 

\slavan

Отводящии мя отсюду находят, содержаще мя всюду: душа же моя отлагает и страшится, многа исполнена мятежа, юже утеши, Чистая, явлением Твоим. 

\inynen

 Спечалующаго ниединаго в скорби моей, ниже утешающаго обретох, Владычице, ибо друзи мои и знаемии вкупе оставиша мя ныне; но, Надежде моя, никакоже да не оставиши мя.

\mysubtitle{Песнь 8}

\irmos{Из пламене преподобным росу источил еси, и праведнаго жертву водою попалил еси, вся бо твориши, Христе, токмо еже хотети. Тя превозносим во вся веки. }

\pripev{Пресвятая Богородице, спаси нас.}

Яко Бога Человеколюбца Мати Человеколюбивая, тихим и милостивым вонми оком, егда от тела душа моя отлучается, да Тя во вся веки славлю, Святая Богородице. 

\pripev{Пресвятая Богородице, спаси нас.}

Убегнути ми варвар безплотных полки и воздушныя бездны возникнути, и к Небеси взыти мя сподоби, да Тя во веки славлю, Святая Богородице. 

\pripev{Пресвятая Богородице, спаси нас.}

Рождшая Господа Вседержителя, горьких мытарств начальника миродержца отжени далече от мене, внегда скончатися хощу, да Тя во веки славлю, Святая Богородице. 

\slavan

Великой последней гласящей трубе, в страшное и грозное воскрешение суда, воскресающим всем, помяни мя тогда, Святая Богородице. 

\inynen

 Высокая Владыки Христа палато, Твою благодать свыше пославши, предвари мя ныне в день озлобления, да Тя славлю во веки вся, Святая Богородице.

\mysubtitle{Песнь 9}

\irmos{Бога человеком невозможно видети, на Негоже не смеют чини ангельстии взирати. Тобою же, Всечистая, явися человеком Слово Воплощенно. Егоже величающе, с небесными вои Тя ублажаем.} 

\pripev{Пресвятая Богородице, спаси нас.}

О, како узрю Невидимаго; како ужасное оно претерплю видение? Како дерзну отверсти очи? Како моего Владыку смею видети, Егоже не престаях от юности огорчевая присно? 

\pripev{Пресвятая Богородице, спаси нас.}

Святая Отроковице, Богородительнице, на мое смирение милосердно призри, умиленное мое и последнее моление сие приимши, и мучащаго вечнующаго огня потщися избавити мя. 

\pripev{Пресвятая Богородице, спаси нас.}

Храмы святыя осквернившая, скверный и телесный храм оставивши, Тебе, Божий Всечестный Храме, молит, Отроковице Дево Мати, душа моя тьмы кромешния убежати и лютаго геенскаго жжения. 

\slavan

Зря конец близу жития моего, и помышляя безместных мыслей, деяний же, Всечистая, душу мою делательницу, люте уязвляюся стрелами совести, но преклоньшися милостивно, буди ми Предстательница. 

\inynen

 Сын дастся нам за милость, Сын Божий, и Ангельский Царь, Превечный от чистых кровей Твоих Человек прошед. Егоже умилостиви, Отроковице, страстней моей душе, исторгаемей люте от окаяннаго моего телесе. 

\myemph{Таже:} Достойно есть яко воистинну блажити Тя Богородицу, Присноблаженную и Пренепорочную и Матерь Бога нашего. Честнейшую Херувим и славнейшую без сравнения Серафим, без истления Бога Слова рождшую, сущую Богородицу Тя величаем.

\mysubtitle{Молитва, от иерея глаголемая на исход души:}

Владыко Господи Вседержителю, Отче Господа нашего Иисуса Христа, Иже всем человеком хотяй спастися и в разум истины приити, не хотяй смерти грешному, но обращения и живота; молимся и мили ся Ти деем, душу раба Твоего (имярек) от всякия узы разреши и от всякия клятвы свободи, остави прегрешения ему, яже от юности, ведомая и неведомая, в деле и слове, и чисто исповеданная, или забвением, или студом утаеная. Ты бо Един еси разрешаяй связанныя и исправляяй сокрушенныя, надежда неначаемым, могий оставляти грехи всякому человеку, на Тя упование имущему. Ей, человеколюбивый Господи, повели, да отпустится от уз плотских и греховных, и приими в мир душу раба Твоего сего (имярек), и покой ю в вечных обителех со святыми Твоими, благодатию Единороднаго Сына Твоего, Господа Бога и Спаса нашего Иисуса Христа, с Нимже благословен еси, с Пресвятым и Благим, и Животворящим Твоим Духом, ныне и присно, и во веки веков. Аминь.

\end{mymulticols}

\mychapterending

\mychapter{Канон молебный ко Господу Иисусу Христу и Пречистой Богородице Матери Господни  при разлучении души от тела всякого правоверного}\begin{mymulticols}
%http://www.molitvoslov.org/text174.htm 
 

\myemph{ Приходит игумен, к мирскому же отец его духовный, и вопрошает, аще есть кое слово, или дело забвения ради, или студа, или кая злоба к коему брату неисповедана, или непрощена есть, вся должен есть изыскивати и вопрошати по единому умирающаго.}

\myemph{Посем начинает иерей:} Благословен Бог наш, всегда, ныне и присно, и во веки веков. (\myemph{аще ли мирский:} \MolitvamiSviatyhOtecNashih)

\TrisviatoePoOtcheNash

Господи, помилуй \myemph{(12)}.

\priiditepoklonimsia

\mysubtitle{Псалом 50}

\PsalmFifty

\mysubtitle{Песнь 1}

\irmos{Яко по суху пешешествовав Израиль, по бездне стопами, гонителя фараона видя потопляема, Богу победную песнь поим, вопияше.}

\pripev{Пресвятая Богородице, спаси нас.}

Каплям подобно дождевным, злии и малии дние мои, летним обхождением оскудевающе, помалу исчезают уже, Владычице, спаси мя.

\pripev{Пресвятая Богородице, спаси нас.}

Твоим благоутробием и многими щедротами Твоими, Владычице, преклоняема естественно, в час сей ужасный предстани ми, Помощнице Непоборимая.

\pripev{Пресвятая Богородице, спаси нас.}

Содержит ныне душу мою страх велик, трепет неисповедим и болезнен есть, внегда изыти ей от телесе, Пречистая, юже утеши.

\slava

Грешным и смиренным известное прибежище, о мне извести Твою милость, Чистая, и бесовския избави руки, якоже бо пси мнози обступиша мя.

\inyne

 Се время помощи, се время Твоего заступления, се, Владычице, время, о немже день и нощь припадах тепле и моляхся Тебе.

\mysubtitle{Песнь 3}

\irmos{Несть свят, якоже Ты, Господи Боже мой, вознесый рог верных Твоих, Блаже, и утвердивый нас на камени исповедания Твоего.}

\pripev[]{Пресвятая Богородице, спаси нас.}

Издалеча сего дне, Владычице, провидя, и того яко пришедша помышляяй присно, слезами теплыми моляхся не забыти мене.

\pripev[]{Пресвятая Богородице, спаси нас.}

Обыдоша мя мысленнии рыкающе скимны, и ищут восхитити и растерзати мя горце, ихже зубы, Чистая, и челюсти сокруши и спаси мя.

\pripev[]{Пресвятая Богородице, спаси нас.}

Угасшу убо отнюд органу словесному, и связавшуся языку, и затворившуся гласу, в сокрушении сердца молю Тя, Спасительнице моя, спаси мя.

\slava

Приклони ухо Твое ко мне, Христа Бога моего Мати, от высоты многия славы Твоея, Благая, и услыши стенание конечное, и руку ми подаждь.

\inyne

 Не отврати от мене многия щедроты Твоя, не затвори утробу Твою человеколюбную, Чистая; но предстани ми ныне, и в час судный помяни мя.

\mysubtitle{Песнь 4}

\irmos{Христос моя сила, Бог и Господь, честная Церковь боголепно поет, взывающи, от смысла чиста о Господе празднующи.}

\pripev[]{Пресвятая Богородице, спаси нас.}

Умовение согрешением, ток слезный ныне положи, Благая, сердца моего сокрушение приемлющи: о Тебе утверждающу упование, Благая, егда како страшнаго мя избавиши огненнаго мучения, яко Сама благодати еси источник, Богородительнице.

\pripev[]{Пресвятая Богородице, спаси нас.}

Непостыдное и непогрешительное всем, иже в нуждах, прибежище, Владычице Пренепорочная, Ты ми буди Заступница в час испытания.

\pripev[]{Пресвятая Богородице, спаси нас.}

Простерши пречистеи Твои и всечестнеи руце, яко священнии голубине криле, под кровом и сению тех покрый мя, Владычице.

\slava

Воздушнаго князя насильника, мучителя, страшных путей стоятеля и напраснаго сих словоиспытателя, сподоби мя прейти невозбранно отходяща от земли.

\inyne

 Се мя, Владычице, страх усрете, егоже и бояхся: се подвиг велик объят мя, в немже буди ми Помощница, Надежде спасения моего.

\mysubtitle{Песнь 5}

\irmos{Божиим светом Твоим, Блаже, утренюющих Ти души любовию озари, молюся, Тя ведети, Слове Божий, Истиннаго Бога, от мрака греховнаго взывающа.}

\pripev[]{Пресвятая Богородице, спаси нас.}

Не забуди мя, Благая, ниже отврати от мене, Твоего отрока, лице Твое, но услыши мя, яко скорблю, и вонми души моей, и сию избави.

\pripev[]{Пресвятая Богородице, спаси нас.}

Иже по плоти сродницы мои, и иже по духу братие, и друзи, и обычнии знаемии, плачите, воздохните, сетуйте, се бо от вас ныне разлучаюся.

\pripev[]{Пресвятая Богородице, спаси нас.}

Ныне избавляяй никако и помогаяй воистинну никтоже; Ты помози ми, Владычице, да не яко человек безпомощен, в руках враг моих затворен буду.

\slava

Вшедше, святии мои Ангели, предстаните судищу Христову, колене свои мысленнии преклоньше, плачевне возопийте Ему: помилуй, Творче всех, дело рук Твоих, Блаже, и не отрини его.

\inyne

 Поклоньшеся Владычице и Пречистей Матери Бога моего помолитеся, яко да преклонит колена с вами, и преклонит Его на милость: Мати бо сущи и Питательница услышана будет.

\mysubtitle{Песнь 6}

\irmos{Житейское море, воздвизаемое зря напастей бурею, к тихому пристанищу Твоему притек, вопию Ти: возведи от тли живот мой, Многомилостиве.}

\pripev[]{Пресвятая Богородице, спаси нас.}

Устне мои молчат, и язык не глаголет, но сердце вещает: огнь бо сокрушения сие снедая внутрь возгарается, и гласы неизглаголанными Тебе, Дево, призывает.

\pripev[]{Пресвятая Богородице, спаси нас.}

Призри на мя свыше, Мати Божия, и милостивно вонми ныне на мое посещение снити, яко да видев Тя, от телесе изыду радуяся.

\pripev[]{Пресвятая Богородице, спаси нас.}

Растерзаеми соузы, раздираеми закони естественнаго сгущения, и составления всего телеснаго, нужду нестерпимую и тесноту сотворяют ми.

\slava

Святых Ангел священным и честным рукам преложи мя, Владычице, яко да тех крилы покрывся, не вижу безчестнаго и смраднаго и мрачнаго бесов образа.

\inyne

 Чертоже Божий Всечестне, небесному разумному чертогу сподоби мя, мою угасшую и несияющую свещу вжегши, святым елеем милости Твоея.

\mysubtitle{Кондак, глас 6}

Душе моя, душе моя, востани, что спиши, конец приближается, и нужда ти молвити: воспряни убо, да пощадит тя Христос Бог, Иже везде сый, и вся исполняяй.

\mysubtitle{Икос}

Христово врачевство видя отверсто, и от сего Адаму источающе здравие, пострадав уязвися диавол, яко беды приемля рыдаше, и своим другом возопи: что сотворю Сыну Мариину? Убивает мя Вифлеемлянин, Иже везде сый и вся исполняяй.

\mysubtitle{Песнь 7}

\irmos{Росодательну убо пещь содела Ангел, преподобным отроком, халдеи же опаляющее веление Божие, мучителя увеща вопити: благословен еси, Боже отец наших.}

\pripev[]{Пресвятая Богородице, спаси нас.}

Нощь смертная мя постиже неготова, мрачна же и безлунна, препущающи неприготовлена к долгому оному пути страшному: да спутешествует ми Твоя милость, Владычице.

\pripev[]{Пресвятая Богородице, спаси нас.}

Се вси дние мои исчезоша воистинну в суете, якоже пишется, и лета моя со тщанием, сети же смертныя воистинну и горькия предвариша мою душу, яже мя обдержат.

\pripev[]{Пресвятая Богородице, спаси нас.}

Множество грехов моих да не возможет победити Твоего многаго благоутробия, Владычице; но да обыдет мя Твоя милость и вся да покрыет беззакония моя.

\slava

Отводящии мя отсюду находят, содержаще мя всюду: душа же моя отлагает и страшится, многа исполнена мятежа, юже утеши, Чистая, явлением Твоим.

\inyne

 Спечалующаго ниединаго в скорби моей, ниже утешающаго обретох, Владычице, ибо друзи мои и знаемии вкупе оставиша мя ныне; но, Надежде моя, никакоже да не оставиши мя.

\mysubtitle{Песнь 8}

\irmos{Из пламене преподобным росу источил еси, и праведнаго жертву водою попалил еси, вся бо твориши, Христе, токмо еже хотети. Тя превозносим во вся веки.}

\pripev[]{Пресвятая Богородице, спаси нас.}

Яко Бога Человеколюбца Мати Человеколюбивая, тихим и милостивым вонми оком, егда от тела душа моя отлучается, да Тя во вся веки славлю, Святая Богородице.

\pripev[]{Пресвятая Богородице, спаси нас.}

Убегнути ми варвар безплотных полки и воздушныя бездны возникнути, и к Небеси взыти мя сподоби, да Тя во веки славлю, Святая Богородице.

\pripev[]{Пресвятая Богородице, спаси нас.}

Рождшая Господа Вседержителя, горьких мытарств начальника миродержца отжени далече от мене, внегда скончатися хощу, да Тя во веки славлю, Святая Богородице.

\slava

Великой последней гласящей трубе, в страшное и грозное воскрешение суда, воскресающим всем, помяни мя тогда, Святая Богородице.

\inyne

 Высокая Владыки Христа палато, Твою благодать свыше пославши, предвари мя ныне в день озлобления, да Тя славлю во веки вся, Святая Богородице.

\mysubtitle{Песнь 9}

\irmos{Бога человеком невозможно видети, на Негоже не смеют чини ангельстии взирати. Тобою же, Всечистая, явися человеком Слово Воплощенно. Егоже величающе, с небесными вои Тя ублажаем.}

\pripev[]{Пресвятая Богородице, спаси нас.}

О, како узрю Невидимаго; како ужасное оно претерплю видение? Како дерзну отверсти очи? Како моего Владыку смею видети, Егоже не престаях от юности огорчевая присно?

\pripev[]{Пресвятая Богородице, спаси нас.}

Святая Отроковице, Богородительнице, на мое смирение милосердно призри, умиленное мое и последнее моление сие приимши, и мучащаго вечнующаго огня потщися избавити мя.

\pripev[]{Пресвятая Богородице, спаси нас.}

Храмы святыя осквернившая, скверный и телесный храм оставивши, Тебе, Божий Всечестный Храме, молит, Отроковице Дево Мати, душа моя тьмы кромешния убежати и лютаго геенскаго жжения.

\slava

Зря конец близу жития моего, и помышляя безместных мыслей, деяний же, Всечистая, душу мою делательницу, люте уязвляюся стрелами совести, но преклоньшися милостивно, буди ми Предстательница.

\inyne

 Сын дастся нам за милость, Сын Божий, и Ангельский Царь, Превечный от чистых кровей Твоих Человек прошед. Егоже умилостиви, Отроковице, страстней моей душе, исторгаемей люте от окаяннаго моего телесе.

\myemph{Таже:} \Chestneyshuyu

\end{mymulticols}

\mychapterending


%

\mypart{БОЖЕСТВЕННАЯ ЛИТУРГИЯ}\label{_content_Liturgiya}
%http://www.molitvoslov.org/content/Liturgiya

\mychapter{О Литургии}

В разделе \textbf{БОЖЕСТВЕННАЯ ЛИТУРГИЯ} приведены песнопения Божественных литургий святителей Иоанна Златоуста и Василия Великого. Избранные песнопения литургии Преждеосвященных Даров помещены в главе «Песнопения из служб Триоди постной».

Литургия святителя Иоанна Златоуста совершается в Православной Церкви в течении всего года, кроме Великого Поста, когда она совершается по субботам, в Благовещение Пресвятой Богородицы и в Неделю ваий.

Литургия святителя Василия Великого совершается 10 раз в году: в воскресные дни Великого поста, в Великие четверг и субботу, в день памяти святителя Василия Великого (1 января), в Рождественский и Крещенский сочельники (если же эти дни выпадают на субботу или воскресенье, то совершается в самые праздники Рождества Христова и Крещения Господня).

Литургия состоит из трех частей: проскомидии, литургии оглашенных и литургии верных.

\mychapter{Проском\'{и}дия}
%http://www.molitvoslov.org/content/Proskomidiya 

До чтения третьего и шестого часов, или во время чтения их, в алтаре совершаются священнодействия \textbf{проском\'{и}дии}, через которые из принесенных хлеба и вина приготовляется вещество для Св. Евхаристии и при этом совершается предварительное поминовение членов Церкви Христовой "--- небесной и земной.

От древнего обычая приносить в храм хлеб и вино для таинства Св. Евхаристии первая часть литургии и называется проском\'{и}дией "--- приношением.

\mychapterending

\mychapter{Литургия оглашенных}
%http://www.molitvoslov.org/text906.htm 

Вторая часть литургии называется \textbf{литургией оглашенных}. Такое название эта часть службы получила от содержания в ее составе молитвословий, песнопений, священнодействий и поучений, имеющих учительный, огласительный характер. В древней Церкви во время ее совершения могли присутствовать, вместе с верными, и оглашенные, т.~е. лица, готовящиеся ко Св. Крещению, а также кающиеся, отлученные от Св. Причастия.

\begin{mymulticols}

\myemph{Диакон:} Благослови, владыко.

\myemph{Иерей:} Благословено Царство Отца и Сына и Святаго Духа, ныне и присно и во веки веков. 

\myemph{Хор:} Аминь.

\mysubtitle{Великая ектения}

\myemph{Диакон:} Миром Господу помолимся. 

\myemph{Хор:} Господи, помилуй.

\myemph{Диакон:} О свышнем мире и спасении душ наших, Господу помолимся. 

\myemph{Хор:} Господи, помилуй.

\myemph{Диакон:} О мире всего мира, благостоянии Святых Божиих Церквей и соединении всех, Господу помолимся.

\myemph{Хор:} Господи, помилуй.

\myemph{Диакон:} О святем храме сем и с верою, благоговением и страхом Божиим входящих в онь, Господу помолимся. 

\myemph{Хор:} Господи, помилуй.

\myemph{Диакон:} О Великом Господине и Отце нашем Святейшем Патриархе Кирилле, и о Господине нашем Преосвященнейшем митрополите ( \myemph{или:} архиепископе, \myemph{или:} епископе) \myemph{(имярек)}, честнем пресвитерстве, во Христе диаконстве, о всем причте и людех, Господу помолимся. 

\myemph{Хор:} Господи, помилуй.

\myemph{Диакон:} О Богохранимей стране нашей, властех и воинстве ея, Господу помолимся.

\myemph{Хор:} Господи, помилуй.

\myemph{Диакон:} О граде сем ( \myemph{или:} o веси сей, \myemph{если в монастыре, то:} о святей обители сей), всяком граде, стране и верою живущих в них, Господу помолимся. 

\myemph{Хор:} Господи, помилуй.

\myemph{Диакон:} О благорастворении воздухов, о изобилии плодов земных и временех мирных, Господу помолимся.

\myemph{Хор:} Господи, помилуй.

\myemph{Диакон:} О плавающих, путешествующих, недугующих, страждущих, плененных и о спасении их, Господу помолимся. 

\myemph{Хор:} Господи, помилуй.

\myemph{Диакон:} О избавитися нам от всякия скорби, гнева и нужды, Господу помолимся. 

\myemph{Хор:} Господи, помилуй.

\myemph{Диакон:} Заступи, спаси, помилуй и сохрани нас, Боже, Твоею благодатию. 

\myemph{Хор:} Господи, помилуй.

\myemph{Диакон:} Пресвятую, Пречистую, Преблагословенную, Славную Владычицу нашу Богородицу и Приснодеву Марию, со всеми святыми помянувше, сами себе, и друг друга, и весь живот наш Христу Богу предадим. 

\myemph{Хор:} Тебе, Господи. 

\myemph{Иерей:} Яко подобает Тебе всякая слава, честь и поклонение, Отцу и Сыну и Святому Духу, ныне и присно и во веки веков. 

\myemph{Хор:} Аминь.

\mysubtitle{Антифоны}

\myemph{Антифоны на литургии бывают трех родов: праздничные, изобразительные и вседневные (будничные). Какие из них поются, определяется на каждый день церковным Уставом. Праздничные антифоны поются в праздники Господни, за исключением Сретения (праздничные антифоны в Неделю ваий, на Пасху, на Вознесение и в День Святой Троицы приведены в главе «Песнопения из служб Триоди цветной»).}

\myemph{Вседневные антифоны положено петь в будни. Чаще всего в воскресные и праздничные дни, поются антифоны изобразительные (псалмы 102, 145 и Блаженны "--- Мф. 5, 3-12)}

\mysubtitle{Первый антифон} 

\myemph{Хор 1-й:} Благослови, душе моя, Господа. Благословен еси, Господи. Благослови, душе моя, Господа, и вся внутренняя моя Имя святое Его. 

\myemph{Хор 2-й:} Благослови, душе моя, Господа, и не забывай всех воздаяний Его. 

\myemph{Хор 1-й:} Очищающаго вся беззакония твоя, исцеляющаго вся недуги твоя. 

\myemph{Хор 2-й:} Избавляющаго от истления живот твой, венчающаго тя милостию и щедротами. 

\myemph{Хор 1-й:} Исполняющаго во благих желание твое: обновится, яко орля, юность твоя. 

\myemph{Хор 2-й:} Творяй милостыни Господь, и судьбу всем обидимым. 

\myemph{Хор 1-й:} Сказа пути Своя Моисеови, сыновом Израилевым хотения Своя. 

\myemph{Хор 2-й:} Щедр и милостив Господь, долготерпелив и многомилостив.

\myemph{Хор 1-й:} Не до конца прогневается, ниже в век враждует. 

\myemph{Хор 2-й:} Не по беззаконием нашим сотворил есть нам, ниже по грехом нашим воздал есть нам. 

\myemph{Хор 1-й:} Яко по высоте небесней от земли, утвердил есть Господь милость Свою на боящихся Его. 

\myemph{Хор 2-й:} Елико отстоят востоцы от запад, удалил есть от нас беззакония наша. 

\myemph{Хор 1-й:} Якоже щедрит отец сыны, ущедри Господь боящихся Его. 

\myemph{Хор 2-й:} Яко Той позна создание наше, помяну, яко персть есмы. 

\myemph{Хор 1-й:} Человек, яко трава дние его, яко цвет сельный, тако оцветет. 

\myemph{Хор 2-й:} Яко дух пройде в нем, и не будет, и не познает ктому места своего. 

\myemph{Хор 1-й:} Милость же Господня от века и до века на боящихся Его. 

\myemph{Хор 2-й:} И правда Его на сынех сынов, хранящих завет Его, и помнящих заповеди Его творити я. 

\myemph{Хор 1-й:} Господь на небеси уготова Престол Свой, и Царство Его всеми обладает. 

\myemph{Хор 2-й:} Благословите Господа, ангели Его, сильнии крепостию, творящии слово Его, услышати глас словес Его. 

\myemph{Хор 1-й:} Благословите Господа, вся силы Его, слуги Его, творящии волю Его. 

\myemph{Хор 2-й:} Благословите Господа, вся дела Его, на всяком месте владычества Его. 

\myemph{Хор 1-й:} Слава Отцу и Сыну и Святому Духу. 

\myemph{Хор 2-й:} И ныне и присно и во веки веков. Аминь. 

\myemph{Хор 1-й:} Благослови, душе моя, Господа, и вся внутренняя моя, имя святое Его. Благословен еси, Господи.

\mysubtitle{Ектения малая} 

\myemph{Диакон:} Паки и паки миром Господу помолимся. 

\myemph{Хор:} Господи, помилуй. 

\myemph{Диакон:} Заступи, спаси, помилуй и сохрани нас, Боже, Твоею благодатию. 

\myemph{Хор:} Господи, помилуй. 

\myemph{Диакон:} Пресвятую, Пречистую, Преблагословенную, Славную Владычицу нашу Богородицу и Приснодеву Марию, со всеми святыми помянувше, сами себе, и друг друга, и весь живот наш Христу Богу предадим. 

\myemph{Хор:} Тебе, Господи. 

\myemph{Иерей:} Яко Твоя держава, и Твое есть Царство, и сила, и слава, Отца и Сына и Святаго Духа, ныне и присно и во веки веков. 

\myemph{Хор:} Аминь. 

\mysubtitle{Второй антифон} 

\myemph{Во время второго антифона зажигается пономарская свеча. Алтарник берёт свечу во время «Единородный Сыне…» и становится с ней на горнем месте.}

\myemph{Хор:} Слава Отцу, и Сыну, и Святому Духу. 

\myemph{Хор 1-й:} Хвали, душе моя, Господа. Восхвалю Господа в животе моем, пою Богу моему, дондеже есмь. 

\myemph{Хор 2-й:} Не надейтеся на князи, на сыны человеческия, в них же несть спасения. 

\myemph{Хор 1-й:} Изыдет дух его, и возвратится в землю свою: в той день погибнут вся помышления его. 

\myemph{Хор 2-й:} Блажен, емуже Бог Иаковль помощник его, упование его на Господа Бога своего. 

\myemph{Хор 1-й:} Сотворшаго небо и землю, море и вся, яже в них. 

\myemph{Хор 2-й:} Хранящаго истину в век, творящаго суд обидимым, дающаго пищу алчущим. 

\myemph{Хор 1-й:} Господь решит окованныя, Господь умудряет слепцы. 

\myemph{Хор 2-й:} Господь возводит низверженныя, Господь любит праведники. 

\myemph{Хор 1-й:} Господь хранит пришельцы, сира и вдову приимет, и путь грешных погубит. 

\myemph{Хор 2-й:} Воцарится Господь во век, Бог твой, Сионе, в род и род. И ныне и присно и во веки веков. Аминь. 

\myemph{Хор:} И ныне и присно и во веки веков. Аминь. 

\mysubtitle{Песнь Господу Иисусу Христу} 

\myemph{Хор:} Единородный Сыне и Слове Божий, Безсмертен Сый, и изволивый спасения нашего ради воплотитися от Святыя Богородицы и Приснодевы Марии, непреложно вочеловечивыйся; распныйся же, Христе Боже, смертию смерть поправый, един Сый Святыя Троицы, спрославляемый Отцу и Святому Духу, спаси нас. 

\mysubtitle{Ектения малая} 

\myemph{Диакон:} Паки и паки миром Господу помолимся. 

\myemph{Хор:} Господи, помилуй. 

\myemph{Диакон:} Заступи, спаси, помилуй и сохрани нас, Боже, Твоею благодатию. 

\myemph{Хор:} Господи, помилуй. 

\myemph{Диакон:} Пресвятую, Пречистую, Преблагословенную, Славную Владычицу нашу Богородицу и Приснодеву Марию, со всеми святыми помянувше, сами себе, и друг друга, и весь живот наш Христу Богу предадим.

\myemph{Хор:} Тебе, Господи. 

\myemph{Иерей:} Яко Твоя держава, и Твое есть Царство, и сила, и слава, Отца и Сына и Святаго Духа, ныне и присно и во веки веков. 

\myemph{Хор:} Аминь.

\mysubtitle{Третий антифон; Блаженны} 

\myemph{«Блаженны» положено петь с тропарями, назначенными в этот день церковным Уставом: особыми тропарями на «Блаженных», или тропарями из песней утреннего канон празднику или святому.}

\myemph{Хор 1-й:} Во Царствии Твоем помяни нас, Господи, егда приидеши во Царствии Твоем. 

\myemph{Хор 2-й, стих 12-й:} Блажени нищий духом, яко тех есть Царство Небесное. 

\myemph{Хор 1-й:} Блажени плачущии, яко тии утешатся.

\myemph{Хор 2-й, стих 10-й:} Блажени кротции, яко тии наследят землю. 

\myemph{Хор 1-й:} Блажени алчущии и жаждущии правды, яко тии насытятся.

\myemph{Хор 2-й, стих 8-й:} Блажени милостивии, яко тии помиловани будут. 

\myemph{Хор 1-й:} Блажени чистии сердцем, яко тии Бога узрят.

\myemph{Хор 2-й, стих 6-й:} Блажени миротворцы, яко тии сынове Божии нарекутся. 

\myemph{Хор 1-й:} Блажени изгнани правды ради, яко тех есть Царство Небесное.

\myemph{Хор 2-й, стих 4-й:} Блажени есте, егда поносят вам, и изженут, и рекут всяк зол глагол на вы, лжуще Мене ради. 

\myemph{Хор 1-й:} Радуйтеся и веселитеся, яко мзда ваша многа на небесех. 

\myemph{Хор:} Слава Отцу и Сыну и Святому Духу. 

\myemph{Хор:} И ныне и присно и во веки веков. Аминь. 

\mysubtitle{Антифоны вседневные (будничные)} 

\myparagraph{Антифон 1-й}

\myemph{Хор 1-й:} Благо есть исповедатися Господеви. Молитвами Богородицы, Спасе, спаси нас. 

\myemph{Хор 2-й:} Благо есть исповедатися Господеви, и пети имени Твоему, Вышний. Молитвами Богородицы, Спасе, спаси нас. 

\myemph{Хор 1-й:} Возвещати заутра милость Твою, и истину Твою на всяку нощь. Молитвами Богородицы, Спасе, спаси нас. 

\myemph{Хор 2-й:} Яко прав Господь Бог наш, и несть неправды в Нем. Молитвами Богородицы, Спасе, спаси нас.

\myemph{Хор 1-й:} Слава Отцу и Сыну и Святому Духу: Молитвами Богородицы, Спасе, спаси нас.

\myemph{Хор 2-й:} И ныне и присно и во веки веков. Аминь. Молитвами Богородицы, Спасе, спаси нас.

\myparagraph{Антифон 2-й}

\myemph{Хор 1-й:} Господь воцарися, в лепоту облечеся. Молитвами святых Твоих, Спасе, спаси нас. 

\myemph{Хор 2-й:} Господь воцарися, в лепоту облечеся, облечеся Господь в силу, и препоясася. Молитвами святых Твоих, Спасе, спаси нас.

\myemph{Хор 1-й:} Ибо утверди вселенную, яже не подвижится. Молитвами святых Твоих, Спасе, спаси нас.

\myemph{Хор 2-й:} Свидения Твоя уверишася зело: дому Твоему подобает святыня, Господи, в долготу дний. Молитвами святых Твоих, Спасе, спаси нас.

\myemph{Хор:} Слава Отцу и Сыну и Святому Духу. 

\myemph{Хор:} И ныне и присно и во веки веков. Аминь. 

\myparagraph{ Песнь Господу Иисусу Христу}

\myemph{Хор:} Единородный Сыне и Слове Божий, Безсмертен Сый, и изволивый спасения нашего ради воплотитися от Святыя Богородицы и Приснодевы Марии, непреложно вочеловечивыйся; распныйся же, Христе Боже, смертию смерть поправый, един Сый Святыя Троицы, спрославляемый Отцу и Святому Духу, спаси нас. 

\myparagraph{Антифон 3-й}

\myemph{Хор 1-й:} Приидите возрадуемся Господеви, воскликнем Богу Спасителю нашему. Спаси ны, Сыне Божий, во святых дивен сый, поющия Ти: аллилуиа. 

\myemph{Хор 2-й:} Предварим лице Его во исповедании, и во псалмех воскликнем Ему: Спаси ны, Сыне Божий, во святых дивен сый, поющия Ти: аллилуиа. 

\myemph{Хор 1-й:} Яко Бог Велий Господь, и Царь Велий по всей земли. Спаси ны, Сыне Божий, во святых дивен сый, поющия Ти: аллилуиа. 

\myemph{Хор 2-й:} Яко в руце Его вси концы земли, и высоты гор Того суть. Спаси ны, Сыне Божий, во святых дивен сый, поющия Ти: аллилуиа. 

\myemph{Хор 1-й:} Яко Того есть море, и Той сотвори е, и сушу руце Его создаете. Спаси ны, Сыне Божий, во святых дивен сый, поющия Ти: аллилуиа. 

\mysubtitle{Вход с Евангелием} 

\myemph{Вход с Евангелием. Дьякон заходит в алтарь, открывает Царские врата, вместе со священником крестится и целует престол и берёт евангелие, алтарник в этот момент крестится с ними синхронно, кланяется, горнему месту, священнику и в момент перехода священника от престола к горнему месту идёт к северным вратам. Когда священник с диаконом тоже направятся к вратам, открывает дверь и по амвону проходит до царских врат, затем сворачивает к аналою и становится перед ним спиной к народу, когда священник зайдёт в алтарь, алтарник заходит через южные врата. В алтаре понамарь проходит до горнего места, крестится, кланяется горнему месту, священнику и проходит, чтобы поставить свечу на место.}

\myemph{Диакон:} Премудрость, пр\'{о}сти.

\myemph{Хор:} Приидите, поклонимся и припадем ко Христу. Спаси Сыне Божий, воскресый из мертвых, поющия Ти: аллилуиа.

\myemph{[B таком виде это песнопение поется во все обычные воскресенья, на Пасху и все дни пасхальной седмицы. Вместо слов «воскресый из мертвых» в будни поется «во святых Дивен сын», а на праздники "--- по смыслу праздника, как указано в Церковном уставе: на Рождество "--- «рождейся от Девы»; на Крещение "--- «во Иордане крестивыйся»; на Вознесение "--- «вознесыйся ко славе» в Пятидесятницу и в День Святого Духа "--- «Спаси ны. Утешителю Благий»; на Преображение "--- «преобразивыйся на горе»; на Воздвижение "--- «плотию распныйся»; в Неделю ваий "--- «возседый на жребя». В праздники Богородицы "--- «молитвами Богородицы». Праздничные входные стихи поются и в дни попразднства, до отдания.]}

\mysubtitle{Тропари и кондаки «по входе»} 

\myemph{Хор поет тропари и кондаки «по входе», назначенные в этот день церковным Уставом (воскресные тропари и кондаки приведены в главе «Песнопения из служб воскресных», дневные "--- в главе «Песнопения из служб будничных», общие ликам святых "--- в главе «Песнопения из служб общих ликам святых», праздничные "--- в главе «Песнопения из служб праздничных»).}

\myemph{Иерей:} Яко Свят еси, Боже наш, и Тебе славу возсылаем, Отцу и Сыну и Святому Духу, ныне и присно.

\myemph{Диакон:} И во веки веков.

\myemph{Хор:} Аминь.

\mysubtitle{Трисвятое} 

\myemph{[В праздники Рождества Христова, Богоявления, в Лазареву и Великую субботы, во все дни пасхальной седмицы и в период Пятидесятницы вместо Трисвятого поется: «Елицы во Христа крестистеся, во Христа облекостеся. Аллилуя». В праздник Воздвижения Креста Господня и в Неделю крестопоклонную поется: «Кресту Твоему поклоняемся Владыко, и Святое Воскресение Твое славим»]}

\myemph{Хор:} Святый Боже, Святый Крепкий, Святый Безсмертный, помилуй нас. \myemph{(Трижды)}

\myemph{Хор:} Слава Отцу и Сыну и Святому Духу, и ныне и присно и во веки веков. Аминь.

\myemph{Хор:} Святый Безсмертный, помилуй нас.

\myemph{Хор:} Святый Боже, Святый Крепкий, Святый Бессмертный, помилуй нас.

\mysubtitle{Прокимен} 

\myemph{Дьякону подаётся кадило}

\myemph{Диакон:} Вонмем.

\myemph{Иерей:} Мир всем.

\myemph{Чтец Апостола:} И духови твоему. Прокимен. Псалом Давидов, глас..

\myemph{[В Богородичные праздники: «Прокимен, песнь Богородицы: Величит душа Моя Господа и возрадовася дух Мой о Бозе Спасе Моем».]}

\myemph{Произносится один или два прокимна, назначенные в этот день на литургии церковным Уставом (воскресные прокимны со своими стихами приведены в главе «Песнопения из служб воскресных восьми гласов», дневные (будничные) "--- в главе «Песнопения из служб будничных», из служб Триодей постной и цветной "--- в главах «Песнопения из служб Триоди постной» и «Песнопения из служб Триоди цветной».}

\myemph{Чтец произносит Прокимен, называя глас его, хор поет прокимен, чтец произносит стих, хор повторяет прокимен, чтец произносит первую половину прокимна, хор поет вторую половину его. Когда Устав назначает два прокимна, первый поется дважды, т.~е. чтец: прокимен, хор: прокимен, чтец: стих, хор: прокимен, затем чтец произносит второй прокимен, и хор поет его один раз.}

\mysubtitle{Прокимны и аллилуиарии воскресные на литургии}

\myemph{Глас 1-й:} Буди, Господи, милость Твоя на нас, якоже уповахом на Тя.

\myemph{Стих:} Радуйтеся, праведнии, о Господе, правым подобает похвала.

\myemph{Аллилуиа:} Бог даяй отмщение мне и покоривый люди под мя.

\myemph{Стих:} Величай спасения царева и творяй милость Христу Своему Давиду и семени его до века.

\myemph{Глас 2-й:} Крепость моя и пение мое Господь. и бысть мне во спасение.

\myemph{Стих:} Наказуя наказа мя Господь, смерти же не предаде мя.

\myemph{Аллилуиа:} Услышит тя Господь в день печали, защитит тя имя Бога Иаковля.

\myemph{Стих:} Господи, спаси царя и услыши ны, в оньже аще день призовем Тя.

\myemph{Глас 3-й:} Пойте Богу нашему, пойте пойте Цареви нашему, пойте.

\myemph{Стих:} Вси языцы, восплещите руками, воскликните Богу гласом радования.

\myemph{Аллилуиа:} На Тя, Господи, уповах, да не постыжуся во век.

\myemph{Стих:} Буди ми в Бога Защитителя и в дом прибежища, еже спасти мя.

\myemph{Глас 4-й:} Яко возвеличишася дела Твоя, Господи, вся премудростию сотворил еси.

\myemph{Стих:} Благослови, душе моя, Господа, Господи Боже мой, возвеличился еси зело.

\myemph{Аллилуиа:} Наляцы и успевай и царствуй, истины ради и кротости, и правды.

\myemph{Стих:} Возлюбил еси правду и возненавидел еси беззконие.

\myemph{Глас 5-й:} Ты, Господи, сохраниши ны и соблюдеши ны от рода сего и во век.

\myemph{Стих:} Спаси мя, Господи, яко оскуде преподобный.

\myemph{Аллилуиа:} Милости Твоя, Господи, во век воспою, в род и род возвещу истину Твою усты моими.

\myemph{Стих:} Зане рекл еси: в век милость созиждется, на небесех уготовится истина Твоя.

\myemph{Глас 6-п:} Спаси, Господи, люди Твоя и благослови достояние Твое.

\myemph{Стих:} К Тебе, Господи, воззову, Боже мой, да не премолчиши от мене.

\myemph{Аллилуиа:} Живый в помощи Вышняго, в крове Бога Небеснаго водворится.

\myemph{Стих:} Речет Господеви: Заступник мой еси и Прибежище мое, Бог мой, и уповаю на Него.

\myemph{Глас 7-й:} Господь крепость людем Своим даст Господь благословит люди Своя миром.

\myemph{Стих:} Принесите Господеви сынове Божии, принесите Господеви сыны овни,

\myemph{Аллилуиа:} Благо есть исповедатися Господеви и пети Имени Твоему, Вышний.

\myemph{Стих:} Возвещати заутра милость Твою, и истину Твою на всяку нощь.

\myemph{Глас 8-й:} Помолитеся и воздадите Господеви Богу нашему.

\myemph{Стих:} Ведом во Иудеи Бог, во Израили велие Имя Его.

\myemph{Аллилуиа:} Приидите, возрадуемся Господеви, воскликнем Богу Спасителю нашему.

\myemph{Стих:} Предварим лице Его во исповедании, и во псалмех воскликнем Ему.

\mysubtitle{Прокимны и аллилуиарии дневные (будничные)}

\myemph{В понедельник, гл. 4-й:} Творяй ангелы Своя духи, и слуги Своя пламень огненный.

\myemph{Стих:} Благослови, душе моя. Господа, Господи Боже мой, возвеличился еси зело.

\myemph{Аллилуиа, гл. 5-й:} Хвалите Господа, вси ангели Его, хвалите Его, вся силы Его.

\myemph{Стих:} Яко Той рече, и быша; Той повеле, и создашася.

\myemph{Во вторник, гл. 7-й:} Возвеселится праведник о Господе и уповает на Него.

\myemph{Стих:} Услыши, Боже, глас мой, внегда молитися ми к Тебе.

\myemph{Аллилуиа, гл. 4-й:} Праведник яко финикс процветет, яко кедр, иже в Ливане, умножится

\myemph{Стих:} Насаждени в дому Господни, во дворех Бога нашего процветут.

\myemph{В среду, гл. 3-й:} Величит душа Моя Господа, и возрадовася дух Мой о Бозе Спасе Моем

\myemph{Стих:} Яко призре на смирение рабы Своея, се бо отныне ублажат Мя вси роди.

\myemph{Аллилуиа, гл. 8-й:} Слыши, Дщи, и виждь, и приклони ухо Твое.

\myemph{Стих:} Лицу Твоему помолятся богатии людстии.

\myemph{В четверг, гл. 8-й:} Во всю землю изыде вещание их, и в концы вселенныя глаголы их.

\myemph{Стих:} Небеса поведают славу Божию, творение же руку Его возвещает твердь.

\myemph{Аллилуиа, гл. 1-й;} Исповедят небеса чудеса, Господи, ибо истину Твою в Церкви святых

\myemph{Стих:} Бог прославляем в совете святых.

\myemph{В пятницу, гл. 7-й:} Возносите Господа Бога нашего, и покланяйтеся подножию ногу Его, яко свято есть.

\myemph{Стих:} Господь воцарися, да гневаются людие.

\myemph{Аллилуиа, гл. 1-й:} Помяни сонм Твой, егоже стяжал еси исперва.

\myemph{Стих:} Бог же Царь наш прежде века, содела спасение посреди земли. 

\myemph{В субботу, гл. 8-й:} Веселитеся о Господе, и радуйтеся, праведнии.

\myemph{Стих:} Блажени, ихже оставишася беззакония и ихже прикрышася греси.

\myemph{Заупокойный, гл. 6-й:} Души их во благих водворятся.

\myemph{Аллилуиа, гл. 4-й:} Воззваша праведнии, и Господь услыша их, и от всех скорбей их избави их.

\myemph{Стих:} Многи скорби праведным, и от всех их избавит я Господь.

\myemph{Стих:} Блажени, яже избрал и приял еси, Господи, и память их в род и род.

\myemph{Диакон:} Премудрость.

\myemph{Чтец:} Деяний святых апостол чтение. \myemph{Или:} Соборнаго послания Петрова \myemph{[или:} Иоаннова\myemph{, причем не принято говорить, какое это послание "--- первое, или второе, или третье}] чтение. \myemph{Или:} К римляном [К коринфяном; К галатом; К Тимофею\myemph{и т.п.}] послания святаго апостола Павла чтение.

\myemph{Диакон:} Вонмем.

\mysubtitle{Чтение Апостола} 

\myemph{Во время чтения апостола, ставится аналой на амвоне для Евангелия. Когда чтение закончится, иерей говорит чтецу:} Мир ти.

\myemph{Чтец:} И духови твоему.

\mysubtitle{Аллилуиа} 

\myemph{Диакон:} Премудрость.

\myemph{Чтец:} Аллилуиа, глас… \myemph{Если прислуживает один алтарник, то выносится пономарская свеча и ставится перед аналоем (с Евангелием), если два алтарника то во время пения аллилуя вдвоём подходят на горнее место со свечами, синхронно крестятся, кланяются горнему месту, священнику, друг другу, и выходят на амвон северными и южными вратами, до чтения евангелия стоят лицом к иконостасу, не кланяясь и не крестясь, в начале чтения поворачиваются лицом к евангелию, по окончанию кланяются иконам и заходят теми же вратами в алтарь, также крестятся и кланяются на горнем месте и проходят, чтобы поставить свечи на место. Не забудьте убрать аналой.}

\myemph{Хор поет «Аллилуиа» "--- трижды на указанный глас, чтец произносит первый стих аллилуиария, хор: «Аллилуиа», чтец произносит второй стих аллилуиария, хор поет в третий раз «Аллилуиа». В богослужебных книгах перед первым стихом аллилуария пишется «Аллилуиа, глас…», а перед вторым "--- «Стих» (воскресные аллилуиарии приведены в главе «Песнопения из служб воскресных восьми гласов», дневные (будничные) "--- в главе «Песнопения из служб будничных», аллилуиарии из служб Триодей постной и цветной "--- в главах «Песнопения из служб Триоди постной» и «Песнопения из служб Триоди цветной».)} 

\myemph{Диакон:} Благослови, владыко, благовестителя святаго апостола и евангелиста \myemph{(имя евангелиста)}.

\myemph{Иерей, благословляя его, произносит:} Бог, молитвами святаго, славнаго, всехвальнаго апостола и евангелиста \myemph{(имярек)}, да даст тебе глагол, благовествующему силою многою, во исполнение Евангелиа Возлюбленнаго Сына Своего, Господа нашего Иисуса Христа.

\myemph{Диакон:} Аминь.

\myemph{Иерей:} Премудрость, пр\'{о}сти, услышим святаго Евангелиа. Мир всем.

\myemph{Хор:} И духови твоему.

\myemph{Диакон:} От \myemph{(имя)} святаго Евангелиа чтение.

\myemph{Хор}: Слава Тебе, Господи, слава Тебе.

\myemph{Иерей:} Вонмем.

\mysubtitle{Чтение Евангелия}

\myemph{Читается Евангелие. Церковный Устав назначает определенные евангельские чтения на каждый день (евангельские чтения Пресвятой Богородице о общие ликам святых приведены в главе «Песнопения из служб общих ликам святых»).}

\myemph{По окончании чтения хор:} Слава Тебе, Господи, слава Тебе.

\myemph{Выносятся записки о здравии и о упокоении.}

\mysubtitle{Ектения сугубая}

\myemph{Диакон:} Рцем вси от всея души, и от всего помышления нашего рцем.

\myemph{Хор:} Господи, помилуй.

\myemph{Диакон:} Господи Вседержителю, Боже отец наших, молим Ти ся, услыши и помилуй.

\myemph{Хор:} Господи, помилуй.

\myemph{Диакон:} Помилуй нас, Боже, по велицей милости Твоей, молим Ти ся, услыши и помилуй.

\myemph{Хор:} Господи, помилуй. \myemph{(Трижды)} 

\myemph{Диакон:} Еще молимся о Великом Господине и Отце нашем Святейшем Патриархе \myemph{(имярек)}, и о Господине нашем Преосвященнейшем митрополите \myemph{(или:} архиепископе, \myemph{или:} епископе \myemph{) (имярек)}, и всей во Христе братии нашей. 

\myemph{Хор:} Господи, помилуй. \myemph{(Трижды)} 

\myemph{Диакон:} Еще молимся о Богохранимей стране нашей, властех и воинстве ея, да тихое и безмолвное житие поживем во всяком благочестии и чистоте. 

\myemph{Хор:} Господи, помилуй. \myemph{(Трижды)} 

\myemph{Диакон:} Еще молимся о братиях наших, священицех, священномонасех и всем во Христе братстве нашем.

\myemph{Хор:} Господи, помилуй. \myemph{(Трижды)} 

\myemph{Диакон:} Еще молимся о блаженных и приснопамятных создателех святаго храма сего \myemph{(если в монастыре:} святыя обители сея), и о всех преждепочивших отцех и братиях, зде лежащих и повсюду, православных. 

\myemph{Хор:} Господи, помилуй. \myemph{(Трижды)} 

\myemph{Диакон:} Еще молимся о милости, жизни, мире, здравии, спасении, посещении, прощении и оставлении грехов рабов Божиих, братии святаго храма сего ( \myemph{если в монастыре:} святыя обители сея). 

\myemph{Хор:} Господи, помилуй. \myemph{(Трижды)} 

\myemph{Диакон:} Еще молимся о плодоносящих и добродеющих во святем и всечестнем храме сем, труждающихся, поющих и предстоящих людех, ожидающих от Тебе великия и богатыя милости.

\myemph{Хор:} Господи, помилуй. \myemph{(Трижды)} 

\myemph{Иерей:} Яко Милостив и Человеколюбец Бог еси, и Тебе славу возсылаем, Отцу и Сыну и Святому Духу, ныне и присно и во веки веков.

\myemph{Хор}: Аминь.

\mysubtitle{Ектения заупокойная} 

\myemph{[В некоторые дни церковного года (кроме двунадесятых и храмовых праздников) за сугубой ектенией читается следующая ектения об усопших, при открытых царских вратах и с кадильницей:]}

\myemph{Диакон:} Помилуй нас, Боже, по велицей милости Твоей, молим Ти ся, услыши и помилуй.

\myemph{Хор:} Господи помилуй. \myemph{(трижды)}. 

\myemph{Диакон:} Еще молимся о упокоении душ усопших рабов Божиих \myemph{(имена)} и о еже проститися им всякому прегрешению, вольному же и невольному.

\myemph{Хор:} Господи помилуй. \myemph{(трижды)}. 

\myemph{Диакон:} Яко да Господь Бог учинит души их, идеже праведнии упокояются.

\myemph{Хор:} Господи помилуй. \myemph{(трижды)}. 

\myemph{Диакон:} Милости Божия, Царства Небеснаго и оставления грехов их у Христа, Безсмертнаго Царя и Бога нашего, просим.

\myemph{Хор:} Подай, Господи.

\myemph{Диакон:} Господу помолимся.

\myemph{Хор:} Господи, помилуй.

\myemph{Иерей:} Яко Ты еси воскресение, и живот, и покой усопших раб Твоих \myemph{(имена рек)}, Христе Боже наш, и Тебе славу возсылаем, со Безначальным Твоим Отцем и Пресвятым и Благим и Животворящим Твоим Духом, ныне и присно и во веки веков.

\myemph{Хор}: Аминь. \myemph{Царские врата закрываются.}

\mysubtitle{Ектения об оглашенных} 

\myemph{Диакон:} Помолитеся, оглашеннии, Господеви.

\myemph{Хор:} Господи, помилуй.

\myemph{Диакон:} Вернии, о оглашенных помолимся, да Господь помилует их.

\myemph{Хор:} Господи, помилуй.

\myemph{Диакон:} Огласит их словом истины.

\myemph{Хор:} Господи, помилуй.

\myemph{Диакон:} Открыет им Евангелие правды.

\myemph{Хор:} Господи, помилуй.

\myemph{Диакон:} Соединит их Святей Своей, Соборней и Апостольстей Церкви.

\myemph{Хор:} Господи, помилуй.

\myemph{Диакон:} Спаси, помилуй, заступи и сохрани их, Боже, Твоею благодатию.

\myemph{Хор:} Господи, помилуй.

\myemph{Диакон:} Оглашеннии, главы ваша Господеви приклоните.

\myemph{Хор:} Тебе, Господи.

\myemph{Иерей:} Да и тии с нами славят пречестное и великолепое Имя Твое, Отца и Сына и Святаго Духа, ныне и присно и во веки веков.

\myemph{Хор:} Аминь.

\myemph{Диакон:} Елицы оглашеннии, изыдите, оглашеннии, изыдите; елицы оглашеннии, изыдите. Да никто от оглашенных, елицы вернии, паки и паки миром Господу помолимся.

\myemph{Хор:} Господи, помилуй.

\myemph{Возглашением диакона: «Оглашеннии, изыдите…» заканчивается вторая часть литургии.}

\end{mymulticols}

\mychapterending

\mychapter{Литургия верных}
%http://www.molitvoslov.org/node/456 

\textbf{Литургия верных} "--- третья, самая важная часть литургии, на которой Св. Дары, приготовленные на проскомидии, силою и действием Святого Духа пресуществляются в Тело и Кровь Христовы и возносятся в спасительную для людей жертву Богу Отцу, а затем преподаются верующим для причащения. Эта часть литургии получила название оттого, что присутствовать при ее совершении и приступать к причащению Св. Тайн могут только верные, то есть лица, принявшие православную веру через Св. Крещение и оставшиеся верными обетам, данным при Св. Крещении.

На литургии верных воспоминаются страдания Господа Иисуса Христа, Его смерть, погребение, Воскресение, Вознесение на небо, седение одесную Бога Отца и второе славное пришествие на землю.

В состав этой части литургии входят важнейшие священнодействия:

\begin{enumerate}

\item Перенесение Честных Даров с жертвенника на престол, приготовление верующих молитвенному участию при совершении Бескровной Жертвы.

\item Самое совершение Таинства, с молитвенным воспоминанием членов Церкви Небесной и земной. 

\item Приготовление к причащению и причащение священнослужителей и мирян.

\item Благодарение за причащение и благословение на исход из храма (отпуст).

\end{enumerate}

\begin{mymulticols}

\mysubtitle{Ектении}

\myemph{Диакон от лица верных произносит две ектении:}

\myemph{Диакон:} Заступи, спаси, помилуй и сохрани нас, Боже, Твоею благодатию.

\myemph{Хор:} Господи, помилуй.

\myemph{Диакон:} Премудрость.

\myemph{Иерей:} Яко подобает Тебе всякая слава, честь и поклонение, Отцу и Сыну и Святому Духу, ныне и присно и во веки веков.

\myemph{Хор:} Аминь.

\myemph{Готовится кадило, зажигаются понамарская свеча.}

\myemph{Диакон:} Паки и паки миром Господу помолимся.

\myemph{Хор:} Господи, помилуй.

\myemph{Диакон:} О свышнем мире и спасении душ наших Господу помолимся. 

\myemph{Хор:} Господи, помилуй.

\myemph{Диакон:} О мире всего мира, благостоянии святых Божиих Церквей и соединении всех Господу помолимся. 

\myemph{Хор:} Господи, помилуй.

\myemph{Диакон:} О святем храме сем и с верою, благоговением и страхом Божиим входящих в онь Господу помолимся. 

\myemph{Хор:} Господи, помилуй.

\myemph{Диакон:} О избавитися нам от всякия скорби, гнева и нужды Господу помолимся. 

\myemph{Хор:} Господи, помилуй.

\myemph{Диакон:} Заступи, спаси, помилуй и сохрани нас, Боже, Твоею благодатию. 

\myemph{Диакон:} Премудрость. 

\myemph{Иерей:} Яко да под державою Твоею всегда храними, Тебе славу возсылаем, Отцу и Сыну и Святому Духу, ныне и присно и во веки веков. 

\myemph{(отверзаются царские врата.)}

\myemph{Подаётся кадило. Алтарник становится на горнем месте так, чтобы не мешать каждению. Когда каждение закончится и дьякон зайдёт в алтарь, синхронно крестится и кланяется со священнослужителями, на третий раз кланяется как обычно (горнее место, священник) и проходит к северным вратам. По сигналу священника открывает дверь и выходит как обычно к аналою. Стоит перед аналоем до тех пор пока не закроются Царские врата. По обычаю заходит в Алтарь}. 

\myemph{Хор:} Аминь, \myemph{и поет Херувимскую песнь} 

\mysubtitle{Херувимская песнь} 

Иже Херувимы тайно образующе и Животворящей Троице Трисвятую песнь припевающе, всякое ныне житейское отложим попечение… 

\mysubtitle{ Великий вход} 

\myemph{Диакон и иерей, взяв Святые Дары, выходят из алтаря на солею.}

\myemph{Диакон:} Великаго Господина и Отца нашего \myemph{(имярек),} Святейшаго Патриарха Московскаго и всея Руси, и Господина нашего Преосвященнейшаго \myemph{(имя епархиального архиерея),} да помянет Господь Бог во Царствии Своем, всегда, ныне и присно и во веки веков.

\myemph{Иерей:} Вас и всех православных христиан да помянет Господь Бог во Царствии Своем, всегда, ныне и присно и во веки веков.

\myemph{Хор:} Аминь. Яко да Царя всех подымем, ангельскими невидимо дароносима чинми. Аллилуиа, аллилуиа, аллилуиа.

\myemph{[ Вместо Херувимской на литургии в Великий четверг поется} «Вечери Твоея Тайныя…», \myemph{а в Великую субботу "---} «Да молчит всякая плоть…» \myemph{(эти песнопения приведены в главе «Песнопения из служб Триоди постной»).]}

\myemph{Режутся просфоры для причастников}. 

\mysubtitle{Просительная ектения} 

\myemph{Диакон:} Исполним молитву нашу Господеви. 

\myemph{Хор:} Господи, помилуй.

\myemph{Диакон:} О предложенных честных Дарех Господу помолимся. 

\myemph{Хор:} Господи, помилуй.

\myemph{Диакон:} О святем храме сем и с верою, благоговением и страхом Божиим входящих в онь Господу помолимся. 

\myemph{Хор:} Господи, помилуй.

\myemph{Диакон:} О избавитися нам от всякия скорби, гнева и нужды Господу помолимся. 

\myemph{Хор:} Господи, помилуй.

\myemph{Диакон:} Заступи, спаси, помилуй и сохрани нас, Боже, Твоею благодатию. 

\myemph{Хор:} Господи, помилуй.

\myemph{Диакон:} Дне всего совершенна, свята, мирна и безгрешна у Господа просим. 

\myemph{Хор:} Подай, Господи. 

\myemph{Диакон:} Ангела мирна, верна наставника, хранителя душ и телес наших у Господа просим. 

\myemph{Хор:} Подай, Господи. 

\myemph{Диакон:} Прощения и оставления грехов и прегрешений наших у Господа просим. 

\myemph{Хор:} Подай, Господи. 

\myemph{Диакон:} Добрых и полезных душам нашим и мира мирови у Господа просим. 

\myemph{Хор:} Подай, Господи. 

\myemph{Диакон:} Прочее время живота нашего в мире и покаянии скончати, у Господа просим. 

\myemph{Хор:} Подай, Господи. 

\myemph{Диакон:} Христианския кончины живота нашего, безболезнены, непостыдны, мирны, и добраго ответа на Страшном Судищи Христове просим. 

\myemph{Хор:} Подай, Господи. 

\myemph{Диакон:} Пресвятую, Пречистую, Преблагословенную, Славную Владычицу нашу Богородицу и Приснодеву Марию, со всеми святыми помянувше, сами себе и друг друга, и весь живот наш Христу Богу предадим.

\myemph{Хор:} Тебе, Господи.

\myemph{Иерей:} Щедротами Единороднаго Сына Твоего, с Нимже благословен еси, со Пресвятым и Благим и Животворящим Твоим Духом, ныне и присно, и во реки веков.

\myemph{Хор:} Аминь.

\myemph{Иерей:} Мир всем.

\myemph{Хор:} И духови твоему.

\myemph{Диакон:} Возлюбим друг друга, да единомыслием исповемы.

\myemph{Хор:} Отца и Сына и Святаго Духа, Троицу Единосущую и Нераздельную.

\myemph{Диакон:} Двери, двери, премудростию вонмем. \myemph{(Открывается завеса царских врат.)}

\myemph{Ставится кипятиться чайник.} 

\mysubtitle{ Символ веры} 

\myemph{Хор (или все молящиеся):}

\begin{enumerate}

\item Верую во Единаго Бога Отца Вседержителя, Творца небу и земли, видимым же всем и невидимым. 

\item И во Единаго Господа Иисуса Христа, Сына Божия, Единороднаго, Иже от Отца рожденнаго прежде всех век. Света от Света, Бога истинна от Бога истинна, рожденна, несотворенна, единосущна Отцу, Имже вся быша.

\item Нас ради, человек, и нашего ради спасения сшедшаго с Небес, и воплотившагося от Духа Свята и Марии Девы, и вочеловечшася.

\item Распятаго же за ны при Понтийстем Пилате, и страдавша, и погребенна.

\item И воскресшаго в третий день по Писанием.

\item И восшедшаго на Небеса, и седяща одесную Отца.

\item И паки грядущаго со славою судити живым и мертвым, Его же Царствию не будет конца.

\item И в Духа Святаго, Господа Животворящаго, Иже от Отца исходящаго, Иже со Отцем и Сыном споклоняема и сславима, глаголавшаго пророки.

\item Во едину Святую Соборную и Апостольскую Церковь.

\item Исповедую едино Крещение во оставление грехов.

\item Чаю воскресения мертвых,

\item и жизни будущаго века. Аминь. 

\end{enumerate}

\mysubtitle{ Евхаристический канон.} 

\myemph{Диакон:} Станем добре, станем со страхом, вонмем, Святое Возношение в мире приносити.

\myemph{Хор:} Милость мира, Жертву хваления. 

\myemph{Иерей:} Благодать Господа нашего Иисуса Христа, и любы Бога и Отца, и причастие Святаго Духа, буди со всеми вами.

\myemph{Хор:} И со духом твоим. 

\myemph{Иерей:} Горе имеим сердца. 

\myemph{Хор:} Имамы ко Господу. 

\myemph{Иерей:} Благодарим Господа. 

\myemph{Хор:} Достойно и праведно есть покланятися Отцу и Сыну и Святому Духу, Троице Единосущной и Нераздельней.

\myemph{Иерей:} Победную песнь поюще, вопиюще, взывающе и глаголюще:

\myemph{Хор:} Свят, Свят, Свят Господь Саваоф, исполнь Небо и земля славы Твоея; осанна в вышних, благословен Грядый во Имя Господне, осанна в вышних.

\myemph{Иерей:} Приимите, ядите, Сие есть Тело Мое, еже за вы ломимое во оставление грехов.

\myemph{Хор:} Аминь.

\myemph{Иерей:} Пийте от нея вси, сия есть Кровь Моя Новаго Завета, яже за вы и за многи изливаемая во оставление грехов.

\myemph{Хор:} Аминь. 

\myemph{(На литургии св. Василия Великого последние возгласа иерея начинаются словами: «Даде святым Своим учеником и апостолом, рек:».)}

\myemph{Иерей:} Твоя от Твоих Тебе приносяще о всех и за вся. \myemph{Готовится кадило}.

\myemph{Хор:} Тебе поем, Тебе благословим, Тебе благодарим, Господи, и молим Ти ся. Боже наш. \myemph{Подаётся кадило во время «Тебе поем…», после слов священника в алтаре «Приложив Духом Твоим Святым. Аминь. Аминь. Аминь.»}

\myemph{Иерей:} Изрядно о Пресвятей, Пречистей, Преблагословенней, Славней Владычице нашей Богородице и Приснодеве Марии.

\myemph{Хор:} Достойно есть, яко воистинну блажити Тя, Богородицу, Присноблаженную и Пренепорочную и Матерь Бога нашего. Честнейшую Херувим и Славнейшую без сравнения Серафим, без истления Бога Слова рождшую, сущую Богородицу Тя величаем. 

\myemph{[В двунадесятые праздники и их попразднства вместо «Достойно…» поется припев и ирмос 9-й песни праздничного канона, так называемый «задостойник». В Великий четверг поется ирмос 9-й песни «Странствия Владычня…», в Великую субботу "--- «Не рыдай Мене, Мати…», в Неделю ваий "--- «Бог Господь…» (эти песнопения приведены в главах «Песнопения из служб Триоди постной» и «Песнопения из служб Триоди цветной»).}

\myemph{Если же литургия св. Василия Великого, вместо «Достойно…» поем:} 

О Тебе радуется, Благодатная, всякая тварь, ангельский собор и человеческий род, освященный храме и раю словесный, девственная похвало, из Неяже Бог воплотися и Младенец бысть, прежде век сый Бог наш; ложесна бо Твоя престол сотвори и чрево Твое пространнее небес содела. О Тебе радуется, Благодатная, всякая тварь, слава Тебе. \myemph{]}

\myemph{Иерей:} В первых помяни Господи, Великаго Господина и Отца нашего \myemph{(имярек),} Святейшаго Патриарха Московскаго и всея Руси, и Господина нашего Преосвященнейшаго \myemph{(имя епархиального епископа)}, ихже даруй святым Твоим Церквам в мире, целых, честных, здравых, долгоденствующих, право правящих слово Твоей истины. 

\myemph{Хор:} И всех и вся.

\myemph{Иерей:} И даждь нам единеми усты и единем сердцем славити и воспевати Пречестное и Великолепное Имя Твое, Отца и Сына и Святаго Духа, ныне и присно и во веки веков. 

\myemph{Хор:} Аминь.

\myemph{Иерей:} И да будут милости Великаго Бога и Спаса нашего Иисуса Христа со всеми вами.

\myemph{Хор:} И со духом твоим.

\myemph{Готовится чаша для теплоты и плат для причастия}. 

\mysubtitle{ Ектения просительная} 

\myemph{Диакон:} Вся святыя помянувше, паки и паки миром Господу помолимся. 

\myemph{Хор:} Господи, помилуй.

\myemph{Диакон:} О принесенных и освященных Честных Дарех, Господу помолимся. 

\myemph{Хор:} Господи, помилуй.

\myemph{Диакон:} Яко да Человеколюбец Бог наш, приемь я во святый, и пренебесный, и мысленный Свой Жертвенник, в воню благоухания духовнаго, возниспослет нам Божественную благодать и дар Святаго Духа, помолимся. 

\myemph{Хор:} Господи, помилуй.

\myemph{Диакон:} О избавится нам от всякия скорби, гнева и нужды, Господу помолимся. 

\myemph{Хор:} Господи, помилуй.

\myemph{Диакон:} Заступи, спаси, помилуй и сохрани нас, Твоею благодатию. 

\myemph{Хор:} Господи, помилуй.

\myemph{Диакон:} Дне всего совершенна, свята, мирна и безгрешна, у Господа просим. 

\myemph{Хор:} Подай, Господи.

\myemph{Диакон:} Ангела мирна, верна наставника, хранителя душ и телес наших, у Господа просим. 

\myemph{Хор:} Подай, Господи.

\myemph{Диакон:} Прощения и оставления грехов и прегрешений наших, у Господа просим. 

\myemph{Хор:} Подай, Господи.

\myemph{Диакон:} Добрых и полезных душам нашим и мира мирови, у Господа просим. 

\myemph{Хор:} Подай, Господи.

\myemph{Диакон:} Прочее время живота нашего в мире и покаянии скончати, у Господа просим. 

\myemph{Хор:} Подай, Господи.

\myemph{Диакон:} Христианския кончины живота нашего, безболезнены, непостыдны, мирны, и добраго ответа на Страшнем Судищи Христове, просим. 

\myemph{Хор:} Подай, Господи.

\myemph{Диакон:} Соединение веры и причастие Святаго Духа испросивше, сами себе, и друг друга, и весь живот наш Христу Богу предадим. 

\myemph{Хор:} Тебе, Господи. 

\myemph{Иерей:} И сподоби нас, Владыко, со дерзновением, неосужденно смети призывати Тебе, Небеснаго Бога Отца, и глаголати: 

\mysubtitle{ Отче наш} 

\myemph{Хор (или все молящиеся):} Отче наш, Иже еси на Небесех! Да святится Имя Твое, да приидет Царствие Твое, да будет воля Твоя, яко на небеси и на земли. Хлеб наш насущный даждь нам днесь, и остави нам долги наша, якоже и мы оставляем должником нашим; и не введи нас во искушение, но избави нас от лукаваго. 

\myemph{Иерей:} Яко Твое есть Царство, и сила, и слава. Отца и Сына и Святаго Духа, ныне и присно и во веки веков. 

\myemph{Хор:} Аминь. 

\myemph{Иерей:} Мир всем. 

\myemph{Хор:} И духови твоему.

\myemph{Диакон:} Главы ваша Господеви приклоните.

\myemph{Хор:} Тебе, Господи.

\myemph{Иерей:} Благодатию, и щедротами, и человеколюбием Единороднаго Сына Твоего, с Нимже благословен еси, со Пресвятым и Благим и Животворящим Твоим Духом, ныне и присно и во веки веков. 

\myemph{Хор:} Аминь.

\myemph{(Закрываются царские врата и завеса)}

\myemph{Диакон:} Вонмем. 

\myemph{Подносится теплота.}

\myemph{Иерей:} Святая святым. 

\myemph{Хор:} Един Свят, един Господь Иисус Христос, во славу Бога Отца. Аминь. 

\mysubtitle{ Причащение священнослужителей} 

\myemph{В алтаре причащаются священнослужители.} 

\myemph{Хор поет назначенный в этот день церковным Уставом причастен "--- стих, оканчивающихся троекратным «Аллилуиа». Причастнов может быть назначено два, однако «Аллилуия» поется только после второго.}

\mysubtitle{ Причастны} 

\myemph{Во время причастных выносится понамарская свеча и ставится перед Царскими вратами. Затем выносятся запивка и просфоры для причастников.}

\myemph{В воскресенье:} Хвалите Господа с небес, хвалите Его в вышних. Аллилуиа, аллилуиа, аллилуиа. 

\myemph{В понедельник:} Творяй ангелы Своя духи, и слуги Своя пламень огненный. 

\myemph{Во вторник:} В память вечную будет праведник, от слуха зла не убоится. 

\myemph{В среду:} Чашу спасения прииму и Имя Господне призову. 

\myemph{В четверг:} Во всю землю изыде вещание их, и в концы вселенныя глаголы их. 

\myemph{В пятницу:} Спасение соделал еси посреде земли, Боже. 

\myemph{В субботу:} Радуйтеся, праведнии, о Господе, правым подобает похвала. 

\myemph{Заупокойный:} Блажени, яже избрал и приял еси, Господи, и память их в род и род. 

\myemph{В праздники Богородицы:} Чашу спасения прииму и Имя Господне призову. 

\myemph{В праздники апостолов:} Во всю землю изыде вещание их, и в концы вселенныя глаголы их. 

\myemph{В дни памяти святых:} В память вечную будет праведник, от слуха зла не убоится. 

\myemph{Забирается свеча}.

\myemph{Открываются царские врата. Диакон, вынося Святую Чашу, возглашает:} Со страхом Божиим и верою приступите! 

\myemph{(Передает Чашу иерею.)} 

\myemph{Хор:} Благословен Грядый во Имя Господне, Бог Господь и явися нам.

\myemph{[В пасхальную седмицу вместо этого поется «Христос воскресе…».]}

\myemph{Иерей (и с ним все, желающие причаститься):} Верую, Господи, и исповедую, яко Ты еси воистину Христос, Сын Бога живаго, пришедый в мир грешныя спасти, от них же первый есмь аз. Еще верую, яко Cие есть самое Пречистое Тело Твое, и сия есть самая Честная Кровь Твоя. Молюся убо Тебе: помилуй мя, и прости ми прегрешения моя, вольная и невольная, яже словом, яже делом, яже ведением и неведением, и сподоби мя неосужденно причаститися Пречистых Твоих таинств, во оставление грехов, и в Жизнь Вечную. Аминь.

Вечери Твоея Тайныя днесь, Сыне Божий, причастника мя приими; не бо врагом Твом тайну повем, ми лобзания Ти дам, яко Ииуда, на яко разбойник исповедаю Тя: помяни мя, Господи, во Царствии Твом.

Да не в суд или во осуждение будет мне причащение Святых Таин, Господи, но во исцеление души и тела. Аминь.

\mysubtitle{Причащение мирян}

\myemph{Причащая мирян, иерей говорит:} Причащается раб Божий \myemph{(имя)} Честнаго и Святаго Тела и Крове Господа и Бога и Спаса нашего Иисуса Христа, во оставление грехов своих и в Жизнь Вечную. 

\myemph{Хор (во время причащения)}: Тело Христово приимите, Источника безсмертнаго вкусите.

\myemph{[В Великий четверг вместо этого поется «Вечери Твоея тайныя…» (это песнопение приведено в главе «Песнопения из служб Триоди постной»); в пасхальную седмицу "--- «Христос воскресе…».]}

\myemph{Иерей:} Спаси, Боже, люди Твоя и благослови достояние Твое.

\myemph{Подносится кадило в алтаре.} 

\myemph{Хор:} Видехом Свет истинный, прияхом Духа Небеснаго, обретохом веру истинную, Нераздельней Троице покланяемся: Та бо нас спасла есть.

\myemph{[Вместо «Видехом свет истинный…» от Пасхи до отдания поется «Христос воскресе…»; от Вознесения до отдания "--- тропарь Вознесения (эти песнопения приведены в главе «Песнопения из служб Триоди цветной»); в Троицкую родительскую субботу "--- «Глубиною мудрости…» (этот тропарь приведен в главе «Песнопения из служб Триоди цветной», в службе мясопустной родительской субботы).]}

\myemph{Иерей:} Всегда, ныне и присно и во веки веков.

\myemph{Хор:} Аминь. Да исполнятся уста наша хваления Твоего, Господи, яко да поем славу Твою, яко сподобил еси нас причаститися Святым Твоим, Божественным, Безсмертным и Животворящим Тайнам; соблюди нас во Твоей святыни, весь день поучатися правде Твоей. Аллилуиа, аллилуиа, аллилуиа. 

\myemph{[В Великий четверг вместо «Да исполнятся…» поется «Вечери Твоея тайныя…» (это песнопение приведено в главе «Песнопения из служб Триоди постной»); в пасхальную седмицу "--- «Христос воскресе…».]}

\mysubtitle{Ектения} 

\myemph{Диакон:} Пр\'{о}сти приимше Божественных, Святых, Пречистых, Безсмертных, Небесных и Животворящих, Страшных Христовых Тайн, достойно благодарим Господа. 

\myemph{Хор:} Господи, помилуй. 

\myemph{Диакон:} Заступи, спаси, помилуй и сохрани нас, Боже, Твоею благодатию. 

\myemph{Хор:} Господи, помилуй. 

\myemph{Диакон:} День весь совершен, свят, мирен и безгрешен испросивше, сами себе, и друг друга, и весь живот наш Христу Богу предадим. 

\myemph{Хор:} Тебе, Господи. 

\myemph{Иерей:} Яко Ты еси Освящение наше, и Тебе славу возсылаем, Отцу и Сыну и Святому Духу, ныне и присно и во веки веков,

\myemph{Хор:} Аминь.

\myemph{Иерей:} С миром изыдем.

\myemph{Хор:} О имени Господни. 

\myemph{Диакон:} Господу помолимся.

\myemph{Хор:} Господи, помилуй. 

\mysubtitle{Молитва заамвонная} 

\myemph{Иерей (стоя пред амвоном):} Благословляяй благословящия Тя, Господи, и освящаяй на Тя уповающия, спаси люди Твоя и благослови достояние Твое, исполнение Церкве Твоея сохрани, освяти любящия благолепие дому Твоему. Ты тех возпрослави Божественною Твоею силою и не остави нас, уповающих на Тя. Мир мирови Твоему даруй, Церквам Твоим, священником, воинству и всем людем Твоим. Яко всякое даяние благо и всяк дар совершен свыше есть, сходяй от Тебе, Отца Светов. И Тебе славу, и благодарение, и поклонение возсылаем, Отцу и Сыну и Святому Духу, ныне и присно, и во веки веков. 

\myemph{Хор:} Аминь. Буди Имя Господне благословено отныне и до века. \myemph{(Tрижды)}

\myemph{[На пасхальной седмице вместо этого поется «Христос воскресе…».]}

\mysubtitle{ Псалом 33} 

\myemph{Хор:} Благословлю Господа на всякое время, выну хвала Его во устех моих. О Господе похвалится душа моя. Да услышат кротции, и возвеселятся. Возвеличите Господа со мною, и вознесем Имя Его вкупе. Взысках Господа, и услыша мя, и от всех скорбей моих избави мя. Приступите к Нему и просветитеся, и лица ваша не постыдятся. Сей нищий воззва, и Господь услыша и, и от всех скорбей его спасе и. Ополчится ангел Господень окрест боящихся Его, и избавит их. Вкусите, и видите, яко благ Господь; блажен муж, иже уповает Нань. Бойтеся Господа вси святии Его, яко несть лишения боящимся Его. Богатии обнищаша и взалкаша: взыскающии же Господа не лишатся всякаго блага. Приидийте, чада, послушайте мене, страху Господню научу вас. Кто есть человек хотяй живот, любяй дни видети благи? Удержи язык твой от зла, и устне твои, еже не глаголати льсти. Уклонися от зла, и сотвори благо, взыщи мира, и пожени и. Очи Господни на праведныя и уши Его в молитву их. Лице же Господне на творящия злая, еже потребити от земли память их. Воззваша праведнии, и Господь услыша их, и от всех скорбей их избави их. Близ Господь сокрушенных сердцем, и смиренныя духом спасет. Многи скорби праведным, и от всех их избавит я Господь. Хранит Господь вся кости их, ни едина от них сокрушится. Смерть грешников люта, и ненавидящии праведнаго прегрешат. Избавит Господь души раб Своих, и не прегрешат вси уповающий на Него.

\myemph{[На пасхальной седмице вместо этого поется «Христос воскресе…».]}

\myemph{Иерей:} Благословение Господне на вас. Того благодатию и человеколюбием, всегда, ныне и присно и во веки веков.

\myemph{Хор:} Аминь.

\myemph{Иерей:} Слава Тебе, Христе Боже, Упование наше, слава Тебе.

\myemph{[На Пасху, в пасхальную седмицу и в отдание Пасхи вместо «Слава Тебе, Христе Боже…» священнослужители поют «Христос воскресе из мертвых, смертию смерть поправ», а хор заканчивает: «и сущим во гробех живот даровав».} 

\myemph{От Недели о Фоме до отдания Пасхи священник произносит: «Слава Тебе, Христе Боже, Упование наше, Слава Тебе», а хор поет «Христос воскресе…» (Трижды).]}

\myemph{Хор:} Слава Отцу и Сыну и Святому Духу. И ныне и присно и во веки веков. Аминь. 

\myemph{Хор:} Господи, помилуй \myemph{(Трижды).}

\myemph{Хор:} Благослови. 

\mysubtitle{ Отпуст} 

\myemph{Иерей произносит отпуст. В воскресенье:} Воскресый из мертвых, Христос, истинный Бог наш, молитвами Пречистыя Своея Матере, святых славных и всехвальных Апостол, иже во святых отца нашего Иоанна, архиепископа Константина града, Златоустаго ( \myemph{или:} св. Василия Великаго, архиепископа Кесарии Каппадокийския), и святаго \myemph{(храма и святого, которого память в этот день),} святых и праведных Богоотец Иоакима и Анны и всех святых, помилует и спасет нас, яко Благ и Человеколюбец.

\mysubtitle{ Многолетие} 

\myemph{Хор:} Великаго Господина и Отца нашего \myemph{(имярек)}, Святейшаго Патриарха Московскаго и всея Руси, и Господина нашего Преосвященнейшаго \myemph{(имя)} митрополита ( \myemph{или:} архиепископа, \myemph{или:} епископа) \myemph{(епархиальный титул его),} богохранимую державу нашу Российскую, настоятеля, братию и прихожан святаго храма сего и вся православныя христианы, Господи, сохрани их на многая лета.

\myemph{По обычаю, перед отпустом священник берет крест с престола и после отпуста, осенив крестом народ и сам поцеловав крест, дает его для целования молящимся, а чтец читает благодарственные молитвы; затем священник опять осеняет крестом народ и возвращается в алтарь, причем царские врата и завеса закрываются.}

\myemph{Алтарники убираются в алтаре, чистят кадило и готовятся к вечернему Богослужению.}

\end{mymulticols}

\mychapterending

%\renewcommand{\ornament}{uzor_begin_9}
\mychapter{Список изменений}
\markright{Список изменений}

Этот раздел содержит информацию об изменениях и исправлениях в материалах сайта \url{www.molitvoslov.com}, вошедших в эту и предыдущие публикации. Информация о более новых версиях располагается в начале, а о более старых "--- в конце раздела.

\mysubtitle{Версия 20160406}

Обновление текстов от 4-го апреля 2016 г.

Исправлены опечатки в следующих молитвах и разделах:

\begin{itemize}

\item 115-м псалом в Последовании ко Святому Причащению.
\item Благодарственные молитвы по Святому Причащению.
\item Величание преподобному Александру Свирскому.
\item Канон за болящего.
\item Канон преподобному Паисию Великому об избавлении от мук умерших без покаяния.
\item Моление о том, чтобы Бог даровал нам усердие к молитве за усопших и принял бы её.
\item Молитва cвятителю Нектарию Эгинскому.
\item Молитва «Вопль к Богоматери».
\item Молитва Пресвятой Богородице в честь Ее иконы «Казанская».
\item Молитва архистратигу Михаилу.
\item Молитва за всякого усопшего.
\item Молитва за усопших внезапною (скоропостижною) смертию.
\item Молитва мученику Уару.
\item Молитва преподобному Моисею Мурину.
\item Молитва св. праведного Иоанна Кронштадского об укреплении в Православной вере и единстве.
\item Молитва святителю Тихону, Патриарху Московскому и Всея Руси.
\item Молитва священномученику Ипатию, епископу Гангрскому.
\item Правило от осквернения.
\item Пяточисленные молитвы.
\item Раздел БОЖЕСТВЕННАЯ ЛИТУРГИЯ: О Литургии, Литургия оглашенных и Литургия верных.
\item Раздел ЕДИНОВЕРЧЕСКИЕ МОЛИТВЫ: Келейное правило, Павечерица великая и малая.
\item Тропарь преподобному Иоанну Дамаскину.
\item Тропарь святителю Спиридону, Тримифунтскому чудотворцу.
\item Тропарь священномученику Ипатию, епископу Гангрскому.

\end{itemize}

Добавлены тексты:

\begin{itemize}

\item Последование по исходе души от тела.
\item В раздел «Молитвы» добавлена отдельно молитва Честному Кресту (Да воскреснет Бог\ldots).
\item Молитва «При немилосердии и раздражении на ближнего» отделена от текста предыдущей молитвы.
\item В Последование ко Святому Причащению добавлена 4-я молитва святого Симеона Метафраста («Яко на Страшнем Твоем и нелицеприемнем предстояй Судилищи\ldots»).

\end{itemize}

\mysubtitle{Версия 20130227}

Обновление текстов от 27-го февраля 2013

\begin{itemize}

\item Незначительные изменения в оформлении после переезда сайта на новый домен \texttt{molitvoslov.org}.
\item Исправлена опечатка в Благодарственных молитвах по Святом Причащении.
\item Исправлена опечатка в Литургии верных.
\item Добавлено житие блаженной Ксении Петербургской.
\item Исправление в Акафисте за единоумершего.

\end{itemize}

\mysubtitle{Версия 20121214}

\begin{itemize}

\item Обновление текстов с сайта от 14-го декабря 2012 г.
\item Исправлен русско-язычный поиск, копирование и вставка из PDF в программе Adobe Acrobat Reader.
\item В текст Великого Канона включены припевы и добавлен русский перевод.

\item В текстах Божественной Литургии, канонах и молитвах добавлены положенные между стихами припевы, которые довольно часто из экономии места опускаются в печатных или интернетных изданиях, что может сбить с толку начинающих. По этой же причине часто употребляемые молитвы "--- «Слава:», «И ныне:», «Трисвятое по Отче наш», «Приидите, поклонимся...» и другие "--- приводятся каждый раз полностью, а не обозначаются сокращённо.

\item Исправлены опечатки (большая часть которых была обнаружена посетителями сайта) либо добавлены недостающие фрагменты молитв или изображения икон в следующих текстах:
\begin{itemize}
\item Молитва пред иконою «Неупиваемая Чаша».
\item Молитва пред иконою «Почаевская».
\item Пресвятой Богородице в честь Ее иконы «Казанская».
\item Молитвы задержания.
\item Молитва против антихриста.
\item Молитва о спасении державы Российской и утолении в ней раздоров и нестроений.
\item Тропарь Кресту и молитва за отечество (Молитвы утренние).
\item Молитвы на сон грядущим.
\item Последование ко Святому Причащению.
\item Единоверческие молитвы.
\item Молитва по исходе души из тела.
\item Правило от осквернения.
\item Молитва перед учением.
\item Молитва после учения.
\item Апостолу Симону Зилоту.
\item Святому праведному Иоанну, Кронштадтскому чудотворцу.
\item Молитвы ко Господу Св. Иоанна Кронштадского.
\item Преподобному Серафиму Саровскому.
\item Святым равноапостольным Кириллу и Мефодию.
\item Святой равноапостольной княгине Ольге.
\item Преподобному Даниилу Московскому.
\item Святым равноапостольным Константину и Елене.
\item Святителю Иоанну Златоустому, архиепископу Константинопольскому.
\item Поучение святителя Игнатия Брянчанинова о молитвенном правиле.
\item Молитвы ко Господу Иисусу Христу.
\item Молитва на сон.
\end{itemize}

\end{itemize}

\mysubtitle{Версия 20111108}

Обновление текстов с сайта от 8-го ноября 2011 г.

\begin{itemize}

\item Добавлен раздел «Утешение сродникам самоубийц».
\item Исправлены опечатки в следующих текстах:

\begin{itemize}

\item «Богородичное правило».
\item «Канон за болящего, глас 3-й».
\item Молитва «Пророку, Предтече и Крестителю Господню Иоанну».
\item Молитва «Святителю Димитрию, Ростовскому чудотворцу».
\item Молитва «Святой блаженной Ксении Петербургской».
\item «Молитва покаянная, читаемая в церквах России во дни смуты».
\item «КАНОН ПРЕПОДОБНОМУ ПАИСИЮ ВЕЛИКОМУ ОБ ИЗБАВЛЕНИИ ОТ МУК УМЕРШИХ БЕЗ ПОКАЯНИЯ».
\item Псалом 50-й, в разделе «ЧАСЫ ТРЕТИЙ И ШЕСТОЙ».
\item Псалом 50-й, в «Молитвы утренние».
\item Псалом 50-й и 26-й, в разделе «ЕДИНОВЕРЧЕСКИЕ МОЛИТВЫ».
\item Псалом 26-й, в разделе «МОЛИТВЕННОЕ ПРАВИЛО ПРЕПОДОБНОГО АМВРОСИЯ OПТИНСКОГО».
\item Молитва «Святому преподобному Иосифу Волоцкому».
\item «Молитва от осквернения».
\item Молитва «Великомученику Георгию Победоносцу».
\end{itemize}

\end{itemize}

\mysubtitle{Версия 20110709}

Обновление текстов с сайта от 9 июля 2011 г.

\begin{itemize}

\item Добавлен новый раздел «О Молитве», а также «Словарь и термины» (около 2500 слов и выражений).

\item Небольшое изменение в разделе «Божественная литургия».

\item Добавлены заглавия тропарей, Богородичных, канонов и песен в

\begin{itemize}

\item «Канон Ангелу Хранителю» («Молитвы»).

\item «Канон молебный ко Господу Иисусу Христу и Пречистой Богородице Матери Господни при разлучении души от тела всякого правоверного» («НАПУТСТВИЕ ХРИСТИАНИНА ПЕРЕД СМЕРТЬЮ И ЗАУПОКОЙНЫЕ МОЛИТВЫ»).

\item в разделе «КАНОН ПРЕПОДОБНОМУ ПАИСИЮ ВЕЛИКОМУ ОБ ИЗБАВЛЕНИИ ОТ МУК УМЕРШИХ БЕЗ ПОКАЯНИЯ».

\end{itemize}
\end{itemize}

\mysubtitle{Версия 20110616}

Обновление текстов с сайта от 16 июня 2011 г.
\begin{itemize}

\item В разделе «РАЗЛИЧНЫЕ МОЛИТВЫ»
\begin{itemize}
\item Исправлены опечатки в «Молитвах ко Господу Св. Иоанна Кронштадского».
\end{itemize}

\item В разделе «МОЛИТВЫ СВЯТЫМ»
\begin{itemize}
\item Убраны повторы молитв прп. Александру Свирскому и прп. Серафиму Саровскому.
\item Добавлены молитвы св. первомученику архидиакону Стефану; св. царице Грузинской Тамаре; святителю Иоанну Шанхайскому и Сан-Францисскому.
\item Исправлены опечатки в молитвах св. Царю Страстотерпцу Николаю; св. прп. Силуану Афонскому; cвятителю Тихону, Патриарху Московскому и Всея Руси.
\end{itemize}

\item В разделе «МОЛИТВЫ В СКОРБЯХ И ИСКУШЕНИЯХ ТВОРИМЫЕ»
\begin{itemize}
\item Добавлена «Молитва преследуемого человеками (свт. Игнатия Брянчанинова)».
\end{itemize}

\item В разделе «МОЛИТВЫ ЗА БОЛЯЩИХ»
\begin{itemize}
\item Исправлена опечатка в «Каноне за болящего, глас 3-й», Песнь 9.
\end{itemize}

\item В разделе «НАПУТСТВИЕ ХРИСТИАНИНА ПЕРЕД СМЕРТЬЮ И ЗАУПОКОЙНЫЕ МОЛИТВЫ»
\begin{itemize}
\item Исправлен «Чин литии, совершаемой мирянином дома и на кладбище».
\item Исправлен тропарь прп. Паисию Великому в молитве «Об ослаблении вечных мук умерших без покаяния», а также в разделе «КАНОН ПРЕПОДОБНОМУ ПАИСИЮ ВЕЛИКОМУ ОБ ИЗБАВЛЕНИИ ОТ МУК УМЕРШИХ БЕЗ ПОКАЯНИЯ».
\item Исправлен «Канон молебный ко Господу Иисусу Христу и Пречистой Богородице Матери Господни при разлучении души от тела всякого правоверного».
\end{itemize}

\item В разделе «В ПОМОЩЬ КАЮЩЕМУСЯ»
\begin{itemize}
\item Исправлены опечатки в «Десяти заповедях».
\end{itemize}


\end{itemize}


\mysubtitle{Версия 20110503.02}

Содержимое сайта от 03 мая 2011 г.

\mychapterending
%\normalfont\nopagebreak\bigskip\bigskip\begin{center}\includegraphics[width=0.15\textwidth]{uzor_end_4}\end{center}



\end{document}
